\documentclass[11pt]{article}
\usepackage[T1]{fontenc}
\usepackage{lmodern,amsmath,amssymb,amsthm,mathtools}
\usepackage[margin=1in]{geometry}
\usepackage{booktabs,microtype,hyperref}
\hypersetup{colorlinks=true,linkcolor=blue,citecolor=blue,urlcolor=blue,
pdftitle={Non-classification methods for five to eight points}}
\setlength{\emergencystretch}{2em}
\newtheorem{theorem}{Theorem}[section]
\newtheorem{proposition}[theorem]{Proposition}
\newtheorem{lemma}[theorem]{Lemma}
\newtheorem{corollary}[theorem]{Corollary}
\theoremstyle{remark}
\newtheorem{remark}[theorem]{Remark}
\newcommand{\C}{\mathbb C}
\newcommand{\R}{\mathbb R}
\newcommand{\Q}{\mathbb Q}
\newcommand{\one}{\mathbf 1}
\newcommand{\norm}[1]{\lVert#1\rVert}
\newcommand{\ip}[2]{\langle#1,#2\rangle}
\newcommand{\pos}[1]{[#1]_+}
\newcommand{\op}{\mathrm{op}}
\newcommand{\Rea}{\operatorname{Re}}
\newcommand{\diag}{\operatorname{diag}}
\numberwithin{equation}{section}
\title{Non-classification methods for five to eight points\\with unrestricted exclusion regions and sharp origin bounds}
\author{Research addendum to the supplied manuscript}
\date{October 8, 2026}
\begin{document}
\maketitle
\begin{abstract}
The primary five-point proof follows the supplied second-order moment
manuscript, using a real projection and a logarithm majorant. We extend
its finite-dimensional projection identity to six, seven, and eight
points, and derive Fourier-frame and higher second-order Gram constraints
without a stationary-shape catalogue. A direct strong-concavity proof
gives the sharp six- and seven-point inequalities whenever an origin
is present. For eight points, a seven-row spectral deficit and a complete
two-variable rational interval certificate prove the sharp bound for
every configuration containing the origin, with equality only for the
central regular heptagon. There is no restriction on the number of
radial levels or on angular symmetry. An effective small-harmonic
rigidity theorem further excludes actual eight-point local maxima in
a specified coefficient region. Finite-point conjugate closure and
distance Hankel tests explain why positive-measure relaxations can lose
essential constraints. New whole radial and angular regions are excluded
for unrestricted six, seven, and eight points. Iterated real projections
give sixth- and seventh-moment circle constraints without the unwanted
higher holomorphic moments, and enlarge eight-point small-harmonic
rigidity to a threshold of $10^{-14}$.
A new two-dimensional symmetric trace completion reconstructs the
inverse-fourth-distance norm for general $n$, while a four-column
skew completion is sharp at the six-, seven-, and eight-point candidates.
Global radial polynomial majorants and moment duality enlarge the
whole-angle exclusion regions. A prescribed-column trace envelope
couples the radial variance to the second and third holomorphic moments.
Pseudoinverse projections include the eight-point singular boundaries.
Earlier complete six- and seven-point proofs,
with their stated published classification dependencies, are retained.
The unrestricted eight-point problem still requires a bound for
configurations with no origin; an equivalent centroid-penalty inequality
states that remaining step precisely. No unrestricted eight-point
value or general finite-$n$ formula is claimed here.
\end{abstract}

\section{Scope, normalization and the existing coefficient moments}
For $n\ge2$, write
\[
\Delta_n(z)=\prod_{i<j}|z_i-z_j|^2,\qquad
M_n=\max_{\sum_i|z_i|^2=n}\Delta_n(z),\qquad N=n(n-1)/2.
\]
Compactness gives existence. A unit regular polygon gives a positive
value, so maximizing points are distinct. Centering and scaling show
that a maximizer has zero centroid. At $\lambda=(n-1)/2$ it also
globally maximizes $\Delta_n(z)e^{-\lambda\sum_i|z_i|^2}$ on $\C^n$.

The supplied manuscript determines $M_2=4$, $M_3=27$, $M_4=256$,
$M_5=3125$, $M_6=6^{15}/5^{10}$ and $M_7=7^{21}/6^{15}$.
The primary five-point proof follows the supplied second-order and moment
manuscript~\cite{JiangFive}. The distinct classification method of
Wang~\cite{Wang} is reconstructed in the algebraic cross-check;
neither proof invokes Wang's sharp theorem as a premise.
The six- and
seven-point global coverage relies on the published Faug\`ere--Svartz
classification~\cite{FS}, followed by independent rational certificates.
The six-point $e_2=0$ charts are independently covered below; only the
general $e_2\ne0$ region still uses external decomposition coverage.
For seven points the entire $e_2=e_3=0$ region is independently
closed by discriminant-combination identities.
The omitted degree-$356$ seven-point branch is excluded by the published
physical-emptiness statement, not by the $174$ supplied-record verifier.
Our rerun verified all $174$ parameter candidates and $1218$ point roots.
No eight-point completeness statement is imported from that paper.


The new proofs that avoid stationary classification have a different
scope from the retained complete classification proofs. For six and
seven points they cover every configuration containing an origin
(Theorem~\ref{thm:noclass-origin67}); for eight points they cover
the entire origin branch (Theorem~\ref{thm:noclass-origin8}).
The enlarged small-harmonic result for eight points
(Theorem~\ref{thm:noorigin-oct7-eight-improvedharmonic}) applies to actual local
maxima but does not show that every global maximum has small harmonics.
These scope distinctions are essential to the global claim.

The present unrestricted arguments for six, seven, and eight points
do not invoke the Faug\`ere--Svartz decomposition. They give whole-region
value bounds and exact finite-real-rank interfaces; they do not yet
close every no-origin configuration. The older classification-dependent
proofs are retained later as historical cross-checks and are not premises
of any new theorem.

The existing critical-value method supplies a rational annihilator with
degree at most $n^{2n-4}$ by saturating the entire stationary graph ideal.
Positive-dimensional stationary components are allowed because the
rational objective is constant on each irreducible critical component.
The present addendum does not compute that annihilator for a new point count.
An exact algebraic output consists of a primitive irreducible integer
polynomial with positive leading coefficient and a rational interval
isolating the selected real root.

For completeness, the moment enclosure used below has the following
self-contained formulation. Expand the integer polynomial
\[
V_n(z)^k=\sum_{|\alpha|=kN}c_\alpha z^\alpha,
\qquad V_n(z)=\prod_{i<j}(z_i-z_j),
\qquad \alpha!=\prod_i\alpha_i!.
\]
Set
\begin{equation}\label{eq:coefficient-moment}
C_{n,k}=\sum_{|\alpha|=kN}c_\alpha^2\alpha!,\qquad
A_{n,k}=\frac{n^{kN}(n-2)!C_{n,k}}{(kN+n-2)!}.
\end{equation}
All coefficients, including cancellations, must be computed before
using this formula as an upper certificate.

\begin{proposition}[Coefficient moment enclosure]
For every $k\ge1$,
\begin{equation}\label{eq:moment-enclosure}
A_{n,k}^{1/k}\le M_n\le
\left[\binom{kN+n-2}{n-2}A_{n,k}\right]^{1/k}.
\end{equation}
\end{proposition}
\begin{proof}
For independent standard complex Gaussians $Z_i$, monomial orthogonality
gives $\mathbb E|V_n(Z)|^{2k}=C_{n,k}$. Project onto the centered
hyperplane, of complex dimension $n-1$. Translation invariance
preserves $V_n$, while the squared radius has law $\Gamma(n-1,1)$
and is independent of its spherical direction. Homogeneity shows that
$A_{n,k}$ is the spherical average of $\Delta_n^k$ at moment $n$,
proving the lower bound.

For a homogeneous holomorphic polynomial $p(u)=\sum b_\alpha u^\alpha$
of degree $s$ in orthonormal centered coordinates, Cauchy--Schwarz and
the multinomial formula give
\[
|p(u)|^2\le
\left(\sum_{|\alpha|=s}|b_\alpha|^2\frac{\alpha!}{s!}\right)
\|u\|^{2s}.
\]
For $p=V_n^k$, Gaussian integration identifies the first factor with
$C_{n,k}/(kN)!$. At $\|u\|^2=n$ this proves
$M_n^k\le n^{kN}C_{n,k}/(kN)!$, which is the stated upper endpoint.
\end{proof}

For fixed $n$ the endpoint ratio tends to one. The rational moment
ratios $A_{n,k+1}/A_{n,k}$ are nondecreasing lower bounds converging
to $M_n$. These convergence statements are already in the supplied
work; the fixed finite stopping specification in Section~\ref{sec:fixed-moment-reconstruction}
adds an effective height-based selection bound.


\section{The primary five-point proof: second-order conditions and moments}
\label{sec:five-primary-moments}
This section follows the user-supplied manuscript by Fengcan Jiang,
\emph{A second-order and moment proof of the five-point Mordell inequality},
dated 3 October 2026~\cite{JiangFive}. Its complete analytic reduction is
retained: second-order optimality, radial moments, a two-dimensional real
projection, angular comparison, and a quartic logarithm majorant.
It does not use a classification of stationary configurations.
The independent algebraic proof is retained later as a cross-check.
The local manuscript attributes the alternative classification proof to
Wang~\cite{Wang}; no external priority assertion is made here.

\begin{theorem}[Five-point inequality by the second-order moment method]
\label{ref5-thm:main}
For $z_1,\ldots,z_5\in\C$ with $\sum_i|z_i|^2=5$,
\[
\Delta_5(z)\le5^5=3125.
\]
Equality holds precisely for a rotated unit regular pentagon, up to
permutation. The proof below uses only the analytic reductions and the
exact rational scalar certificate.
\end{theorem}
\subsection{Second-order estimates}
\subsubsection{A general radial bound}
We use the Hermitian inner product $\ip{x}{y}=\sum_i\overline{x_i}y_i$;
a superscript $\mathsf T$ denotes transpose, and $*$ denotes conjugate
transpose. Powers of vectors below are taken coordinatewise.

\begin{proposition}\label{ref5-prop:general}
Let $n\ge2$, and let $z$ be a collision-free local maximizer of
$\Delta_n$ on $\mathcal S_n$. Set
\[
 t_i=|z_i|^2,\quad s=\sum_i(t_i-1)^2,\quad
 p_2=\sum_i z_i^2,\quad \lambda=\frac{n-1}{2}.
\]
Then
\begin{equation}\label{ref5-eq:stationary-general}
 \sum_{j\ne i}\frac1{z_i-z_j}=\lambda\overline{z_i},
 \qquad \sum_i z_i=0.
\end{equation}
The complex symmetric matrix
\begin{equation}\label{ref5-eq:B}
 B_{ij}=-\frac1{(z_i-z_j)^2}\quad(i\ne j),\qquad
 B_{ii}=\sum_{j\ne i}\frac1{(z_i-z_j)^2}
\end{equation}
satisfies
\begin{equation}\label{ref5-eq:B-general}
 B\one=0,\qquad Bz=\lambda\overline z,\qquad
 \norm{B}_{\op}\le\lambda.
\end{equation}
Moreover,
\begin{equation}\label{ref5-eq:radial-general}
 3ns+|p_2|^2\le n^2,\qquad
 |z_i-z_j|^2\ge\frac{2n}{(n-1)^2}.
\end{equation}
If $n\ge3$, the second power sum also satisfies
\begin{equation}\label{ref5-eq:power-radial}
 |p_2|^2\le
 \frac{ns(n+s)}{s+n(n-2)^2/(n-1)^2}.
\end{equation}
For the real weighted Laplacian with weights
$|z_i-z_j|^{-2}$, one has
\begin{equation}\label{ref5-eq:L-general}
 Lt=(n-1)(t-\one).
\end{equation}
\end{proposition}

\begin{proof}
Lagrange multipliers for $\log\Delta_n$ give the first equation
in \eqref{ref5-eq:stationary-general} with a real multiplier.
Multiplying by $z_i$ and summing yields
$\lambda n=n(n-1)/2$; summing the equations gives the zero centroid.
The first two identities in \eqref{ref5-eq:B-general} follow directly.

For $\ip{z}{h}=0$, every phase multiple of $h$ is a real tangent vector
to $\mathcal S_n$. The second derivative of the Lagrangian is
\[
 -2\Rea(h^{\mathsf T}Bh)-2\lambda\norm{h}^2.
\]
Its nonpositivity for every phase implies
$|h^{\mathsf T}Bh|\le\lambda\norm h^2$.
Write any vector as $h+\alpha z$ with $\ip z h=0$.
Since $h^{\mathsf T}Bz=0$ and $z^{\mathsf T}Bz=\lambda n$,
the same quadratic bound holds for all vectors.
For unit vectors $x,y$, the identity
\[
4x^{\mathsf T}By=(x+y)^{\mathsf T}B(x+y)
 -(x-y)^{\mathsf T}B(x-y)
\]
and the parallelogram law give $|x^{\mathsf T}By|\le\lambda$.
Taking the supremum proves the operator norm bound.

Pairwise expansion gives
$Bz^{\circ2}=(n-1)(t-\one)$. Applying \eqref{ref5-eq:B-general} to
$z^{\circ2}-p_2\one/n$ yields
\[
 (n-1)^2s\le\lambda^2
 \left(\sum_i t_i^2-\frac{|p_2|^2}{n}\right).
\]
As $\sum_i t_i^2=n+s$, this is the first bound in
\eqref{ref5-eq:radial-general}. To sharpen the entry bound, let
$P_0=I-\one\one^{\mathsf T}/n$. Since $B=P_0BP_0$ and
$\norm{P_0e_i}^2=(n-1)/n$, we have
$|B_{ij}|\le\lambda(n-1)/n$.
This proves the second inequality.
To prove \eqref{ref5-eq:power-radial}, multiply the stationary equations
by $z_i^3$ and sum. Since $\sum_i z_i=0$, pairwise expansion gives
$\sum_i t_i z_i^2=(2n-3)p_2/(n-1)$. Thus
\[
 \ip{t-\one}{z^{\circ2}-p_2\one/n}
   =\frac{n-2}{n-1}p_2.
\]
Cauchy--Schwarz yields
$((n-2)/(n-1))^2|p_2|^2
 \le s(n+s-|p_2|^2/n)$.
Rearranging proves \eqref{ref5-eq:power-radial}, including $s=0$.
Finally, multiply \eqref{ref5-eq:stationary-general} by $z_i$, take real
parts, and use
$2\Rea(z_i/(z_i-z_j))=1+(t_i-t_j)/|z_i-z_j|^2$.
This gives \eqref{ref5-eq:L-general}.
\end{proof}

\subsubsection{Five-point moments and a Gram matrix}
From now on $n=5$, and we work at a collision-free local maximizer.
In addition to $t_i$ and $s$, put
\begin{equation}\label{ref5-eq:notation}
\begin{gathered}
 T=\sum_i t_i^2=5+s,\qquad U=\sum_i t_i^3,\qquad V=\sum_i t_i^4,\\
 p_j=\sum_i z_i^j,\quad q=|p_2|^2,\quad p=|p_3|^2,\quad
 w=|p_4|,\quad m=\sum_i t_i^2z_i^2.
\end{gathered}
\end{equation}
Proposition \ref{ref5-prop:general} gives
\begin{equation}\label{ref5-eq:domain}
 0\le s\le\frac53,\qquad
 0\le q\le \min\left\{25-15s,\frac{80s(5+s)}{45+16s}\right\}.
\end{equation}
Multiplication of the stationary equations by $z_i^2,z_i^3,z_i^4$,
followed by summation and pairing, gives
\begin{equation}\label{ref5-eq:mixed}
 \sum_i t_i z_i=0,\qquad
 \sum_i t_i z_i^2=\frac74p_2,\qquad
 \sum_i t_i z_i^3=\frac32p_3.
\end{equation}
For example, for the middle identity the pair sum is
$4p_2+\sum_{i<j}z_i z_j=7p_2/2$.
Direct expansion also gives
\begin{equation}\label{ref5-eq:B-powers}
 Bz^{\circ2}=4(t-\one),\qquad
 Bz^{\circ4}=(8t-9\one)\circ z^{\circ2}-p_2\one.
\end{equation}

The matrix $4I-B^*B$ is positive semidefinite.
Let $D$ be its Gram matrix on the two vectors
$z^{\circ2}-p_2\one/5$ and $z^{\circ4}-p_4\one/5$.
Using \eqref{ref5-eq:mixed}--\eqref{ref5-eq:B-powers}, its entries are
\begin{equation}\label{ref5-eq:D}
\begin{aligned}
 D_{22}&=20-12s-\frac45q,\\
 D_{44}&=-60V+144U-81T+5q-\frac45w^2,\\
 D_{24}&=-28m+83p_2-\frac45\overline{p_2}p_4.
\end{aligned}
\end{equation}
In particular,
\begin{equation}\label{ref5-eq:D-psd}
 D_{22},D_{44}\ge0,\qquad |D_{24}|^2\le D_{22}D_{44}.
\end{equation}

\subsection{A two-dimensional projection constraint}
\begin{lemma}\label{ref5-lem:projection}
For $s>0$ at a collision-free local maximizer,
\begin{equation}\label{ref5-eq:w-lower}
 w\ge T-\frac{2q}{5}-\frac{9q}{8s}-\frac{p}{5-\sqrt q}.
\end{equation}
\end{lemma}
\begin{proof}
By \eqref{ref5-eq:domain}, $q<25$.
Thus $E=\operatorname{span}_{\C}\{\one,z,\overline z\}$ has
dimension three: $\one$ is orthogonal to the other two vectors, whose
Gram determinant is $25-q>0$.
The vector $u=t-\one$ is real, lies in $E^\perp$, and has norm squared $s$.
Let $v$ be the orthogonal projection of $z^{\circ2}$ onto $E^\perp$.
The Gram matrix of $z,\overline z$ is
$\left(\begin{smallmatrix}5&\overline{p_2}\\p_2&5\end{smallmatrix}\right)$.
Since $\ip z {z^{\circ2}}=0$ and
$\ip{\overline z}{z^{\circ2}}=p_3$, the linear part of the
projection onto $E$ is
\[
 \frac{-\overline{p_2}p_3z+5p_3\overline z}{25-q}.
\]
Consequently,
\begin{equation}\label{ref5-eq:projection-identities}
\begin{aligned}
 \norm v^2&=T-\frac q5-\frac{5p}{25-q},&
 \ip u v&=\frac34p_2,\\
 v^{\mathsf T}v&=p_4-\frac{p_2^2}{5}
    +\frac{\overline{p_2}p_3^2}{25-q}.
\end{aligned}
\end{equation}
The conjugation-invariant space $E^\perp$ has a real orthonormal basis
$u/\sqrt s,e$. Write $v=a\,u/\sqrt s+b e$.
Then $|a|^2=9q/(16s)$ and
\[
 |v^{\mathsf T}v|=|a^2+b^2|
 \ge\norm v^2-2|a|^2
 =\norm v^2-\frac{9q}{8s}.
\]
Combine this with \eqref{ref5-eq:projection-identities} and the triangle
inequality. The coefficient of $p$ becomes
$(5+\sqrt q)/(25-q)=1/(5-\sqrt q)$, proving the assertion.
\end{proof}

\subsection{Excluding small nonzero radial variance}
\begin{proposition}[Local radial gap]\label{ref5-prop:gap}
A collision-free local maximizer satisfies either $s=0$ or $s>1/20$.
\end{proposition}
\begin{proof}
Suppose $0<s\le1/20$, and write $z_i=r_i e^{i\theta_i}$.
For $\delta_i=t_i-1$, we have $\sum_i\delta_i=0$ and
$\sum_i\delta_i^2=s$. Hence
\begin{equation}\label{ref5-eq:local-radii}
 t_i\ge1-\sqrt{4s/5}\ge\frac45,\qquad
 \frac78<r_ir_j<\frac98.
\end{equation}
For completeness, set $a=(t_i+t_j)/2$, $b=(t_i-t_j)/2$.
Cauchy--Schwarz for the remaining three deviations gives
$s\ge(10/3)(a-1)^2+2b^2$.
It follows that
$t_i t_j\ge a^2-s/2+(5/3)(a-1)^2$.
This last expression is increasing for
$a\ge1-\sqrt{3s/10}>7/8$, since its derivative is $(16a-10)/3$.
At the left endpoint it equals $a^2$, so
$r_ir_j\ge1-\sqrt{3s/10}>7/8$.
The upper bound follows from $r_ir_j\le a\le1+\sqrt{3s/10}<9/8$;
here $(1/8)^2>3/200$.

Let $\overline r=5^{-1}\sum_i r_i$ and
$\sigma=\sum_i(r_i-\overline r)^2$.
Using $(r_i+r_j)^2\ge4r_ir_j>7/2$,
\begin{equation}\label{ref5-eq:sigma}
 \sigma=\frac15\sum_{i<j}
       \frac{(t_i-t_j)^2}{(r_i+r_j)^2}
 <\frac{2s}{7}\le\frac1{70},\qquad
 (r_i-r_j)^2<\frac{4s}{7}\le\frac1{35}.
\end{equation}
All squared distances are at least $1/2$ by
Proposition \ref{ref5-prop:general}. Thus
\[
 \sin^2\frac{\theta_i-\theta_j}{2}
 =\frac{|z_i-z_j|^2-(r_i-r_j)^2}{4r_ir_j}
 >\frac{11}{105}>\frac1{10}.
\]
The angles are distinct and all the corresponding half-angle cotangents
have absolute value less than $3$.

Lift and order the angles so that
$\theta_1<\cdots<\theta_5<\theta_1+2\pi$, and put
\[
 \mathcal C(\theta)=\sum_{i<j}
       \log|e^{i\theta_i}-e^{i\theta_j}|^2.
\]
For $\sum_i h_i=0$,
\[
 h^{\mathsf T}\nabla^2\mathcal C(\theta)h
 =-\frac12\sum_{i<j}
   \csc^2\frac{\theta_i-\theta_j}{2}(h_i-h_j)^2
 \le-\frac52\norm h^2.
\]
Let $\theta^*$ be the equally spaced angles with the same mean, and let
$\eta=\norm{\theta-\theta^*}$. The ordered region of span less than
$2\pi$, intersected with the fixed-mean hyperplane, is convex.
Integration of the preceding Hessian bound along the connecting segment
gives
\begin{equation}\label{ref5-eq:eta-gradient}
 \eta\le\frac25\norm{\nabla\mathcal C(\theta)}.
\end{equation}
The angular stationary equations for the actual radii imply
\begin{equation}\label{ref5-eq:angular-stationary}
 \frac{\partial\mathcal C}{\partial\theta_i}
 =\sum_{j\ne i}\cot\frac{\theta_i-\theta_j}{2}
       \frac{(r_i-r_j)^2}{|z_i-z_j|^2}.
\end{equation}
The incidence matrix of $K_5$ has operator norm $\sqrt5$.
For $v_i=r_i-\overline r$, the zero sum gives
$\max_i v_i^2\le4\sigma/5$, and expansion yields
$\sum_{i<j}(v_i-v_j)^4=5\sum_i v_i^4+3\sigma^2\le7\sigma^2$.
Combining these observations with
\eqref{ref5-eq:eta-gradient}--\eqref{ref5-eq:angular-stationary}, if squared
distances are at least $d$ and cotangents at most $h$ in absolute value,
then
\begin{equation}\label{ref5-eq:eta-basic}
 \eta\le\gamma\sigma h/d,\qquad \gamma=\frac{2\sqrt{35}}5.
\end{equation}
Initially $d=1/2,h=3$ and \eqref{ref5-eq:sigma} give
$\eta<3\gamma/35<21/100$.

We next use only algebraic constants. Put
\begin{equation}\label{ref5-eq:algebraic-trig}
\begin{gathered}
 S_1=\sqrt{\frac{5-\sqrt5}{8}},\quad C_1=\frac{1+\sqrt5}{4},
 \qquad
 S_2=\sqrt{\frac{5+\sqrt5}{8}},\quad C_2=\frac{\sqrt5-1}{4},\\
 P_j(u)=S_j(1-u^2/4)-C_j u/\sqrt2,\qquad
 N(u)=C_1+S_1u/\sqrt2.
\end{gathered}
\end{equation}
Here $S_j=\sin(j\pi/5)$, $C_j=\cos(j\pi/5)$.
The inequalities $\cos x\ge1-x^2/2$ and $\sin x\le x$ for $x\ge0$
give
\[
 \sin(j\pi/5-u/\sqrt2)\ge P_j(u),\qquad
 \cos(\pi/5-u/\sqrt2)\le N(u).
\]
If $\eta\le u\le21/100$, each pairwise angular error is at most
$\sqrt2u$. Consequently every absolute half-angle sine is at least
$P_1(u)>0$ and the absolute cotangent is at most $N(u)/P_1(u)$.
Equations \eqref{ref5-eq:local-radii}, \eqref{ref5-eq:sigma}, and
\eqref{ref5-eq:eta-basic} then give
\begin{equation}\label{ref5-eq:algebraic-iteration}
 \eta<\frac{\gamma N(u)}{245P_1(u)^3}.
\end{equation}
The exact algebraic checks in the ancillary program prove
\begin{equation}\label{ref5-eq:two-steps}
 \frac{\gamma N(21/100)}{245P_1(21/100)^3}<\frac9{100},
 \qquad
 \frac{\gamma N(9/100)}{245P_1(9/100)^3}<\frac3{50}.
\end{equation}
Thus $\eta<3/50$. The same algebraic checks give
\begin{equation}\label{ref5-eq:weights}
 \frac{2}{7P_1(3/50)^2}<\frac{19}{20},\qquad
 \frac{2}{7P_2(3/50)^2}<\frac13.
\end{equation}
These bound the inverse squared distances on consecutive pairs
and diagonals in the cyclic order, respectively.

Let $L$ be the real Laplacian with weights $|z_i-z_j|^{-2}$.
Comparison in the positive semidefinite order with a pentagon of
edge weight $19/20$ and diagonal weight $1/3$ gives
\begin{equation}\label{ref5-eq:spectral-gap}
 \lambda_{\max}(L)\le
 \frac{19}{20}\frac{5+\sqrt5}{2}
 +\frac13\frac{5-\sqrt5}{2}<4.
\end{equation}
The two Laplacians of the pentagon and its complement are simultaneously
diagonalized by the discrete Fourier basis; their nonzero eigenvalues
are $(5\pm\sqrt5)/2$ in complementary order.
However, \eqref{ref5-eq:L-general} gives $L(t-\one)=4(t-\one)$.
As $s>0$, this contradicts \eqref{ref5-eq:spectral-gap}.
\end{proof}

\subsection{A quartic logarithm majorant}
Put
\begin{equation}\label{ref5-eq:majorant-constants}
 a=\frac{5-\sqrt5}{2},\quad b=\frac{5+\sqrt5}{2},\quad
 c=\frac15-\frac{2}{5\sqrt5}\log\frac ba,\quad k=\frac9{2500}.
\end{equation}
The following bounds are certified in the ancillary program:
\begin{equation}\label{ref5-eq:c-bounds}
 \frac{22}{3}k<c<C:=\frac{278365}{10000000}<\frac{17}{2}k.
\end{equation}
Let $H$ be the cubic interpolating $\log x$ and its derivative at $a,b$.
Solving the four interpolation conditions gives
\begin{equation}\label{ref5-eq:H}
 H(x)=\frac12\log5-\frac74-\frac{25c}{4}
 +(1+15c)x-\left(\frac1{10}+\frac{15c}{2}\right)x^2+cx^3.
\end{equation}

\begin{lemma}\label{ref5-lem:majorant}
For $0<x\le R:=41/6$,
\begin{equation}\label{ref5-eq:majorant}
 \log x\le H(x)-k(x^2-5x+5)^2.
\end{equation}
\end{lemma}
\begin{proof}
Add the node $R$ to the double nodes $a,b$.
The quartic Hermite interpolant is
$H(x)-k_R(x^2-5x+5)^2$, where
\[
 k_R=\frac{H(R)-\log R}{(R^2-5R+5)^2}.
\]
The fifth derivative of $\log x$ is positive.
Its Hermite remainder has the sign of
$(x-a)^2(x-b)^2(x-R)$, hence is nonpositive for $0<x\le R$.
The exact series check in Section~\ref{ref5-app:certificate} proves
$k_R>k$. This yields \eqref{ref5-eq:majorant}.
\end{proof}

All squared distances at a local maximizer belong to this range:
\[
 |z_i-z_j|^2\le2(t_i+t_j)
 \le4+2\sqrt{6s/5}\le4+2\sqrt2<41/6.
\]
Write $d_{ij}=|z_i-z_j|^2$ and
$\mathcal Q=\sum_{i<j}(d_{ij}^2-5d_{ij}+5)^2$.
Expansion using \eqref{ref5-eq:mixed} gives
\begin{equation}\label{ref5-eq:distance-moments}
\begin{aligned}
 \sum_{i<j}d_{ij}&=25,\\
 \sum_{i<j}d_{ij}^2&=5T+50+q,\\
 \sum_{i<j}d_{ij}^3&=5U+45T+\frac{21}{2}q-p,\\
 \sum_{i<j}d_{ij}^4&=5V+80U+18T^2+49q-12p+w^2
                 +12\Rea(m\overline{p_2}).
\end{aligned}
\end{equation}
These identities can be checked uniformly by putting
$M_{\alpha\beta}=\sum_i z_i^\alpha\overline{z_i}^{\,\beta}$ and using
\begin{equation}\label{ref5-eq:moment-expansion}
 \sum_{i<j}d_{ij}^r=
 \frac12\sum_{\alpha,\beta=0}^r(-1)^{\alpha+\beta}
 \binom r\alpha\binom r\beta
 M_{r-\alpha,r-\beta}M_{\alpha\beta}.
\end{equation}

The regular pentagon has five squared distances $a$ and five $b$.
Subtract its value in \eqref{ref5-eq:majorant} and use
\eqref{ref5-eq:distance-moments} to obtain
\begin{equation}\label{ref5-eq:log-bound}
 \log\Delta_5-5\log5\le
 \left(-\frac12+\frac{15c}{2}\right)s+5c(U-5)
 +\left(-\frac1{10}+3c\right)q-cp-k\mathcal Q.
\end{equation}
The moment expansion of $\mathcal Q$ is
\begin{equation}\label{ref5-eq:Q-expand}
\begin{aligned}
 \mathcal Q={}&5V+30U+18T^2-275T+750-21q-2p+w^2\\
             &+12\Rea(m\overline{p_2}).
\end{aligned}
\end{equation}
Set
\[
 R_0=5V+30U+18T^2-275T+750,\qquad
 D_0=-60V+144U-81T+5q.
\]
Solving \eqref{ref5-eq:D} for $m$ and applying \eqref{ref5-eq:D-psd} gives
\begin{equation}\label{ref5-eq:Q-lower}
 \mathcal Q\ge R_0+\frac{102}{7}q-2p+w^2-\frac{12}{35}qw
       -\frac37\sqrt{qD_{22}D_0}.
\end{equation}
Indeed, the contribution of $m$ is
\[
12\Rea(m\overline{p_2})
 =\frac{249}{7}q-\frac{12}{35}\Rea(\overline{p_2}^{\,2}p_4)
 -\frac37\Rea(D_{24}\overline{p_2}).
\]
Also $D_0=D_{44}+4w^2/5\ge D_{44}\ge0$, so replacing $D_{44}$
by $D_0$ in the square root weakens the lower bound in the required
direction.

\subsection{Reduction to two variables}
\subsubsection{Eliminating the complex power sums}
Assume $s>0$. Put $L_0=T-2q/5-9q/(8s)$ and $w_0=6q/35$.
Lemma \ref{ref5-lem:projection} implies
$p\ge(5-\sqrt q)\pos{L_0-w}$.
The expression to be bounded below after combining
\eqref{ref5-eq:log-bound} and \eqref{ref5-eq:Q-lower} is
$k(w^2-12qw/35)+(c-2k)p$.
It satisfies
\begin{equation}\label{ref5-eq:eliminate-pw}
 k(w^2-12qw/35)+(c-2k)p
 \ge k\left(\pos{T-\left(\frac47+\frac9{8s}\right)q}^{\,2}
                  -\frac{36q^2}{1225}\right).
\end{equation}
To see this, complete the square at $w_0$.
If $L_0\le w_0$, the unconstrained minimum at $w_0$ is feasible.
If $L_0>w_0$, the minimum is at $L_0$ provided
\[
 L_0-w_0\le\frac{c-2k}{2k}(5-\sqrt q).
\]
Here $L_0>0$ implies
$q<T/(2/5+9/(8s))\le800/129<25/4$.
By \eqref{ref5-eq:c-bounds}, the last right-hand side is greater than
$(8/3)(5/2)=20/3\ge T\ge L_0-w_0$.
This proves the asserted piecewise minimum.

\subsubsection{Eliminating the radial moments}
Cauchy--Schwarz and projection of $t^{\circ2}$ onto the orthogonal
vectors $\one,t-\one$ give
\begin{equation}\label{ref5-eq:UV}
 U\ge\frac{T^2}{5},\qquad
 V\ge V_0(U,s):=\frac{T^2}{5}+\frac{(U-T)^2}{s}.
\end{equation}
The zero-sum deviations $\delta_i=t_i-1$ also satisfy
\[
 \sum_i\delta_i^3\ge-\frac3{\sqrt{20}}s^{3/2}.
\]
For a proof, at an extremum of the third moment on
$\sum\delta_i=0$, $\sum\delta_i^2=s$, the multiplier equation is a
quadratic in $\delta_i$. Thus at most two values occur; checking
multiplicities $1+4$ and $2+3$ gives the stated sharp lower bound.
Consequently,
\begin{equation}\label{ref5-eq:U0}
 U\ge U_0(s)=
 \begin{cases}
 5+3s-\dfrac3{\sqrt{20}}s^{3/2},&0<s\le5/4,\\[2pt]
 (5+s)^2/5,&5/4\le s\le5/3.
 \end{cases}
\end{equation}
The two candidates meet at $s=5/4$, as follows by setting
$x=\sqrt{s/20}$ and factoring
$1-3x-4x^2=(1-4x)(1+x)$.

After \eqref{ref5-eq:eliminate-pw}, the right-hand side of the logarithm
bound decreases with $V$: its nonradical derivative is $-5k$ and
the radicand decreases with $V$.
We may therefore set $V=V_0(U,s)$.
For $U\ge U_0(s)$,
\[
 \partial_U V_0=\frac{2(U-T)}s
 \ge\begin{cases}
 4-3\sqrt{s/5}\ge5/2,&s\le5/4,\\
 2(5+s)/5\ge5/2,&s\ge5/4.
 \end{cases}
\]
Hence $\partial_U D_0\le-6$, while the remaining derivative is at most
$5c-(85/2)k<0$. We may also set $U=U_0(s)$.
Both downward substitutions increase $D_0$, preserving its
nonnegativity along the substitution path.
Finally the resulting upper bound increases with $c$, with derivative
$(15/2)s+5(U_0-5)+3q>0$, so we may replace $c$ by the rational upper
bound $C$ in \eqref{ref5-eq:c-bounds}.

Define
\begin{equation}\label{ref5-eq:scalar-definitions}
\begin{gathered}
 U_*=U_0(s),\quad V_*=\frac{(5+s)^2}{5}
                  +\frac{(U_*-5-s)^2}{s},\\
 A_q=\frac1{10}-3C+\frac{102k}{7},\quad
 B_q=\frac{36k}{1225},\quad D_q=\frac{3k}{7},\\
 f(s)=\left(-\frac12+\frac{15C}{2}\right)s+5C(U_*-5)\\
 \hspace{15mm}-k\{5V_*+30U_*+18(5+s)^2-275(5+s)+750\},\\
 D_*(s)=-60V_*+144U_*-81(5+s).
\end{gathered}
\end{equation}
We have proved the following scalar reduction.

\begin{proposition}\label{ref5-prop:scalar}
At a collision-free local maximizer with $s>0$,
\begin{equation}\label{ref5-eq:F-bound}
 \log\Delta_5-5\log5\le F(s,q),
\end{equation}
where
\begin{equation}\label{ref5-eq:F}
\begin{aligned}
 F(s,q)={}&f(s)-A_q q+B_q q^2\\
 &+D_q\sqrt{q(20-12s-4q/5)(D_*(s)+5q)}\\
 &-k\pos{5+s-\left(\frac47+\frac9{8s}\right)q}^{\,2}.
\end{aligned}
\end{equation}
The parameters satisfy \eqref{ref5-eq:domain}, and both nontrivial factors
inside the square root are nonnegative.
\end{proposition}

\subsection{Exact certificate and completion of the proof}
\begin{lemma}[Scalar certificate]\label{ref5-lem:certificate}
With the constants in \eqref{ref5-eq:scalar-definitions},
\begin{equation}\label{ref5-eq:certificate}
 F(s,q)<0\qquad
 \left(\frac1{20}\le s\le\frac53,\quad
 0\le q\le\min\left\{25-15s,\frac{80s(5+s)}{45+16s}\right\}\right).
\end{equation}
\end{lemma}

\begin{proof}
We treat \eqref{ref5-eq:certificate} in two regions:
\[
 \mathrm A:\ s=20x^2,\quad\frac1{20}\le x\le\frac14;
 \qquad
 \mathrm B:\ \frac54\le s\le\frac53.
\]
In A, $U_*=5+60x^2-60x^3$ and
$V_*=5+120x^2-240x^3+260x^4$.
In B, $U_*=(5+s)^2/5$ and $V_*=(5+s)^3/25$.
The endpoint $20(1/20)^2=1/20$ exactly matches the local gap,
and the regions meet at $s=5/4$.

Region A starts with its containing rectangle and $0\le q\le25$.
For a rectangle $[x_l,x_h]\times[q_l,q_h]$, put
$s_{\min}=20x_l^2$, $s_{\max}=20x_h^2$, and replace $q_h$ by
\[
 \min\left\{q_h,25-15s_{\min},
            \frac{80s_{\max}(5+s_{\max})}{45+16s_{\max}}\right\}.
\]
The last expression is justified by \eqref{ref5-eq:power-radial} and the
monotonicity of $80s(5+s)/(45+16s)$ for $s\ge0$.
Discard an empty rectangle. Evaluate the interval form of
\eqref{ref5-eq:F}, clipping its radicand factors to their nonnegative parts.
If the upper endpoint is negative, accept the rectangle.
Otherwise bisect its longer normalized side, using scales $1/5$ and
$1$. The depth-first search exhausts all rectangles.

\begin{center}
\begin{tabular}{lrrr}
\toprule
Region & Visited nodes & Accepted rectangles & Maximum depth\\
\midrule
A & 183 & 92 & 21\\
\bottomrule
\end{tabular}
\end{center}
No rectangles remain unresolved. Every accepted upper endpoint is less
than the following convenient rational bound:
\begin{equation}\label{ref5-eq:certificate-margins}
 -\frac1{200000}.
\end{equation}
All comparisons use exact fractions. The interval arithmetic and the
constant checks are detailed in Section~\ref{ref5-app:certificate}.
Pruning removes only infeasible points and each bisection covers its
parent, so the finite certified cover proves the assertion.

On the second branch no subdivision is necessary. In fact,
\begin{equation}\label{ref5-eq:branch-b-analytic}
 F(s,q)\le-\frac{56657645489}{1713169664000}<-\frac3{100}
 \qquad\left(\frac54\le s\le\frac53,\quad 0\le q\le25-15s\right).
\end{equation}
To see this, write
\[
 f(s)=a_0s+b_0s^2-\frac{k}{5}s^3,\quad
 a_0=-\frac12+\frac{35}{2}C+20k,\quad b_0=C-27k.
\]
The rational inequalities $b_0<0$ and $a_0+(5/2)b_0<0$
imply that $f$ is decreasing for $s\ge5/4$. Likewise,
\[
 D_*(s)=15+27s-\frac{36}{5}s^2-\frac{12}{5}s^3
\]
is decreasing on this interval, since its derivative is at most $-9/4$.
Consequently $f(s)\le-9181/256000$, $D_*(s)\le525/16$,
$q\le25/4$, $20-12s-(4/5)q\le5$, and
$D_*(s)+5q\le1025/16$. Put
\[
 \alpha=A_q-\frac{25}{4}B_q=\frac{6692069}{98000000}>0.
\]
Dropping the nonpositive cutoff term in $F$ and completing the square
in $\sqrt q$ gives
\[
 F(s,q)\le-\frac{9181}{256000}-\alpha q
       +D_q\sqrt{\frac{5125}{16}q}
 \le-\frac{9181}{256000}
       +\frac{5125D_q^2}{64\alpha}
 =-\frac{56657645489}{1713169664000}.
\]
Every constant comparison in this argument is checked exactly in the
ancillary program.
\end{proof}

\begin{proof}[Proof of Theorem \ref{ref5-thm:main}]
The constraint sphere is compact, and a regular pentagon gives
$\Delta_5=5^5>0$. A global maximizer exists and is collision-free.
If $0<s\le1/20$, Proposition \ref{ref5-prop:gap} excludes it.
If $s\ge1/20$, Proposition \ref{ref5-prop:scalar} and
Lemma \ref{ref5-lem:certificate} give $\Delta_5<5^5$, a contradiction.
Thus $s=0$ at every global maximizer, so $|z_i|=1$.

For completeness, the classical unit-circle conclusion can be read
directly from the root polynomial, as in \cite{Mordell}.
Let $P(\zeta)=\prod_i(\zeta-z_i)$.
The stationary equations give
\[
 \frac{P''(z_i)}{2P'(z_i)}
 =\sum_{j\ne i}\frac1{z_i-z_j}=2\overline{z_i}=\frac2{z_i}.
\]
The polynomial $\zeta P''(\zeta)-4P'(\zeta)$ has degree at most three
and vanishes at five distinct points, so it is zero.
Coefficient comparison gives $P(\zeta)=\zeta^5-\alpha$, $|\alpha|=1$.
The roots are precisely a unit regular pentagon, and
$\Delta_5=\prod_i|P'(z_i)|=5^5$.
This establishes both the inequality and all equality cases.
\end{proof}

\begin{remark}
The local gap is a statement about every local maximizer, whereas
\eqref{ref5-eq:F-bound} need not exclude a local maximizer whose value is
strictly below $5^5$. The global comparison in the final proof uses the
known regular-pentagon competitor.
\end{remark}


\section{A general moment hierarchy and its limitations}
\label{ref5-sec:extensions}
The second-order method extends to arbitrary $n$, although the target
$\Delta_n\le n^n$ does not: Chalk's counterexample already disproves this
bound for every $n\ge6$ \cite{Chalk}. We record constraints that remain
valid at all collision-free local maximizers and a quantitative six-point
consequence. The polygonal spectral calculation below records a
limitation of the method; it is not presented as a new stability result.

\subsection{Higher-degree second-order constraints}
Work at a collision-free local maximizer on $\mathcal S_n$, and put
$p_k=\sum_i z_i^k$, $p_0=n$, and $\lambda=(n-1)/2$.
For every $k\ge1$, the stationary equations imply
\begin{equation}\label{ref5-eq:general-mixed-hierarchy}
 \sum_i t_i z_i^k=
 \frac{(2n-k-1)p_k+\sum_{j=1}^{k-1}p_jp_{k-j}}{n-1}.
\end{equation}
Indeed, multiply the stationary equation by $z_i^{k+1}$, sum, and pair
indices; the result before division by $\lambda$ is
$\frac12\{\sum_{j=0}^k p_jp_{k-j}-(k+1)p_k\}$.
For $k\ge2$ a related coordinatewise identity is
\begin{equation}\label{ref5-eq:general-B-hierarchy}
 Bz^{\circ k}=
 \left(k\lambda t-(k-1)(n-k/2)\one\right)\circ z^{\circ(k-2)}
 -\sum_{j=1}^{k-2}(k-1-j)p_jz^{\circ(k-2-j)}.
\end{equation}
To verify it, expand
$z_i^k-z_j^k=(z_i-z_j)\sum_{r=0}^{k-1}z_i^{k-1-r}z_j^r$
and then use, for $r\ge1$,
\[
 \sum_{j\ne i}\frac{z_j^r}{z_i-z_j}
 =\lambda z_i^r\overline z_i
 -\sum_{\ell=0}^{r-1}p_\ell z_i^{r-1-\ell}+r z_i^{r-1}.
\]

\begin{proposition}[Moment Gram hierarchy]\label{ref5-prop:moment-hierarchy}
For any finite set of positive integers $\mathcal K$, set
$x_k=z^{\circ k}-p_k\one/n$ and define
\begin{equation}\label{ref5-eq:general-D-hierarchy}
 D_{k\ell}=\lambda^2\ip{x_k}{x_\ell}
       -\ip{Bz^{\circ k}}{Bz^{\circ\ell}},\qquad k,\ell\in\mathcal K.
\end{equation}
Then $D$ is Hermitian positive semidefinite. Its entries are polynomials
in mixed moments, with the identities
\eqref{ref5-eq:general-mixed-hierarchy}--\eqref{ref5-eq:general-B-hierarchy}
providing exact reductions.
\end{proposition}
\begin{proof}
The matrix $\lambda^2I-B^*B$ is positive semidefinite by
Proposition~\ref{ref5-prop:general}, and $B\one=0$. Thus $D$ is its Gram
matrix on the vectors $x_k$. Constant projection removes a nonnegative
term that would otherwise weaken the constraints.
\end{proof}

For an explicit reusable instance, write
\[
 T=\sum_i t_i^2=n+s,\quad U=\sum_i t_i^3,\quad V=\sum_i t_i^4,
 \quad q=|p_2|^2,\quad w=|p_4|,\quad m=\sum_i t_i^2z_i^2.
\]
The submatrix on $\mathcal K=\{2,4\}$ has entries
\begin{equation}\label{ref5-eq:general-D24}
\begin{aligned}
 D_{22}&=\frac{(n-1)^2}{4}\left(n-3s-\frac qn\right),\\
 D_{44}&=-\frac{15(n-1)^2}{4}V+12(n-1)(n-2)U
       -9(n-2)^2T+nq-\frac{(n-1)^2}{4n}w^2,\\
 D_{24}&=-\frac{7(n-1)^2}{4}m+(7n^2-22n+18)p_2
       -\frac{(n-1)^2}{4n}\overline{p_2}p_4.
\end{aligned}
\end{equation}
They satisfy $D_{22},D_{44}\ge0$ and
$|D_{24}|^2\le D_{22}D_{44}$.
These formulas follow from
$Bz^{\circ4}=(2(n-1)t-3(n-2)\one)\circ z^{\circ2}-p_2\one$
and \eqref{ref5-eq:general-mixed-hierarchy}, and reduce exactly to
\eqref{ref5-eq:D} when $n=5$. Increasing $\mathcal K$ supplies additional
necessary constraints; it does not by itself prove a global upper bound.

\subsection{The dimensional obstruction in the projection step}
Assume $n\ge3$ and $s>0$, set $p=|p_3|^2$, and let
$E=\operatorname{span}_{\C}\{\one,z,\overline z\}$.
The radial bound gives $q<n^2$, so $\dim E=3$.
Let $u=t-\one$ and let $v$ be the projection of $z^{\circ2}$ onto
$E^\perp$. The computation in Lemma~\ref{ref5-lem:projection} now gives
\begin{equation}\label{ref5-eq:general-projection-identities}
\begin{aligned}
 \norm v^2&=T-\frac qn-\frac{np}{n^2-q},&
 \ip u v&=\frac{n-2}{n-1}p_2,\\
 v^{\mathsf T}v&=p_4-\frac{p_2^2}{n}
                  +\frac{\overline{p_2}p_3^2}{n^2-q}.
\end{aligned}
\end{equation}
Since $u\in E^\perp$ and $\norm u^2=s$, Cauchy--Schwarz proves the
dimension-independent bound
\begin{equation}\label{ref5-eq:general-projection-CS}
 \frac{(n-2)^2}{(n-1)^2}q
 \le s\left(T-\frac qn-\frac{np}{n^2-q}\right).
\end{equation}
Dropping the last nonpositive term recovers \eqref{ref5-eq:power-radial}.
For $n=5$, $E^\perp$ admits an orthonormal basis of two real vectors,
and the further
bound on $|v^{\mathsf T}v|$ used earlier is available.
For $n\ge6$ that particular argument fails: after fixing the real
direction $u$, two real orthonormal directions $e_1,e_2$ remain, and
$b(e_1+i e_2)$ has positive Hermitian norm but zero transpose square.
Thus no analogous lower bound follows solely from the norm and the
component along $u$. Additional moment constraints would be needed.

\subsection{A quantitative six-point radial gap}
\begin{proposition}\label{ref5-prop:six-gap}
At a collision-free local maximizer of $\Delta_6$ on $\mathcal S_6$,
either $s=0$ or $s>1/200$.
\end{proposition}
\begin{proof}
Suppose $0<s\le1/200$, and write $z_i=r_i e^{i\theta_i}$.
For $a=(t_i+t_j)/2$ and $b=(t_i-t_j)/2$, the zero-sum deviations give
$s\ge3(a-1)^2+2b^2$.
The function $a^2-s/2+\frac32(a-1)^2$ is increasing for $a>3/5$,
so its value at $a=1-\sqrt{s/3}$ proves
\[
 \frac{23}{24}<1-\sqrt{s/3}\le r_ir_j
 \le1+\sqrt{s/3}<\frac{25}{24}.
\]
Here $s/3\le1/600<1/576$.
With $\sigma=\sum_i(r_i-\overline r)^2$, the same pairwise variance
identities as in the five-point proof give
\[
 \sigma=\frac16\sum_{i<j}
       \frac{(t_i-t_j)^2}{(r_i+r_j)^2}<\frac{6s}{23},
 \qquad (r_i-r_j)^2<\frac{12s}{23}\le\frac3{1150}.
\]
The strengthened general distance bound gives $d_{ij}\ge12/25$.
Consequently,
\[
 \sin^2\frac{\theta_i-\theta_j}{2}
 >\frac{12/25-3/1150}{25/6}>\frac1{10};
 \qquad \left|\cot\frac{\theta_i-\theta_j}{2}\right|<3.
\]
Order the lifted angles and let $\theta^*$ be the equally spaced
six-point angles with the same mean. The circle logarithmic product
has Hessian at most $-3I$ on the mean-zero subspace. Its gradient
at the actual angles is given by \eqref{ref5-eq:angular-stationary}, with
six indices in place of five. The incidence matrix of $K_6$ has norm
$\sqrt6$, while, for $v_i=r_i-\overline r$,
\[
 \sum_{i<j}(v_i-v_j)^4
 =6\sum_i v_i^4+3\sigma^2\le8\sigma^2.
\]
Thus the angular error satisfies
\[
 \eta:=\norm{\theta-\theta^*}
 \le\frac4{\sqrt3}\frac{\sigma h}{d}
 <\frac73\frac{6s}{23}\,3\,\frac{25}{12}
 =\frac{175s}{46}\le\frac7{368}<\frac1{50}.
\]
Every half-angle error is at most $\eta/\sqrt2<1/70$, since
$\sqrt2>7/5$.
As sine is $1$-Lipschitz, the absolute half-angle sines in the three
cyclic separation classes are greater than $17/35$, $59/70$, and
$69/70$, respectively; here $\sqrt3/2>6/7$.
Multiplication by $4r_ir_j>23/6$ gives
\[
 d_{ij}>\frac{9}{10},\qquad d_{ij}>\frac83,\qquad
 d_{ij}>\frac{100}{27}
\]
in those three classes. Hence their Laplacian weights are less than
$a_0=10/9$, $b_0=3/8$, $c_0=27/100$.
The comparison circulant Laplacian has nonzero eigenvalues
$a_0+3b_0+2c_0$, $3a_0+3b_0$, and $4a_0+2c_0$, so
\[
 \lambda_{\max}(L)\le\frac{2243}{450}<5.
\]
This contradicts $L(t-\one)=5(t-\one)$ and $s>0$.
All numerical comparisons are rational inequalities or follow by
squaring positive rational numbers.
\end{proof}


\subsection{Why a six-point global argument needs a different target}
The inverse-distance Laplacian of the unit regular polygon similarly
has eigenvalues $k(n-k)/2$, $0\le k\le n-1$.
Indeed, its $k$th eigenvalue is
$\frac12\sum_{r=1}^{n-1}
|\sin(k\pi r/n)/\sin(\pi r/n)|^2$; expansion of
$|\sum_{j=0}^{k-1}e^{2\pi i jr/n}|^2$ and Fourier orthogonality
give $k(n-k)/2$.
The largest eigenvalue is $\frac12\lfloor n^2/4\rfloor$.
The strict spectral margin below $n-1$ holds for $2\le n\le6$, is zero
at $n=7$, and is absent for $n\ge8$. Thus the local exclusion by strict
Laplacian comparison cannot be copied unchanged for $n\ge7$.

For completeness, Chalk's competitor consists of the origin and a
regular $(n-1)$-gon of radius $\sqrt{n/(n-1)}$. It has
\begin{equation}\label{ref5-eq:Chalk-competitor}
 \Delta_n=(n-1)^{n-1}
       \left(\frac n{n-1}\right)^{n(n-1)/2}.
\end{equation}
At $n=6$ its ratio to $6^6$ is
$6^9/5^{10}=10077696/9765625>1$.
For real $x\ge6$, twice the logarithm of this ratio is
$g(x)=x(x-3)\log x-(x-1)(x-2)\log(x-1)$, whose derivative is
\[
 g'(x)=(2x-3)\log\frac{x}{x-1}-1
 \ge\frac{2x-5}{2x-1}>0.
\]
The lower bound uses
$\log(1+y)\ge2y/(2+y)$ for $y\ge0$, which follows by
differentiation. Hence \eqref{ref5-eq:Chalk-competitor} exceeds $n^n$
for every integer $n\ge6$.
Proposition~\ref{ref5-prop:six-gap} therefore separates the regular hexagon
from possible global maximizers, rather than proving that it is
globally optimal. A full six-point sharp bound would require a new
global comparison and additional constraints beyond the
two-dimensional projection.



\section{Extending the projection and logarithm-majorant methods}
\subsection{Recovering the projection bound with more real directions}
The following lemma supplies a direct extension of the main
projection idea. Its rank condition is a condition to check, rather
than an automatic consequence of larger point count.

\begin{lemma}\label{lem:alignment-real-projection}
Let $E\subset\C^n$ be invariant under conjugation, let $v\in E^\perp$,
and let $u_1,\ldots,u_r$ be linearly independent real vectors in
$E^\perp$. Suppose
$\dim_{\C}E^\perp=r+1$. Set
\[
 G_{ij}=u_i^{\mathsf T}u_j,\qquad b_j=u_j^{\mathsf T}v.
\]
Then $G$ is real positive definite and
\[
 |v^{\mathsf T}v|\ge
 \lVert v\rVert^2-b^*G^{-1}b-|b^{\mathsf T}G^{-1}b|.
\]
\end{lemma}
\begin{proof}
Let $a=\sum_j u_j(G^{-1}b)_j$ be the orthogonal projection of $v$
onto the complex span of the real vectors. Its squared norm is
$b^*G^{-1}b$ and its transpose square is $b^{\mathsf T}G^{-1}b$.
The remaining conjugation-invariant one-dimensional space has a real
unit vector $e$, so $v=a+\gamma e$. Orthogonality gives
$|\gamma|^2=\lVert v\rVert^2-b^*G^{-1}b$ and
$v^{\mathsf T}v=b^{\mathsf T}G^{-1}b+\gamma^2$.
The reverse triangle inequality proves the claim.
\end{proof}

At a stationary local maximizer for general $n$, let
$E=\operatorname{span}_{\C}\{\mathbf1,z,\overline z\}$,
$q=|p_2|^2$, $p=|p_3|^2$, $T=n+s$, and
$v=\Pi_{E^\perp}z^{\circ2}$. For $s>0$ the radial bound gives
$q<n^2$, so $\dim E=3$. The general projection identities are
\[
 \lVert v\rVert^2=T-\frac qn-\frac{np}{n^2-q},\qquad
 v^{\mathsf T}v=p_4-\frac{p_2^2}{n}
                    +\frac{\overline{p_2}p_3^2}{n^2-q}.
\]
Choose $n-4$ independent real projected radial vectors
$u_j=\Pi_{E^\perp}t^{\circ k_j}$. The lemma gives
\begin{equation}\label{eq:alignment-general-projection}
 |p_4|\ge T-\frac{2q}{n}-\frac{p}{n-\sqrt q}
       -b^*G^{-1}b-|b^{\mathsf T}G^{-1}b|.
\end{equation}
For $n=5$, taking $u_1=t-\mathbf1$ gives $G=s$ and
$b_1=3p_2/4$, exactly recovering
Lemma~\ref{ref5-lem:projection}. For $n=6,7,8$ the required
numbers of independent real directions are two, three, and four.

All entries are explicit mixed moments. Put
$T_j=\sum_i t_i^j$, $R_j=\sum_i t_i^jz_i$,
$W_j=\sum_i t_i^jz_i^2$, and
\[
 \alpha_j=\frac{n\overline{R_j}-\overline{p_2}R_j}{n^2-q}.
\]
For $u_j=\Pi_{E^\perp}t^{\circ j}$, one has
\begin{align*}
u_j={}&t^{\circ j}-\frac{T_j}{n}\mathbf1
                  -\alpha_jz-\overline{\alpha_j}\overline z,\\
G_{ij}={}&T_{i+j}-\frac{T_iT_j}{n}
 -\frac{2n\operatorname{Re}(\overline{R_i}R_j)
        -2\operatorname{Re}(\overline{p_2}R_iR_j)}{n^2-q},\\
b_j={}&W_j-\frac{T_jp_2}{n}-\alpha_jp_3.
\end{align*}
The last formula uses $R_1=0$, a stationary mixed moment identity.
The determinant and adjugate of $G$ give a polynomial formulation
after clearing its positive determinant.

If every radius layer has zero centroid and there are at least
$n-3$ distinct radii, the default choices $k_j=j$, $1\le j\le n-4$,
are automatically independent. Indeed, their dependence would make
a nonconstant polynomial $f(t)$ of degree at most $n-4$ equal to
$a+2\operatorname{Re}(\beta z_i)$ at every point. Averaging over each
centered layer makes $f(t)$ equal to $a$ at at least $n-3$ distinct
values, which is impossible. Without this layer-centroid hypothesis,
rank must be checked. Rank failure still provides the exact relation
$f(t_i)=a+2\operatorname{Re}(\beta z_i)$ for a separate geometric
branch analysis.

\subsection{A logarithm majorant adapted to a candidate configuration}
\begin{lemma}\label{lem:alignment-node-majorant}
Let $0<d_1<\cdots<d_k<R$. Let $H$ be the degree at most $2k-1$
polynomial interpolating $\log x$ and its derivative at the $d_j$.
Define
\[
 \kappa_R=\frac{H(R)-\log R}{\prod_j(R-d_j)^2}>0.
\]
For $0<\kappa\le\kappa_R$ and $0<x\le R$,
\[
 \log x\le H(x)-\kappa\prod_j(x-d_j)^2.
\]
\end{lemma}
\begin{proof}
The even derivative $(\log x)^{(2k)}$ is negative, so the double-node
Hermite remainder gives $H(R)>\log R$. Adding the node $R$ gives
the interpolant $H(x)-\kappa_R\prod_j(x-d_j)^2$. Its remainder
has the sign of $(x-R)\prod_j(x-d_j)^2$, because
$(\log x)^{(2k+1)}>0$. It is nonpositive on $0<x\le R$.
Decreasing $\kappa_R$ preserves the upper bound.
\end{proof}

For a central regular $(n-1)$-gon at moment $n$, the distance nodes
are $n/(n-1)$ and
$4n\sin^2(j\pi/(n-1))/(n-1)$. There are three distinct nodes for
$n=6$ and for $n=7$; for $n=7$ the first chord node coincides with
the center-to-ring distance. There are four distinct nodes for $n=8$.
The corresponding majorant degrees are therefore six, six, and eight.
The uniform distance-moment expansion
\[
 \sum_{i<j}|z_i-z_j|^{2r}
 =\frac12\sum_{\alpha,\beta=0}^r(-1)^{\alpha+\beta}
  \binom r\alpha\binom r\beta
  M_{r-\alpha,r-\beta}M_{\alpha\beta},\qquad
  M_{ab}=\sum_i z_i^a\overline{z_i}^{\,b},
\]
then supplies a finite moment comparison through $r=6$ or $8$.
Together with the second-order Gram hierarchy and
\eqref{eq:alignment-general-projection}, this defines a concrete
extension of the five-point strategy. A certified nonpositive final
scalar bound remains to be constructed; the lemma alone does not
prove optimality of the central polygon.

\subsection{The local angular spectral argument has a precise limit}
For a unit regular $n$-gon, the weighted Laplacian with weights
$1/(4\sin^2(\pi r/n))$ has eigenvalues
\[
                         \frac{k(n-k)}2,
                         \qquad 0\le k<n.
\]
To verify the formula, sum
$|\sum_{j=0}^{k-1}e^{2\pi irj/n}|^2$ over $r=0,\ldots,n-1$;
orthogonality gives $nk$, and removal of the $r=0$ term gives
$k(n-k)$. Thus the largest eigenvalues for $n=5,6,7,8$ are
$3,9/2,6,8$. The strict comparison with $n-1$ used for the local
radial gap holds at five and six, becomes an equality at seven,
and fails at eight. This limits that specific proof mechanism.

The original coordinatewise $B$ recurrence also gives, at a unit
regular polygon,
\[
 Bz^{\circ k}=-\frac{(k-2)(n-k)}2z^{\circ(k-2)},
 \qquad 2\le k\le n-1.
\]
For $n=8$, $k=5$ forces $\lVert B\rVert_{\mathrm{op}}\ge9/2>7/2$,
so the regular octagon violates the second-order necessary condition
for a local maximum. This is a valid eight-point exclusion, while the
seven-point equality case requires further analysis. Neither statement
exhausts the remaining eight-point configurations.

\subsection{A uniform generating identity for the dual force}
% Optional general-n lemma replacing the expanded mixed-trace necessity proof.
\begin{lemma}[Generating identity for the double-force equations]
\label{lem:general-dual-resultant-identity}
Let $p$ be a monic polynomial of degree $n\ge2$ with distinct roots $x_i$,
let $\delta=\operatorname{Res}(p,p')$, and put
$g_i=p''(x_i)/p'(x_i)$. Assume the $g_i$ are also distinct. Define
\[
 R_0(Y)=\operatorname{Res}_X(p,Yp'-p''),\qquad
 R_t(Y)=\left.\partial_T\operatorname{Res}_X
 (p,(Y+TX)p'-p'')\right|_{T=0}.
\]
Then $x_i=2\sum_{j\ne i}(g_i-g_j)^{-1}$ for every $i$ if and only if
\begin{equation}\label{eq:general-dual-resultant-identity}
 R_t(Y)=R_0''(Y).
\end{equation}
For centered $p$, its nontrivial coefficient equations are
\[
 E_r=[Y^{n-2-r}]R_t-(n-r)(n-r-1)[Y^{n-r}]R_0=0,
 \qquad 1\le r\le n-2.
\]
Under root scaling $x_i\mapsto\alpha x_i$, $E_r$ has weight
$n(n-1)-r$. Every physical stationary shape satisfies these equations,
and its $g_i$ are distinct.
\end{lemma}
\begin{proof}
Let $H(Y)=\prod_i(Y-g_i)$. The product formula gives
$R_0=\delta H$ and $R_t=\delta\sum_i x_iH/(Y-g_i)$.
For the forward implication, summing the two terms of each pair yields
\[
 \frac{R_t}\delta
 =2\sum_{i<j}\frac{H}{(Y-g_i)(Y-g_j)}=H''.
\]
Conversely, evaluate the identity at $Y=g_i$ and divide by
$\delta H'(g_i)\ne0$: one gets
$x_i=H''(g_i)/H'(g_i)=2\sum_{j\ne i}(g_i-g_j)^{-1}$.

Always $\sum_i g_i=0$ and $\sum_i x_ig_i=n(n-1)$ by pairwise
summation. If $p$ is centered, $\sum_i x_i=0$ as well. These are
the leading coefficient identities; the displayed $E_r$ give all
remaining coefficients. Scaling sends $\delta$ to
$\alpha^{n(n-1)}\delta$, $g_i$ to $\alpha^{-1}g_i$, and each relevant
mixed sum $\sum_i x_i g_i^{r+1}$ to $\alpha^{-r}$ times itself,
proving the weight claim.

For a physical stationary shape of moment $Q>0$, write
$\kappa=n(n-1)/Q$. Then $g_i=\kappa\overline{x_i}$, so the $g_i$
are distinct and
$\sum_{j\ne i}(g_i-g_j)^{-1}
=\kappa^{-1}\overline{\sum_{j\ne i}(x_i-x_j)^{-1}}=x_i/2$.
\end{proof}

The generating identity encodes the double-force equations. Physical feasibility still requires the positive conjugate pairing of the original and dual roots.


\paragraph{Computing the first variation without its higher powers.}
Modulo monic $p$, replace the second resultant polynomial by
$Yp'-p''+T(Xp'-np)$, of degree at most $n-1$.
Its $(2n-1)$-square Sylvester matrix has $n$ rows that depend on $T$.
Writing it as $M_0+TM_1$, determinant multilinearity gives
\[
 R_t(Y)=\sum_{j\in J}\det M_0[j\leftarrow(M_1)_{j,:}],
 \qquad |J|=n.
\]
Thus only $n$ row-replacement determinants are needed; no higher
$T$ coefficient is expanded. For the next eight-point chart
$p=X^8+X^4+bX^3+cX^2+dX+u$, this exact method generated $R_t$
and all six necessary coefficient equations in about fourteen seconds.
The full symbolic data and generator are supplied in the verification
package. Sixteen independent rational sample checks are diagnostic,
not a complete interpolation certificate. No physical exclusion or
objective bound for this four-parameter chart follows from this
generation step.

% Standalone fragment; usual amsmath/amssymb/amsthm preamble suffices.
\section{Finite projection and Fourier frames without a shape classification}
\label{sec:noclass-frames}

Throughout this section the collision-free configuration is centered
and normalized by $\sum_i|z_i|^2=n$. Write
\[
 p_k=\sum_i z_i^k,\quad t_i=|z_i|^2,\quad
 T_j=\sum_i t_i^j,\quad s=T_2-n,\quad
 R_j=\sum_i t_i^jz_i,\quad u=(t_i-1)_i .
\]
At a stationary configuration put $\lambda=(n-1)/2$ and define the
complex symmetric Laplacian
\[
 (Bv)_i=\sum_{j\ne i}\frac{v_i-v_j}{(z_i-z_j)^2}.
\]
The force equations give $Bz=\lambda\overline z$ and, for $k\ge1$,
\begin{equation}\label{eq:noclass-mixed-force}
 \sum_i t_i z_i^k
 =\frac{(2n-k-1)p_k+\sum_{r=1}^{k-1}p_rp_{k-r}}{n-1}.
\end{equation}
In particular, $R_1=0$. Formula \eqref{eq:noclass-mixed-force}
follows by multiplying the force equation by $z_i^{k+1}$ and
summing the divided differences over unordered pairs.

\begin{lemma}\label{lem:noclass-finite-projection}
Suppose $n\ge5$, $s>0$, and $|p_2|<n$. Let
$E=\operatorname{span}_{\mathbb C}\{\mathbf1,z,\overline z\}$ and
$W=(E\oplus\mathbb C u)^\perp$. For $2\le k\le n-3$, put
$v_k=\Pi_W z^{\circ k}$ and define
$H_{kl}=v_k^*v_l$, $C_{kl}=v_k^Tv_l$. Then
\[
                         \det H=|\det C|.
\]
All entries are finite moment expressions. More explicitly, put
\[
 K=\begin{pmatrix}n&\overline{p_2}\\p_2&n\end{pmatrix},
 \quad J=\begin{pmatrix}0&1\\1&0\end{pmatrix},
 \quad h_k=\binom{a_k}{p_{k+1}},
\]
\[
 a_k=\frac{(2n-k)p_{k-1}+
                 \sum_{r=1}^{k-2}p_rp_{k-1-r}}{n-1},\qquad
 \beta_k=\frac{(n-k)p_k+
                 \sum_{r=1}^{k-1}p_rp_{k-r}}{n-1}.
\]
With $M_{ab}=\sum_i z_i^a\overline z_i^{\,b}$, the entries are
\[
 H_{kl}=M_{l,k}-\frac{\overline{p_k}p_l}{n}
          -h_k^*K^{-1}h_l-\frac{\overline{\beta_k}\beta_l}{s},
\]
\[
 C_{kl}=p_{k+l}-\frac{p_kp_l}{n}
          -h_k^TJK^{-1}h_l-\frac{\beta_k\beta_l}{s}.
\]
\end{lemma}
\begin{proof}
Equation \eqref{eq:noclass-mixed-force} gives $u\perp E$.
The assumptions make $\dim E=3$ and $\dim W=n-4$.
The displayed formulas are the ordinary orthogonal projection
formulas; the bilinear formula uses $K^{-T}J=JK^{-1}$.
The space $W$ is invariant under complex conjugation, so it has a
real orthonormal basis. In that basis the $n-4$ vectors $v_k$ form a
square coordinate matrix $A$, possibly singular. Thus $H=A^*A$ and
$C=A^TA$, and both sides of the determinant identity equal
$|\det A|^2$. No invertibility of $A$ is assumed.
Moments $p_k$ with $k>n$ can be eliminated by Newton's identities;
this is a finite algebraic interface, not a global upper bound by itself.
\end{proof}

For actual local maxima the condition $|p_2|<n$ is automatic when
$n\ge3$. Equality in $|p_2|\le\sum_i|z_i|^2=n$ would make the
configuration collinear after rotation. Then $u=z^2-\mathbf1\ne0$
(there cannot be three distinct points among $\{-1,1\}$), and
$Bu=(n-1)u$, contradicting $\|B\|_{\rm op}\le(n-1)/2$.
The general operator bound follows by the same horizontal
second-variation and polarization argument given explicitly in
the six-point proof below.

\begin{lemma}\label{lem:noclass-fourier-frame}
For $n=6,7,8$ and $s>0$, respectively use the normalized columns
\[
 \begin{array}{c|l}
6&\mathbf1/\sqrt6,\ z/\sqrt6,\ \overline z/\sqrt6,\
 z^2/\sqrt{T_2},\ \overline z^2/\sqrt{T_2},\ u/\sqrt s\\
7&\mathbf1/\sqrt7,\ z/\sqrt7,\ \overline z/\sqrt7,\
 z^2/\sqrt{T_2},\ \overline z^2/\sqrt{T_2},\
 z^3/\sqrt{T_3},\ u/\sqrt s\\
8&\mathbf1/\sqrt8,\ z/\sqrt8,\ \overline z/\sqrt8,\
 z^2/\sqrt{T_2},\ \overline z^2/\sqrt{T_2},\
 z^3/\sqrt{T_3},\ \overline z^3/\sqrt{T_3},\ u/\sqrt s .
 \end{array}
\]
Let $V$ denote this square matrix and
$\varepsilon_n^2=\|V^*V-I\|_F^2$. Put $T=T_2$, $U=T_3$, $q=|p_2|^2$.
Then
\[
\begin{aligned}
\varepsilon_6^2={}&q\left(\frac1{18}+\frac2{3T}+\frac{64}{25sT}\right)
 +\frac{2|p_3|^2}{3T}+\frac{2|p_4|^2}{T^2},\\
\varepsilon_7^2={}&q\left(\frac4{7T}+\frac2{49}
                  +\frac{121}{126U}+\frac{25}{9sT}\right)\\
 &+|p_3|^2\left(\frac2{7U}+\frac4{7T}+\frac8{9sU}\right)
  +|p_4|^2\left(\frac2{T^2}+\frac2{7U}\right)
  +\frac{2(|R_2|^2+|p_5|^2)}{TU},\\
\varepsilon_8^2={}&q\left(\frac1{2T}+\frac1{32}
                  +\frac{169}{98U}+\frac{144}{49sT}\right)\\
 &+|p_3|^2\left(\frac1{2U}+\frac1{2T}+\frac{100}{49sU}\right)
  +|p_4|^2\left(\frac2{T^2}+\frac1{2U}\right)\\
 &+\frac{4(|R_2|^2+|p_5|^2)}{TU}
  +\frac{2|p_6|^2}{U^2}.
\end{aligned}
\]
The squared row norms are
\[
\begin{aligned}
 P_6(t)&=\frac16+\frac t3+\frac{2t^2}{T}+\frac{(t-1)^2}{s},\\
 P_7(t)&=\frac17+\frac{2t}7+\frac{2t^2}{T}
                   +\frac{t^3}{U}+\frac{(t-1)^2}{s},\\
 P_8(t)&=\frac18+\frac t4+\frac{2t^2}{T}
                   +\frac{2t^3}{U}+\frac{(t-1)^2}{s}.
\end{aligned}
\]
They obey
\[
                    \sum_i(P_n(t_i)-1)^2\le\varepsilon_n^2.
\]
If $\rho=n/(n-1)$ and $s>\rho$, also
\[
                 (s-\rho)^2\le n^2\rho^2s\,\varepsilon_n^2.
\]
\end{lemma}
\begin{proof}
Every column has norm one. Summing twice the squared absolute
values of the off-diagonal Gram entries, and using
\eqref{eq:noclass-mixed-force}, gives the formulas above.
For a square matrix,
$\|VV^*-I\|_F^2=\|V^*V-I\|_F^2$; retaining its diagonal proves
the discrepancy inequality.
For the last assertion, $P_n(t)-1$ has negative constant term
$1/s-1/\rho$, positive coefficients in degrees two and three,
and exactly one positive root $r<n$ (use strict convexity and
$P_n(n)>1$). Division by $t-r$ gives a polynomial whose coefficients
are positive, and whose constant term is $(1/\rho-1/s)/r$.
Consequently
$|P_n(t)-1|\ge(1/\rho-1/s)|t-r|/n$ for $t\ge0$.
Finally $\sum_i(t_i-r)^2\ge s$ proves the stated bound.
\end{proof}

\begin{lemma}\label{lem:noclass-zero-harmonics}
Let $n\ge4$. If a collision-free centered stationary configuration
has $p_2=\cdots=p_{n-2}=0$, it is either a unit regular $n$-gon or
the origin together with a regular $(n-1)$-gon of squared radius
$n/(n-1)$.
\end{lemma}
\begin{proof}
Newton's identities give its monic root polynomial
$P(z)=z^n-a z-b$. The force equation is
$P''(z_i)=(n-1)\overline z_i P'(z_i)$. Multiplying by $z_i^2$ and
using $P(z_i)=0$ gives, without dividing by a root,
\[
                  a((n-1)t_i-n)z_i+n b(t_i-1)=0.
\]
If $a,b\ne0$, the coefficient $(n-1)t_i-n$ cannot vanish:
at its zero, $t_i-1=1/(n-1)$ and the second term is nonzero.
Thus all $z_i$ lie on the line of $b/a$. This contradicts
$p_2=0$ and $\sum_i|z_i|^2=n$. If $a=0$, collision-freeness forces
$b\ne0$ and gives a regular $n$-gon. If $b=0$, it forces $a\ne0$
and gives a simple origin plus a regular $(n-1)$-gon. Normalization
fixes the radii in both cases.
\end{proof}

\subsection{A quantitative six-point region}

\begin{theorem}\label{thm:noclass-small6}
Let $z$ be a collision-free local maximum for $n=6$ and put
$h=1/14000$. If
\[
                   \max(|p_2|,|p_3|,|p_4|)\le h,
\]
then it is either a unit regular hexagon or the origin together
with a regular pentagon of squared radius $6/5$.
In particular, a noncircular local maximum in this region is the
central pentagon. A global maximum in this region is the central
pentagon and has value $C_6=6^{15}/5^{10}$.
No restriction on the initial shape or radial population is imposed.
\end{theorem}

\begin{proof}
The local second-order condition gives $\|B\|_{\rm op}\le5/2$.
Indeed, the horizontal tangent space is
$\{\mathbf1,z\}^{\perp}$; applying the spherical Hessian to a
direction and to every phase multiple of it bounds its complex
symmetric quadratic form by $(5/2)\|v\|^2$.
The relation $Bz=(5/2)\overline z$ and symmetry make this
quadratic form orthogonal to the $z$ direction. Polarization gives
the same operator bound on $\mathbf1^\perp$, and $B\mathbf1=0$.
Writing $g_i=e_i-\mathbf1/6$ gives
$|B_{ij}|=|g_i^*Bg_j|\le(5/2)(5/6)$ for $i\ne j$, hence
\begin{equation}\label{eq:noclass-six-basic}
 |z_i-z_j|^2\ge\frac{12}{25}.
\end{equation}
Also $Bz^2=5u$. Apply the operator bound to
$z^2-(p_2/6)\mathbf1$ to obtain
$4s\le T-|p_2|^2/6$, where $T=6+s$; thus $s\le2$.

We record a common uniqueness tube for the two candidate orbits.
For a candidate $c$, the smallest pair distance is at least one,
$Bc=(5/2)\overline c$, and the operator norm on $c^\perp$
is less than $21/10$. For the unit hexagon this follows from the
orthogonal Fourier basis and
$Bc^k=-(k-2)(6-k)c^{k-2}/2$ for $2\le k\le5$;
the largest transverse singular value is $2$.
For the central pentagon use the orthogonal basis
$\mathbf1,c,\overline c,c^2,\overline c^2,u$:
\[
 B\mathbf1=B\overline c=B\overline c^2=0,\quad
 Bc=(5/2)\overline c,\quad Bc^2=5u,\quad
 Bu=(5/6)\overline c^2 .
\]
The largest transverse singular value is $5/\sqrt6<21/10$.
If $\|x-c\|\le\eta=1/100$, every pair segment stays at distance
at least $1-2\eta$, and the derivative of the inverse-square
function has modulus at most $2/(1-2\eta)^3<5/2$.
Each edge weight changes by less than $4\eta$. The identity
\[
 v^T(B(x)-B(c))v
  =\sum_{i<j}a_{ij}(v_i-v_j)^2,\qquad
 \sum_{i<j}|v_i-v_j|^2\le6\|v\|^2
\]
and polarization imply $\|B(x)-B(c)\|_{\rm op}<24\eta$.
For a horizontal vector $x^*v=0$, its component in $\mathbb Cc$
has squared norm at most $\eta^2\|v\|^2/6$. Thus
\[
 |v^TB(x)v|
 <\left(\frac{21}{10}+\frac{5\eta^2}{12}+24\eta\right)\|v\|^2
 <\frac{12}{5}\|v\|^2 .
\]
For $f=\log\Delta_6$ the degree of homogeneity is $30$; its
horizontal spherical Hessian therefore satisfies
\[
       f''\le 2|v^TB(x)v|-5\|v\|^2\le-\frac15\|v\|^2.
\]
After choosing the phase so $c^*x$ is positive real, the short
great circle from $c$ to $x$ is horizontal and stays in this tube.
Since $c$ is stationary, its strictly decreasing directional
derivative excludes any different stationary endpoint. Moreover,
\begin{equation}\label{eq:noclass-six-tube}
 \log\Delta_6(x)\le\log\Delta_6(c)
             -\frac1{10}\operatorname{dist}(x,\mathbb S^1c)^2
\end{equation}
within the tube.

First handle $s\le1/200$; this supplies a self-contained small-circle
step rather than assuming a previously proved radial gap.
Then $13/14<t_i<15/14$ and $24/25<|z_i|<21/20$.
Newton's identities write
\[
 P(z)=z^6-a z-b+E(z),\qquad
 E(z)=c_4z^4+c_3z^3+c_2z^2,
\]
with $|c_4|\le h/2$, $|c_3|\le h/3$, $|c_2|<h/3$.
The exact force equation gives
\[
 a(5t_i-6)z_i+6b(t_i-1)
 =\sum_{k=2}^4c_kz_i^k
                \left((6-k)t_i-6+\frac{k(k-1)}5\right).
\]
The right side has modulus less than
\[
 h\left[\frac12(21/20)^4\frac74
       +\frac13(21/20)^3\frac{21}{10}
       +\frac13(21/20)^2\,2\right]<3h.
\]
If $a,b\ne0$, rotate along $b/a$. Since $|p_2|\le h$, some point
has imaginary part of modulus greater than $2/3$.
As $|5t_i-6|\ge9/14$, the imaginary part of the preceding identity
gives $|a|<7h$. If $b=0$ the modulus and $|z_i|>24/25$ give an
even stronger bound; if $a=0$ the same conclusion is immediate.
Also $|E(z_i)|<3h/2$, so $|z_i^6-b|<9h$ and
$|b|>(24/25)^6-9h>3/4$.
The six roots of $z^6-b$ have common radius $r_b>19/20$.
For the nearest such root let $\delta$ be its distance from $z_i$.
Initially $\delta^6\le9h<(1/3)^6$, so $\delta<1/3$.
The product of the other five target separations is $6r_b^5$.
The reverse triangle inequality therefore gives
\[
 \delta<
 \frac{9h}{6(19/20)^5(37/57)^5}<18h<1/500.
\]
Repeating with the improved distance gives
\[
 \delta<\frac{9h}{6(19/20)^5(474/475)^5}<2h.
\]
By \eqref{eq:noclass-six-basic} the matching is injective, hence
a bijection. The Euclidean norm comparison with $\sum_i|z_i|^2=6$
now gives $|r_b-1|<2h$. Each point is within $4h$ of a rotated
unit hexagon and the total distance is less than $4\sqrt6h<1/100$.
The uniqueness tube forces exactly that hexagon, so $s=0$.
Consequently a noncircular local maximum under the stated harmonic
condition has $s>1/200$.

Now suppose $s>1/200$. Use the six-column frame of
Lemma~\ref{lem:noclass-fourier-frame}; write $\varepsilon=\varepsilon_6$.
Since $6\le T\le8$,
$\varepsilon^2<(257/3)h^2<86h^2$.
The discrepancy gives
\[
 \left(\frac2T+\frac1s\right)R_2
          =\sum_i(P_6(t_i)-1)z_i,\qquad |R_2|<31h.
\]
It also gives $t_i<3$: otherwise
$P_6(t_i)\ge t_i/3+t_i^2/4\ge13/4>1+\varepsilon$.
Put $\gamma=p_5/T$, $g=|\gamma|$, and
$r=z^3-\gamma\overline z^2$. Its squared norm is
\[
                  d=\|r\|^2=U-|p_5|^2/T,\qquad U=T_3.
\]
Here $g<\sqrt3$. The six components of $V^*r$ are
\[
 \frac{p_3-\gamma\overline{p_2}}{\sqrt6},\
 \frac{(9/5)p_2-\gamma\overline{p_3}}{\sqrt6},\
 \frac{p_4}{\sqrt6},\
 \frac{R_2-\gamma\overline{p_4}}{\sqrt T},\
 0,\
 \frac{(3/5)p_3-(4/5)\gamma\overline{p_2}}{\sqrt s}.
\]
The harmless larger bound $g<2$ gives
\[
 \|V^*r\|^2\le
 \left[\frac96+\frac{(19/5)^2}{6}+\frac16
       +\frac{33^2}{6}+200(11/5)^2\right]h^2.
\]
Since $VV^*\ge(1-\varepsilon)I$, $d<1156h^2$.
The inequalities $U\ge T^2/6$ and $T\ge6$ imply $g>99/100$.
Importantly, reverse triangle gives an inequality, not an equality:
\begin{equation}\label{eq:noclass-six-residual}
                \sum_i t_i^2(\sqrt{t_i}-g)^2\le d.
\end{equation}
If every $t_i\ge a_0=3/25$, then
\[
 s\le\sum_i(t_i-g^2)^2
 \le12\sum_i(\sqrt{t_i}-g)^2
 \le12d/a_0^2
 <\frac{2890000}{3}h^2<1/200,
\]
a contradiction. Equation \eqref{eq:noclass-six-basic} allows
at most one point with $t_i<a_0$. Thus exactly one such point
exists; denote it by $z_0$.
Its radius is less than $7/20$, so
$t_0<34h/(99/100-7/20)<54h$.
For the other five points,
\[
 |\sqrt{t_i}-g|<34h/a_0<284h,\qquad
 |t_i-g^2|<1000h,\qquad |g^2-6/5|<1011h.
\]
The last bound gives $21/20<g<23/20$, and the five radii belong
to $(1,6/5)$. The radial moment alone also gives
\[
 s=t_0^2+\frac{(6-t_0)^2}{5}-6
       +\sum_{\rm ring}(t_i-\overline t_{\rm ring})^2
 \ge\frac65-\frac{12}{5}t_0>1,
\]
and $g^2-t_0>1$.

Repeat the frame estimates using $s>1$. They give
$\varepsilon<h$, $|R_2|<4h$, and $d<16h^2$.
Since each ring $t_i>1$, \eqref{eq:noclass-six-residual} yields
$|\sqrt{t_i}-g|<4h$ and $|t_i-g^2|<10h$.
Using $R_1=0$ in the form
\[
       (g^2-t_0)z_0=\sum_{\rm ring}(t_i-g^2)z_i
\]
gives $|z_0|<60h$.
The moment now sharpens to
$|g^2-6/5|\le10h+720h^2<11h$; each ring radius differs
from $\sqrt{6/5}$ by less than $10h$.

Let $w_1,\ldots,w_5$ be the ring points and $e_k$ their elementary
symmetric functions. Their first power sum has modulus less than
$60h$, and the next three have modulus less than $2h$.
Newton's identities give
$|e_1|<60h$, $|e_2|<2h$, $|e_3|,|e_4|<h$.
Thus $|w_i^5-e_5|<131h$, while $e_5\ne0$.
The target roots $\alpha_j^5=e_5$ have radius $r_a>1$,
minimum separation greater than one, and derivative product
$\prod_{j\ne k}|\alpha_j-\alpha_k|=5r_a^4>5$.
For the nearest target root, $\delta^5\le131h<(2/5)^5$.
The reverse triangle inequality then gives
\[
 \delta<\frac{131h}{5(3/5)^4}<203h<1/50,\qquad
 \delta<\frac{131h}{5(49/50)^4}<29h.
\]
The matching is injective by \eqref{eq:noclass-six-basic}.
The geometric mean radius $r_a$ differs from $\sqrt{6/5}$
by less than $10h$. Hence each ring point is within $39h$ of
the corresponding central-pentagon candidate, and the center is
within $60h$. The total squared distance is less than
\[
                    (60^2+5\cdot39^2)h^2<1/10000.
\]
The uniqueness tube again forces the exact candidate.
Finally $C_6/6^6=6^9/5^{10}>1$, so the hexagon cannot be
the global maximum when the harmonic condition holds.
\end{proof}

All numerical inequalities in this proof are rational or integer
comparisons replayed by \path{noclass_checks.py}.
The theorem leaves the complementary region
$\max(|p_2|,|p_3|,|p_4|)>1/14000$ open for this method.

% Standalone fragment; uses only the usual amsmath/amssymb/amsthm preamble.
\section{A direct global bound for six and seven points containing the origin}
\label{sec:noclass-origin67}

The following argument treats every radius and angle simultaneously.
It does not assume stationarity, a reflection, or a prescribed number
of radial layers.

\begin{theorem}\label{thm:noclass-origin67}
Let $m\in\{5,6\}$ and $w_1,\ldots,w_m\in\mathbb C$ satisfy
$\sum_i|w_i|^2=m$. Put
\[
 W(w)=\prod_i|w_i|^2\prod_{i<j}|w_i-w_j|^2,\qquad
 D=\sum_i(|w_i|^2-1)^2.
\]
Then
\[
 W(w)\le m^m\exp\!\left(-\frac{D}{2m^3}\right)\le m^m.
\]
Equality in the final bound holds exactly for a rotated unit regular
$m$-gon. Consequently, for $n=6,7$, every configuration containing
the origin and satisfying $\sum_i|z_i|^2=n$ obeys
\[
 \Delta_n(z)\le
 C_n:=\frac{n^{n(n-1)/2}}{(n-1)^{(n-1)(n-2)/2}}.
\]
If the other $n-1$ points are denoted by $z_i$, the sharper bound is
\[
 \Delta_n(z)\le C_n\exp\!\left[
 -\frac{1}{2(n-1)^3}
 \sum_i\left(\frac{n-1}{n}|z_i|^2-1\right)^2\right].
\]
The equality configuration is the origin together with a regular
$(n-1)$-gon of squared radius $n/(n-1)$.
\end{theorem}

\begin{proof}
If a point is zero or two points coincide, $W=0$, so suppose otherwise.
Write $t_i=|w_i|^2$ and
\[
 S_m(t)=\sum_{k=1}^m t^k,\qquad
 P_m(t)=S_m'(t)^2-S_m(t)S_m''(t).
\]
We first prove
\begin{equation}\label{eq:noclass-origin-curvature}
 -(\log S_m)''(t)=\frac{P_m(t)}{S_m(t)^2}
 \ge\frac1{m^3}\qquad(0<t\le m).
\end{equation}
The ascending coefficient lists of $P_m$ are
\[
\begin{array}{c|l}
m&[t^0,t^1,\ldots]P_m\\ \hline
5&(1,2,2,0,-5,4,8,8,5)\\
6&(1,2,2,0,-5,-14,0,8,11,10,6).
\end{array}
\]
For $0\le t\le1$, express
$P_m(t)=\sum_{k=0}^{d}b_k\binom dk t^k(1-t)^{d-k}$.
The exact Bernstein coefficient lists are
\[
\begin{array}{c|c|l}
m&d&(b_0,\ldots,b_d)\\ \hline
5&8&
\left(1,\frac54,\frac{11}7,\frac{55}{28},\frac{33}{14},
\frac{75}{28},\frac{45}{14},\frac{25}4,25\right)\\[2pt]
6&12&
\left(1,\frac76,\frac{15}{11},\frac{35}{22},\frac{182}{99},
\frac{91}{44},\frac{145}{66},\frac{827}{396},
\frac{262}{165},\frac7{10},\frac7{33},\frac72,21\right).
\end{array}
\]
Every coefficient is at least $1/m$. Since the Bernstein basis sums
to one and $S_m(t)\le m$ on this interval, the desired curvature
bound follows there. These are polynomial identities, verifiable by
the elementary formula
\[
 b_k=\sum_{i=0}^{\min(k,\deg P_m)}
 [t^i]P_m\,\frac{\binom ki}{\binom di}.
\]
For $t\ge1$, put $x=t-1$. The ascending coefficient lists of
$P_m(1+x)-m(1+x)^{2m-2}$ are
\[
\begin{array}{c|l}
5&(20,110,300,460,415,220,64,8,0)\\
6&(15,150,668,1576,2235,2030,1204,456,101,10,0).
\end{array}
\]
Thus $P_m(t)\ge m t^{2m-2}$ for $t\ge1$. On $1\le t\le m$,
$S_m(t)\le m t^m$, giving
$P_m/S_m^2\ge1/(m t^2)\ge1/m^3$ and completing
\eqref{eq:noclass-origin-curvature}.

Taylor's inequality for this strongly concave function gives
\[
 \log S_m(t)\le\log m+\frac{m+1}{2}(t-1)
                   -\frac{(t-1)^2}{2m^3}.
\]
All $t_i$ lie in $(0,m]$, and $\sum_i(t_i-1)=0$. Summing therefore
gives $\prod_iS_m(t_i)\le m^m e^{-D/(2m^3)}$.
For the square Vandermonde matrix $V_{ik}=w_i^k$, $1\le i,k\le m$,
\[
 |\det V|^2=W(w).
\]
Hadamard's inequality on its rows gives
$W\le\prod_iS_m(t_i)$, proving the claimed bounds.

Final equality requires $D=0$ and orthogonal rows. With $t_i=1$,
orthogonality says
$\sum_{k=1}^m(w_i\overline{w_j})^k=0$ for $i\ne j$.
Hence every ratio $w_i/w_j$ is a nontrivial $m$th root of unity,
so the $m$ distinct points form a rotated regular $m$-gon.
Finally use $w_i=\sqrt{(n-1)/n}\,z_i$. There are
$n(n-1)/2$ squared distances in $\Delta_n$, so the scaling factor
is $(n/(n-1))^{n(n-1)/2}$, yielding the stated $C_n$.
\end{proof}

The source \path{noclass_checks.py} replays both derivative
polynomials, both Bernstein identities, and both shifted identities
coefficient by coefficient with standard-library rational arithmetic.
This origin-containing theorem does not assert that every global
maximizer has an origin point.

\section{New rational constraints from subset scaling}
\label{sec:subset-scaling-new}

The following arguments supply necessary conditions for global maximizers
without assuming any polygonal or reflection symmetry. They improve the
separation and radius constants in the existing compact formulation, but
do not by themselves classify the maximizing shapes for arbitrary $n$.
All proofs below are self-contained; no priority assertion is intended.
Put $N=n(n-1)/2$ and $\lambda=(n-1)/2$.

\begin{lemma}[A boundary maximum derivative]\label{lem:boundary-derivative-new}
Let $q$ be a nonconstant polynomial with
$|q(1)|=\max_{|t|=1}|q(t)|$. Then $q(1)\ne0$ and
\[
                 \Re\frac{q'(1)}{q(1)}>0.
\]
\end{lemma}
\begin{proof}
Set $f=q/q(1)$ and $b=f(0)$. The maximum-modulus principle gives
$|f|\le1$ in the unit disk, and $|b|<1$ since $f$ is nonconstant.
Schwarz--Pick, applied to $f$ at zero, gives for $0<r<1$
\[
 \left|\frac{f(r)-b}{1-\overline b f(r)}\right|\le r.
\]
Using the elementary identity for this disk automorphism yields
\[
 1-|f(r)|^2\ge
 (1-r^2)\frac{|1-\overline b f(r)|^2}{1-|b|^2}.
\]
Divide by $2(1-r)$ and let $r\uparrow1$. Since $f(1)=1$ and $f$
is analytic across the boundary, this gives
\[
 \Re f'(1)\ge\frac{|1-b|^2}{1-|b|^2}>0.
\]
\end{proof}

\begin{theorem}[Every-subset internal moment]\label{thm:subset-spread-new}
Let $z_1,\ldots,z_n$ be a global maximizer with
$\sum_i|z_i|^2=n$. For a subset $S$ of size $m$, put
\[
 a_S=\frac1m\sum_{i\in S}z_i,\qquad
 Q_S=\sum_{i\in S}|z_i-a_S|^2.
\]
For every proper subset with $2\le m<n$,
\begin{equation}\label{eq:subset-spread-new}
             Q_S>\frac{m(m-1)}{n-1}.
\end{equation}
For the full set $Q_S=n$, and for a singleton $Q_S=0$.
\end{theorem}
\begin{proof}
A maximizer is centered and collision-free. Homogeneity shows that it
also globally maximizes
\[
 \Phi(z)=\Delta_n(z)\exp\left(-\lambda\sum_i|z_i|^2\right)
\]
over all of $\C^n$: for a fixed nonzero shape of moment $Q$, the best
scale satisfies $s^2Q=N/\lambda=n$.
Fix a proper subset $S$, all outside points, and its centroid $a_S$.
Write $y_i=z_i-a_S$, so $\sum_{i\in S}y_i=0$, and replace the subset
by $a_S+t y_i$. Up to a positive factor independent of $t$, the
conditional objective is
\[
 |P(t)|^2e^{-\alpha|t|^2},\qquad
 P(t)=t^k q(t),\qquad k=\frac{m(m-1)}2,\qquad
 \alpha=\lambda Q_S,
\]
where
\[
 q(t)=\prod_{i\in S,\ j\notin S}(a_S+t y_i-z_j).
\]
The original $t=1$ globally maximizes this expression. On the unit
circle the Gaussian factor and $|t|^{2k}$ are constant, so
$|q(1)|=\max_{|t|=1}|q(t)|$. Stationarity at $t=1$ gives
\[
 \frac{P'(1)}{P(1)}=\alpha,
 \qquad \frac{q'(1)}{q(1)}=\alpha-k.
\]
The polynomial $q$ is nonconstant: some $y_i\ne0$ because $m\ge2$
and the points are distinct; its outside factors have positive degree.
If a $y_i=0$, its outside factors are nonzero constants by
collision-freeness. Thus the degree cannot cancel in the product.
Lemma~\ref{lem:boundary-derivative-new} yields $\alpha>k$, proving
\eqref{eq:subset-spread-new}. The full-set and singleton statements
follow from centering and the prescribed moment.
\end{proof}

\begin{remark}
For closed rational optimization constraints one may relax the strict
inequality to $Q_S\ge m(m-1)/(n-1)$. This weaker version also follows
using only the usual maximum-modulus principle: if
$M_q(r)=\max_{|t|=r}|q(t)|$, then $M_q$ is nondecreasing, while
$r^{2k}e^{-\alpha r^2}$ is strictly increasing for
$r<\sqrt{k/\alpha}$. Consequently a global radial maximizer at $r=1$
requires $\alpha\ge k$.
\end{remark}

\begin{corollary}[Pair and triple separation]\label{cor:pair-triple-new}
For $n\ge3$ every pair in a global maximizer satisfies
\begin{equation}\label{eq:pair-new-rational}
                    |z_i-z_j|^2>\frac4{n-1}.
\end{equation}
For $n\ge4$ every triple $S$ satisfies
\begin{equation}\label{eq:triple-new-rational}
 Q_S>\frac6{n-1},\qquad
 \sum_{\{i,j\}\subset S}|z_i-z_j|^2>\frac{18}{n-1}.
\end{equation}
At force-one scale $w_i=\sqrt\lambda\,z_i$, these become
$|w_i-w_j|^2>2$, and respectively $Q_S(w)>3$ and a triple
squared-distance sum greater than $9$.
\end{corollary}
\begin{proof}
For any $m$ points,
$mQ_S=\sum_{\{i,j\}\subset S}|z_i-z_j|^2$.
Apply Theorem~\ref{thm:subset-spread-new} for $m=2,3$ and rescale.
For $n=2$ the pair squared distance is $4$ in the original scale and
$2$ at force-one scale. Since $e>2$, \eqref{eq:pair-new-rational}
is strictly stronger than $8/[e(n-1)]$ whenever $n\ge3$.
\end{proof}

\begin{corollary}[Radius and every-subset centroid]\label{cor:subset-centroid-new}
For $n\ge3$ and every nonempty proper subset $S$ of size $m$,
\begin{equation}\label{eq:subset-centroid-new}
 \left|\sum_{i\in S}z_i\right|^2
      <\frac{2m^2(n-m)^2}{n(n-1)}.
\end{equation}
In particular,
\begin{equation}\label{eq:radius-new-rational}
                    |z_i|^2<\frac{2(n-1)}n.
\end{equation}
At force-one scale the bounds are
\[
 \left|\sum_{i\in S}w_i\right|^2
       <\frac{m^2(n-m)^2}{n},\qquad
 |w_i|^2<\rho_n:=\frac{(n-1)^2}{n}.
\]
All displayed strict inequalities may be relaxed to nonstrict rational
constraints retaining every global maximizer. For $n=2$ the radius
inequality is equality.
\end{corollary}
\begin{proof}
Let $T$ be the complement. Centering gives
\[
 n=Q_S+Q_T+\frac{n}{m(n-m)}
                    \left|\sum_{i\in S}z_i\right|^2.
\]
Theorem~\ref{thm:subset-spread-new} gives
\[
 Q_S+Q_T>
 \frac{m(m-1)+(n-m)(n-m-1)}{n-1}
 =n-\frac{2m(n-m)}{n-1}.
\]
Singleton blocks have internal moment zero; at least one block has
size at least two when $n\ge3$, so the sum is strict. Rearrangement
proves \eqref{eq:subset-centroid-new}, and $m=1$ gives
\eqref{eq:radius-new-rational}.
\end{proof}

\begin{proposition}[Directional majorization]\label{prop:projection-new}
Fix an angle $\theta$, and sort
$x_i=\Re(e^{-i\theta}z_i)$ in decreasing order. Put
$c=\sqrt{2/[n(n-1)]}$ and $b_i=c(n+1-2i)$. Then $(x_i)$ is
majorized by $(b_i)$: for $1\le m<n$,
\[
 \sum_{i=1}^m x_i\le c m(n-m)=\sum_{i=1}^m b_i,
 \qquad \sum_i x_i=\sum_i b_i=0.
\]
For $n\ge3$ every proper partial-sum inequality is strict. Consequently,
for every convex real function $\varphi$,
\[
                   \sum_i\varphi(x_i)\le\sum_i\varphi(b_i).
\]
In particular, for $n\ge3$,
\begin{equation}\label{eq:quadrupole-new}
 \sum_i x_i^2<\frac{2(n+1)}3,\qquad
 \left|\sum_i z_i^2\right|<\frac{n+4}3.
\end{equation}
For real $t\ne0$,
\[
 \sum_i e^{t x_i}\le\frac{\sinh(nct)}{\sinh(ct)}.
\]
The ratio has continuous value $n$ at $t=0$.
\end{proposition}
\begin{proof}
Apply Corollary~\ref{cor:subset-centroid-new} to the subset giving the
largest $m$ projections; the projection of its sum is bounded by the
modulus of that sum. The arithmetic progression $(b_i)$ has the stated
partial sums, proving majorization. The convex-function inequality is
the elementary majorization inequality. For the square it is strict
because strict proper partial sums exclude a permutation of $(b_i)$.
Now
\[
 \sum_i b_i^2=\frac{2(n+1)}3,\qquad
 \max_\theta\sum_i\Re(e^{-i\theta}z_i)^2
       =\frac12\left(n+\left|\sum_i z_i^2\right|\right),
\]
which proves \eqref{eq:quadrupole-new}. The exponential expression is
just the finite geometric sum over the arithmetic progression.
\end{proof}

\begin{remark}[Population cut]
If a disk of radius $r$ contains $m$ points, then its internal moment
is at most $mr^2$. For a proper subset with $m\ge2$, the new theorem
therefore gives $m-1<(n-1)r^2$. This is a useful finite population cut,
but its coefficient is still twice the limiting bulk count in the
uniform disk of radius $\sqrt2$. No density classification or exact
all-$n$ formula follows from this cut alone.
\end{remark}

\subsection{Rational trace bounds using coefficient Euclidean norms}
Write the force-one polynomial as
$P(t)=\prod_i(t-w_i)=t^n+c_1t^{n-1}+\cdots+c_n$, and let
$V_{ji}=w_i^j$ for $0\le j<n$. On physical roots the Hermitian trace
matrix is $H=VV^*$. Put $d=n-1$, $\rho_n=(n-1)^2/n$, and define
\[
 T_{n,k}=\binom dk^2\rho_n^k\quad(0\le k\le d),\qquad
 A_n^{\mathrm{coef}}=\sum_{k=0}^dT_{n,k},\qquad T_n=\max_{0\le k\le d}T_{n,k}.
\]
All these constants are positive rationals.

\begin{proposition}[Improved compact coefficient constants]\label{prop:trace-new}
Every global maximizer at force-one scale satisfies
\begin{equation}\label{eq:trace-new}
 H\succeq\mu_n I,\qquad
 \mu_n=\frac{2^d}{nA_n^{\mathrm{coef}}},\qquad
 |c_k|^2\le\binom nk^2\rho_n^k.
\end{equation}
Let $R(t)=\sum_{\ell=0}^{n-1}r_\ell t^\ell$ be the unique polynomial
of degree below $n$ satisfying $R(w_i)=\overline w_i$. Then
\begin{equation}\label{eq:conjugation-coeff-new}
 |r_\ell|^2\le B_n^{\mathrm{interp}}:=\frac{NnT_n}{2^d}.
\end{equation}
At $n=6$ the constants are
\[
 \rho_6=\frac{25}6,\qquad A_6=\frac{138927151}{7776},\qquad
 \mu_6=\frac{41472}{138927151}>\frac1{3350},\qquad
 B_6=\frac{48828125}{2304}<21193.
\]
\end{proposition}
\begin{proof}
The radius bound is $|w_i|^2\le\rho_n$, and the pair bound is
$|w_i-w_j|^2\ge2$. Hence
\[
 |P'(w_i)|^2=\prod_{j\ne i}|w_i-w_j|^2\ge2^d.
\]
For the Lagrange polynomial
\[
 L_i(t)=\frac{\prod_{j\ne i}(t-w_j)}{P'(w_i)},
\]
the coefficient of $t^{d-k}$ has squared modulus at most
$T_{n,k}/2^d$. Thus its squared coefficient Euclidean norm is at most
$A_n^{\mathrm{coef}}/2^d$. These coefficient vectors form the inverse of $V^T$;
therefore
\[
 \|V^{-1}\|_F^2\le\frac{nA_n^{\mathrm{coef}}}{2^d},\qquad
 \lambda_{\min}(H)\ge\|V^{-1}\|_F^{-2}\ge\mu_n.
\]
The elementary symmetric functions give the bound on $c_k$.
Finally, interpolation gives $R=\sum_i\overline w_i L_i$.
For each coefficient, Cauchy--Schwarz and $\sum_i|w_i|^2=N$ give
\[
 |r_\ell|^2\le
 \left(\sum_i|w_i|^2\right)
 \left(\sum_i|[t^\ell]L_i|^2\right)
 \le\frac{NnT_n}{2^d}.
\]
The six-point values follow by rational arithmetic.
\end{proof}

The disk localizing matrix can now use
\[
 J_{jk}=\rho_n H_{jk}-H_{j+1,k+1}
       =\bigl[V\operatorname{diag}(\rho_n-|w_i|^2)V^*\bigr]_{jk}.
\]
For the separation-product cut, the quotient-algebra polynomial
$D(t)=P'(t)\overline{P'}(R(t))$ satisfies
\[
 K_{jk}=\operatorname{Tr}\bigl(t^jR(t)^k[D(t)-2^d]\bigr)
       =\bigl[V\operatorname{diag}(|P'(w_i)|^2-2^d)V^*\bigr]_{jk}.
\]
Hence both $J\succeq0$ and $K\succeq0$ are closed rational necessary
conditions for every global maximizer. As with the existing model,
$K$ bounds products of separations and does not independently impose
every pair inequality equivalently. These strengthened cuts preserve
the exact maximum when added to a model whose physical points are
actual moment-$n$ configurations; they do not turn necessary conditions
into a classification theorem.

% Self-contained section fragment; needs amsmath, amssymb, amsthm
% and theorem/lemma environments. No original-main theorem is invoked.
\section{A sharp eight-point bound with a fixed origin}
\label{sec:noclass-origin8}

This section proves the entire fixed-origin eight-point problem, with
no angular, reflection, or radial-population classification. One
two-variable inequality is verified by a finite certificate using exact
rational arithmetic. All other steps are analytic. The result does not
assert that an unrestricted eight-point maximizer contains the origin.

For seven points define
\[
 W(u)=\prod_{i=1}^7|u_i|^2\prod_{i<j}|u_i-u_j|^2,\qquad
 S(t)=\sum_{k=1}^7t^k,\quad L(t)=\sum_{k=0}^6t^k,
\]
\[
 \ell(t)=\log\frac{S(t)}7-4(t-1).
\]

\begin{lemma}[A strict boundary derivative]
\label{lem:noclass-boundary}
Let $H$ be a nonconstant polynomial, with
$|H(1)|=\max_{|\tau|=1}|H(\tau)|$. Then
\[
 \Re\frac{H'(1)}{H(1)}>0.
\]
\end{lemma}
\begin{proof}
The maximum is positive. Set $g=H/H(1)$ and
$h=1-\Re g$. The maximum-modulus principle gives $|g|\le1$
in the closed unit disk, so $h$ is nonnegative on the boundary,
harmonic inside, and $h(0)>0$ because $g$ is nonconstant.
The Poisson formula and the lower bound for its kernel give
\[
 h(r)\ge\frac{1-r}{1+r}h(0)\qquad(0<r<1).
\]
Since $g(1)=1$, divide by $1-r$ and let $r\uparrow1$:
\[
 \Re g'(1)\ge\frac12h(0)>0.
\]
\end{proof}

\begin{lemma}[Radial cuts for the rooted maximum]
\label{lem:noclass-rooted-radii}
At a global maximizer of $W$ under $\sum_i|u_i|^2=7$,
every proper subset $A$ of $k$ free points satisfies
\[
 \sum_{i\in A}|u_i|^2>\frac{k(k+1)}8,\qquad 1\le k\le6.
\]
In particular,
\[
 \frac14<|u_i|^2<\frac74\qquad(1\le i\le7).
\]
No force equation at the fixed origin is imposed.
\end{lemma}
\begin{proof}
A unit regular heptagon has $W=7^7>0$. Hence a global
maximizer exists, has distinct nonzero free points, and has positive
objective. Homogeneity implies that it also globally maximizes
\[
 \log W(u)-4\sum_i|u_i|^2
\]
over all free-point configurations: radial scaling of a configuration
has its Gaussian maximum at total moment seven.

Scale just the points of $A$ by $\tau$. The holomorphic factor
whose squared modulus is $W$ takes the form
\[
 C\,\tau^d H_A(\tau),\qquad
 d=\frac{k(k+1)}2,\qquad
 H_A(\tau)=\prod_{\substack{i\in A\\j\notin A}}(\tau u_i-u_j),
\]
with $C\ne0$. The polynomial $H_A$ has degree $k(7-k)>0$
and $H_A(1)\ne0$. Unit-modulus scaling preserves the total
moment, so $|H_A|$ has its circle maximum at one.
Radial stationarity at $\tau=r=1$ gives
\[
 0=2d+2\Re\frac{H_A'(1)}{H_A(1)}
       -8\sum_{i\in A}|u_i|^2.
\]
Lemma~\ref{lem:noclass-boundary} proves the strict subset cut.
A singleton gives the lower bound $1/4$; its six-point
complement gives the upper bound $7-42/8=7/4$.
\end{proof}

\begin{lemma}[The seven-row spectral deficit]
\label{lem:noclass-seven-gram}
Let $K$ be a positive definite $7\times7$ Hermitian matrix
with diagonal entries one. If
\[
 X=\sum_{i<j}|K_{ij}|^2,
\]
then
\[
 \log\det K\le-\frac12X,
\]
with equality only for $K=I$.
\end{lemma}
\begin{proof}
The positive eigenvalues $\lambda_i$ sum to seven and
$\sum_i(\lambda_i-1)^2=2X$. Maximize
\[
 \Phi(\lambda)=\sum_{i=1}^7
 \left(\log\lambda_i+\frac14(\lambda_i-1)^2\right)
\]
on the positive trace-seven simplex. Its value tends to minus
infinity at the boundary, so its maximum is attained inside.
The scalar critical equation has at most two roots, whose product
is two. If both roots occur, the smaller is below one and the
larger is above two. The larger root cannot have multiplicity
two or more at a maximum: varying two equal larger entries by
$(+\varepsilon,-\varepsilon)$ gives a strictly positive second
derivative, since $-1/\lambda^2+1/2>0$ there.

There is therefore at most one larger entry. Writing the other
six entries as $a$, the equations are $6a+b=7$ and $ab=2$.
They give $(b,a)=(4,1/2)$ or $(3,2/3)$. Their values are
\[
 \frac{21}8-4\log2<-\frac1{24},\qquad
 \frac76-\log\frac{243}{64}<-\frac1{30}.
\]
Indeed $\log x>2(x-1)/(x+1)$ for $x>1$,
so $\log2>2/3$ and
$\log(243/64)>\log(27/8)=3\log(3/2)>6/5$.
The remaining critical possibility is that all entries are one,
where $\Phi=0$. Thus $\Phi\le0$, with equality only there.
\end{proof}

\begin{lemma}[An exact two-variable radial compensation certificate]
\label{lem:noclass-pair8}
For $1/4\le t,s\le16/9$, put
\[
 f(t,s)=
 \frac{L(\sqrt{ts})^2(ts-1)^2}
 {(1+\sqrt{ts})^4L(t)L(s)}.
\]
Then
\begin{equation}\label{eq:noclass-pair8}
 f(t,s)\ge
 \frac{\ell(t)+\ell(s)}3+
 \frac{(t-1)^2+(s-1)^2}{96}
 +\frac{(t-1)(s-1)}8.
\end{equation}
Equality holds only at $t=s=1$.
\end{lemma}
\begin{proof}[Exact-certificate proof]
Write $t=x^2,s=y^2$ and subtract the right-hand side from
$f(x^2,y^2)$; call the difference $\Psi(x,y)$.
The square to be certified is
\[
 (x,y)\in[1/2,4/3]^2.
\]
Its rational part is evaluated using
\[
 f(x^2,y^2)=
 \frac{(xy-1)^2L(xy)^2}
 {(xy+1)^2L(x^2)L(y^2)}.
\]
Exact differentiation gives $\Psi(1,1)=0$,
$\nabla\Psi(1,1)=0$, and
$\nabla^2\Psi(1,1)=(5/12)I$.

On the central square $J=[1-1/4000,1+1/4000]^2$,
rational interval evaluation certifies
\[
 \Psi_{xx}>\frac14,\qquad
 \Psi_{yy}>\frac14,\qquad
 |\Psi_{xy}|<\frac1{80}.
\]
Thus the Hessian has eigenvalues greater than $19/80$
throughout $J$, and $\Psi>0$ there except at $(1,1)$.

Divide each coordinate interval at the endpoints of $J$.
The other eight rectangles are bisected along their longer
side until their Taylor lower bound is strictly positive.
For a box with rational center $c$, half-widths $a,b$, center
value interval $V$, center gradient intervals $G_1,G_2$,
and Hessian intervals $H_{11},H_{12},H_{22}$ over the box,
the accepted lower bound is exactly
\[
 \inf V-a\sup|G_1|-b\sup|G_2|
 -\frac12\left(
 a^2\sup|H_{11}|+2ab\sup|H_{12}|+b^2\sup|H_{22}|
 \right)>0.
\]
Every arithmetic operation and this final comparison use
rational numbers.

For the logarithms, first write $u=2^k v$ with
$3/4\le v\le3/2$, and use $\log u=k\log2+\log v$.
Each logarithm is enclosed by the 25-term expansion
\[
 \log v=2\sum_{j=0}^{24}\frac{z^{2j+1}}{2j+1}+R,
 \qquad z=\frac{v-1}{v+1}.
\]
Both this $z$ and the value $z=1/3$ used for $\log2$
have modulus at most $17/50$. The exact tail check is
\[
 |R|\le
 \frac{2(17/50)^{51}}{51(1-(17/50)^2)}<2^{-70}.
\]

The complete rational checker is
\path{noclass_origin8.py}. It produces 1102 accepted leaf
rectangles, with maximum bisection depth 22 and no pending
rectangles. The leaf file
\path{noclass_origin8_leaves.json} records each sector,
binary splitting path, exact rational rectangle, and depth.
Its independent coverage check verifies full binary coverage
in every sector and the exact area identity
\[
 \sum_{\text{leaves}}\operatorname{area}(B)
 +(2/4000)^2=(4/3-1/2)^2=\frac{25}{36}.
\]
A second mode independently recomputes the positive Taylor
bound for every saved leaf.

The original exact run completed in batches of 862 and 240
leaves. An independent exact replay completed in batches of
342, 366, 363, and 31 leaves, with the same 1102 leaves and
maximum depth. Its additional checks include the center
value, gradient and Hessian, and the polynomial derivatives.
The saved results and reproduction instructions are
\path{noclass_origin8_results.json} and
\path{noclass_origin8_README.md}. Displayed decimal
margins do not enter any acceptance decision.
\end{proof}

\begin{theorem}[The sharp weighted seven-point inequality]
\label{thm:noclass-weighted7}
For arbitrary $u_1,\ldots,u_7\in\mathbb C$ with
$Q=\sum_i|u_i|^2>0$,
\[
 \boxed{\displaystyle
 \prod_i|u_i|^2\prod_{i<j}|u_i-u_j|^2
 \le\frac{Q^{28}}{7^{21}}.}
\]
Equality holds precisely for a regular heptagon of radius
$\sqrt{Q/7}$, up to rotation and permutation.
\end{theorem}
\begin{proof}
By homogeneity it is enough to consider $Q=7$.
It is enough to bound a global maximum, and
Lemma~\ref{lem:noclass-rooted-radii} places all its squared
radii in $(1/4,7/4)$. It has distinct nonzero points.

For $V_{ik}=u_i^k$, $1\le i,k\le7$, one has
$W=|\det V|^2$. Set $t_i=|u_i|^2$, $H=\prod_iS(t_i)$
and normalize the row Gram matrix to
\[
 K_{ij}=
 \frac{S(u_i\overline{u_j})}{\sqrt{S(t_i)S(t_j)}}.
\]
Then $K$ is positive definite, has diagonal one, and
$W=H\det K$. Put $R=\sum_i\ell(t_i)=\log H-7\log7$.

The geometric-sum identity and reverse triangle inequality,
with $q=u_i\overline{u_j}$ and $r=|q|=\sqrt{t_it_j}$, give
\[
 |S(q)|\ge\frac{r|r^7-1|}{r+1}.
\]
The bound is also valid at $q=1$, where its required
right-hand side is zero. Dividing its square by
$S(t_i)S(t_j)$ gives
\[
 |K_{ij}|^2\ge f(t_i,t_j).
\]
Apply Lemma~\ref{lem:noclass-pair8} and sum over the 21
unordered pairs. Each $\ell(t_i)$ occurs six times.
For $d_i=t_i-1$, $\sum_i d_i=0$, so the two quadratic
contributions cancel:
\[
 \frac6{96}\sum_i d_i^2+
 \frac18\sum_{i<j}d_id_j=0.
\]
Consequently, $X=\sum_{i<j}|K_{ij}|^2\ge2R$.
Lemma~\ref{lem:noclass-seven-gram} now gives
\[
 \log W=7\log7+R+\log\det K
 \le7\log7+R-\frac12X\le7\log7.
\]

If equality holds, the strict equality statement in the
spectral lemma forces $K=I$. For distinct nonzero points,
$S(u_i\overline{u_j})=0$ implies
$(u_i\overline{u_j})^7=1$ and
$u_i\overline{u_j}\ne1$. Thus $t_it_j=1$ for every pair;
since there are at least three points, all $t_i=1$.
The seven points are therefore a rotated set of seventh
roots of unity. Conversely that configuration has $W=7^7$.
\end{proof}

\begin{theorem}[Every eight-point configuration containing the origin]
\label{thm:noclass-origin8}
If $z_1=0$ and $\sum_{i=1}^8|z_i|^2=8$, then
\[
 \boxed{\displaystyle
 \Delta_8(z)\le C_8=\frac{8^{28}}{7^{21}}.}
\]
Equality holds precisely for the origin together with a
regular heptagon of radius $\sqrt{8/7}$, up to rotation
and permutation. No centering, origin force equation,
angular symmetry, or restriction on radial levels is required.
\end{theorem}
\begin{proof}
Apply Theorem~\ref{thm:noclass-weighted7} to the seven
nonzero entries, whose total moment is eight. Cases with
an additional zero or a collision have zero objective.
\end{proof}

For the unrestricted eight-point problem, the remaining step
is to exclude a global maximizer with no point at the origin.
Translating an existing point to zero changes the prescribed
moment, so the fixed-origin theorem does not supply that step.

\begin{lemma}[An explicit spectral deficit in every dimension]
\label{lem:noclass-general-gram}
Let $m\ge2$ and let $K$ be positive definite Hermitian with
diagonal entries one. Then
\[
 \boxed{\displaystyle
 \log\det K\le-\frac{4(m-1)}{m^2}
       \sum_{i<j}|K_{ij}|^2.}
\]
Equality holds only for $K=I$. For $m=7$, the preceding
specialized constant $1/2$ is stronger than $24/49$.
\end{lemma}
\begin{proof}
For $m=2$, $\det K=1-X$ and $\log(1-X)\le-X$.
For $m>2$, put $c=2(m-1)/m^2$ and maximize
\[
 \Phi_m=\sum_i\left(\log\lambda_i+c(\lambda_i-1)^2\right),
 \qquad \lambda_i>0,\quad \sum_i\lambda_i=m.
\]
The boundary values tend to minus infinity. The scalar
critical equation has at most two roots, with product
$q=1/(2c)=m^2/[4(m-1)]$. If both occur, the larger root
has square greater than $q$; its multiplicity cannot exceed
one because $-1/\lambda^2+2c>0$ there.
With one larger entry $b$ and $m-1$ smaller entries $a$,
the equations $ab=q$ and $(m-1)a+b=m$ give
\[
 (b-m/2)^2=0,\qquad
 b=m/2,\qquad a=\frac{m}{2(m-1)}.
\]
Set $y=(m-2)/m$. At this critical spectrum,
\[
 (1-y)\Phi_m=
 y^2-\left[(1-y)\log(1-y)+(1+y)\log(1+y)\right]<0.
\]
Indeed the bracket minus $y^2$ has value and derivative
zero at zero, and second derivative
$2y^2/(1-y^2)>0$ for $0<y<1$.
The all-one spectrum has value zero and is the unique maximum.
Finally $\sum_i(\lambda_i-1)^2=2\sum_{i<j}|K_{ij}|^2$.
\end{proof}


% Fragment for the existing manuscript; not a standalone document.
\section{Higher angular constraints without classification}
\label{sec:noclass-n8-higher}
The seventh holomorphic moment is retained by a higher second-order
Gram matrix. This supplies necessary inequalities for every actual
eight-point local maximum, without a stationary-shape catalogue.

Let $z$ be a collision-free local maximum with $\sum_i|z_i|^2=8$.
First-order optimality gives $\sum_i z_i=0$. Put
\[
 t_i=|z_i|^2,\quad T_j=\sum_i t_i^j,\quad
 p_k=\sum_i z_i^k,\quad N_{r,k}=\sum_i t_i^r z_i^k,\quad
 M_{j,l}=\sum_i z_i^j\overline{z_i}^{\,l},\quad s=T_2-8.
\]
The force multiplier is $\lambda=7/2$. The inverse-square Laplacian
$B$ satisfies $B\mathbf1=0$, $Bz=\lambda\overline z$, and
$\|B\|_{\rm op}\le\lambda$ by second-order optimality. The pairwise
divided-difference identity gives, for $k\ge2$,
\begin{equation}\label{eq:noclass-n8-B-recurrence}
 Bz^{\circ k}=
 \left\{k\lambda t-(k-1)(8-k/2)\right\}z^{\circ(k-2)}
 -\sum_{r=1}^{k-2}(k-1-r)p_rz^{\circ(k-2-r)}.
\end{equation}
The paired force equations also give
\begin{equation}\label{eq:noclass-n8-mixed}
 N_{1,k}=\frac{(15-k)p_k+\sum_{r=1}^{k-1}p_rp_{k-r}}7
 \qquad(k\ge1),
\end{equation}
where $p_0=8$ and $p_1=0$. In particular $N_{1,1}=0$.

\begin{lemma}[The complete $\{2,9\}$ Gram matrix]
\label{lem:noclass-n8-29}
Define
\[
 u_2=z^{\circ2}-\frac{p_2}{8}\mathbf1,\qquad
 u_9=z^{\circ9}-\frac{p_9}{8}\mathbf1,\qquad
 L=\sum_{j=1}^5c_jz^{\circ j},\quad c_j=(j+1)p_{7-j},
\]
\[
 Q_j=\frac{63}{2}N_{j+1,7-j}-28N_{j,7-j}.
\]
The Hermitian matrix
$D_{jk}=\lambda^2\langle u_j,u_k\rangle-
\langle Bu_j,Bu_k\rangle$, $j,k\in\{2,9\}$, is positive
semidefinite, with
\begin{align}
 D_{22}={}&\frac{49}{32}\{64-24s-|p_2|^2\},
 \label{eq:noclass-n8-D22}\\
 D_{2,9}={}&-\frac{833}{4}N_{2,7}+280p_7+155p_2p_5
 +159p_3p_4+17p_2^2p_3-\frac{49}{32}\overline{p_2}p_9,
 \label{eq:noclass-n8-D29}\\
 D_{99}={}&8|p_7|^2-980T_9+1764T_8-784T_7
 -\frac{49}{32}|p_9|^2
 +2\operatorname{Re}\sum_{j=1}^5\overline{c_j}Q_j-\|L\|^2.
 \label{eq:noclass-n8-D99}
\end{align}
Here $\|L\|^2=\sum_{j,l=1}^5c_j\overline{c_l}M_{j,l}$.
\end{lemma}
\begin{proof}
Since $B\mathbf1=0$, subtraction of the means affects only the
input Gram matrix. Its entries are
\[
 \|u_2\|^2=T_2-|p_2|^2/8,\quad
 \langle u_2,u_9\rangle=N_{2,7}-\overline{p_2}p_9/8,
 \quad\|u_9\|^2=T_9-|p_9|^2/8.
\]
Equation~\eqref{eq:noclass-n8-B-recurrence} gives
\[
 Bz^{\circ2}=7(t-\mathbf1),\qquad
 Bz^{\circ9}=\left(\frac{63}{2}t-28\right)z^{\circ7}
              -p_7\mathbf1-L.
\]
Its first output norm is $49s$. For the off-diagonal entry,
\eqref{eq:noclass-n8-mixed} gives exactly
\[
 \langle Bz^2,Bz^9\rangle
 =\frac{441}{2}N_{2,7}-280p_7-155p_2p_5-159p_3p_4-17p_2^2p_3.
\]
Indeed, before eliminating $N_{1,7}$ this is
\[
 \frac{441}{2}N_{2,7}-\frac{833}{2}N_{1,7}+196p_7
 -36p_2p_5-40p_3p_4-17p_2^2p_3.
\]
Use $N_{1,7}=(8p_7+2p_2p_5+2p_3p_4)/7$ to obtain the
displayed formula. For the second output norm, set $A(t)=63t/2-28$.
Then
\[
 \|Az^7\|^2=\frac{3969}{4}T_9-1764T_8+784T_7.
\]
The $p_7$--$L$ cross terms cancel because
\[
 \sum_iL(z_i)=9(p_2p_5+p_3p_4),\qquad
 \sum_iA(t_i)z_i^7=8p_7+9(p_2p_5+p_3p_4).
\]
Expanding $\|Az^7-p_7\mathbf1-L\|^2$ proves
\eqref{eq:noclass-n8-D99}. Finally $\|B\|\le\lambda$ proves
$D\succeq0$ for every linear combination of $u_2,u_9$.
\end{proof}

\begin{lemma}[A quantitative seventh-moment bound]
\label{lem:noclass-n8-29-error}
Put $a=8/7$, $P=|p_7|^2$, and write
\[
 D_{22}=56+e_A,\quad D_{2,9}=8p_7+e_B,
 \quad D_{99}=8P-48T_7+e_C,
\]
where
\begin{align*}
 e_A={}&-\frac{147}{4}(s-8/7)-\frac{49}{32}|p_2|^2,\\
 e_B={}&-\frac{833}{4}(N_{2,7}-a^2p_7)+155p_2p_5+159p_3p_4
 +17p_2^2p_3-\frac{49}{32}\overline{p_2}p_9,\\
 e_C={}&-980(T_9-a^2T_7)+1764(T_8-aT_7)-\frac{49}{32}|p_9|^2
 +2\operatorname{Re}\sum_j\overline{c_j}Q_j-\|L\|^2.
\end{align*}
If $|e_A|\le A$, $|e_B|\le B_0$, and $|e_C|\le C$, then
\begin{equation}\label{eq:noclass-n8-p7-error}
 P\ge7T_7-\mathcal E,\qquad
 \mathcal E=\frac{48T_7A+(56+A)C+16\sqrt{8T_7}B_0}{384}.
\end{equation}
Moreover
\begin{equation}\label{eq:noclass-n8-Fischer}
 \Delta_8\le(8T_7-P)\prod_{k=1}^6T_k
 \le(T_7+\mathcal E)\prod_{k=1}^6T_k.
\end{equation}
\end{lemma}
\begin{proof}
Expansion gives
\[
 \det D=384(P-7T_7)+e_A(8P-48T_7)+(56+e_A)e_C
 -16\operatorname{Re}(\overline{p_7}e_B)-|e_B|^2.
\]
Use $0\le P\le8T_7$, hence $|8P-48T_7|\le48T_7$, and drop
the final nonpositive term. Positivity of the determinant proves
\eqref{eq:noclass-n8-p7-error}. Apply Hadamard--Fischer to the
monomial Gram matrix, keeping the $\{0,7\}$ block together and
all other diagonal entries separately, to prove
\eqref{eq:noclass-n8-Fischer}.
\end{proof}
All errors vanish at the central heptagon. Small nonzero errors only
give an enclosure: a global proof still needs a radial product deficit
that compensates $\mathcal E$. No such uniform deficit is asserted here.

\subsection{A common high-order Gram saturation for central polygons}
\begin{proposition}\label{prop:noclass-central-profile}
Let $n>4$, $m=n-1$, $a=n/m$, and $\lambda=m/2$. Let $\nu$ be a
positive measure of total mass $n$, consisting of mass one at zero
and mass $m$ on $|z|^2=a$. Assume $p_1=\cdots=p_{m-1}=0$, where
$p_j=\int z^j\,d\nu$. Put
\[
 R_2=m(|z|^2-1),\qquad
 R_{n+1}=\{(n+1)\lambda|z|^2-n\lambda\}z^m-p_m.
\]
If the unprojected matrix
\[
 \lambda^2\operatorname{Gram}_\nu(z^2,z^{n+1})
       -\operatorname{Gram}_\nu(R_2,R_{n+1})
\]
is positive semidefinite, then $\nu$ is precisely the mass-one
origin and a rotated set of $m$th roots on $|z|^2=a$, each with mass one.
These are necessary Gram conditions at an actual local maximum
having the stated radial and moment profile.
\end{proposition}
\begin{proof}
Write $M_m=ma^m$ and $P=|p_m|^2$. On the circle
$R_{n+1}=nz^m-p_m$, and at zero it equals $-p_m$.
The input norms and cross term are $n^2/m$, $a^2M_m$, and $a^2p_m$.
The output norms and cross term are $mn$, $n^2M_m-nP$, and $np_m$.
Thus the defect matrix is
\[
 \begin{pmatrix}
 mn(n-4)/4&n(n-4)p_m/4\\
 n(n-4)\overline{p_m}/4&nP-3n^2M_m/4
 \end{pmatrix},
\]
with determinant $3n^3(n-4)(P-mM_m)/16$.
Positivity and Cauchy on the circle force $P=mM_m$. Equality in
Cauchy means $z^m$ is constant almost everywhere on the circle.
The zero Fourier moments $p_1,\ldots,p_{m-1}$ and total circle
mass $m$ then give mass one at each of the $m$ roots.
\end{proof}
The relevant pairs are $\{2,7\}$, $\{2,8\}$, and $\{2,9\}$ for
six, seven, and eight points. The proposition assumes the candidate
radial profile; it does not prove that all local maxima have it.

\subsection{An effective small-harmonic exclusion for actual eight-point maxima}
\begin{theorem}\label{thm:noclass-n8-small-harmonic}
Let $z$ be a collision-free eight-point local maximum with
$\sum_i|z_i|^2=8$. If
\[
 \epsilon:=\max_{2\le k\le6}|p_k|\le10^{-38},
\]
then $z$ is a rotated and permuted central regular heptagon.
\end{theorem}
\begin{proof}
The configuration is centered. Newton's identities put its monic
root polynomial in the form
\[
 \Pi(X)=X^8-\alpha X-\beta+E(X),\quad
 E(X)=\sum_{j=2}^6a_jX^j,\quad |a_j|\le\epsilon/2.
\]
For $\epsilon\le1$, the bounds for the corresponding elementary
symmetric coefficients of orders $2,\ldots,6$ are respectively
$\epsilon/2$, $\epsilon/3$, $3\epsilon/8$, $11\epsilon/30$,
and $53\epsilon/144$. Since $|z_i|\le\sqrt8<3$,
\[
 |E(z_i)|\le545\epsilon,\quad |E'(z_i)|\le1002\epsilon,
 \quad |E''(z_i)|\le1549\epsilon.
\]
The physical force equations read
$\Pi''(z_i)=7\overline{z_i}\Pi'(z_i)$. Multiplying their tail
part by $z_i^2/7$ and using $\Pi(z_i)=0$ gives
\begin{equation}\label{eq:noclass-tail-real-line}
 |\alpha(7t_i-8)z_i+8\beta(t_i-1)|\le60000\epsilon.
\end{equation}
In detail, the perturbation in the force equation has modulus at
most $(1549+21\cdot1002)\epsilon=22591\epsilon$; the resulting
linear-relation error is at most
$\{8\cdot22591/7+56\cdot545\}\epsilon<60000\epsilon$.

If $|\alpha|\le300000\epsilon$, use the model $\Pi_0=X^8-\beta$.
Otherwise $\alpha\ne0$. If $\beta=0$, use $\Pi_0=X^8-\alpha X$.
If $\beta\ne0$, rotate the coordinates along $\beta/\alpha$.
Because $|p_2|\le\epsilon$,
\[
 \sum_i(\operatorname{Im}z_i)^2\ge(8-\epsilon)/2>2.
\]
Some point has $|\operatorname{Im}z_i|>1/2$. The imaginary part
of \eqref{eq:noclass-tail-real-line}, divided by $\alpha$, gives
\[
 |7t_i-8|\le120000\epsilon/|\alpha|<2/5,
 \qquad |t_i-1|>1/14.
\]
Taking moduli in that equation and using $|z_i|<3$ gives
$8|\beta||t_i-1|\le420000\epsilon$, hence
$|\beta|\le735000\epsilon$. Again use $\Pi_0=X^8-\alpha X$.
In either case every actual root satisfies
\[
 |\Pi_0(z_i)|\le10^6\epsilon,
 \qquad \operatorname{dist}(z_i,\{\Pi_0=0\})
       \le\eta:=(10^6\epsilon)^{1/8}\le10^{-4}.
\]
The distance assertion follows by factoring the monic degree-eight
polynomial and bounding its smallest factor by the geometric mean.

There is no multiplicity loophole in matching the roots. With
$P_0=I-\mathbf1\mathbf1^*/8$, one has $B=P_0BP_0$; consequently
$|B_{ij}|\le(7/2)(7/8)=49/16$, and
\[
 |z_i-z_j|\ge4/7\qquad(i\ne j).
\]
For the regular model $\beta$ cannot be zero: otherwise all eight
actual roots have modulus at most $\eta$, contradicting their
second moment. For the central model $\alpha\ne0$ by its branch
definition. Both models therefore have eight distinct roots. Two
actual points cannot match the same model root since
$2\eta<4/7$. Thus the matching is a bijection.

Let $u$ be the matched model root vector. Then
$\|z-u\|\le\sqrt8\eta$, and
$|\|u\|-\sqrt8|\le\sqrt8\eta$. After normalizing $u$ to moment
eight, the resulting regular octagon or central heptagon $w$ obeys
\[
 \|z-w\|\le2\sqrt8\eta<6\eta\le6\cdot10^{-4}.
\]
For the regular octagon, $B(w)w^{\circ5}=-(9/2)w^{\circ3}$, so
$\|B(w)\|\ge9/2$. Its minimum distance is greater than $3/4$.
On the stated neighborhood all actual pair distances exceed $7/10$,
and
\[
 \left|\frac1{(z_i-z_j)^2}-\frac1{(w_i-w_j)^2}\right|
 <21|(z_i-z_j)-(w_i-w_j)|.
\]
The diagonal and off-diagonal row sums, and symmetry of absolute
entries, give
\[
 \|B(z)-B(w)\|\le14\cdot21\sqrt2\|z-w\|
 <441\|z-w\|\le1323/5000<1.
\]
This contradicts $\|B(z)\|\le7/2$.

It remains to prove stationary uniqueness near the central heptagon,
rather than only compare its value. Choose the closest phase and
permutation of $z^*$ and write
\[
 z=\frac{\sqrt8(z^*+h)}{\sqrt{8+\rho^2}},\quad
 \sum_i h_i=\langle z^*,h\rangle=0,\quad \rho=\|h\|.
\]
If $d=\|z-z^*\|\le6\cdot10^{-4}$, then
$a_0=1-d^2/16\ge399999991/400000000$ and
\[
 \rho^2=8(a_0^{-2}-1)
 \le57599999352/159999992800000081<10^{-6}.
\]
The central-polygon Fourier computation gives slice singular values
$7/\sqrt8,7/\sqrt8,7/8,7/4,7/4,7/8$. Thus its quadratic logarithmic
term $Q_2(h)$ is at most $-\rho^2$. The minimum squared distance
satisfies $d_*^2>4/5$. Indeed, $x=\cos(2\pi/7)$ is the unique
positive root of $8x^3+4x^2-4x-1$, as follows from summing the
seventh roots of unity and Descartes' rule. The polynomial is
negative at $1/2$, equal to $287/1000>0$ at $13/20$, and strictly increasing
on $[1/2,1]$, since its derivative there equals
$6+32(x-1/2)+24(x-1/2)^2$. Hence $x<13/20$ and the shortest
chord has squared length $2(8/7)(1-x)>4/5$; the central distances
are still larger. Consequently, for
$v_{ij}=(h_i-h_j)/(z_i^*-z_j^*)$,
\[
 \max|v_{ij}|\le b=8\rho/5,\qquad
 \sum_{i<j}|v_{ij}|^2\le10\rho^2.
\]
The derivative in the radial chart direction is exactly
\[
 DF(h)[h]=2Q_2(h)+2\operatorname{Re}\sum_{i<j}
       \frac{v_{ij}^3}{1+v_{ij}}+\frac{7\rho^4}{8+\rho^2}.
\]
For $\rho\le10^{-3}$ this is at most
\[
 \left\{-2+\frac{20b}{1-b}+\frac78\rho^2\right\}\rho^2
 \le\left\{-2+\frac5{156}+\frac7{8000000}\right\}\rho^2
 <-\rho^2.
\]
At a stationary point this derivative must vanish. Hence $h=0$.
This completes the proof without a stationary classification.
\end{proof}

\begin{corollary}\label{cor:noclass-n8-frame-gap}
For the frame quantity $\varepsilon_8$ defined by
\begin{align*}
 \varepsilon_8^2={}&|p_2|^2\left\{\frac1{2T_2}+\frac1{32}
 +\frac{169}{98T_3}+\frac{144}{49sT_2}\right\}\\
 &+|p_3|^2\left\{\frac1{2T_3}+\frac1{2T_2}+\frac{100}{49sT_3}\right\}
 +|p_4|^2\left\{\frac2{T_2^2}+\frac1{2T_3}\right\}\\
 &+\frac{4(|N_{2,1}|^2+|p_5|^2)}{T_2T_3}
 +\frac{2|p_6|^2}{T_3^2},
\end{align*}
an actual eight-point local maximum with $\varepsilon_8\le10^{-41}$
is the central heptagon.
\end{corollary}
\begin{proof}
Since $T_2\le64$ and $T_3\le512$, positive terms in the definition give
$|p_2|<6\varepsilon_8$, $|p_3|<12\varepsilon_8$,
$|p_4|\le32\varepsilon_8$, $|p_5|<91\varepsilon_8$, and
$|p_6|<363\varepsilon_8$. Apply the theorem. The quantity is defined
for $s>0$. The case $s=0$ cannot be an eight-point local maximum:
stationarity on a common unit circle forces a regular octagon, which
fails the $Bz^5$ second-order test above.
\end{proof}
This is an effective exclusion of a small-harmonic region. It does
not prove that a global maximum has small $\varepsilon_8$; the region
$\varepsilon_8>10^{-41}$ remains unexcluded by this argument.

\subsection{Finite-point closure and a precise relaxation obstruction}
\begin{lemma}[Conjugate closure at finite point count]
\label{lem:noclass-conjugate-closure}
For distinct $n$ points, let $V_{ij}=z_i^j$, $0\le j<n$, and
$G=V^*V$, $H=V^{\mathsf T}V=(p_{j+k})$. Then
\[
 H G^{-1}\overline H=\overline G.
\]
Thus the joint Gram matrix
$\left(\begin{smallmatrix}G&\overline H\\H&\overline G\end{smallmatrix}\right)$
has zero Schur complement, not merely a positive semidefinite one.
\end{lemma}
\begin{proof}
The Vandermonde is invertible. Substitution gives
$HG^{-1}\overline H=V^{\mathsf T}V(V^*V)^{-1}V^*\overline V
=V^{\mathsf T}\overline V=\overline G$.
\end{proof}

Consider instead the positive measure of total mass eight
\[
 \nu=\delta_0+\frac12\sum_{j=0}^{13}\delta_{r e^{2\pi ij/14}},
 \qquad r^2=8/7=:a.
\]
It has second moment eight, $p_1=\cdots=p_{13}=0$, and $T_k=7a^k$.
Its monomial Gram matrix on $1,z,\ldots,z^7$ is diagonal with
determinant $8\cdot7^7a^{28}=8C_8$. All positive single-point
moment Gram tests hold, but the conjugate closure fails: $H$ on
these columns has rank two while $G$ has rank eight.

Even the formal recurrence outputs $R_k$ for $0\le k\le8$ do not
remove this measure. Both input and output columns are mutually
orthogonal. Their squared norm ratios, in order $k=0,\ldots,8$,
are
\[
 0,\quad49/4,\quad49/8,\quad49/64,\quad49/16,\quad
 49/16,\quad49/64,\quad49/64,\quad49/4.
\]
Thus the complete formal $\lambda^2G-R^*R$ matrix is positive
semidefinite through degree eight. The necessary degree-nine
output has norm ratio $7>\lambda$, and the $\{2,9\}$ test excludes
it. In the indexing of \eqref{eq:noclass-n8-mixed}, the identities
hold for $1\le k\le13$ and first fail at $k=14$:
$N_{1,14}=a p_{14}=8p_{14}/7$, while the displayed right side is
$p_{14}/7$. Here the force equation is multiplied by $z^{k+1}$;
thus the first failure uses $z^{15}$, not the mixed-moment index $15$.
Also $\operatorname{Gram}(1,z,\ldots,z^8)$ has rank nine, violating
the rank bound for eight actual points. This is a counterexample
only to the stated relaxed moment conditions; its outputs are not
the inverse-square Laplacian of an eight-point configuration.

\subsection{Positive distance moments also retain point count}
Let the formal mixed moments be $M_{jj}=(n-1)(n/(n-1))^j$ for $j>0$,
$M_{00}=n$, and $M_{jk}=0$ for $j\ne k$. Put $m=n-1$ and $a=n/m$.
Substitution into the squared-distance expansion gives the formal
functional
\[
 \mathcal L(1)=n(n-1)/2,\qquad
 \mathcal L(x^k)=a^k\left\{m+\frac{m^2}{2}\binom{2k}{k}\right\}
 \quad(k\ge1).
\]
It cannot be assumed positive merely because the single-point mixed
moment matrices are positive. Indeed, define
\begin{align*}
 Q_6(x)&=(x-6/5)(x^2-6x+36/5),\\
 Q_7(x)&=(x-7/6)(x-7/2)(x-14/3),\\
 Q_8(x)&=(x-8/7)(x^3-8x^2+128x/7-512/49).
\end{align*}
Exact arithmetic gives
\[
 \mathcal L_6(Q_6^2)=0,\quad
 \mathcal L_6(x^2Q_6)=7776/125\ne0,
\]
\[
 \mathcal L_7(Q_7^2)=-117649/1296<0,
\]
\[
 \mathcal L_8(Q_8^2)=0,\quad
 \mathcal L_8(x^3Q_8)=2097152/16807\ne0.
\]
For a positive distance functional, Cauchy forces every cross term
with a zero-norm polynomial to vanish. Therefore distance Hankel
positivity of orders three, three, and four excludes these formal
examples respectively. These are not positive distance measures.
In particular the fourteen-gon measure above has actual weighted
pair mass $119/4$, whereas the formal eight-point distance functional
fixes $\mathcal L(1)=28$. This self-counting constraint must not be
discarded when passing from eight points to a positive measure.

The accompanying checker verifies polynomial identities, moment
values, and all rational constants using only the Python standard
library. Its finite tests audit these formulas; the continuum
inequalities and matching argument are proved above. No general
exclusion of origin-free eight-point configurations, nor an
unrestricted proof of $M_8=C_8$, is claimed.

\section{Exact obstructions to dropping finite-point information}
\label{sec:noclass-relaxation}

The examples here are formal relaxations, not configurations that
contradict the central-polygon bound. They identify information that
must be retained in a nonclassification proof.

Put $m=n-1$, $a=n/m$, and consider the positive measure
\[
 \mu=\delta_0+m\,\sigma_{\sqrt a},
\]
where $\sigma_{\sqrt a}$ is uniform probability on the circle of
radius $\sqrt a$. Its total mass is $n$ and its second moment is $n$.
All positive analytic moments vanish, and
\[
 T_k=m a^k\quad(k\ge1),\qquad
 M_{rs}=0\quad(r\ne s).
\]
Consequently all ordinary point-moment matrices are positive
semidefinite, and the analytic mixed-force identities
\eqref{eq:noclass-mixed-force} hold formally. The radial bounds agree
with the central polygon. The formal second-order matrix obtained
from the centered monomials $z,\ldots,z^n$ is diagonal. Its diagonal
lists for $n=6,7,8$ are
\[
\begin{array}{c|l}
6&(0,15,54,324/5,1728/25,0)\\
7&(0,63/2,84,735/8,343/3,12005/96,0)\\
8&(0,56,120,768/7,6144/49,61440/343,491520/2401,0).
\end{array}
\]
For $k\ge3$ these entries follow from
\[
 Bz^k=\left[k\lambda t-(k-1)(n-k/2)\right]z^{k-2},
\]
with $\lambda=(n-1)/2$; the $k=2$ case uses
$Bz^2=(n-1)(t-1)$, including the origin atom.
Thus these finite second-order tests alone do not restore the
finite-point structure.

The exact projection identity does distinguish the relaxation.
For $2\le k\le n-3$, projection away from
$\mathbf1,z,\overline z,u$ leaves $z^k$ unchanged.
Hence $H=\operatorname{diag}(T_2,\ldots,T_{n-3})$ is positive
definite while $C=0$. The identity
$\det H=|\det C|$ fails because this measure lives in an
infinite-dimensional space, rather than at $n$ points.

\subsection{A separate continuous radial-count obstruction at six}
For the six-point frame, setting its harmonic errors to zero
imposes the radial recurrence
\[
 t^2=\alpha t+\beta,\quad
 \alpha=\frac{(6+s)(6-s)}{3(6+3s)},\quad
 \beta=\frac{(6+s)(5s-6)}{6(6+3s)}.
\]
Consider the further scalar relaxation that substitutes the lower
bound $|p_5|^2= T_2T_3$ into the degree-five objective, without
enforcing its vector Cauchy equality case or a common positive
joint-moment realization. This scalar core admits $s=1$. Then
\[
 t^2=\frac{35}{27}t-\frac7{54},\qquad
 T_2=7,\quad T_3=\frac{224}{27},\quad
 T_4=\frac{14357}{1458},\quad T_5=\frac{460159}{39366}.
\]
The degree-five Hermite interpolant of $\log x$ at
$6/5,3-3/\sqrt5,3+3/\sqrt5$, with double interpolation at each
node, has formal objective strictly above $\log C_6$ by more than
$1/150$ on this radial core. This sign is verified with outward
rational interval arithmetic, including the logarithm tails.
The corresponding distance moments are
\[
 (D_0,\ldots,D_5)
 =\left(15,36,114,\frac{3850}{9},
               \frac{422219}{243},\frac{47872510}{6561}\right).
\]
Nevertheless this recurrence cannot describe six actual radii:
its two roots are $(35\pm\sqrt{847})/54$. If $k$ of six radii
take the smaller root, their sum is six only if
$(6-2k)\sqrt{847}=114$, impossible because $\sqrt{847}$ is irrational.
The integer radial population is one explicit obstruction to
realizing the scalar core. The core is not asserted to satisfy every
joint-moment, projection, or distance-positivity constraint.
It is not a counterexample configuration.

All computations in this section are replayed by
\path{noclass_relaxation_check.py}, with no floating acceptance
criterion and no assertion removed by optimized Python.

\section{A subset dilation test that retains complex phase}
\label{sec:noclass-subset-dilation}
The internal-spread inequality in Theorem~\ref{thm:subset-spread-new}
is a necessary condition for a global maximum. The following additional
condition follows already from the full complex second variation. It
retains phase information that a scalar spread bound discards.

\begin{proposition}[Complex subset cut]\label{prop:noclass-complex-cut}
Let $z$ be a centered, collision-free local maximum of $\Delta_n$ on
$\sum_i|z_i|^2=Q>0$, and put $\kappa=n(n-1)/(2Q)$. For a proper subset
$S$ with $0<Q_S:=\sum_{i\in S}|z_i|^2<Q$, define
\[
 C_S=\sum_{i\in S,\,j\notin S}
       \frac{z_i z_j}{(z_i-z_j)^2},\qquad
 K_S=\kappa Q_S(1-Q_S/Q).
\]
Then
\[
 |K_S+C_S|\le K_S,
 \qquad |C_S|^2\le-2K_S\Re C_S.
\]
In particular $\Re C_S\le0$.
\end{proposition}
\begin{proof}
Write $B_{ij}=-(z_i-z_j)^{-2}$ for $i\ne j$ and choose the diagonal
to make every row sum zero. The force equation gives
$Bz=\kappa\bar z$. For every complex tangent vector $h$ satisfying
$\sum_i\bar z_i h_i=0$, the two real directions $h$ and $ih$ give
\[
 |h^{\mathsf T}Bh|\le\kappa\|h\|^2.
\]
Take $v_i=z_i$ on $S$ and $v_i=0$ elsewhere, and
$h=v-(Q_S/Q)z$. Expanding by pairs gives
\[
 v^{\mathsf T}Bv=\kappa Q_S+C_S,
 \quad h^{\mathsf T}Bh=K_S+C_S,
 \quad \kappa\|h\|^2=K_S.
\]
The asserted disk condition follows; its squared form gives the last
inequality. All pairs are taken exactly once, and no reflection or
radial-group assumption is used.
\end{proof}

For the unit regular octagon and the alternating four-point subset,
$Q_S=4$, $K_S=7$ and $C_S=-16$. Thus $|K_S+C_S|=9>7$.
This reproduces its instability by a concrete subset direction.
The cut is a reusable necessary condition; it does not by itself force
an origin point in an arbitrary eight-point global maximum.

\paragraph{The exact remaining global inequality.}
The origin theorem~\ref{thm:noclass-origin8} can be translated into a
precise formulation of the remaining difficulty. For arbitrary seven
points $w_1,\ldots,w_7$, put $Q_w=\sum|w_i|^2$ and $b=\sum w_i$.
After adjoining $0$ and centering the eight points, their moment is
$Q_w-|b|^2/8$. Hence the unrestricted assertion $M_8=8^{28}/7^{21}$
is equivalent to
\begin{equation}\label{eq:noclass-centroid-target}
 \prod_i|w_i|^2\prod_{i<j}|w_i-w_j|^2
 \le \frac{(Q_w-|b|^2/8)^{28}}{7^{21}}.
\end{equation}
The proved origin inequality has $Q_w^{28}$ in the numerator. It gives
\eqref{eq:noclass-centroid-target} when $b=0$, but an additional
centroid penalty is needed for general $b$. Merely translating one
point to the origin increases its normalized moment and cannot finish
the global argument.

\section{A stronger two-dimensional closure for the six-point problem}
\label{sec:noorigin-oct7-six}

This section gives new necessary inequalities for a nonclassification
proof at six points. It also records an exact obstruction to using the
three-node logarithm majorant as a global certificate by itself. The
large-harmonic region is not closed by the results below.

\subsection{Whole radial and angular regions, without stationarity}

\begin{lemma}\label{lem:noorigin-six-radial-majorant}
For $S(t)=1+t+\cdots+t^5$ and every $t\ge0$,
\[
 \log S(t)\le\log6+\frac52(t-1)+\frac{17}{24}(t-1)^2,
\]
with equality exactly at $t=1$.
\end{lemma}
\begin{proof}
For the difference $g$ of the right and left sides,
\[
 g'(t)=\frac{(t-1)Q_*(t)}{S(t)},\qquad
 Q_*(t)=\frac{17t^5+47t^4+17t^3-t^2-7t-1}{12}.
\]
The coefficients of $Q_*$ have exactly one sign change.
Descartes' rule, $Q_*(0)<0$, and $Q_*(1)>0$ therefore give
exactly one positive root $r\in(0,1)$. Thus $g$ increases on
$[0,r]$, decreases on $[r,1]$, and increases on $[1,\infty)$.
Its only possible global minima are $0$ and $1$.
We have $g(1)=0$ and $g(0)=\log6-43/24>0$; the latter strict
inequality is certified by the rational logarithm enclosure in the
checker. This proves the assertion.
\end{proof}

For actual six points, put
\[
 K_{ij}=\sum_{k=0}^5(z_i\overline z_j)^k,
 \qquad X=\sum_{i<j}\frac{|K_{ij}|^2}{S(t_i)S(t_j)}.
\]
No centering assumption is needed for the next value bound.

\begin{theorem}\label{thm:noorigin-six-whole-regions}
Every six-point configuration with $\sum_i|z_i|^2=6$ obeys
\begin{equation}\label{eq:noorigin-six-all-angle-value}
 \Delta_6\le6^6\exp\left(\frac{17}{24}s-\frac59X\right).
\end{equation}
Consequently each of the following entire regions has value
strictly less than $C_6$:
\[
 s\le\frac3{68},\qquad
 \frac{17}{24}s-\frac59X\le\frac1{32}.
\]
Moreover a collision-free local maximum with $X\ge5/2$ has
value strictly below $C_6$. In particular every global maximum
satisfies the new necessary conditions
\[
 s>\frac3{68},\qquad X<\frac52,\qquad
 X<\frac{51}{40}s-\frac9{160}.
\]
These statements impose no restriction on low analytic harmonics.
\end{theorem}
\begin{proof}
The row-normalized Gram matrix of the square Vandermonde matrix
has diagonal one, trace six, and off-diagonal square sum $X$.
Lemma~\ref{lem:noclass-general-gram}, applied with $m=6$, gives
$\log\det K_{\rm normalized}\le-5X/9$.
Multiplying its determinant by $\prod_iS(t_i)$ and applying
Lemma~\ref{lem:noorigin-six-radial-majorant} proves the value bound.
The singular case has $\Delta_6=0$ and is immediate.

The strict comparison needed for both all-angle regions is
\[
                   e^{1/32}<\frac{6^9}{5^{10}}.
\]
For rational $0<x<K+2$ the positive exponential series, truncated after
$K$ terms, has upper bound
\[
 E_K(x)=\sum_{j=0}^K\frac{x^j}{j!}
       +\frac{x^{K+1}}{(K+1)!}\frac1{1-x/(K+2)}.
\]
All tail ratios are at most $x/(K+2)$, which proves this bound.
The checker proves $E_{12}(1/32)<6^9/5^{10}$ by exact rational
subtraction. Finally at a local maximum $s\le2$, so $X\ge5/2$
forces the exponent at most $17/12-25/18=1/36<1/32$.
A central pentagon competitor has value $C_6$; contraposition
gives the three necessary conditions.
\end{proof}

\subsection{A universal Takagi trace restriction}

\begin{lemma}\label{lem:noorigin-takagi-trace}
At a centered collision-free local maximum with $n\ge3$, normalized
by $\sum_i|z_i|^2=n$, put $\lambda=(n-1)/2$ and $q=|p_2|^2$. Then
\[
                      |p_2|\le\frac{2n}{n+1}.
\]
In particular six-point local maxima satisfy $q\le144/49$.
If $b_i=B_{ii}$ and $T=T_2$, then also
\[
 \sum_i|b_i|^2\ge
 \frac{(\lambda^2T-2\binom n3)^2}{T}
 +\frac{4\binom n3^2q}{T(nT-q)},
\]
and hence, with $D_{-2}=\sum_{i<j}|z_i-z_j|^{-4}$,
\[
 D_{-2}\le\frac{\lambda^2(n-1)}2
 -\frac{(\lambda^2T-2\binom n3)^2}{2T}
 -\frac{2\binom n3^2q}{T(nT-q)}.
\]
\end{lemma}
\begin{proof}
For three distinct complex numbers the elementary divided-difference
identities give
\[
 \sum_{\rm cyc}\frac1{(a-b)(a-c)}=0,\quad
 \sum_{\rm cyc}\frac a{(a-b)(a-c)}=0,\quad
 \sum_{\rm cyc}\frac {a^2}{(a-b)(a-c)}=1.
\]
Squaring the force $F_i=\sum_{j\ne i}(z_i-z_j)^{-1}
=\lambda\overline z_i$ and summing therefore gives
\[
 \operatorname{tr}B=\lambda^2\overline{p_2},\quad
 \sum_i z_i b_i=0,\quad
 \sum_i z_i^2 b_i=\lambda^2T-2\binom n3.
\]
Extract the known Takagi block
$B_0=(\lambda/n)\overline z\,\overline z^{\,T}$.
Because $Bz=\lambda\overline z$ and
$B^*\overline z=\lambda z$, this block is orthogonal to the
remaining operator $B_t=B-B_0$ in both domain and image.
Thus $\|B_t\|_{\rm op}\le\lambda$ and
$B_t\mathbf1=B_tz=0$, giving $\operatorname{rank}B_t\le n-2$.
The nuclear-norm trace bound yields
\[
 (\lambda^2-\lambda/n)|p_2|
       =|\operatorname{tr}B_t|\le(n-2)\lambda.
\]
For $n\ge3$, division simplifies to $2n/(n+1)$.

To prove the diagonal bound, project $b=(b_i)$ onto the span of
$\mathbf1,\overline z^{\circ2}$. Their Gram matrix is
$\bigl(\begin{smallmatrix}n&\overline{p_2}\\p_2&T\end{smallmatrix}\bigr)$,
which is positive definite since $q<n^2\le nT$.
The two known inner products are $\lambda^2\overline{p_2}$ and
$r=\lambda^2T-2\binom n3$. Its projection norm is
\[
 \frac{\lambda^4Tq+nr^2-2\lambda^2rq}{nT-q}
 =\frac{r^2}{T}+\frac{4\binom n3^2q}{T(nT-q)}.
\]
Finally $\|B\|_F^2=2D_{-2}+\sum_i|b_i|^2$ and
$\|B\|_F^2\le\lambda^2(n-1)$ since $B\mathbf1=0$.
\end{proof}

\begin{corollary}\label{cor:noorigin-principal-angle-trace}
Under the local-maximum assumptions of
Lemma~\ref{lem:noorigin-takagi-trace}, the stronger bound is
\[
 |p_2|\le\frac{2n(n-3)}{n(n-1)-4}.
\]
For $n=3$ this is $p_2=0$; the respective bounds for $n=6,7,8$
are $18/13$, $28/19$, and $20/13$.
In particular the six-point squared bound is $q\le324/169$.
\end{corollary}
\begin{proof}
The remainder $B_t$ in the preceding proof has domain supported in
$D=\{\mathbf1,z\}^\perp$ and image supported in
$I=\{\mathbf1,\overline z\}^\perp$. If $P_D,P_I$ are their
orthogonal projectors, trace duality gives
\[
 |\operatorname{tr}B_t|
       =|\operatorname{tr}(P_DP_I B_t)|
       \le\lambda\|P_DP_I\|_*.
\]
Inside $\mathbf1^\perp$, the two subspaces have codimension one.
The singular values of $P_DP_I$ are $1$ with multiplicity $n-3$
and the one remaining principal-angle cosine $|p_2|/n$.
Thus the original trace inequality improves to
\[
 (\lambda^2-\lambda/n)|p_2|
       \le\lambda(n-3+|p_2|/n).
\]
Since $\lambda-2/n>0$ for $n\ge3$, rearranging proves the bound.
\end{proof}

\begin{lemma}[A finite-column inverse-moment lower bound]
\label{lem:noorigin-skew-inverse}
At a centered collision-free stationary configuration with
$\sum_i|z_i|^2=n$, $\lambda=(n-1)/2$, and $q<nT$, write
$T=T_2$, $U=T_3$, and
$F=\lambda^2U-2\lambda(n-2)T+(n-2)^2n$. Then
\begin{equation}\label{eq:noorigin-skew-inverse}
 D_{-1}\ge\frac{\lambda^2T}{n}+
                          \frac{nF-q/4}{nT-q}.
\end{equation}
For six points the bound is
\[
 D_{-1}\ge\frac{25T}{24}
       +\frac{(75/2)U-120T+576-q/4}{6T-q}.
\]
It is sharp at the central pentagon: the right side there is $25/3$.
At the unit regular hexagon it gives $17/2$.
\end{lemma}
\begin{proof}
Define the complex skew-symmetric matrix
$A_{ij}=1/(z_i-z_j)$ for $i\ne j$ and $A_{ii}=0$.
Its squared Frobenius norm is $2D_{-1}$.
For $X=[\mathbf1,z,z^2]$ and $Y=AX$, force gives
\[
 Y=[\lambda\overline z,\ \lambda(t-2\mathbf1),\
                             (\lambda t-(n-2))z].
\]
Rotate so that $p_2\ge0$ is real. Since $R_1=0$,
\[
 G=X^*X=\begin{pmatrix}n&0&p_2\\0&n&0\\p_2&0&T\end{pmatrix},
 \quad
 R=X^TY=\begin{pmatrix}0&-\lambda n&0\\
             \lambda n&0&p_2/2\\0&-p_2/2&0\end{pmatrix}.
\]
The squared norms of the columns of $Y$ are
$\lambda^2n$, $\lambda^2T$, and $F$, and
$(Y^*Y)_{02}=\lambda p_2/2$.

For completeness, the general skew-symmetric completion inequality
is
\[
 \|A\|_F^2\ge
 2\operatorname{tr}(Y^*YG^{-1})
 -\operatorname{tr}(G^{-1}R^*G^{-T}R).
\]
To prove it, orthonormalize $X$ to $Q=XG^{-1/2}$ and extend $Q$
to a unitary matrix. Under unitary congruence the first three columns
of the skew matrix are fixed by $Y$. If they are written as the
blocks $K,L$, the full matrix is
$\bigl(\begin{smallmatrix}K&-L^T\\L&M\end{smallmatrix}\bigr)$.
Its squared norm is $\|K\|_F^2+2\|L\|_F^2+\|M\|_F^2$,
at least $2\|Y G^{-1/2}\|_F^2-\|K\|_F^2$.
This gives the displayed trace formula, with no bound on the
unknown bottom-right block.

Using the displayed $G,R$, the two traces simplify to
\[
 \operatorname{tr}(Y^*YG^{-1})
 =\frac{\lambda^2T}{n}
       +\frac{\lambda^2nT+nF-\lambda q}{nT-q},
\]
\[
 \operatorname{tr}(G^{-1}R^*G^{-T}R)
       =\frac{2\lambda^2nT-2\lambda q+q/2}{nT-q}.
\]
Their combination proves \eqref{eq:noorigin-skew-inverse}.
At the central pentagon $T=36/5$, $U=216/25$, $q=0$;
direct substitution gives $25/3$, its exact inverse-distance sum.
The regular-hexagon substitution is $T=U=6$, $q=0$.
\end{proof}

Continue to use the normalization $p_1=0$, $T_1=6$, put
$s=T_2-6$, and write $q=|p_2|^2$. At an actual local maximum
we have $q<36$, $0\le s\le2$ and $\|B\|_{\rm op}\le5/2$.
For $s>0$, let $W=(\operatorname{span}_{\mathbb C}
\{\mathbf1,z,\overline z,t-\mathbf1\})^\perp$ and
$v_k=\Pi_W z^{\circ k}$, $k=2,3$. Thus $W$ is a
two-dimensional conjugation-invariant space. Define
\[
 A=v_2^*v_2,\quad B_0=v_3^*v_3,\quad G=v_2^*v_3,
 \quad K_0=v_2^Tv_2,\quad L=v_2^Tv_3.
\]

\begin{lemma}\label{lem:noorigin-six-two-dimensional}
If $A>0$, then
\begin{equation}\label{eq:noorigin-six-exact-closure}
 |A L-GK_0|^2=(A^2-|K_0|^2)(AB_0-|G|^2).
\end{equation}
Equivalently,
\begin{equation}\label{eq:noorigin-six-row-closure}
 |L|^2+|G|^2-\frac{2}{A}\operatorname{Re}
       (L\overline G\,\overline{K_0})
       =AB_0-\frac{B_0}{A}|K_0|^2.
\end{equation}
For every $w>1$ this implies the phase-free bound
\begin{equation}\label{eq:noorigin-six-weighted-closure}
 |L|^2+w|G|^2\ge AB_0-
             \frac{w}{w-1}\frac{B_0}{A}|K_0|^2.
\end{equation}
The case $A=0$ has $v_2=0$ and $G=K_0=L=0$; no division by $A$
is used in that case.
\end{lemma}
\begin{proof}
Choose a real orthonormal basis of $W$ and write $v_2=(a,b)$.
The vector $e=(-\overline b,\overline a)$ is Hermitian-orthogonal
to $v_2$ and has squared norm $A$. Write
$v_3=(G/A)v_2+\xi e$. Then
$AB_0-|G|^2=A^2|\xi|^2$ and
$|v_2^Te|^2=A^2-|K_0|^2$, proving
\eqref{eq:noorigin-six-exact-closure}. Expanding and cancelling the
common $|G|^2|K_0|^2$ term proves
\eqref{eq:noorigin-six-row-closure}.
By the ordinary bilinear Cauchy inequality $|L|^2\le AB_0$.
Put $r=|K_0|/A$. The identity gives
\[
 |L|^2+w|G|^2\ge AB_0-\frac{B_0}{A}|K_0|^2
             -2r|L||G|+(w-1)|G|^2.
\]
Completing the square in $|G|$ bounds the last two terms below by
$-r^2|L|^2/(w-1)\ge-r^2AB_0/(w-1)$, proving the claim.
\end{proof}

Unlike the determinant-only identity, this closure does not contain
$p_6$. Its complete moment interface is as follows. Put $d=36-q$,
$T=T_2$, $U=T_3$, $R_2=\sum_i t_i^2z_i$, and
$h=(9p_2/5,p_4)^T$. Then
\begin{align*}
 A={}&T-\frac q6-\frac{6|p_3|^2}{d}-\frac{16q}{25s},\\
 B_0={}&U-\frac{|p_3|^2}{6}
       -h^*\frac1d\begin{pmatrix}6&-\overline{p_2}\\-p_2&6\end{pmatrix}h
       -\frac{9|p_3|^2}{25s},\\
 G={}&R_2-\frac{\overline{p_2}p_3}{6}
       -\frac{\overline{p_3}(6p_4-9p_2^2/5)}{d}
       -\frac{12\overline{p_2}p_3}{25s},\\
 K_0={}&p_4-\frac{p_2^2}{6}
       +\frac{\overline{p_2}p_3^2}{d}-\frac{16p_2^2}{25s},\\
 L={}&p_5-\frac{p_2p_3}{6}
       -\frac{p_3(54p_2/5-\overline{p_2}p_4)}d
       -\frac{12p_2p_3}{25s}.
\end{align*}
These follow directly from Lemma~\ref{lem:noclass-finite-projection}.
In particular $A\ge0$ yields the additional necessary bounds
\[
 q\le\frac{150s(6+s)}{96+25s},\qquad
 |p_3|^2\le\frac{36-q}{6}
       \left(6+s-\frac q6-\frac{16q}{25s}\right).
\]
The weight $w=100$ in \eqref{eq:noorigin-six-weighted-closure}
matches the pair of negative terms
$-|p_5|^2-100|R_2|^2$ in the degree-five distance moment.
The projection corrections displayed above must be retained when
using this statement; replacing $L,G$ by $p_5,R_2$ is not an identity.

\subsection{A distance-positivity cut through the continuous radial core}

Write $D_k=\sum_{i<j}|z_i-z_j|^{2k}$, with $D_0=15$.
The zero-error continuous radial core of the six-point frame has
\[
 \alpha=\frac{36-s^2}{9(s+2)},\qquad
 \beta=\frac{(6+s)(5s-6)}{18(s+2)},\qquad
 T_0=T_1=6,\ T_2=6+s,\quad
 T_j=\alpha T_{j-1}+\beta T_{j-2}.
\]
In this formal core set $p_2=p_3=p_4=R_2=0$,
$|p_5|^2=T_2T_3$, and $p_6=0$. These assignments satisfy
\eqref{eq:noorigin-six-exact-closure} but are not asserted to be
moments of actual points. In this core the formal relation
$t^2=\alpha t+\beta$ is also imposed after multiplication by the
mixed monomials needed below. In particular
$W_2=\sum_i t_i^2z_i^2=\alpha(9p_2/5)+\beta p_2=0$;
the analogous recurrences give the other vanishing mixed terms.
The resulting formal distance moments are
\begin{align*}
 D_1={}&36,&D_2={}&6T_2+72,\\
 D_3={}&6T_3+54T_2,&D_4={}&6T_4+96T_3+18T_2^2,\\
 D_5={}&6T_5+150T_4+99T_2T_3,\\
 D_6={}&6T_6+216T_5+225T_2T_4+200T_3^2
                   -\frac{72}{5}T_2T_3.
\end{align*}
Let $\mathcal L(x^k)=D_k$ and put
\[
 Q(x)=(x-6/5)(x^2-6x+36/5).
\]
Direct rational expansion gives
\begin{equation}\label{eq:noorigin-six-core-hankel}
 \mathcal L(Q^2)=\frac{(5s-6)\,P(s)}{273375(s+2)^4},
\end{equation}
where
\begin{align*}
 P(s)={}&50s^8+280140s^7+9286425s^6+76111596s^5\\
       &+183146508s^4+157689504s^3+67919472s^2
         -419904s-31492800.
\end{align*}
For $s\ge1/2$, $P'(s)>0$ follows already from
$135838944s-419904>0$, with all remaining derivative terms
nonnegative. Also $P(1/2)=2426974375/128>0$.
Thus the entire formal core $1/2\le s<6/5$ violates
$\mathcal L(Q^2)\ge0$, a necessary condition for actual pair distances.
At $s=1$ the exact violation is
\[
                    \mathcal L(Q^2)=-\frac{462520991}{22143375}.
\]
This improves the diagnosis of the core that exceeds the Hermite
objective: distance positivity already excludes its exhibited positive
value at $s=1$,
independently of the integer radial-count objection. It does not
exclude a neighborhood with uncontrolled nonzero harmonic errors.

\subsection{Why the three-node Hermite family is insufficient by itself}

Let $H_5$ be the polynomial interpolating $\log x$ and its first
derivative at $6/5,3-3/\sqrt5,3+3/\sqrt5$. For the actual unit
regular hexagon its squared distances $1,3,4$ have multiplicities
$6,6,3$. Outward rational interval arithmetic proves
\[
 \frac{81}{1000}<6H_5(1)+6H_5(3)+3H_5(4)-\log C_6
                       <\frac{83}{1000}.
\]
Its $Q^2$ distance sum is $396/5$.
If $H_6=H_5-\kappa Q^2$ is a logarithm majorant on $(0,8]$,
evaluation at $8$ necessarily gives
\[
 \kappa\le\frac{H_5(8)-\log8}{Q(8)^2}<\frac3{5000}.
\]
Consequently the regular hexagon still has
\[
                 \sum_{i<j}H_6(|z_i-z_j|^2)-\log C_6
                  >\frac{81}{1000}-\frac3{5000}\frac{396}{5}
                  >\frac{33}{1000}.
\]
For a majorant on $(0,12]$ the analogous endpoint bound is even
smaller: $\kappa<1/2500$. Hence this entire three-node degree-six
family cannot supply an unconditional global upper bound on its own,
even though the hexagon has actual objective less than $C_6$.
This is an obstruction to that proof relaxation, not a counterexample
to the desired six-point inequality. A successful use must retain
stationary constraints and the already proved small-harmonic rigidity,
or incorporate additional distance nodes.

The accompanying standard-library checker verifies the closure over
exact Gaussian rationals, the full rational identity
\eqref{eq:noorigin-six-core-hankel}, and the logarithm enclosures.
It is safe under optimized Python. These new constraints are suitable
for a finite scalar certificate, but no such full six-point certificate
is claimed here.

\subsection{A Laurent majorant with a lower moment requirement}

There is an alternative logarithm majorant that removes all positive
distance moments of degree four and above. Let $P_5$ interpolate
$f(x)=x^2\log x$ and $f'$ at the same three candidate distance nodes.
Since $f^{(6)}(x)=-12/x^4<0$, Hermite's remainder proves the global
majorant
\[
 \log x\le\frac{P_5(x)}{x^2}
        =\sum_{j=-2}^3 a_jx^j\qquad(x>0).
\]
Equality holds at precisely the three interpolation nodes.
Consequently
\[
             \log\Delta_6\le\sum_{j=-2}^3a_jD_j.
\]
The coefficient signs $a_{-2}>0$, $a_{-1}<0$, $a_2<0$, $a_3>0$
are verified by outward rational enclosures in the checker.
At the actual unit regular hexagon the Laurent upper value is
strictly below $\log C_6$ by more than $1/50$; thus the obstruction
to the polynomial Hermite family above does not apply here.

At a stationary local maximum the degree-three force defect is
\[
 0\le\frac{25}{4}\left(T_3-\frac{|p_3|^2}{6}\right)
       -\sum_i t_i\left(\frac{15}{2}t_i-9\right)^2
 =-50T_3+135T_2-486-\frac{25}{24}|p_3|^2.
\]
Thus $T_3\le162/25+27s/10-|p_3|^2/48$.
Substitution, using $a_3>0$ and the exact identities for $D_2,D_3$,
gives the rigorous inverse-moment interface
\begin{align*}
 \log\Delta_6-\log C_6\le{}& c_0+c_s s+c_q q
       -\frac98a_3|p_3|^2+a_{-2}D_{-2}+a_{-1}D_{-1},\\
 c_0={}&15a_0+36a_1+108a_2+\frac{9072}{25}a_3-\log C_6,\\
 c_s={}&6a_2+\frac{351}{5}a_3<0,\qquad
 c_q=a_2+\frac{54}{5}a_3<0.
\end{align*}
Here $D_{-1}=\sum_{i<j}|z_i-z_j|^{-2}$ and $D_{-2}$ is as above.
The new diagonal and Takagi trace bounds constrain these quantities,
but are not yet sharp enough to prove the desired inequality on the
remaining region. The unresolved issue is a joint inverse-moment
bound, rather than positive distance moments $D_4,D_5,D_6$.

% Standalone fragment; standard amsmath/amssymb/amsthm preamble suffices.
\section{Deleting one vertex: an exact force defect and a quantitative gap}
\label{sec:noorigin-six-centroid}

Let a collision-free six-point stationary configuration satisfy
$\sum_i z_i=0$, $\sum_i|z_i|^2=6$, and
\[
 \sum_{j\ne i}\frac1{z_i-z_j}=\frac52\overline{z_i}.
\]
Choose a vertex of smallest squared radius $t=|z_i|^2\le1$.
For $j\ne i$ put
\[
 v_j=z_j-z_i,\quad w_j=z_j+z_i/5,
 \quad Q=\sum_{j\ne i}|w_j|^2=6-6t/5,
 \quad R=\sum_{j\ne i}|v_j|^{-2}.
\]
Thus the five $w_j$ are centered. Define their defect from the
five-point force at the same moment by
\[
 g_j=\sum_{\substack{k\ne i,j}}\frac1{w_j-w_k}
                   -\frac{10}{Q}\overline{w_j}.
\]

\begin{proposition}\label{prop:noorigin-six-deletion}
With the notation above, put
\[
 c=\frac{5-3t}{Q},\qquad
 \kappa=\frac{5(1-t)}{2(5-t)}.
\]
Then
\begin{align}
 -g_j&=\frac1{v_j}-c\overline{v_j}
                         -\kappa\overline{z_i},\label{eq:six-delete-residual}\\
 \|g\|^2
 &=R-\frac54t-\frac{(5-3t)^2}{Q}
   =\frac{25}2+\frac{25}4t+R-\frac{100}{Q}.\label{eq:six-delete-norm}
\end{align}
Consequently
\begin{equation}\label{eq:six-delete-reciprocal}
 R\ge\frac{25}{6(1+t)}
       +\frac{25t(1-t)^2}{4(5-t)(1+t)}.
\end{equation}
If $0<t<1$ and equality holds in
\eqref{eq:six-delete-reciprocal}, all six points are collinear.
If $t=1$, equality cannot hold for a collision-free stationary
configuration. If $t=0$, equality is equivalent to the origin
together with a regular pentagon of squared radius $6/5$.

If the configuration is an actual local maximum, then for every
chosen nonzero vertex of smallest radius,
\[
 \|g\|^2\ge
 \frac{25t(1-t)^2}{64(5-t)^2}
 \left[6-\operatorname{Re}\left(p_2\frac{\overline{z_i}^{\,2}}t\right)\right],
 \qquad p_2=\sum_k z_k^2.
\]
In particular, with $s=T_2-6$ its right side is at least
$25t(1-t)^2[6-\sqrt{36-18s}]/[64(5-t)^2]$.
For this same nonzero chosen vertex, the uniform estimate is
\begin{equation}\label{eq:six-delete-uniform-gap}
                 \|g\|^2>\frac{t}{20000000}.
\end{equation}
This is a force defect bound, rather than a sharp objective bound.
\end{proposition}

\begin{proof}
The full force equation gives
\[
 \sum_{j\ne i}v_j=-6z_i,\qquad
 \sum_{j\ne i}\frac1{v_j}=-\frac52\overline{z_i}.
\]
Project the vectors $(\overline{v_j})_j$ and $(1/v_j)_j$
off the constant vector. Their squared norms are $Q$ and
$R-5t/4$, and their Hermitian inner product is $5-3t$.
Their orthogonal residual therefore has squared norm
$R-5t/4-(5-3t)^2/Q$.
Also, using the full force equation at each remaining point,
\[
 g_j=\left(\frac52-\frac{10}{Q}\right)\overline{w_j}
                     -\frac12\overline{z_i}-\frac1{v_j}.
\]
The identities $5/2-10/Q=c$ and $6c/5-1/2=\kappa$
prove \eqref{eq:six-delete-residual} and
\eqref{eq:six-delete-norm}. Elementary reduction of
$5t/4+(5-3t)^2/Q$ gives the right side of
\eqref{eq:six-delete-reciprocal}.

If $g=0$, multiplication of \eqref{eq:six-delete-residual}
by $v_j$ yields
\[
       1-c|v_j|^2-\kappa\overline{z_i}v_j=0.
\]
For $0<t<1$, its imaginary part puts every $v_j$ on the line
through $z_i$, and hence all six points on a line.
For $t=1$, it gives $|v_j|^2=1/c=12/5$ for all five points.
For points on this circle,
\[
 \operatorname{Re}\!\left[
 v_j\sum_{k\ne i,j}\frac1{v_j-v_k}\right]=2,
\]
because each of the four pair terms has real part $1/2$.
The full force at $z_j$ now gives
$3=(5/2)(12/5+\operatorname{Re}(\overline{z_i}v_j))$.
Thus all five points have the same projection onto the line of
$z_i$, and the same distance from its origin; at most two distinct
points can satisfy these conditions, a contradiction.
For $t=0$, equality says $|v_j|^2=6/5$.
The monic polynomial $P(X)=\prod_{j\ne i}(X-v_j)$ obeys
$v_jP''(v_j)-4P'(v_j)=0$. This polynomial has degree at most
four and vanishes at five distinct points, so it is identically
zero. Its coefficients force $P(X)=X^5-b$, $b\ne0$,
which proves the asserted equality configuration and its converse.

We next prove the quantitative local-maximum assertion.
We use only the already established local-maximum consequences
\[
 |z_j-z_k|^2\ge\frac{12}{25},\qquad T_2=\sum_k|z_k|^4\le8.
\]
The latter gives $|z_j|^2<3$, so for $t\le1$ we have
$|v_j|^2<8$.
For any choice of axis, write $z_k=x_k+\mathrm i y_k$ and
$b=\sum_k y_k^2$. Then
\begin{equation}\label{eq:six-delete-strip}
                             b\ge\frac6{85}.
\end{equation}
Indeed, if $b<6/85$, every pair satisfies
$(y_j-y_k)^2\le2b<12/25$. Sorting the six real coordinates
therefore gives adjacent gaps at least
$\sqrt{12/25-2b}$. The variance identity implies
\[
 \sum_kx_k^2\ge\frac{35}{2}\left(\frac{12}{25}-2b\right).
\]
Hence $6=\sum_kx_k^2+b\ge42/5-34b$, which contradicts
$b<6/85$.

Rotate $z_i=\sqrt t>0$. The imaginary part of the residual
identity implies
\[
 |g_j|^2|v_j|^2\ge\kappa^2t\,(\operatorname{Im}z_j)^2,
 \qquad
 \|g\|^2>\frac{3\kappa^2t}{340}.
\]
Since $\sum_j(\operatorname{Im}z_j)^2
=(6-\operatorname{Re}(p_2\overline{z_i}^{\,2}/t))/2$
in this rotated coordinate system, $|v_j|^2<8$ gives the displayed
moment-dependent estimate. The established local-maximum bound
$|p_2|^2\le36-18s$ gives its stated scalar consequence.
If $t\le199/200$, then $\kappa\ge1/400$, and thus
$\|g\|^2>3t/54400000>t/20000000$.

It remains to consider $199/200\le t\le1$.
Suppose for contradiction that $\|g\|\le1/1000$.
Put $d_j=|v_j|^2$ and $a=1/c$. On this interval
\[
 c\ge5/12,\quad 19/8<a\le12/5,
 \quad 0\le\kappa\le1/320.
\]
Equation \eqref{eq:six-delete-residual}, $\sqrt{d_j}<3$,
and $\sqrt t\le1$ give
\[
 |d_j-a|\le\frac{\sqrt{d_j}}c
                (|g_j|+\kappa\sqrt t)
 <\frac{36}5\left(\frac1{1000}+\frac1{320}\right)
 =\frac{297}{10000}<\frac3{100}.
\]
Write $v_j=x_j+\mathrm i y_j$. The exact pair identity
\[
 \operatorname{Re}\frac{v_j}{v_j-v_k}
       =\frac12+\frac{d_j-d_k}{2|v_j-v_k|^2}
\]
and the full force at $z_j$ give
\[
 x_j=\frac{6/5-d_j}{\sqrt t}
       +\frac1{5\sqrt t}\sum_{k\ne i,j}
                                  \frac{d_j-d_k}{|v_j-v_k|^2}.
\]
With
$x_0=(6/5-a)/\sqrt t=-(12/5)\sqrt t/(5-3t)$,
the separation bound therefore gives
\[
 |x_j-x_0|<\frac{13}{3\sqrt t}\frac3{100}
              \le\frac{26}{199},\qquad
 -6/5\le x_0<-47/40.
\]
In particular $-11/8<x_j<-1$. Also
$469/200<d_j<243/100$, so
\[
 \frac49<\frac{469}{200}-\frac{121}{64}
        <y_j^2<\frac{243}{100}-1<\frac{36}{25}.
\]
Thus each of the five points belongs to one of the two rectangles
\[
 (-11/8,-1)\times(2/3,6/5),\qquad
 (-11/8,-1)\times(-6/5,-2/3).
\]
Each rectangle has squared diameter
$9/64+64/225<12/25$. The pair separation permits at most
one point in each rectangle, contradicting the existence of five
distinct points. We conclude $\|g\|>1/1000$, which also
implies \eqref{eq:six-delete-uniform-gap}.
\end{proof}

\paragraph{What the direct translation penalty can and cannot prove.}
Write $T=\sum_k|z_k|^4$ and $p_2=\sum_k z_k^2$.
The stationary identity $\sum_k|z_k|^2z_k=0$ gives
\[
 \sum_{j\ne i}|z_j-z_i|^4
       =T+6t^2+24t+2\operatorname{Re}(p_2\overline{z_i}^{\,2}).
\]
Translate $z_i$ to the origin and rescale the remaining points by
$\sqrt{5/[6(1+t)]}$. The established origin-containing bound
therefore gives the exact necessary upper bound
\[
 \log\frac{\Delta_6}{C_6}
 \le15\log(1+t)-\frac1{250}
 \left[
 \frac{25\bigl(T+6t^2+24t+
       2\operatorname{Re}(p_2\overline{z_i}^{\,2})\bigr)}
      {36(1+t)^2}-5\right].
\]
The bracket is the variance of five nonnegative numbers of sum five,
so it lies in $[0,20]$. This particular penalty is consequently at
most $0.08$; it cannot alone establish the desired upper bound
when $t\ge e^{2/375}-1$.
For the unit regular hexagon the bracket is exactly $5/4$, and
the resulting upper-bound ratio is $2^{15}e^{-1/200}>1$,
although its actual ratio is $5^{10}/6^9<1$.

\paragraph{A stronger deletion bound with an explicit dependency.}
If the established sharp five-point inequality
$\Delta_5(w)\le5^5(Q/5)^{10}$ is used as an input, then
the arithmetic-geometric mean inequality on the five star distances
gives
\[
 \frac{\Delta_6}{C_6}
             \le(1+t)^5(1-t/5)^{10}.
\]
Multiplying the six deletion inequalities before applying any
star-distance relaxation instead gives
\[
 \frac{\Delta_6}{C_6}
       \le\left[5\prod_{i=1}^6(1-t_i/5)\right]^{5/2}
       \le\left[5(4/5)^6\right]^{5/2}<2.
\]
Neither bound is sharp at a regular hexagon: the first ratio is
$33554432/9765625>3$, and the second is greater than one.
These bounds use the five-point theorem, not a classification of
six-point stationary shapes. The force defect above is a new
constraint that such continuous deletion bounds do not retain;
turning its strict positive gap into a sufficiently strong
objective deficit remains open.

% New no-classification seven-point regions. Standard theorem macros suffice.
\section{Seven-point regions excluded without a stationary classification}
\label{sec:noorigin-oct7-seven}

Write $t_i=|z_i|^2$, $\sum_{i=1}^7t_i=7$, and
$s=\sum_i(t_i-1)^2$. The central hexagon competitor has value
\[
 C_7=\frac{7^{21}}{6^{15}},\qquad
 R:=\frac{C_7}{7^7}=\frac{7^{14}}{6^{15}}.
\]
The conclusions below concern actual seven points and retain their
pairwise angular correlations. No Faug\`ere--Svartz component,
stationary-shape catalogue, or numerical coverage assumption is used.

\begin{lemma}[A seven-column radial majorant]
\label{lem:noorigin-oct7-radial-majorant}
For $S(t)=1+t+\cdots+t^6$ and every $t\ge0$,
\begin{equation}\label{eq:noorigin-oct7-radial-majorant}
 \log S(t)\le\log7+3(t-1)+\frac87(t-1)^2.
\end{equation}
Equality holds exactly at $t=1$.
\end{lemma}
\begin{proof}
Put $g(t)=\log7+3(t-1)+8(t-1)^2/7-\log S(t)$.
Direct multiplication gives
\[
 g'(t)=\frac{(t-1)Q(t)}{S(t)},\qquad
 Q(t)=\frac{2-5t-5t^2+2t^3+16t^4+37t^5+16t^6}{7}.
\]
We give a finite rational proof that $Q(t)>0$ for $t\ge0$.
On $[a,b]$, write $t=a+(b-a)y$ and let $A_h$ be the coefficient
of $y^h$ in $Q(a+(b-a)y)$. Its degree-six Bernstein coefficients are
\[
 B_k=\sum_{h=0}^k A_h\frac{\binom kh}{\binom6h},\qquad
 A_h=(b-a)^h\sum_{j=h}^6[t^j]Q(t)\binom jh a^{j-h}.
\]
For the eight intervals $[j/8,(j+1)/8]$, $0\le j\le7$,
the respective minimum Bernstein coefficients are
\[
 \frac{6113}{32768},\quad\frac{585}{7168},\quad
 \frac{5329}{229376},\quad\frac{31315}{1376256},\quad
 \frac{29}{224},\quad\frac{142011}{229376},\quad
 \frac{13283}{7168},\quad\frac{1005125}{229376}.
\]
They are all positive. The Bernstein basis is nonnegative and sums
to one, so $Q>0$ on $[0,1]$. The ascending coefficient list of
$Q(1+x)$ is
\[
       (9,48,101,108,63,19,16/7),
\]
which is positive for $x\ge0$. Thus $g$ strictly decreases to
$g(1)=0$ and strictly increases after $1$, proving the result.
All displayed coefficient identities and signs are replayed using
standard-library rational arithmetic by the accompanying checker.
\end{proof}

\begin{theorem}[A whole radial-variance region]
\label{thm:noorigin-oct7-nearcircle}
Every seven-point configuration with $\sum_i|z_i|^2=7$ obeys
\begin{equation}\label{eq:noorigin-oct7-variance-bound}
                    \Delta_7(z)\le7^7\exp(8s/7).
\end{equation}
Consequently,
\[
                   s\le\frac8{25}\quad\Longrightarrow\quad
                   \Delta_7(z)<C_7.
\]
This implication imposes no restriction on $p_2,\ldots,p_5$ or on
the angles. In particular every global maximizer has $s>8/25$.
\end{theorem}
\begin{proof}
Let $V_{ik}=z_i^k$, $0\le k\le6$. Its squared determinant is
$\Delta_7$ and its squared row norms are $S(t_i)$. Hadamard and
the preceding lemma give
\[
 \log\Delta_7\le\sum_i\log S(t_i)
             \le7\log7+\frac87s,
\]
with the zero-determinant case immediate. It remains to justify the
strict value comparison without a floating-point logarithm.
For rational $x>0$ and $K+2>x$, the positive exponential series gives
\begin{equation}\label{eq:noorigin-oct7-exp-tail}
 e^x\le E_K(x):=\sum_{j=0}^K\frac{x^j}{j!}
             +\frac{x^{K+1}}{(K+1)!}\frac1{1-x/(K+2)}.
\end{equation}
This follows because the ratio of two successive tail terms is at
most $x/(K+2)$. Exact subtraction at $x=64/175$, $K=12$, gives
\[
 R-E_{12}(64/175)=
 \frac{10616970982804289077504712500519377022727929}
 {11580834889123260173059130859375000000000000000}>0.
\]
Since $8s/7\le64/175$, the strict claim follows. A central hexagon
attains $C_7$, so the final necessary condition follows as well.
\end{proof}

\begin{theorem}[An angular-correlation exclusion]
\label{thm:noorigin-oct7-correlations}
For actual seven points define
\[
 K_{ij}=\sum_{k=0}^6(z_i\overline z_j)^k,\qquad
 X=\sum_{i<j}\frac{|K_{ij}|^2}{S(t_i)S(t_j)}.
\]
Then $0\le X\le21$ and
\begin{align}
 \Delta_7&\le\left\{\prod_iS(t_i)\right\}
                    (1-X/21)^{7/2},
                    \label{eq:noorigin-oct7-correlation-bound}\\
 \Delta_7&\le7^7\exp(8s/7-X/2).
                    \label{eq:noorigin-oct7-correlation-linear}
\end{align}
Thus every configuration satisfying
\[
                      \frac87s-\frac X2\le\frac{64}{175}
\]
has $\Delta_7<C_7$. Every collision-free seven-point local maximum
with $X\ge14/3$ also has $\Delta_7<C_7$.
In particular, every global maximizer necessarily satisfies
\begin{equation}\label{eq:noorigin-oct7-necessary-region}
 s>\frac8{25},\qquad
 X<\frac{14}{3},\qquad
 X<\frac{16}{7}s-\frac{128}{175}.
\end{equation}
\end{theorem}
\begin{proof}
The row-normalized Gram matrix $A_{ij}=K_{ij}/\sqrt{S(t_i)S(t_j)}$
is positive semidefinite with diagonal one and trace seven.
If $\mu_1,\ldots,\mu_7$ are its nonnegative eigenvalues, then
\[
 e_2(\mu)=\frac{(\operatorname{tr}A)^2-\operatorname{tr}(A^2)}2
        =21-X.
\]
Maclaurin's inequality for nonnegative numbers gives
$\det A\le(e_2/21)^{7/2}$. Multiplying by $\prod_iS(t_i)$
proves \eqref{eq:noorigin-oct7-correlation-bound}; it also proves
$X\le21$. The seven-row spectral deficit
Lemma~\ref{lem:noclass-seven-gram} gives the stronger linear
estimate $\log\det A\le-X/2$. This applies whenever the points
are distinct, since $A$ is then positive definite; otherwise
$\Delta_7=0$. Combining that estimate with the radial majorant proves
\eqref{eq:noorigin-oct7-correlation-linear}. The exponential
comparison in the preceding proof gives the first exclusion.

For completeness, second-order optimality at a local maximum gives
$\|B\|_{\rm op}\le3$, where $B$ is the inverse-square Laplacian.
First-order optimality gives $\sum_i z_i=0$ and
$B(z^{\circ2}-p_2\mathbf1/7)=6(t-\mathbf1)$.
Since $\|z^{\circ2}-p_2\mathbf1/7\|^2=7+s-|p_2|^2/7$,
\[
             36s\le9(7+s-|p_2|^2/7),\qquad s\le7/3.
\]
If $X\ge14/3$, the linear correlation bound is therefore
\[
             \Delta_7\le7^7e^{8/3-7/3}
                 =7^7e^{1/3}<7^7e^{64/175}<C_7.
\]
The first strict inequality uses $1/3<64/175$; the second is
the exact exponential certificate above. A global maximizer is a
collision-free local maximum
with value at least $C_7$, giving
\eqref{eq:noorigin-oct7-necessary-region}.
\end{proof}

These are genuine value bounds on entire regions, including
origin-free configurations with unrestricted low harmonics.
The remaining region in \eqref{eq:noorigin-oct7-necessary-region}
has not been shown empty or bounded sharply by $C_7$ here. In
particular this argument does not prove the full seven-point
global theorem without classification.

% Exact seven-point finite-rank interface, without stationary classification.
\subsection{A sixth-moment circle from the three real remaining directions}
\label{sec:noorigin-oct7-projection}

\begin{lemma}[The seven-point truncated conjugate closure]
\label{lem:noorigin-oct7-truncated-closure}
Let $z$ be a collision-free seven-point local maximum, centered and
normalized to $\sum_i|z_i|^2=7$. Put
\[
 p_j=\sum_i z_i^j,\quad T_j=\sum_i|z_i|^{2j},\quad
 s=T_2-7>0,\quad q=|p_2|^2,\quad D=49-q,\quad
 R_2=\sum_i|z_i|^4z_i.
\]
Then $D>0$. Let $u=(|z_i|^2-1)_i$,
$W=\operatorname{span}_{\mathbb C}(\mathbf1,z,\overline z,u)^\perp$,
and $v_j=\Pi_Wz^{\circ j}$ for $j=2,3$. Write
\[
 A=\|v_2\|^2,\quad K=v_2^{\mathsf T}v_2,\quad
 B=\|v_3\|^2,\quad M=v_3^{\mathsf T}v_3,\quad
 G=v_2^*v_3,\quad L=v_2^{\mathsf T}v_3.
\]
The following formulas use no moment $p_j$ with $j>6$:
\begin{align}
 A={}&T_2-\frac q7-\frac{7|p_3|^2}{D}-\frac{25q}{36s},
       \label{eq:noorigin-oct7-A}\\
 K={}&p_4-\left(\frac17+\frac{25}{36s}\right)p_2^2
                 +\frac{\overline{p_2}p_3^2}{D},
       \label{eq:noorigin-oct7-K}\\
 B={}&T_3-\left(\frac17+\frac4{9s}\right)|p_3|^2
       -\frac{847q/36+7|p_4|^2-(11/3)
                 \operatorname{Re}(\overline{p_2}^{2}p_4)}D,
       \label{eq:noorigin-oct7-B}\\
 M={}&p_6+J,\qquad
 J=-\left(\frac17+\frac4{9s}\right)p_3^2
       +\frac{121p_2^3/36-(77/3)p_2p_4+\overline{p_2}p_4^2}{D},
       \label{eq:noorigin-oct7-J}\\
 G={}&R_2-\left(\frac17+\frac5{9s}\right)\overline{p_2}p_3
       +\frac{\overline{p_3}(11p_2^2/6-7p_4)}D,
       \label{eq:noorigin-oct7-G}\\
 L={}&p_5-\left(\frac17+\frac5{9s}\right)p_2p_3
       -\frac{p_3(77p_2/6-\overline{p_2}p_4)}D.
       \label{eq:noorigin-oct7-L}
\end{align}
In particular $A\ge|K|$. If $A>|K|$, put
\[
 d=A^2-|K|^2,\qquad
 U=\frac{A(|G|^2+|L|^2)-2\operatorname{Re}(KG\overline L)}d,
 \qquad
 V=\frac{-KG^2+2AGL-\overline K L^2}d.
\]
Then $U\ge0$, $B-U\ge0$, $|V|\le U$, and the exact constraint is
\begin{equation}\label{eq:noorigin-oct7-sixth-circle}
                  |p_6+J-V|=B-U.
\end{equation}
Consequently
\begin{align}
 |p_6|&\ge\ell_6:=\bigl|B-U-|J-V|\bigr|,
                    \label{eq:noorigin-oct7-sixth-lower}\\
 |p_6|&\ge\left[B-\frac{2(|G|^2+|L|^2)}{A-|K|}-|J|\right]_+.
                    \label{eq:noorigin-oct7-sixth-simple}
\end{align}
The assumption $A>|K|$ is essential for these divided formulas;
the degenerate case is retained separately below.
\end{lemma}
\begin{proof}
Stationarity gives $p_1=0$ and $\sum_i|z_i|^2z_i=0$.
For the inverse-square Laplacian, second-order optimality gives
$\|B_{\rm lap}\|\le3$. Equality $q=49$ would put the seven
points on a real line after rotation. Then
$u=z^{\circ2}-\mathbf1\ne0$ and $B_{\rm lap}u=6u$, a
contradiction. Thus $D>0$.

Put $E=\operatorname{span}_{\mathbb C}(\mathbf1,z,\overline z)$.
The two nonconstant columns have Gram matrix
\[
 K_E=\begin{pmatrix}7&\overline{p_2}\\p_2&7\end{pmatrix},
 \qquad
 K_E^{-1}=\frac1D
       \begin{pmatrix}7&-\overline{p_2}\\-p_2&7\end{pmatrix}.
\]
Moreover $u\perp E$ and $\|u\|^2=s$. The paired-force identities
give the projection data
\[
 h_2=\binom0{p_3},\quad
 h_3=\binom{11p_2/6}{p_4},\qquad
 \beta_2=5p_2/6,\quad\beta_3=2p_3/3,
\]
where $h_j=(\langle z,z^j\rangle,
\langle\overline z,z^j\rangle)^{\mathsf T}$ and
$\beta_j=\langle u,z^j\rangle$. Orthogonal projection subtracts
$\overline{p_j}p_k/7+h_j^*K_E^{-1}h_k+
\overline{\beta_j}\beta_k/s$ from Hermitian products.
It subtracts
$p_jp_k/7+h_j^{\mathsf T}J_2K_E^{-1}h_k+\beta_j\beta_k/s$
from transpose products, with
$J_2=\left(\begin{smallmatrix}0&1\\1&0\end{smallmatrix}\right)$.
The latter formula follows because the two-column span and $u$
are invariant under conjugation, and
$K_E^{-\mathsf T}J_2=J_2K_E^{-1}$.
Using $\langle z^2,z^3\rangle=R_2$ and expanding gives
\eqref{eq:noorigin-oct7-A}--\eqref{eq:noorigin-oct7-L}.

The space $W$ is conjugation-invariant and has complex dimension
three, hence has an orthonormal basis of three real vectors.
The Hermitian Gram matrix of $v_2,\overline{v_2}$ is
\[
 K_2=\begin{pmatrix}A&\overline K\\K&A\end{pmatrix}.
\]
It is positive definite exactly when $A>|K|$. Project $v_3$ off
their span and call the residual $r$. Its projection vector has
inner products $h=(G,L)^{\mathsf T}$ with that pair. The same
projection formulas give
\[
 \|r\|^2=B-h^*K_2^{-1}h=B-U,\qquad
 r^{\mathsf T}r=M-h^{\mathsf T}J_2K_2^{-1}h=M-V.
\]
The residual space has complex dimension one and is invariant
under conjugation. Thus $r=\gamma e$ for a real unit vector $e$,
and $|r^{\mathsf T}r|=\|r\|^2$. This proves the exact circle.
The removed projection has norm squared $U$ and transpose square
$V$, proving $|V|\le U$. The inverse matrix has norm
$1/(A-|K|)$, so $U\le(|G|^2+|L|^2)/(A-|K|)$. Reverse triangle
inequalities give the first lower bound in both directions;
its weaker positive-part version, together with $|V|\le U$ and
the upper bound on $U$, gives the second.
\end{proof}

\paragraph{A division-free equation and the degenerate branch.}
For all cases, including $A=|K|$, the joint Gram matrix
\begin{equation}\label{eq:noorigin-oct7-joint-gram}
 \mathcal G=
 \begin{pmatrix}
 A&G&\overline K&\overline L\\
 \overline G&B&\overline L&\overline M\\
 K&L&A&\overline G\\
 L&M&G&B
 \end{pmatrix}
 \succeq0,\qquad \det\mathcal G=0
\end{equation}
because four vectors live in $W$ of dimension three.
Put
\[
 U_0=A(|G|^2+|L|^2)-2\operatorname{Re}(KG\overline L),\qquad
 V_0=-KG^2+2AGL-\overline K L^2.
\]
The polynomial identity
\begin{equation}\label{eq:noorigin-oct7-cleared-rank}
 d\det\mathcal G=(Bd-U_0)^2-|Md-V_0|^2
\end{equation}
is valid without division. It gives the squared form of the circle
when $d>0$. If $A=0$, then $v_2=0$ and $K=G=L=0$.
If $A=|K|>0$, then $v_2$ is a complex multiple of a real vector:
equality in $|v_2^{\mathsf T}v_2|\le\|v_2\|^2$ gives a common
phase to its nonzero coordinates. The pair spans only one complex
direction and its removal leaves two real directions. In that
case one has only the ordinary transpose Cauchy inequality, not
\eqref{eq:noorigin-oct7-sixth-circle}. No sixth-moment lower bound
from the one-dimensional residual is asserted on this branch.

\begin{corollary}[A finite-moment value certificate]
\label{cor:noorigin-oct7-value-certificate}
Under the nondegenerate hypotheses of the preceding lemma,
\begin{equation}\label{eq:noorigin-oct7-Fischer-sixth}
 \Delta_7\le
             (7T_6-\ell_6^2)\,7T_2T_3T_4T_5.
\end{equation}
Consequently, the finite-moment condition
\begin{equation}\label{eq:noorigin-oct7-certified-region}
     (7T_6-\ell_6^2)\,7T_2T_3T_4T_5\le C_7
\end{equation}
proves $\Delta_7\le C_7$. Equality for an actual local maximum
satisfying this condition holds exactly at a central regular hexagon.
In particular every origin-free global maximizer in the
nondegenerate branch must violate
\eqref{eq:noorigin-oct7-certified-region} strictly.
\end{corollary}
\begin{proof}
Apply Hadamard--Fischer to the Gram matrix of
$1,z,\ldots,z^6$, with block $\{1,z^6\}$ and the five remaining
singletons. Its block determinant is $7T_6-|p_6|^2$, and the
remaining squared norms multiply to $7T_2T_3T_4T_5$. Substituting
the lower bound on $|p_6|$ proves the value certificate.
The quantity $7T_6-\ell_6^2$ is nonnegative for actual points,
by $\ell_6\le|p_6|\le\sqrt{7T_6}$.

If equality with $C_7$ holds throughout, the positive definite
Vandermonde Gram matrix has equality in Hadamard--Fischer.
All cross-block entries vanish, in particular
$p_2=p_3=p_4=p_5=0$. The purely algebraic zero-harmonics
Lemma~\ref{lem:noclass-zero-harmonics} gives either a regular
heptagon or a central regular hexagon. The former has value
$7^7<C_7$, so only the latter remains. Conversely, at a central
hexagon, $R_2=p_2=\cdots=p_5=0$,
$A=T_2=49/6$, $B=T_3=343/36$, and $G=L=J=K=0$.
Thus $|p_6|=\ell_6=T_3$, the certificate is exactly sharp,
and Hadamard--Fischer is an equality.
\end{proof}

The truncated closure is a new exact interface involving
only $p_2,\ldots,p_6$, $R_2$, and radial moments. It rejects the
continuous central-circle relaxation even at the candidate radial
profile: that relaxation has $A=T_2$, $B=T_3>0$, $G=L=J=K=0$,
but $p_6=0$, contrary to \eqref{eq:noorigin-oct7-sixth-circle}.
It has not been proved that every remaining actual local maximum
satisfies the value condition. The nondegenerate and degenerate
residual regions are therefore both still open in this argument.

\section{Further unrestricted eight-point reductions without classification}
\label{sec:noorigin-oct7-eight}

The results here address the part of the problem for which an origin
point has not been established. They use no stationary-shape list.
They give a substantial exclusion region and a stronger finite-point
projection interface, but do not yet prove the unrestricted sharp bound.

\subsection{A full angular exclusion region from a scalar logarithm bound}
\begin{lemma}\label{lem:noorigin-oct7-eight-tangent}
For every $t\ge0$, with $S(t)=\sum_{k=0}^7t^k$,
\[
 \log S(t)\le\log8+\frac72(t-1)+\frac{11}{7}(t-1)^2.
\]
Equality occurs only at $t=1$.
\end{lemma}
\begin{proof}
Let $F(t)$ be the right-hand side minus the left-hand side.
An exact polynomial division gives
\[
 F'(t)=\frac{(t-1)Q(t)}{S(t)},\qquad
 Q(t)=\frac{1}{14}(9-12t-19t^2-12t^3+9t^4+44t^5+93t^6+44t^7).
\]
For $0\le t\le1$, expand $Q(a+x/32)$ in the Bernstein basis
of degree seven on $0\le x\le1$, for
$a=j/32$, $0\le j<32$. The minimum of all $256$ coefficients is
$27/784>0$. For $t=1+x$, all eight monomial coefficients are strictly
positive. These exact identities are recorded in
\path{noorigin_oct7_eight_result.json}; the independent coefficient
formula used by the standard-library checker is
\[
 b_k=\sum_{i=0}^k[x^i]Q(a+x/32)
                         \frac{\binom ki}{\binom7i}.
\]
Consequently $Q(t)>0$ on the full half-line. Since $F(1)=0$,
$F$ strictly decreases to $1$ and then strictly increases, proving
the assertion. The displayed coefficient list is checked by exact
polynomial multiplication, not by numerical sampling.
\end{proof}

\begin{theorem}\label{thm:noorigin-oct7-eight-lowvariance}
For any eight complex points with $\sum_i|z_i|^2=8$, put
$s=\sum_i(|z_i|^2-1)^2$. Then
\[
 \log\Delta_8(z)\le8\log8+\frac{11}{7}s
\]
when the points are distinct. In particular,
\[
 s\le\frac{23}{50}\quad\Longrightarrow\quad
 \Delta_8(z)<C_8=\frac{8^{28}}{7^{21}}.
\]
This covers all angles and all harmonic moments in the stated radial
region; no force equation is assumed.
\end{theorem}
\begin{proof}
Hadamard's inequality on the rows of the Vandermonde matrix with
columns $1,z,\ldots,z^7$ gives $\Delta_8\le\prod_iS(|z_i|^2)$.
Sum the preceding lemma and use $\sum_i(|z_i|^2-1)=0$.
For the strict comparison, set $r=8^{20}/7^{21}>1$ and
$u=(r-1)/(r+1)$. The first sixteen terms of
\[
 \log r=2\sum_{j\ge0}\frac{u^{2j+1}}{2j+1}
\]
already exceed $181/250$, as verified by rational arithmetic.
But $(11/7)(23/50)=253/350<181/250$.
Thus the resulting upper bound is strictly below $\log C_8$.
Collisions have $\Delta_8=0$ and satisfy the final assertion directly.
\end{proof}

\begin{theorem}\label{thm:noorigin-oct7-eight-correlation}
For eight actual points define
\[
 X=\sum_{i<j}\frac{|\sum_{k=0}^7(z_i\overline z_j)^k|^2}
 {S(t_i)S(t_j)},\qquad t_i=|z_i|^2.
\]
Then
\[
 \log\Delta_8\le8\log8+\frac{11}{7}s-\frac7{16}X.
\]
Consequently $(11/7)s-(7/16)X\le181/250$ implies
$\Delta_8<C_8$. In particular, every collision-free local maximum
with $X\ge8$ has $\Delta_8<C_8$. A global maximizer necessarily
has
\[
 s>\frac{23}{50},\qquad X<8,\qquad
 X<\frac{176}{49}s-\frac{1448}{875}.
\]
\end{theorem}
\begin{proof}
Apply the dimension-eight version of the universal Gram lemma
\ref{lem:noclass-general-gram} to the row-normalized Vandermonde
Gram matrix. Its coefficient is $4(8-1)/8^2=7/16$.
Combine it with Lemma~\ref{lem:noorigin-oct7-eight-tangent} and
the strict logarithm comparison in the preceding proof.
At a local maximum $Bz^2=7(t-\mathbf1)$, $B\mathbf1=0$, and
$\|B\|_{\rm op}\le7/2$, so
\[
 49s\le\frac{49}{4}(8+s-|p_2|^2/8),\qquad s\le8/3.
\]
For $X\ge8$, the excess over $8\log8$ is thus at most
$(11/7)(8/3)-(7/16)8=29/42<181/250$.
A central heptagon supplies a competitor of value $C_8$, giving
the final necessary inequalities for a global maximizer.
\end{proof}

\subsection{A larger rigorously covered harmonic region}
\begin{theorem}\label{thm:noorigin-oct7-eight-improvedharmonic}
Every collision-free local maximum with total moment eight and
\[
                    \max_{2\le k\le6}|p_k|\le10^{-14}
\]
is a rotated central regular heptagon. In particular, no actual
local maximum without an origin point lies in this harmonic region.
\end{theorem}
\begin{proof}
We retain the Newton, force, and candidate-neighborhood estimates
proved in Theorem~\ref{thm:noclass-n8-small-harmonic}, and improve the
root-matching step that caused its very small threshold.
Write $\varepsilon=10^{-14}$. Those estimates produce either
$q(X)=X^8-\beta$ or $q(X)=X^8-\alpha X$, and at each actual root,
\[
                   |q(z_i)|\le10^6\varepsilon=10^{-8}.
\]
They also give $|z_i-z_j|\ge4/7$. Every $z_i$ is therefore
within $\delta_0=1/10$ of some model root. Assign a nearest model
root to every actual root. Two actual roots cannot share an
assignment, since $2\delta_0=1/5<4/7$; the assignment is a
bijection. This also excludes $\alpha=0$ or $\beta=0$, which
would leave all model roots at zero.

Two distinct model roots are separated by at least
$4/7-2\delta_0=13/35$. Thus, for an actual root and every model
root other than its assigned one, the distance is at least
$13/35-\delta_0=19/70$. Factoring $q$ now improves the assigned
root distance to
\[
 \delta\le10^6\varepsilon(70/19)^7
              <10^{10}\varepsilon=10^{-4}.
\]
All roots are matched before this division, so no unproved model
separation or root multiplicity assumption is used.
Normalizing the model moment to eight changes the total orbit
distance by less than $6\delta\le6\cdot10^{-4}$, by the same
moment comparison in the earlier theorem. For the regular octagon,
the spectral perturbation is at most $441\cdot6\cdot10^{-4}<1$,
which cannot repair its second-variation violation. For the central
heptagon the local chart has radius less than $10^{-3}$, where the
already proved derivative bound is strictly negative unless the
chart vector is zero. Only the exact central regular heptagon
remains. The rational root bootstrap and both final numerical
constant comparisons are replayed in the new checker.
\end{proof}

\subsection{A two-dimensional identity that eliminates the eighth holomorphic moment}
The full determinant identity in Lemma~\ref{lem:noclass-finite-projection}
uses $v_2,\ldots,v_5$ for $n=8$. Its highest bilinear entries involve
moments that do not occur with a useful sign in a short logarithm
majorant. We can instead remove a conjugate pair and retain only two
vectors. This is the same finite real-projection idea that underlies
the five-point argument.

\begin{proposition}\label{prop:noorigin-oct7-eight-twoD}
At a centered stationary eight-point configuration with $|p_2|<8$,
use $W,H,C$ from Lemma~\ref{lem:noclass-finite-projection}, and suppose
$s>0$, $A_0:=H_{22}>|K_0|:=|C_{22}|$. Put
\[
 D_0=\begin{pmatrix}A_0&\overline{K_0}\\K_0&A_0\end{pmatrix},
 \quad J=\begin{pmatrix}0&1\\1&0\end{pmatrix},\qquad
 d_k=\binom{H_{2k}}{C_{2k}}\quad(k=3,4),
\]
and
\[
 \widehat H_{jk}=H_{jk}-d_j^*D_0^{-1}d_k,
 \qquad
 \widehat C_{jk}=C_{jk}-d_j^TJD_0^{-1}d_k.
\]
Set $A=\widehat H_{33}$, $B=\widehat H_{44}$,
$G=\widehat H_{34}$, $K=\widehat C_{33}$,
$L=\widehat C_{34}$. Then
\begin{equation}\label{eq:noorigin-oct7-eight-twoD}
 |AL-GK|^2=(A^2-|K|^2)(AB-|G|^2).
\end{equation}
If $A>0$, then for every $w>1$,
\begin{equation}\label{eq:noorigin-oct7-eight-weighted}
 |L|^2+w|G|^2\ge AB-\frac{w}{w-1}\frac BA|K|^2.
\end{equation}
The quantities in these two equations depend only on
$p_2,\ldots,p_7$, $T_2,T_3,T_4$, $R_2,R_3$, and $M_{4,2}$.
They do not involve $p_8$.
\end{proposition}
\begin{proof}
$W$ is a conjugation-invariant complex space of dimension four.
The condition $A_0>|K_0|$ makes $v_2,\overline v_2$ independent.
Their span is conjugation-invariant, so its orthogonal complement
$U$ in $W$ is the complexification of a real two-dimensional space.
The displayed hatted matrices are precisely the Hermitian and
bilinear Gram matrices of $a=\Pi_Uv_3$, $b=\Pi_Uv_4$.
The three vectors $a,b,\overline a$ lie in $U$, so
\[
 \det\begin{pmatrix}
 A&G&\overline K\\\overline G&B&\overline L\\K&L&A
 \end{pmatrix}
 =A^2B-A(|G|^2+|L|^2)-B|K|^2
                 +2\Re(GK\overline L)=0.
\]
Multiplication by $-A$ and rearrangement proves
\eqref{eq:noorigin-oct7-eight-twoD}, including singular $a,b$.
For $A>0$ the same determinant equation gives
\[
 |L|^2+w|G|^2
 =AB-\frac BA|K|^2+\frac2A\Re(GK\overline L)
                                      +(w-1)|G|^2.
\]
Completing the square in $G$ bounds the last two terms below by
$-|K|^2|L|^2/[A^2(w-1)]$. The bilinear Cauchy inequality
$|L|^2\le AB$ proves \eqref{eq:noorigin-oct7-eight-weighted}.
Finally, $C_{34}$ involves $p_7$ and $C_{33}$ involves $p_6$;
$C_{44}$ is never used. All additional projection entries are
$H_{23},H_{24},C_{23},C_{24}$, involving $R_2,M_{4,2},p_5,p_6$.
This verifies the stated moment scope.
\end{proof}

For a useful interpretation, write $L=p_7-\ell_0$ and
$G=R_3-g_0$, with $\ell_0,g_0$ given by the explicit projection
formulas. The last inequality is a lower bound for the two quantities
that occur with negative signs in the seventh distance moment:
\[
              D_7=\text{other moment terms}
                              -|p_7|^2-1225|R_3|^2.
\]
The coefficient $1225=\binom73^2$ follows directly from expanding
$|z_i-z_j|^{14}$. With $w=1225$, the unwanted eighth holomorphic
moment is absent from the resulting interface.
The projection assumptions and the shifts $\ell_0,g_0$ still need
uniform control in the remaining global region; this proposition
alone is not a complete upper-bound certificate.

\begin{corollary}[A scalar eight-point value certificate]
\label{cor:noorigin-oct7-eight-valuecertificate}
Under the assumptions of Proposition~\ref{prop:noorigin-oct7-eight-twoD},
suppose also $A>0$, and define
\[
 \ell_0=p_7-L,\qquad
 c_7=\ell_0+\frac{GK}{A},\qquad
 \rho_7=\frac{\sqrt{(A^2-|K|^2)(AB-|G|^2)}}{A},\qquad
 L_7=\bigl|\rho_7-|c_7|\bigr|.
\]
The shift $\ell_0$ uses only moments below $p_7$: it is the
sum of the two displayed projection corrections to $C_{34}$.
Then
\[
 |p_7-c_7|=\rho_7,\qquad |p_7|\ge L_7,
 \qquad
 \Delta_8\le(8T_7-L_7^2)\prod_{k=1}^6T_k.
\]
Thus any actual stationary configuration for which the final
right-hand side is at most $C_8$ has the desired sharp upper bound,
without looking up its stationary component.
If the configuration is an actual local maximum and attains $C_8$
under this certificate, it must be the central regular heptagon.
\end{corollary}
\begin{proof}
Divide \eqref{eq:noorigin-oct7-eight-twoD} by $A^2$ to obtain
the exact circle constraint. Hermitian and bilinear Cauchy
inequalities make both factors under the square root nonnegative.
The reverse triangle inequality proves $|p_7|\ge L_7$.
For the original Vandermonde column Gram matrix, keep the
$\{1,z^7\}$ block together and apply Fischer to the remaining
columns. The block has determinant $8T_7-|p_7|^2$, yielding
the asserted bound. This proves a checkable sufficient certificate;
we have not proved that every remaining global maximizer passes it.
For equality with $C_8$, equality in Fischer forces all cross-block
entries to vanish, in particular $p_2=\cdots=p_6=0$.
Lemma~\ref{lem:noclass-zero-harmonics} leaves a regular octagon or
a central regular heptagon; the former has value $8^8<C_8$.
\end{proof}

The new unconditional exclusion and the enlarged rigidity theorem
give the following precise reduction. If an eight-point global
maximizer without an origin exists, then necessarily
\[
 \sum_i(|z_i|^2-1)^2>\frac{23}{50},\quad X<8,\qquad
                   \max_{2\le k\le6}|p_k|>10^{-14}.
\]
The remaining task is to close this simultaneous region, including
the degenerate projection boundary, without importing a classification.

\section{A coupled trace envelope using the prescribed quadratic column}
\label{sec:noorigin-oct8-trace}

The earlier principal-angle trace bound used only the saturated column
$Bz=\lambda\overline z$. Retaining the additional prescribed action
$Bz^2=2\lambda(t-\mathbf1)$ makes the trace constraint depend jointly
on the radial variance and the second and third holomorphic moments.
This does not use a stationary-shape classification.

\begin{theorem}\label{thm:noorigin-oct8-coupled-trace}
Let $n\ge5$ and let $z$ be a centered, collision-free local maximum
with $\sum_i|z_i|^2=n$ and $s=T_2-n>0$. Put
\[
 \lambda=\frac{n-1}{2},\quad a=|p_2|,\quad
 b=\frac{n-2}{n-1},\quad w=n+s-\frac{a^2}{n},
\]
\[
 \xi=\frac{a^2}{n^2},\quad
 \gamma=\frac{|p_3|^2}{nw},\quad
 \tau=\frac{4s}{w},\quad
 \zeta=\frac{b^2a^2}{sw}.
\]
Then $w>0$, $0\le\tau\le1$, and the two quantities
\[
 E=1+\xi+\gamma-\tau+\tau\zeta,\qquad
 F=(\xi+\gamma)(1-\tau)+\tau\xi\zeta
\]
are nonnegative. The following necessary inequality holds:
\begin{equation}\label{eq:noorigin-oct8-coupled-trace}
 \left(\lambda-\frac1n-\frac{2b}{w}\right)a
       \le n-5+\sqrt{E+2\sqrt F}.
\end{equation}
\end{theorem}
\begin{proof}
Rotate the points so that $p_2=a\ge0$ is real. Retain the notation
$B_0=(\lambda/n)\overline z\,\overline z^{\,T}$ and
$B_t=B-B_0$ from Corollary~\ref{cor:noorigin-principal-angle-trace}.
Let $D=\{\mathbf1,z\}^{\perp}$ and
$I=\{\mathbf1,\overline z\}^{\perp}$.
The operator $Y=B_t/\lambda$ maps $D$ into $I$, vanishes on
$D^\perp$, and has norm at most one. Moreover,
\[
                 \operatorname{tr}Y=(\lambda-1/n)a.
\]
The vectors
\[
 v=z^{\circ2}-\frac a n\mathbf1,\quad e=v/\sqrt w,\quad
 f=(t-\mathbf1)/\sqrt s
\]
satisfy $e\in D$, $f\in D\cap I$, and $\|e\|=\|f\|=1$.
The force identities give
\[
 Ye=\sqrt\tau f,\qquad
 e^*f=\frac{ba}{\sqrt{sw}}=\sqrt\zeta .
\]
In particular $\tau\le1$ follows from $\|Y\|\le1$.
The strict bound $a<n$ from the earlier trace theorem ensures $w>0$.

Let $H=D\cap e^\perp$, with orthogonal projector $P_H$, and write
$Y=\sqrt\tau f e^*+Z$, where $Z=YP_H$.
Since $YY^*\preceq P_I$ and the two domain columns are orthogonal,
\[
                ZZ^*\preceq C:=P_I-\tau ff^* .
\]
The matrix $C$ is positive semidefinite. This inequality implies
$Z=C^{1/2}U$ for a contraction $U$ with $U=UP_H$:
on the range of $C$, take $U=C^{+1/2}Z$, and the same inequality
gives $UU^*\preceq P_{\operatorname{ran}C}$. Trace duality yields
\[
 |\operatorname{tr}Z|
     \le\|P_H C^{1/2}\|_*.
\]
Thus it remains to compute the singular values in this nuclear norm.

On the $(n-3)$-dimensional space $H$, put
$q=P_H\overline z/\sqrt n$ and $f_H=P_Hf$. Direct projection gives
\[
 P_HCP_H=I_H-qq^*-\tau f_Hf_H^*,
\]
\[
 \|q\|^2=1-\xi-\gamma,\qquad
 \|f_H\|^2=1-\zeta,\qquad
 |q^*f_H|^2=\gamma\zeta .
\]
Here $|\langle\overline z/\sqrt n,e\rangle|^2=\gamma$;
$f$ is orthogonal to both $z$ and $\overline z$, which gives
the last formula after subtracting its component along $e$.
All but two eigenvalues of $P_HCP_H$ are one. Counting a repeated
one in a rank-deficient case, denote the remaining two by
$\eta_1,\eta_2\ge0$. The trace and determinant of the rank-two
update give exactly
\[
 \eta_1+\eta_2=E,\qquad \eta_1\eta_2=F .
\]
Consequently
\[
 \|P_H C^{1/2}\|_*
   =n-5+\sqrt{\eta_1}+\sqrt{\eta_2}
   =n-5+\sqrt{E+2\sqrt F}.
\]
Finally $\operatorname{tr}(\sqrt\tau f e^*)=
\sqrt{\tau\zeta}=2ba/w$. Subtract this real contribution from
the known trace of $Y$ and apply the preceding bound.
This proves the necessary inequality even if its left side is negative.
\end{proof}

\begin{corollary}[A finite radial band with moderate third moment]
\label{cor:noorigin-oct8-trace-band}
For $n=6,7,8$, every collision-free local maximum satisfying
\[
 \left|s-\frac{n}{n-1}\right|\le\frac1{20},
       \qquad |p_3|\le1
\]
obeys, respectively,
\[
              |p_2|<1,\qquad |p_2|<\frac65,\qquad
              |p_2|<\frac43 .
\]
These strengthen the earlier bounds $18/13,28/19,20/13$ on
this entire band. No assumption on $p_4$ or higher moments is made.
\end{corollary}
\begin{proof}
Let $s_\pm=n/(n-1)\pm1/20$ and
$q_*=\{2n(n-3)/(n(n-1)-4)\}^2$. The earlier universal bound gives
$a^2\le q_*$. Define the rational constants
\[
 w_*=n+s_- -q_*/n,\quad
 \xi_*=\frac{q_*}{n^2},\quad \gamma_*=\frac1{nw_*},
\]
\[
 \tau_-=\frac{4s_-}{n+s_+},\quad
 \tau_+=\frac{4s_+}{w_*},\quad
 \zeta_*=\frac{b^2q_*}{s_-w_*}.
\]
All denominators are positive. For the actual configuration,
\[
 E\le E_*:=1+\xi_*+\gamma_*-\tau_-+\tau_+\zeta_*,
\]
\[
 F\le F_*:=(\xi_*+\gamma_*)(1-\tau_-)
                                      +\tau_+\xi_*\zeta_* .
\]
Take $a_*=1,6/5,4/3$ in the respective cases. If $a\ge a_*$,
the left side of \eqref{eq:noorigin-oct8-coupled-trace} is at least
$L_*=(\lambda-1/n-2b/w_*)a_*$; the coefficient in parentheses
is positive. The following strict rational bounds are independently
verified in \path{noorigin_oct8_trace_check.py}:
\[
\begin{array}{c|ccc|c}
 n&E_*<&F_*<&L_*>&c\\ \hline
 6&14/25&7/200&209/100&1\\
 7&63/100&17/500&317/100&1\\
 8&17/25&31/1000&106/25&11/10
\end{array}
\]
For each row, its displayed upper bounds satisfy
$4F_*<(c^2-E_*)^2$ and $E_*<c^2$. Thus
$\sqrt{E_*+2\sqrt{F_*}}<c$, while $L_*>n-5+c$.
This contradicts the necessary trace inequality.
\end{proof}

The new band is a coupled moment exclusion, not a proof that the
entire remaining global region lies in the band or has $|p_3|\le1$.
It therefore does not close the unrestricted six-, seven-, or
eight-point problem.

% Independent fragment; ordinary amsmath/amssymb/amsthm preamble suffices.
\section{Six-point radial improvement and exact trace closure}
\label{sec:noorigin-oct8-six}

\subsection{A quartic radial majorant and a larger all-angle region}

\begin{lemma}\label{lem:oct8-centered-radial-SOS}
For any $n\ge3$ real numbers $u_i$ with $\sum_i u_i=0$, set
$s=\sum_i u_i^2$ and $m_j=\sum_i u_i^j$. Then
\begin{align*}
 &\sum_{i<j}\sum_{k\ne i,j}(u_i-u_j)^2
                    (u_i+u_j+(n-2)u_k)^2\\
 &\hspace{10mm}=(n-2)
          \bigl[(n^2-3n+3)s^2-n(n-1)m_4\bigr].
\end{align*}
In particular
\[
 m_4\le\frac{n^2-3n+3}{n(n-1)}s^2,\qquad
 |m_3|^2\le\frac{(n-2)^2}{n(n-1)}s^3.
\]
\end{lemma}
\begin{proof}
For a fixed pair $i,j$, summing over $k$ uses
$\sum_{k\ne i,j}u_k=-u_i-u_j$ and
$\sum_{k\ne i,j}u_k^2=s-u_i^2-u_j^2$. The remaining pair sums are
\[
 \sum_{i<j}(u_i-u_j)^2(u_i+u_j)^2=nm_4-s^2,
\]
\[
 \sum_{i<j}(u_i-u_j)^2(s-u_i^2-u_j^2)=(n-1)s^2-nm_4.
\]
Substitution gives the identity and hence the fourth-moment bound.
For $s>0$, projecting $u^2$ off $\mathbf1$ and $u$ gives
$m_4-s^2/n-m_3^2/s\ge0$. Combining the two bounds proves the
cubic inequality. The case $s=0$ is immediate.
\end{proof}

\begin{lemma}\label{lem:oct8-six-quartic-log}
For every $t\ge0$, with $u=t-1$,
\[
 \log(1+t+\cdots+t^5)
 \le\log6+\frac52u+\frac5{24}u^2-\frac58u^3+\frac3{10}u^4,
\]
with equality exactly at $t=1$.
\end{lemma}
\begin{proof}
For the difference $g$ of the right and left sides, differentiation
and polynomial division give
\[
 g'(t)=\frac{(t-1)^3Q(t)}{1+t+\cdots+t^5},\qquad
 Q(t)=\frac{144t^5-81t^4-256t^3-81t^2+144t+239}{120}.
\]
The checker verifies that every degree-five Bernstein coefficient of
$Q$ is positive on each of the sixteen intervals
$[j/8,(j+1)/8]$, $0\le j\le15$. On the remaining half-line,
\[
 Q(2+x)=\frac65x^5+\frac{453}{40}x^4+\frac{607}{15}x^3
            +\frac{2653}{40}x^2+\frac{473}{10}x+\frac{489}{40}>0
                  \quad(x\ge0).
\]
Thus $g$ strictly decreases to $g(1)=0$ and then strictly increases.
This proves the global majorant, without an assumption on the radial
population or on $t$ being close to one.
\end{proof}

\begin{theorem}\label{thm:oct8-six-radial-region}
Every six-point configuration with $\sum_i|z_i|^2=6$ satisfies
\[
 \log\Delta_6\le6\log6+\frac5{24}s
             +\frac58\sqrt{\frac8{15}s^3}+\frac{21}{100}s^2
                       -\frac59X,
\]
where $X$ is the complete normalized row-correlation sum from
Theorem~\ref{thm:noorigin-six-whole-regions}. In particular one may
drop the nonpositive last term.
Consequently the entire all-angle region $s\le87/1000$ has value
strictly below $C_6$. Every global maximum satisfies
\[
                             s>\frac{87}{1000}.
\]
More generally the whole region
\[
 \frac5{24}s+\frac58\sqrt{\frac8{15}s^3}
                        +\frac{21}{100}s^2-\frac59X
                         \le\frac{629}{20000}
\]
has value strictly below $C_6$.
\end{theorem}
\begin{proof}
Hadamard's inequality for the square Vandermonde matrix gives
$\Delta_6\le\prod_i(1+t_i+\cdots+t_i^5)$. Set $u_i=t_i-1$;
then $\sum_i u_i=0$. Sum the quartic logarithm bound and apply
Lemma~\ref{lem:oct8-centered-radial-SOS} with $n=6$.
Retaining the normalized Gram determinant, the already established
six-row inequality adds the term $-5X/9$. Collisions give $\Delta_6=0$
and are immediate.
The resulting scalar upper bound is increasing in $s\ge0$.
At $s_*=87/1000$, exact rational comparison gives
\[
 \frac5{24}s_*+\frac58\sqrt{\frac8{15}s_*^3}
                    +\frac{21}{100}s_*^2<\frac{629}{20000}
                       <\log\frac{6^9}{5^{10}}.
\]
For the first comparison, subtract the two rational terms from
$629/20000$, verify positivity, and square; the required remaining
rational square is strictly larger than $(5/24)s_*^3$.
For the second, the positive two-term atanh series at
$r=(6^9-5^{10})/(6^9+5^{10})$ gives
$\log(6^9/5^{10})>2r+2r^3/3>629/20000$.
Both checks are performed with exact fractions. The central pentagon
attains $C_6$, proving the final necessary condition by contraposition.
\end{proof}

\subsection{Trace closure of the inverse-square operator}

The Laurent majorant reduces the six-point question to a joint estimate
of $D_{-1}$ and $D_{-2}$. The previous bound for $D_{-2}$ used
$\|B\|_F^2\le125/4$, whereas its value at the central pentagon is
$175/12$. The following finite-column identity removes this particular
loss. It uses no list of stationary configurations, and its all-row form
has no projection-degeneracy exception.

Let $z_1,\ldots,z_6$ be distinct centered points, put
$P(Z)=\prod_i(Z-z_i)$, and write
\[
 X=[\mathbf1,z,z^2,z^3],\qquad G=X^*X.
\]
Here and below powers of vectors are componentwise. The matrix $G$ is
positive definite. At a stationary configuration with $T_1=6$ and
$\lambda=5/2$, the four columns of $Y=BX$ are known:
\begin{equation}\label{eq:oct8-six-B-columns}
 Y=[0,\lambda\overline z,5(t-\mathbf1),
                                    (15t/2-9)z].
\end{equation}
Set $R=X^TY$ and define the symmetric minimum-Frobenius completion
\begin{equation}\label{eq:oct8-six-B-min}
 B_0=YG^{-1}X^*+\overline XG^{-T}Y^T
               -\overline XG^{-T}RG^{-1}X^*.
\end{equation}
In particular $B_0X=Y$ and $B_0^T=B_0$.
The three matrices in its squared norm have the explicit low-moment
interface, with $T=T_2$, $V=T_3$ and $s=T-6$,
\[
 G=\begin{pmatrix}
 6&0&p_2&p_3\\0&6&0&9p_2/5\\
 \overline{p_2}&0&T&R_2\\
 \overline{p_3}&9\overline{p_2}/5&\overline{R_2}&V
 \end{pmatrix},
\]
\[
 R=\begin{pmatrix}
 0&0&0&0\\0&15&0&9p_2/2\\
 0&0&4p_2&3p_3\\0&9p_2/2&3p_3&3(p_4+p_2^2)/2
 \end{pmatrix},
\]
\[
 Y^*Y=\begin{pmatrix}
 0&0&0&0\\0&75/2&0&45p_2/4\\
 0&0&25s&75R_2/2\\
 0&45\overline{p_2}/4&75\overline{R_2}/2&225V/4-135T+486
 \end{pmatrix}.
\]
These formulas retain the off-diagonal moment corrections.

Choose the two-column complementary matrix
\[
 U_{i0}=\frac1{P'(z_i)},\qquad
 U_{i1}=\frac{z_i}{P'(z_i)},\qquad H=U^*U.
\]
It has full column rank and $X^TU=0$. For $k=0,\ldots,5$ put
\[
 E_{k\cdot}=
 \left(\sum_i z_i^kU_{i0}^2,
       2\sum_i z_i^kU_{i0}U_{i1},
       \sum_i z_i^kU_{i1}^2\right),
 \qquad
 d_k=\ell_k-\sum_i z_i^k(B_0)_{ii},
\]
where, writing $M_{ab}=\sum_i z_i^a\overline z_i^{\,b}$,
\begin{equation}\label{eq:oct8-six-trace-data}
 (\ell_0,\ldots,\ell_5)=
 \left(\frac{25}4\overline{p_2},0,
       \frac{25}4T_2-40,\frac{25}4R_2,
       \frac{25}4M_{42}-16p_2,
       \frac{25}4M_{52}-\frac{38}3p_3\right).
\end{equation}
Equivalently, any full-rank basis of $\ker X^T$ may replace $U$;
the completion and its Frobenius norm are independent of that choice.
The barycentric formula is used to prove that the trace map is
injective, rather than to assume a lower bound on a selected trace minor.
For actual computation, the entirely polynomial-column construction
\[
 U=\overline{(I-XG^{-1}X^*)[z^4,z^5]}
\]
is convenient and uses only finite moment Gram matrices. Its columns
are independent because the full square Vandermonde matrix is invertible.
No reciprocal pair distance or reciprocal derivative $P'(z_i)$ is
needed to compute this version of the completion.

\begin{theorem}\label{thm:oct8-six-trace-completion}
For every distinct stationary six-point configuration above, $E$ has
column rank three. Define
\[
 c=(E^*E)^{-1}E^*d,
 \qquad C=\begin{pmatrix}c_0&c_1\\c_1&c_2\end{pmatrix}.
\]
Then the exact operator identity and exact inverse-moment identity are
\begin{align}
 B&=B_0+UCU^T,\label{eq:oct8-six-B-exact}\\
 2D_{-2}
 &=2\operatorname{tr}(Y^*YG^{-1})
       -\operatorname{tr}(G^{-1}R^*G^{-T}R)
       +\operatorname{tr}(C^*HCH^T)
       -\sum_i|(B_0+UCU^T)_{ii}|^2.
       \label{eq:oct8-six-inverse-exact}
\end{align}
Thus $D_{-2}$ is determined by four polynomial columns and six weighted
diagonal traces. No operator-norm upper bound has been substituted into
\eqref{eq:oct8-six-inverse-exact}.
\end{theorem}
\begin{proof}
Since $B$ is symmetric, $R=X^TY$ is symmetric. Direct multiplication
proves the assertions about $B_0$. Every symmetric matrix annihilating
$X$ is of the form $UCU^T$, with $C$ symmetric: the columns lie in
$\ker X^T$, and symmetry supplies the same support on the other side.
The standard Lagrange coefficient identity
$\sum_i z_i^j/P'(z_i)=0$ for $0\le j\le4$ proves $X^TU=0$;
the two columns of $U$ are independent since the $z_i$ are distinct.

For $q(Z)=c_0+2c_1Z+c_2Z^2$, the equation $Ec=0$ says
\[
               \sum_i z_i^k\frac{q(z_i)}{P'(z_i)^2}=0
                         \quad(0\le k\le5).
\]
The square Vandermonde matrix is invertible, so $q(z_i)=0$ for all
six points. A polynomial of degree at most two with six distinct zeros
is identically zero. Hence $E$ has full column rank, including cases
where any particular three selected rows would be singular.

To verify the trace data, put $F_i=\sum_{j\ne i}(z_i-z_j)^{-1}$
and $b_i=B_{ii}$. Squaring $F_i$ and collecting each unordered triple
gives, for $k\ge2$,
\[
 \sum_i z_i^k b_i
   =\sum_i z_i^k F_i^2
     -2\sum_{\{a,b,c\}}h_{k-2}(z_a,z_b,z_c),
\]
where $h_j$ is the complete homogeneous polynomial of degree $j$ in
three variables. For $k=0,1$ the triple contribution is zero.
For six centered points the triple sums for degrees $0,1,2,3$ are
$20,0,8p_2,19p_3/3$, respectively. Substitution of
$F_i=(5/2)\overline z_i$, together with $R_1=0$, gives
\eqref{eq:oct8-six-trace-data}. The actual remainder therefore obeys
$Ec=d$, whose unique solution is the stated least-squares formula.

Finally $B_0$ is Frobenius-orthogonal to every $UCU^T$.
Orthonormalizing $X$ and using unitary congruence gives
\[
 \|B_0\|_F^2=2\operatorname{tr}(Y^*YG^{-1})
             -\operatorname{tr}(G^{-1}R^*G^{-T}R).
\]
Also $\|UCU^T\|_F^2=\operatorname{tr}(C^*HCH^T)$.
Since $\|B\|_F^2=2D_{-2}+\sum_i|b_i|^2$, these identities
prove \eqref{eq:oct8-six-inverse-exact}.
\end{proof}

At an actual local maximum the earlier diagonal projection bound may
be applied to the last term, giving the directly usable upper bound
\begin{align*}
 D_{-2}\le{}&\operatorname{tr}(Y^*YG^{-1})
 -\frac12\operatorname{tr}(G^{-1}R^*G^{-T}R)
 +\frac12\operatorname{tr}(C^*HCH^T)\\
 &-\frac{(25T/4-40)^2}{2T}
              -\frac{800|p_2|^2}{T(6T-|p_2|^2)}.
\end{align*}
Unlike the earlier scalar norm estimate, this bound is sharp at the
central pentagon. The positive complementary-block term is explicitly
determined by the weighted trace data and must be retained.

\begin{corollary}[The same two-dimensional closure for general $n$]
\label{cor:oct8-general-trace-completion}
For $n\ge3$ distinct centered stationary points normalized by
$\sum_i|z_i|^2=n$, put $\lambda=(n-1)/2$ and use
\[
 X=[\mathbf1,z,\ldots,z^{n-3}],\qquad
 U_{i\cdot}=\frac{(1,z_i)}{P'(z_i)}.
\]
There are always only three complex coefficients in the symmetric
complementary block. Define $Y=BX$, $G=X^*X$, $R=X^TY$, $B_0$ by
\eqref{eq:oct8-six-B-min}, and $E$ with the same three columns as
above and with rows $k=0,\ldots,n-1$. Then $E$ has full column rank
and the formulas \eqref{eq:oct8-six-B-exact} and
\eqref{eq:oct8-six-inverse-exact} still hold. The required polynomial
columns and traces are given by
\[
 (Bz^k)_i=k\lambda t_i z_i^{k-2}
 -\sum_{r=0}^{k-2}(k-1-r)p_r z_i^{k-2-r}
                       +\frac{k(k-1)}2z_i^{k-2}\quad(k\ge2),
\]
\[
 \ell_k=\lambda^2 M_{k,2}
 -2\sum_{\{a,b,c\}}h_{k-2}(z_a,z_b,z_c),
\]
with the triple term zero for $k=0,1$ and $p_0=n$.
Thus the inverse-fourth-moment problem can always be converted to a
finite polynomial-column and weighted-trace problem. The identity
alone is not an upper-value proof for general $n$.
Equivalently one can compute the complementary basis as
\[
 U=\overline{(I-XG^{-1}X^*)[z^{n-2},z^{n-1}]},
\]
whose independence again follows from the square Vandermonde matrix.
\end{corollary}
\begin{proof}
The barycentric identity now gives $X^TU=0$ because the largest
power summed is $n-2$. The same Vandermonde argument says that a
quadratic vanishing at all $n\ge3$ distinct points is zero.
The column formula follows by polynomial division of
$(z_i^k-z_j^k)/(z_i-z_j)^2$ and summation over $j$; the trace
formula is the triple divided-difference identity used above.
All remaining completion arguments are unchanged.
\end{proof}

For the origin plus a regular $(n-1)$-gon with $n=6,7,8$, this
general-$n$ completion has $C=0$. Indeed its two missing domain
directions are $u$ and $\overline z$. Their outputs have no projection
onto the two complementary image directions $u,z$, as follows from
the finite geometric-series sums; the mixed blocks are already forced
by the polynomial columns and symmetry. The diagonal formula at a
ring point is
\[
 B_{ii}=\lambda^2\overline z_i^{\,2}
                -\frac{P'''(z_i)}{3P'(z_i)}
        =\frac{n(8-n)}{12z_i^2},
\]
and at the origin it is zero. In particular at the central heptagon
all eight diagonal entries vanish. The exact inverse-fourth moments
for the three candidate configurations are
\[
            D_{-2}=\frac{50}9,\qquad
                    \frac{1851}{196},\qquad\frac{1029}{64},
\]
respectively. These statements supply candidate-sharp inverse-moment
interfaces for seven and eight as well, without stationary-shape
enumeration.
For completeness, the corresponding norm calculation in the orthogonal
Fourier and radial basis, with $m=n-1$ and $\rho=n/m$, is
\[
 \|B\|_F^2=\frac{m^2}4+\frac{2m^2}n
  +\frac1{\rho^2}\sum_{k=3}^{m-1}
             \left[\frac{(k-2)(k-m)+2}{2}\right]^2,
 \qquad
 \sum_i|B_{ii}|^2=\frac{m^3(8-n)^2}{144}.
\]
The first two terms are the saturated $z$ block and the two radial
blocks. The remaining coefficients follow by applying the displayed
general column formula to each Fourier mode. Subtracting these two
squared norms and dividing by two gives the stated inverse moments.

\paragraph{Sharpness at the central pentagon.}
At the origin plus a pentagon of squared radius $6/5$, one has
$T_2=36/5$, $s=6/5$ and $p_2=p_3=p_4=0$. In the orthogonal Fourier
basis the nonzero actions are
\[
 Bz=(5/2)\overline z,\qquad Bz^2=5u,
                         \qquad Bu=(5/6)\overline z^2.
\]
The minimum completion from the four columns is already the full
operator; consequently its complementary block $C$ is zero. The
identity gives
\[
 \|B\|_F^2=\frac{25}4+2\frac{25s}{T_2}
           =\frac{175}{12},\qquad
 \sum_i|b_i|^2=\frac{125}{36},\qquad D_{-2}=\frac{50}9.
\]
Together with the earlier sharp lower bound $D_{-1}\ge25/3$, this
restores the exact inverse-moment pair at the candidate. The Laurent
interpolant itself is exact there. It is important to retain $T_3$
rather than replacing it by the loose upper bound from $Bz^3$:
$T_3=216/25$ at this candidate, while that upper bound is $243/25$.

\paragraph{Remaining inequality.}
The new identity removes the unexplained inverse-fourth-moment
slack, but it does not prove that the Laurent upper value is at most
$\log C_6$ for every remaining stationary configuration. One still
needs a sign certificate combining this trace closure, the
inverse-second-moment completion, and the finite moment constraints.
All these formulas apply without a six-point shape classification;
the complete non-origin global bound is still open in this method.

For clarity, if the right side of the new $D_{-2}$ bound is denoted
$\mathcal B$, the three-column skew bound already gives the valid
sharp-at-candidate objective interface
\begin{align*}
 \log\Delta_6\le{}&15a_0+36a_1
  +a_2(6T+72+q)
  +a_3(6V+54T+54q/5-|p_3|^2)\\
 &+a_{-2}\mathcal B
  +a_{-1}\left[\frac{25T}{24}
    +\frac{75V/2-120T+576-q/4}{6T-q}\right].
\end{align*}
Here $a_j$ are the Laurent coefficients previously defined and
$q=|p_2|^2$. Since $a_{-2}>0$ and $a_{-1}<0$, both substitutions
have the correct upper-bound direction. No upper estimate on $V$
was used. Showing that this finite expression is at most $\log C_6$
over all remaining actual stationary moment data is the precise
remaining sign problem for this route.

% New exact inverse-distance interfaces; no stationary-shape classification.
\section{A sharp inverse-distance completion and a seven-point Laurent bound}
\label{sec:noorigin-oct8-seven}

\subsection{Four columns recover the missing inverse-distance norm}

Put $p_j=\sum_i z_i^j$, $t_i=|z_i|^2$, $T_j=\sum_i t_i^j$,
$R_j=\sum_i t_i^jz_i$, and
$M_{42}=\sum_i z_i^4\overline z_i^{\,2}=\sum_i t_i^2z_i^2$.
The following lemma applies to any centered, collision-free stationary
configuration normalized by $\sum_i|z_i|^2=n$; local maximality is
not required.

\begin{lemma}[The four-column inverse-distance completion]
\label{lem:noorigin-oct8-skew-four}
Let $n\ge4$, $\lambda=(n-1)/2$, $q=|p_2|^2$, and
$D_{-1}=\sum_{i<j}|z_i-z_j|^{-2}$.
Set $c=(2n-3)/(n-1)$ and form the positive definite Hermitian matrix
\[
 G=\begin{pmatrix}
 n&0&p_2&p_3\\
 0&n&0&cp_2\\
 \overline{p_2}&0&T_2&R_2\\
 \overline{p_3}&c\overline{p_2}&\overline{R_2}&T_3
 \end{pmatrix}.
\]
Let $H$ be the Hermitian matrix whose diagonal and upper triangular
entries are
\begin{gather*}
 H_{00}=\lambda^2n,\qquad H_{11}=\lambda^2T_2,\\
 H_{22}=\lambda^2T_3-2\lambda(n-2)T_2+(n-2)^2n,\\
 H_{33}=\lambda^2T_4-2\lambda(n-3)T_3+(n-3)^2T_2+(n-3)q,\\
 H_{01}=0,\qquad H_{02}=\lambda p_2/2,\qquad H_{03}=\lambda p_3,\\
 H_{12}=\lambda^2R_2,\qquad
 H_{13}=\lambda^2M_{42}-(n-2)(n-3)p_2/2,\\
 H_{23}=\lambda^2R_3-\lambda(2n-5)R_2.
\end{gather*}
Finally let
\[
 R=\begin{pmatrix}
 0&-\lambda n&0&-(2n-3)p_2/2\\
 \lambda n&0&p_2/2&p_3\\
 0&-p_2/2&0&(p_4-p_2^2)/2\\
 (2n-3)p_2/2&-p_3&-(p_4-p_2^2)/2&0
 \end{pmatrix}.
\]
Then the real-valued finite-moment expression
\begin{equation}\label{eq:noorigin-oct8-skew-four}
 \mathcal L_4:=\operatorname{tr}(HG^{-1})
 -\frac12\operatorname{tr}(G^{-1}R^*G^{-\mathsf T}R)
 \quad\hbox{satisfies}\quad D_{-1}\ge\mathcal L_4.
\end{equation}
It contains no holomorphic moment above $p_4$.
At the origin together with a regular $(n-1)$-gon, equality holds
for $n=6,7,8$, and the exact values are respectively
\[
                  D_{-1}=\frac{25}{3},\quad
                  \frac{177}{14},\quad\frac{147}{8}.
\]
Equality also holds at the unit regular hexagon and heptagon.
\end{lemma}
\begin{proof}
Define the complex skew-symmetric matrix
$A_{ij}=1/(z_i-z_j)$ for $i\ne j$ and $A_{ii}=0$.
Its squared Frobenius norm is $2D_{-1}$.
For $X=[\mathbf1,z,z^2,z^3]$, the force equations and polynomial
divided differences give $Y=AX$ with columns
\[
 Y=\left[\lambda\overline z,\quad
          \lambda(t-2\mathbf1),\quad
          (\lambda t-(n-2)\mathbf1)z,\quad
          (\lambda t-(n-3)\mathbf1)z^2-p_2\mathbf1\right].
\]
Here multiplication of vectors is coordinatewise. The identities
$X^*X=G$, $Y^*Y=H$, and $X^{\mathsf T}Y=R$ follow from
the paired-force mixed moments. In particular
\[
 \sum_i t_i z_i^2=cp_2,\quad
 \sum_i t_i z_i^3=\frac{2n-4}{n-1}p_3,\quad
 \sum_i t_i z_i^4=\frac{(2n-5)p_4+p_2^2}{n-1}.
\]
The diagonal identity $R_{33}=0$ uses
$\sum_i t_i z_i^5=((2n-6)p_5+2p_2p_3)/(n-1)$;
the $p_5$ terms cancel completely. Four distinct nodes make the
four columns of $X$ linearly independent, so $G\succ0$ even on
all earlier projection-degenerate branches.

Orthonormalize $X$ and extend its columns to a unitary matrix.
Under unitary congruence, the skew-symmetric matrix $A$ has block form
\[
                 \begin{pmatrix}K&-L^{\mathsf T}\\L&M\end{pmatrix}.
\]
The first four columns are fixed by $Y$. Hence
\[
 \|A\|_F^2=\|K\|_F^2+2\|L\|_F^2+\|M\|_F^2
 \ge2\operatorname{tr}(HG^{-1})
       -\operatorname{tr}(G^{-1}R^*G^{-\mathsf T}R).
\]
This proves the bound and also shows that equality is equivalent
to the remaining skew block being zero.

For the central regular $m$-gon, $m=n-1$, put $\rho=n/m$.
The columns $\mathbf1,z,z^2,z^3$ are mutually orthogonal and have
squared norms $n,n,m\rho^2,m\rho^3$.
Their orthogonal complement consists of one real radial direction
and the outer Fourier modes $4,\ldots,m-1$.
Rotation by $2\pi/m$ shows that a transpose matrix element of $A$
between modes $a,b$ can be nonzero only if
$a+b\equiv1\pmod m$. The radial direction counts as mode zero.
For $m=5,6,7$ no two distinct remaining modes satisfy this condition;
skew symmetry makes the self-pairings zero. The remaining block is
therefore zero. Substituting $T_j=m\rho^j$ in the displayed matrices
gives
\[
 \mathcal L_4=\frac{m^2(m^2+23)}{24n},
\]
which yields the three stated values.
For a unit regular $n$-gon with $n=6,7$, the complement is only
the Fourier modes $4,\ldots,n-1$, and the same test again gives
a zero remaining block. Its exact value is $n(n^2-1)/24$.
\end{proof}

\paragraph{The exact improvement over three columns.}
This new bound is never weaker than the earlier three-column skew
completion. More precisely, let $E=\operatorname{span}\{\mathbf1,z,z^2\}$,
let $\xi=z^3-\Pi_Ez^3$, and put $h=\|\xi\|^2>0$.
If $\mathcal L_3$ denotes the same completion expression for the
first three columns, then
\begin{equation}\label{eq:noorigin-oct8-skew-increment}
 \mathcal L_4-\mathcal L_3
       =\frac{\|\Pi_{\overline E^{\perp}}A\xi\|^2}{h}\ge0.
\end{equation}
Indeed, after the first three columns have been fixed, the remaining
unknown block is skew-symmetric. Fixing its column along
$\xi/\sqrt h$ fixes both that column and its negative transpose row;
its diagonal element is zero. Its minimum additional squared
Frobenius norm is twice the squared norm of that column. Dividing
by two to obtain $D_{-1}$ gives the identity.
At the central pentagon, hexagon, and heptagon the increments are
exactly $0$, $3/14$, and $7/8$, respectively. Thus the old loss at
the seven- and eight-point candidates is now removed.

\subsection{A global Laurent majorant matched to the central hexagon}

\begin{lemma}[The seven-point Laurent interface]
\label{lem:noorigin-oct8-seven-laurent}
Let $\rho=7/6$ and let $P$ be the degree-at-most-five polynomial
interpolating $f(x)=x^2\log x$ and its first derivative at
$\rho,3\rho,4\rho$. Write
\[
       \frac{P(x)}{x^2}=\sum_{j=-2}^{3}a_jx^j.
\]
For every $x>0$,
\[
                   \log x\le\sum_{j=-2}^{3}a_jx^j,
\]
with equality exactly at the three interpolation nodes.
The coefficients satisfy
\[
              a_{-2}>0,\quad a_{-1}<0,\quad a_2<0,\quad a_3>0.
\]
For every centered seven-point stationary configuration normalized
to $\sum_i|z_i|^2=7$,
\begin{align}\label{eq:noorigin-oct8-seven-laurent}
 \log\Delta_7\le{}&21a_0+49a_1
       +a_2(7T_2+98+|p_2|^2)\nonumber\\
       &+a_3(7T_3+63T_2+11|p_2|^2-|p_3|^2)
         +a_{-2}D_{-2}+a_{-1}D_{-1},
\end{align}
where $D_{-2}=\sum_{i<j}|z_i-z_j|^{-4}$.
The right side is exactly $\log C_7$ at the central hexagon.
At the unit regular heptagon, it is less than $\log C_7-6/25$.
\end{lemma}
\begin{proof}
The Hermite remainder has the sign of
$f^{(6)}(x)=-12/x^4$, since the node factor is
$(x-\rho)^2(x-3\rho)^2(x-4\rho)^2$. This proves the global
majorant and its strictness away from the nodes.
Summing it over all $21$ squared distances and using the exact
stationary distance moments
\[
 D_1=49,\quad D_2=7T_2+98+|p_2|^2,\quad
 D_3=7T_3+63T_2+11|p_2|^2-|p_3|^2
\]
proves the interface. Every central-hexagon squared distance is one
of the interpolation nodes, which proves sharpness.

For reproducibility, the exact Hermite system has rational coefficient
matrix and right sides involving only logarithms of rational numbers.
The accompanying rational outward-interval checker proves all four
coefficient signs. At the unit heptagon the distance moments are
\[
        (D_{-2},D_{-1},D_0,D_1,D_2,D_3)
                       =(14,14,21,49,147,490).
\]
The inverse moments follow from the root-of-unity partial-fraction
identities; the positive moments follow by Fourier summation.
The same checker verifies, with strict rational margins,
\[
 14a_{-2}+14a_{-1}+21a_0+49a_1+147a_2+490a_3
              -21\log7+15\log6<-6/25.
\]
No approximate regular-heptagon coordinates are needed.
\end{proof}

At a seven-point local maximum the $Bz^3$ operator bound gives
\[
 T_3\le\frac{539}{72}+\frac{11}{4}s-\frac{|p_3|^2}{56},
                    \qquad s=T_2-7.
\]
Consequently \eqref{eq:noorigin-oct8-seven-laurent} also implies
\begin{align*}
 \log\Delta_7-\log C_7\le{}&c_0+c_s s+c_q|p_2|^2
       -\frac98a_3|p_3|^2+a_{-2}D_{-2}+a_{-1}\mathcal L_4,\\
 c_0={}&21a_0+49a_1+147a_2+\frac{35525}{72}a_3-\log C_7,\\
 c_s={}&7a_2+\frac{329}{4}a_3<0,\qquad
 c_q=a_2+11a_3<0.
\end{align*}
The force equations and the contraction bound for $B$ prove the
$T_3$ estimate exactly as in the six-point cubic-column argument.
The strict signs of $c_s,c_q$ are included in the rational checker.

The four-column inverse bound applies across the degenerate conic
branches of the earlier projection argument, because its Gram matrix
is positive definite whenever the points are distinct. It removes
one specific candidate-level slack and the Laurent bound already
excludes the actual regular heptagon with a substantial exact gap.
The remaining positive $a_{-2}D_{-2}$ term still requires a coupled
upper estimate. These lemmas do not assert continuous coverage of
all remaining local maxima or prove the full unrestricted seven-point
maximum.

\subsection{A larger seven-point radial band with completely unrestricted angles}
\label{sec:noorigin-oct8-seven-radial}

The earlier whole-angle seven-point exclusion $s\le8/25$ can be
enlarged using the variance itself to restrict the admissible radii.
This improvement applies to arbitrary configurations, without
stationarity or any restriction on low holomorphic moments.

\begin{lemma}[A sharper radial tangent on the admissible half-line]
\label{lem:noorigin-oct8-seven-restricted-tangent}
For $S(t)=1+t+\cdots+t^6$ and every $t\ge7/16$,
\begin{equation}\label{eq:noorigin-oct8-seven-restricted-tangent}
 \log S(t)\le\log7+3(t-1)+\frac{127}{128}(t-1)^2.
\end{equation}
Equality holds exactly at $t=1$.
\end{lemma}
\begin{proof}
Let $g(t)$ be the right side minus $\log S(t)$. Multiplication gives
\[
 g'(t)=\frac{(t-1)Q(t)}{S(t)},\qquad
 Q(t)=\frac{-1-65t-65t^2-t^3+127t^4+319t^5+127t^6}{64}.
\]
For $t>0$ the derivative of $Q(t)/t^4$ is strictly positive:
\[
 64\left(\frac{Q(t)}{t^4}\right)'
     =319+254t+\frac4{t^5}+\frac{195}{t^4}
                         +\frac{130}{t^3}+\frac1{t^2}>0.
\]
Since $Q(0)<0$ and $Q(1)>0$, $Q$ has exactly one positive root
$\beta\in(0,1)$. Hence $g$ increases before $\beta$, decreases
from $\beta$ to $1$, and increases after $1$. Its minimum on
$[7/16,\infty)$ is therefore attained at $7/16$ or at $1$.
The latter value is zero. At the former endpoint,
\[
 g(7/16)=\log\frac{117440512}{29734657}-\frac{45009}{32768}.
\]
Put $r=117440512/29734657$ and $v=(r-1)/(r+1)\in(0,1)$.
The positive atanh series and a finite exact rational comparison give
\[
 \log r\ge2\sum_{k=0}^{19}\frac{v^{2k+1}}{2k+1}
       >\frac{45009}{32768}+\frac1{20000}.
\]
Thus the endpoint is strictly positive, and the stated inequality
and equality condition follow. The rational checker reproduces
the derivative identity and this strict endpoint comparison.
\end{proof}

\begin{theorem}[The enlarged whole-angle variance region]
\label{thm:noorigin-oct8-seven-radial-band}
Every seven-point configuration with $\sum_i|z_i|^2=7$ and
$s=\sum_i(|z_i|^2-1)^2$ satisfies
\[
             s\le\frac{189}{512}
           \quad\Longrightarrow\quad \Delta_7<C_7.
\]
In particular every global maximizer necessarily has
$s>189/512$. The conclusion imposes no restriction on angles,
on the presence of the origin, or on any holomorphic moment.
\end{theorem}
\begin{proof}
Put $u_i=|z_i|^2-1$, so $\sum_i u_i=0$.
For each $i$, Cauchy on the other six entries gives
\[
 s=u_i^2+\sum_{j\ne i}u_j^2\ge\frac76u_i^2.
\]
Thus $s\le189/512$ implies
$|u_i|\le\sqrt{6s/7}\le9/16$, and all squared radii are at
least $7/16$. Apply row Hadamard to the seven-column Vandermonde
matrix and the sharper radial tangent to obtain
\[
 \log\Delta_7\le\sum_i\log S(|z_i|^2)
       \le7\log7+\frac{127}{128}s
       \le7\log7+\frac{24003}{65536}.
\]
The zero-determinant case is immediate. For the remaining cases
put $x=24003/65536$. The exact exponential tail bound gives
\[
 e^x\le\sum_{j=0}^{12}\frac{x^j}{j!}
       +\frac{x^{13}}{13!}\frac1{1-x/14}
       <\frac{7^{14}}{6^{15}}.
\]
The last comparison is a strict rational calculation, with margin
greater than $1/10000$. Therefore $\Delta_7<C_7$.
The central hexagon attains $C_7$, proving the final necessary
condition. No stationary-shape classification enters any step.
\end{proof}

\paragraph{A stronger angular certificate when every radius is admissible.}
The same tangent also applies beyond the variance band whenever
every squared radius is at least $7/16$.
With $X$ defined in the earlier seven-point row-correlation theorem,
its spectral deficit estimate gives
\[
 \log\Delta_7\le7\log7+\frac{127}{128}s-\frac X2.
\]
Consequently the joint conditions
\[
 \min_i|z_i|^2\ge\frac7{16},\qquad
 \frac{127}{128}s-\frac X2\le\frac{24003}{65536}
\]
exclude the configuration from global maximality.
Every global maximizer with all squared radii at least $7/16$
must therefore obey the sharper conditional constraint
\[
                    X<\frac{127}{64}s-\frac{24003}{32768}.
\]
These bounds enlarge the rigorously covered continuous regions;
the complementary unrestricted region remains to be controlled.

\subsection{Global quintic radial tangents and a moment-dual exclusion}
\label{sec:noorigin-oct8-quintic-radial}

A fifth-degree radial tangent admits a global sign proof for both
six and seven points. Its sum can then be bounded by a quadratic
moment dual. The quadratic is a pointwise inequality for every
radial deviation, and imposes no two-radius assumption on the
actual configuration.

\begin{lemma}[Global quintic radial majorants]
\label{lem:noorigin-oct8-quintic-majorants}
Put $S_n(t)=1+t+\cdots+t^{n-1}$ and $u=t-1$. For every $t\ge0$,
\begin{align*}
 \log S_6(t)&\le\log6+\frac52u+\frac5{24}u^2-\frac58u^3
                       +\frac{151}{576}u^4+\frac{53}{288}u^5,\\
 \log S_7(t)&\le\log7+3u+\frac12u^2-u^3+\frac14u^4+\frac35u^5.
\end{align*}
Equality in either formula holds exactly at $t=1$.
\end{lemma}
\begin{proof}
Let $P_n(u)$ denote the polynomial after $\log n$ and let
$g_n(t)=\log n+P_n(t-1)-\log S_n(t)$.
Direct exact multiplication gives
\[
 g_6'(t)=\frac{(t-1)^5(265+832t+1161t^2+832t^3+265t^4)}
                      {288S_6(t)},
\]
\[
 g_7'(t)=\frac{t(t-1)^5(3+7t+9t^2+7t^3+3t^4)}{S_7(t)}.
\]
Both derivatives are negative on $0<t<1$ and positive on $t>1$.
Their value at $1$ is zero. Continuity includes the endpoint
$t=0$, proving the global bounds and strictness away from $1$.
\end{proof}

\begin{lemma}[A quadratic moment dual for the summed quintic]
\label{lem:noorigin-oct8-quintic-dual}
For $n=6$ or $7$, put $m=n-1$ and write
$F_n(u)=P_n(u)-(n-1)u/2=a_2u^2+a_3u^3+a_4u^4+a_5u^5$.
Suppose real numbers $u_1,\ldots,u_n$ satisfy
$\sum_i u_i=0$, $\sum_i u_i^2=s$, and
$A=\sqrt{ms/n}\le3/5$. Then
\[
                 \sum_i F_n(u_i)\le\Phi_n(A),
\]
where
\begin{align*}
 \Phi_6(A)&=\frac14A^2+\frac35A^3+\frac{1057}{4000}A^4
                                      -\frac{689}{3750}A^5,\\
 \Phi_7(A)&=\frac7{864}A^2(72+120A+31A^2-74A^3).
\end{align*}
Both functions strictly increase for $0<A\le3/5$.
\end{lemma}
\begin{proof}
The zero-sum condition and Cauchy give $|u_i|\le A$.
For $A>0$, let $q$ be the quadratic interpolating $F_n$ at $-A$
and interpolating its value and derivative at $A/m$.
Polynomial division gives the exact factorization
\[
 q(u)-F_n(u)=(u+A)(u-A/m)^2 R(u),
\]
where
\begin{align*}
 R(u)={}&-a_3-a_4(u-A+2A/m)\\
        &-a_5\{u^2+(-A+2A/m)u+A^2-2A^2/m+3A^2/m^2\}.
\end{align*}
Here $a_5>0$, so $R$ is concave in $u$ and its minimum on
$[-A,A]$ occurs at an endpoint. Those endpoint values are
\begin{align*}
 R(A)&=-a_3-\frac{2a_4}mA-a_5(1+3/m^2)A^2,\\
 R(-A)&=-a_3+2a_4(1-1/m)A-a_5(3-4/m+3/m^2)A^2.
\end{align*}
For both coefficient rows, $A\le3/5$ makes these strictly positive.
For the second endpoint it suffices to discard its positive linear
term; all remaining comparisons are rational and replayed by the
checker. Therefore $q(u)\ge F_n(u)$ throughout the full interval
containing every actual $u_i$.

Because $q$ is quadratic, its sum depends only on
$\sum_i1=n$, $\sum_i u_i=0$, and $\sum_i u_i^2=s$.
The formal list $-A,A/m,\ldots,A/m$ has those same three moments.
Thus
\[
 \sum_iF_n(u_i)\le\sum_iq(u_i)
      =F_n(-A)+mF_n(A/m)=\Phi_n(A).
\]
The formal list is used only to evaluate a quadratic moment
functional; the actual radii have not been classified or replaced.
Expanding gives the two displayed polynomials. If
$\Phi_n=d_2A^2+d_3A^3+d_4A^4-e_5A^5$,
all four coefficients are positive, and
\[
 \Phi_n'(A)=A(2d_2+3d_3A+4d_4A^2-5e_5A^3)>0
\]
on the stated range, since $2d_2-5e_5(3/5)^3>0$.
The case $A=0$ is immediate.
\end{proof}

\begin{theorem}[Stronger whole-angle six- and seven-point bands]
\label{thm:noorigin-oct8-quintic-bands}
For arbitrary configurations normalized by $\sum_i|z_i|^2=n$,
with $s=\sum_i(|z_i|^2-1)^2$, one has
\[
 n=6,\quad s\le\frac{11}{125}\quad\Longrightarrow\quad\Delta_6<C_6,
\]
\[
 n=7,\quad s\le\frac5{13}\quad\Longrightarrow\quad\Delta_7<C_7.
\]
These exclusions hold for all angles and all holomorphic moments.
\end{theorem}
\begin{proof}
Row Hadamard and the preceding lemmas give, on their variance range,
\[
       \log\Delta_n\le n\log n+\Phi_n\bigl(\sqrt{(n-1)s/n}\bigr).
\]
Both band endpoints have $A<3/5$. Monotonicity and exact rational
squaring prove
\[
 \Phi_6\bigl(\sqrt{11/150}\bigr)<\frac{11}{350},\qquad
 \Phi_7\bigl(\sqrt{30/91}\bigr)<\frac{18313}{50000}.
\]
For clarity, after replacing $A^2$ by its rational endpoint value,
each left side has the form $r_0+r_1\sqrt v$ with $r_1>0$.
The displayed comparisons are equivalent to the positive rational
inequalities $((x-r_0)/r_1)^2>v$.
The exact degree-twelve exponential tail bound then gives
\[
 e^{11/350}<\frac{6^9}{5^{10}},\qquad
 e^{18313/50000}<\frac{7^{14}}{6^{15}}.
\]
Their strict rational margins exceed $1/50000$ and $1/10000$,
respectively. Multiplying by $n^n$ proves both value exclusions;
the zero-determinant cases are immediate. All polynomial and
strict-value comparisons are reproduced by the standard-library
checker.
\end{proof}

Thus every six-point global maximizer must have $s>11/125$, and
every seven-point global maximizer must have $s>5/13$.
The moment-dual step is a pointwise polynomial inequality and does
not invoke a catalogue of radial profiles or stationary shapes.
The complementary large-variance regions remain unresolved.

\section{Eight-point projection boundaries and stronger value regions}
\subsection{The first projected quadratic never vanishes at a local maximum}
\label{sec:noorigin-oct8-eight-boundary}

The earlier eight-point circle interface divided by the positive
determinant of the conjugate pair $v_2,\overline v_2$. The following
argument first removes its rank-zero boundary. It applies to every
$n\ge5$ and uses no list of stationary configurations.

\begin{lemma}\label{lem:noorigin-oct8-quadratic-nonzero}
Let $n\ge5$ and let $z$ be a centered, collision-free, noncollinear
stationary configuration with $\sum_i|z_i|^2=n$. If $s>0$, then
\[
 \Pi_{\{\mathbf1,z,\overline z,t-\mathbf1\}^{\perp}}z^{\circ2}
 \ne0.
\]
Consequently the quantity $A_0=H_{22}$ in the finite projection
interface is strictly positive at every actual local maximum to
which that interface applies.
\end{lemma}
\begin{proof}
If the projection vanishes, there are complex constants $a,b,c,d$
such that every point satisfies
\[
             z_i^2=a+bz_i+c\overline z_i+d|z_i|^2.
\]
Let $F(x,y)$ and $G(x,y)$ be the real and imaginary parts of
$z^2-a-bz-c\overline z-d|z|^2$, where $z=x+iy$.
Their homogeneous quadratic parts are linearly independent: their
traceless parts are $x^2-y^2$ and $2xy$, respectively.

Two real quadratic polynomials with independent quadratic parts
and at least five distinct common real zeros must have a common
real linear factor. Here is the elementary intersection argument.
If they have no common factor over $\mathbb C$, choose a generic
real linear coordinate $x$ separating the finitely many common
zeros. Their resultant with respect to the other coordinate has
degree at most four and is nonzero, so there are at most four
common zeros. If a common factor has degree two, both quadratics
are proportional, contrary to the independent quadratic parts.
Thus a common factor has degree one. A nonreal linear factor
would bring its distinct conjugate factor as well, again a common
degree-two factor; hence the common line is real. After factoring
out that line, the two remaining linear factors are independent,
so they have at most one additional common point. Therefore all
but at most one of the $n$ points lie on a real line.

Since the configuration is noncollinear, exactly one point is off
that line. Rotate so the line is $\operatorname{Im}z=b_0$;
write the off-line point as $x_0+ic_0$. Centering gives
$c_0=-(n-1)b_0$, so $b_0\ne0$. For a point $x_i+ib_0$ on
the line the imaginary part of the force equation is
\[
 \frac{b_0-c_0}{(x_i-x_0)^2+(b_0-c_0)^2}=\lambda b_0,
 \qquad \lambda=\frac{n-1}{2}.
\]
Thus every denominator equals $n/\lambda$, and the $x_i$ take
at most two values. This contradicts the at least four distinct
points on the line. The final statement follows because local
maxima are noncollinear, as proved in the finite projection section.
\end{proof}

\subsection{A value certificate on the remaining projection boundary}

The rank-one boundary is a real conic condition; it is not excluded
by the preceding lemma. It can nevertheless be included in a
single sufficient value certificate. One must replace the previous
circle by a disk on this boundary. The circle lower bound cannot
be carried over unchanged.

For eight points retain $W,H,C$ from
Lemma~\ref{lem:noclass-finite-projection}. Put
\[
 A_0=H_{22},\quad K_0=C_{22},\quad
 D_0=\begin{pmatrix}A_0&\overline K_0\\K_0&A_0\end{pmatrix},
 \quad J=\begin{pmatrix}0&1\\1&0\end{pmatrix},\quad
 d_k=\binom{H_{2k}}{C_{2k}}\quad(k=3,4).
\]
Use the following explicit Moore--Penrose inverse:
\[
 D_0^+=
 \begin{cases}
 D_0^{-1},&A_0>|K_0|,\\
 D_0/(4A_0^2),&A_0=|K_0|>0,\\
 0,&A_0=0.
 \end{cases}
\]
Define
\[
 \widehat H_{jk}=H_{jk}-d_j^*D_0^+d_k,
 \qquad
 \widehat C_{jk}=C_{jk}-d_j^TJD_0^+d_k,
\]
and $A=\widehat H_{33}$, $B=\widehat H_{44}$,
$G=\widehat H_{34}$, $K=\widehat C_{33}$,
$L=\widehat C_{34}$. These expressions include the full lower
moment projection corrections.

\begin{theorem}\label{thm:noorigin-oct8-eight-all-rank-certificate}
For every centered stationary eight-point configuration with
$s>0$ and $|p_2|<8$, the above projected quantities satisfy
\begin{equation}\label{eq:noorigin-oct8-eight-disk}
 |AL-GK|^2\le(A^2-|K|^2)(AB-|G|^2).
\end{equation}
The two factors on the right are nonnegative. If
$A_0>|K_0|$, equality holds. If $A>0$, put
\[
 \ell_0=p_7-L,\qquad c_7=\ell_0+GK/A,
 \qquad
 \rho_7=\frac{\sqrt{(A^2-|K|^2)(AB-|G|^2)}}{A}.
\]
Then $\ell_0$ uses only moments below $p_7$, and
\[
 |p_7-c_7|\le\rho_7,
 \qquad |p_7|\ge\mathcal L_7,
\]
where a valid choice, including every projection rank, is
\[
 \mathcal L_7=
 \begin{cases}
 \bigl|\rho_7-|c_7|\bigr|,&A_0>|K_0|,\ A>0,\\
 \max\{0,|c_7|-\rho_7\},&A_0=|K_0|,\ A>0,\\
 |\ell_0|,&A=0.
 \end{cases}
\]
In the last case $p_7=\ell_0$ exactly. Consequently
\begin{equation}\label{eq:noorigin-oct8-eight-all-rank-value}
 \Delta_8\le(8T_7-\mathcal L_7^2)\prod_{k=1}^6T_k.
\end{equation}
Every actual configuration for which this right-hand side is at
most $C_8$ therefore satisfies the sharp value inequality, including
the rank-one projection boundary and the boundary $A=0$.
No eighth holomorphic moment is used.
\end{theorem}
\begin{proof}
The displayed pseudoinverse gives the orthogonal projection onto
$\operatorname{span}\{v_2,\overline v_2\}$ at every rank.
The projection range is invariant under conjugation. The Hermitian
projection formula is the usual one, while the bilinear formula
follows from $(D_0^+)^TJ=JD_0^+$ and
$D_0^+D_0D_0^+=D_0^+$. Thus the hatted quantities are exactly
the Hermitian and bilinear inner products of
$a=\Pi_Uv_3$, $b=\Pi_Uv_4$, where $U$ is the orthogonal
complement of that range in $W$.

The Gram matrix of $a,b,\overline a$ is positive semidefinite.
Its determinant is
\[
 \delta=A^2B-A(|G|^2+|L|^2)-B|K|^2
                 +2\Re(GK\overline L)\ge0.
\]
An exact rearrangement gives
\[
 (A^2-|K|^2)(AB-|G|^2)-|AL-GK|^2=A\delta\ge0,
\]
proving the disk constraint. Cauchy proves nonnegativity of its
factors. If the removed span has rank two, then $U$ has complex
dimension two, so the three-vector determinant vanishes and the
constraint is the exact circle proved earlier.
If $A=0$, then $a=0$ and $G=K=L=0$, which gives
$p_7=\ell_0$ without division. For $A>0$, divide the disk
constraint by $A^2$ and use the reverse triangle inequality.
An exact circle gives the stronger absolute-difference lower
bound stated in its branch.

The only occurrence of $p_7$ is the original $C_{34}$ term
$p_7$; both of its projection corrections involve lower moments.
All the other entries involve at most $p_6$, radial moments through
$T_4$, and the same mixed moments as the earlier interface.
Finally apply Hadamard--Fischer to the original Vandermonde Gram,
keeping its $\{1,z^7\}$ block. Its determinant is
$8T_7-|p_7|^2$, proving the value bound.
\end{proof}

At a rank-one conjugate pair, $K_0/A_0$ has modulus one and the
additional exact consistency conditions are
\[
             C_{23}=(K_0/A_0)H_{23},\qquad
             C_{24}=(K_0/A_0)H_{24}.
\]
They follow from $\overline v_2=(\overline K_0/A_0)v_2$.
In particular they retain constraints on $p_5,p_6$ that would be
lost by treating the conic boundary as an arbitrary disk.

This theorem removes the undefined divisions in the previous
sufficient certificate. It does not prove that all remaining
global maximizers pass the certificate. The disk is genuinely
weaker than a circle: in a real three-dimensional residual space,
$a=(1,i,0)$ and $b=(0,0,1)$ give $A=2$, $B=1$, $G=K=L=0$,
and the disk has radius $\sqrt2$, while the circle assertion is
false. In particular, the earlier weighted two-dimensional
lower inequality must not be applied on that boundary.

\subsection{A short low-degree operator constraint retaining the mixed phase}

\begin{proposition}\label{prop:noorigin-oct8-low-operator-gram}
At any centered collision-free local maximum with moment $n$, put
$q=|p_2|^2$ and $\lambda=(n-1)/2$. The following matrix is positive
semidefinite:
\[
 \begin{pmatrix}
 n-3s-q/n&-5R_2-\overline{p_2}p_3/n\\
 -5\overline R_2-p_2\overline{p_3}/n&
 -8T_3+\dfrac{6(2n-3)}\lambda T_2
       -\dfrac{n(2n-3)^2}{\lambda^2}-|p_3|^2/n
 \end{pmatrix}.
\]
For $n=8$ its lower-right entry is
\[
 -8T_3+\frac{156}{7}T_2-\frac{5408}{49}-\frac{|p_3|^2}{8}.
\]
Besides the individual radial bounds, its determinant retains the
phase-sensitive constraint
\[
 \left|5R_2+\frac{\overline{p_2}p_3}{n}\right|^2
 \le\left(n-3s-\frac qn\right)
 \left\{-8T_3+\frac{6(2n-3)}\lambda T_2
       -\frac{n(2n-3)^2}{\lambda^2}-\frac{|p_3|^2}{n}\right\}.
\]
\end{proposition}
\begin{proof}
Apply $\|B\|\le\lambda$ to all combinations of the two centered
columns $z^2-p_2\mathbf1/n$ and $z^3-p_3\mathbf1/n$, and divide
their defect Gram matrix by $\lambda^2$. The force recurrences give
\[
 Bz^2=2\lambda(t-\mathbf1),\qquad
 Bz^3=\{3\lambda t-(2n-3)\}z.
\]
The identity $R_1=p_1=0$ makes their output inner product
$6\lambda^2R_2$, whereas their centered input inner product is
$R_2-\overline{p_2}p_3/n$. Expanding their two norms gives the
displayed diagonal entries. The determinant inequality follows
from positive semidefiniteness.
\end{proof}

These new assertions are geometric and operator necessary conditions,
or sufficient value certificates, with full proofs. The accompanying
exact checks audit their algebra and singular-rank formulas; they do
not replace the remaining uniform objective comparison.

\subsection{An explicit frame gap on each singular boundary}

\begin{corollary}\label{cor:noorigin-oct8-boundary-frame-gap}
At an eight-point local maximum, use $\varepsilon_8$ from the
finite Fourier-frame lemma. On the conic boundary
$A_0=|K_0|>0$ one necessarily has
\[
                         \varepsilon_8^2\ge\frac87.
\]
If instead $A_0>|K_0|$ and the next residual norm $A$ vanishes,
then
\[
                         \varepsilon_8^2\ge\frac83.
\]
On the full local-maximum domain,
\[
 \varepsilon_8\le\frac{1}{12\cdot10^{14}}
 \quad\Longrightarrow\quad
 z\text{ is the central regular heptagon}.
\]
Thus the singular boundaries are quantitatively separated from
the already proved small-harmonic rigidity region. No upper bound
forcing every global maximum into that region is asserted.
\end{corollary}
\begin{proof}
The normalized eight-column frame has columns
$\mathbf1,z,\overline z,z^2,\overline z^2,z^3,\overline z^3,u$,
each divided by its norm. The first four base directions
$\mathbf1,z,\overline z,u$ are independent. If the projected
quadratic pair has rank one, the whole frame has rank at most
$4+1+2=7$. If that pair has rank two and $A=0$, both the
projected $v_3$ and its conjugate are in its span, so the whole
frame has rank at most six. For a positive matrix of trace eight
and rank at most $r$, Cauchy gives
\[
 \|G-I\|_F^2=\operatorname{tr}(G^2)-8
                         \ge\frac{64}{r}-8.
\]
This proves the two boundary gaps.

The operator bounds give $T_2\le32/3<11$. The lower-right
entry of Proposition~\ref{prop:noorigin-oct8-low-operator-gram}
gives $T_3\le780/49<16$. The positive terms in the exact frame
formula therefore imply
\[
 |p_2|\le6\varepsilon_8,\quad
 |p_3|\le4\varepsilon_8,\quad
 |p_4|\le8\varepsilon_8,\quad
 |p_5|\le7\varepsilon_8,\quad
 |p_6|\le12\varepsilon_8.
\]
For example the last follows from
$2|p_6|^2/T_3^2\le\varepsilon_8^2$, and the other four follow
by retaining respectively the terms $1/32$,
$1/(2T_3)+1/(2T_2)$, $2/T_2^2$ and $4/(T_2T_3)$.
Consequently the displayed frame threshold makes every
$|p_k|$, $2\le k\le6$, at most $10^{-14}$. Apply
Theorem~\ref{thm:noorigin-oct7-eight-improvedharmonic}.
The case $s=0$ is already excluded for eight-point local maxima
by the regular-octagon second-order test.
\end{proof}

\subsection{A sharper global eight-point radial logarithm bound}
\label{sec:noorigin-oct8-eight-radial}

The radial coefficient $11/7$ can be reduced to $3/2$ on the
entire half-line. The certificate below combines one polynomial
Bernstein interval with a positive shifted derivative polynomial;
it makes no restriction on angles or stationary shape.

\begin{lemma}\label{lem:noorigin-oct8-eight-radial}
For $S(t)=1+t+\cdots+t^7$ and every $t\ge0$,
\[
 \log S(t)\le\log8+\frac72(t-1)+\frac32(t-1)^2.
\]
Equality occurs exactly at $t=1$.
\end{lemma}
\begin{proof}
Let $g(t)$ be the right-hand side minus $\log S(t)$.
Exact division gives
\[
 g'(t)=\frac{(t-1)Q(t)}{S(t)},\qquad
 Q(t)=\frac12-t-\frac32t^2-t^3+\frac12t^4+3t^5
                               +\frac{13}{2}t^6+3t^7.
\]
The coefficients of $Q(2/3+x)$, in increasing degree, are
\[
 \left(\frac{161}{1458},\frac{1507}{243},\frac{617}{18},
       \frac{1969}{27},\frac{1529}{18},57,\frac{41}{2},3\right).
\]
They are strictly positive. Hence $g$ strictly decreases on
$[2/3,1]$ and strictly increases on $[1,\infty)$.
Since $g(1)=0$, the assertion follows on $[2/3,\infty)$.

For $0\le t\le2/3$, put
\[
 y(t)=2-\frac t2-\frac32t^2,\qquad
 E(y)=\sum_{j=0}^{12}\frac{y^j}{j!}
                             +\frac{7y^{13}}{6\cdot13!}.
\]
Here $1\le y\le2$. The exponential tail after degree twelve
has successive ratios at most $2/14=1/7$, so $e^y\le E(y)$.
Define the rational polynomial of degree thirty-three
\[
                         P(t)=8-S(t)E(y(t)).
\]
Every Bernstein coefficient of $P(2x/3)$ on $0\le x\le1$
is strictly larger than $1/25$. The exact minimum is
\[
        \frac{1285380596146225650109}{27234837359209351872000}
                                    >\frac1{25}.
\]
This is a single-interval rational certificate. To make the
finite verification reproducible, if
$P(2x/3)=\sum_{i=0}^{33}c_i x^i$, its coefficients are
\[
       b_k=\sum_{i=0}^k c_i\frac{\binom{k}{i}}{\binom{33}{i}}
                         \qquad(0\le k\le33).
\]
The accompanying standard-library checker forms $P$ from the
displayed formula, checks every coefficient by exact arithmetic,
and records all thirty-four values. Thus $P(t)>1/25$ throughout
the interval. Consequently
\[
                S(t)e^{y(t)}\le S(t)E(y(t))<8,
\]
which is precisely $g(t)>0$ on $[0,2/3]$.
This finishes the full half-line proof and the equality statement.
\end{proof}

\begin{theorem}\label{thm:noorigin-oct8-eight-sharper-regions}
For arbitrary eight points with $\sum_i|z_i|^2=8$, retain the
row correlation quantity
\[
 X=\sum_{i<j}\frac{|S(z_i\overline z_j)|^2}{S(t_i)S(t_j)},
 \qquad t_i=|z_i|^2,\qquad s=\sum_i(t_i-1)^2.
\]
For distinct points,
\[
                 \log\Delta_8\le8\log8+\frac32s-\frac7{16}X.
\]
In particular, the entire all-angle region $s\le12/25$ has
value strictly below $C_8$. Every collision-free local maximum
with $X\ge15/2$ also has value strictly below $C_8$.
Thus any global maximizer necessarily satisfies
\[
 s>\frac{12}{25},\qquad X<\frac{15}{2},\qquad
 X<\frac{24}{7}s-\frac{1448}{875}.
\]
These inequalities strengthen the previous exclusions and use no
stationary classification.
\end{theorem}
\begin{proof}
Apply the eight-dimensional normalized-Gram inequality
from Lemma~\ref{lem:noclass-general-gram} and sum the new scalar
bound. Its linear terms cancel because $\sum_i(t_i-1)=0$.
The independent rational logarithm enclosure from the checker gives
\[
 \log(C_8/8^8)=\log(8^{20}/7^{21})>\frac{181}{250}.
\]
For $s\le12/25$, the excess is at most
$(3/2)(12/25)=18/25<181/250$.
At a local maximum $s\le8/3$, so $X\ge15/2$ makes the excess
at most $4-(7/16)(15/2)=23/32<181/250$.
The central heptagon competitor has value $C_8$.
Contraposition gives the first two necessary conditions and,
on rearranging $(3/2)s-(7/16)X>181/250$, the last one.
Collisions have zero value and satisfy both exclusions directly.
\end{proof}

The stronger scalar certificate shrinks the unrestricted remaining
domain. It does not prove that the remaining simultaneous domain is
empty, or establish the full sharp eight-point inequality.

\subsection{An eight-point quintic tangent and a larger whole-angle band}
\label{sec:noorigin-oct8-eight-quintic}

The global quadratic logarithm bound remains useful for large radial
variance and its correlation penalty. A quintic tangent gives a
stronger whole-angle band near the unit circle, using a quadratic
moment dual rather than a radial population assumption.

\begin{lemma}\label{lem:noorigin-oct8-eight-quintic-scalar}
For $t\ge3/10$ and $u=t-1$,
\[
 \log S_8(t)\le\log8+\frac72u+\frac78u^2-\frac{35}{24}u^3
                          +\frac7{64}u^4+\frac{217}{160}u^5.
\]
Equality holds exactly at $t=1$. This scalar majorant is asserted
only on the displayed domain.
\end{lemma}
\begin{proof}
Write the right-hand side as $\log8+P(u)$ and let $g(t)$ be
its difference from $\log S_8(t)$. Exact multiplication gives
\[
 g'(t)=\frac{(t-1)^5Q(t)}{S_8(t)},\qquad
 Q(t)=\frac{-87+217t^2+448t^3+553t^4+448t^5+217t^6}{32}.
\]
The polynomial $Q$ is strictly increasing for $t>0$, is negative
at zero and positive at one. It has exactly one positive root
$r<1$. Thus $g$ increases up to $r$, decreases from $r$ to one,
and increases after one. The minimum on $[3/10,1]$ occurs at an
endpoint, irrespective of whether $r$ lies in that interval.
At $1$ it is zero. At $t_0=3/10$, let
\[
 R=\frac8{S_8(t_0)},\qquad v=\frac{R-1}{R+1}.
\]
All these numbers are rational and $0<v<1$. The first sixteen
positive terms in the logarithm series give the exact inequality
\[
 g(t_0)=\log R+P(-7/10)
 >2\sum_{j=0}^{15}\frac{v^{2j+1}}{2j+1}+P(-7/10)
 >\frac1{10000}.
\]
The last rational comparison is replayed in the checker.
The derivative signs now prove the lemma and its equality statement.
\end{proof}

\begin{lemma}\label{lem:noorigin-oct8-eight-quintic-dual}
Let eight real numbers $u_i$ satisfy $\sum_i u_i=0$ and
$\sum_i u_i^2=s$, and put $A=\sqrt{7s/8}$. If $A\le17/25$, then
\[
 \sum_i\left(\frac78u_i^2-\frac{35}{24}u_i^3
                +\frac7{64}u_i^4+\frac{217}{160}u_i^5\right)
 \le\Phi_8(A),
\]
where
\[
             \Phi_8(A)=A^2+\frac{10}{7}A^3
                        +\frac{43}{392}A^4-\frac{465}{343}A^5.
\]
The function $\Phi_8$ is strictly increasing on $0<A\le17/25$.
\end{lemma}
\begin{proof}
The zero-sum condition and Cauchy give $|u_i|\le A$.
Let $F(u)=a_2u^2+a_3u^3+a_4u^4+a_5u^5$ have the coefficient
row in the statement. For $A>0$, let $q$ be the quadratic
interpolating $F$ at $-A$, and interpolating its value and first
derivative at $A/7$. Exact division gives
\[
 q(u)-F(u)=(u+A)(u-A/7)^2\mathcal R(u),
\]
\[
 \mathcal R(u)=-a_3-a_4(u-A+2A/7)
 -a_5\{u^2+(-A+2A/7)u+A^2-2A^2/7+3A^2/49\}.
\]
This is concave in $u$, and its endpoint values are
\[
 \mathcal R(A)=\frac{35}{24}-\frac A{32}-\frac{403}{280}A^2,
 \qquad
 \mathcal R(-A)=\frac{35}{24}+\frac3{16}A-\frac{1891}{560}A^2.
\]
The first decreases in $A\ge0$. The second is concave in $A$,
so its minimum on $[0,17/25]$ is at an endpoint. At $A=17/25$
their values are respectively
\[
                   \frac{1620271}{2100000},\qquad
                   \frac{6407}{262500},
\]
both positive; the second is also positive at $A=0$.
Thus $\mathcal R$ is positive throughout $[-A,A]$, and
$q(u)\ge F(u)$ there.

Because $q$ is quadratic, $\sum_iq(u_i)$ depends only on the
three specified moments. The formal list $-A,A/7,\ldots,A/7$
has those same moments and gives
\[
 \sum_iF(u_i)\le\sum_iq(u_i)
                   =F(-A)+7F(A/7)=\Phi_8(A).
\]
This evaluates a quadratic moment functional; it makes no
assumption that the actual radial deviations have those two values.
The formal two-variable polynomial identity for $q-F$ and the
displayed expansion of $\Phi_8$ are replayed exactly by the checker.
For monotonicity,
\[
 \Phi_8'(A)=A\left(2+\frac{30}{7}A+\frac{43}{98}A^2
                                      -\frac{2325}{343}A^3\right)>0,
\]
since $30/7-(2325/343)(17/25)^2=9873/8575>0$.
The case $A=0$ is immediate.
\end{proof}

\begin{theorem}\label{thm:noorigin-oct8-eight-quintic-band}
For arbitrary eight points with total moment eight, define $s$ and
$X$ as in Theorem~\ref{thm:noorigin-oct8-eight-sharper-regions}.
Whenever $s\le2312/4375$, distinct points satisfy
\[
       \log\Delta_8\le8\log8+
             \Phi_8\!\left(\sqrt{7s/8}\right)-\frac7{16}X.
\]
In particular the entire all-angle region
\[
                       s\le\frac{31}{60}
\]
has value strictly below $C_8$. Every global maximizer therefore
satisfies $s>31/60$. The correlation restriction $X<15/2$ from
the global quadratic logarithm bound remains valid independently.
\end{theorem}
\begin{proof}
The stated variance domain is exactly the condition
$A=\sqrt{7s/8}\le17/25$. Each deviation $u_i=t_i-1$ has
modulus at most $A$, so $t_i\ge8/25>3/10$. Apply the scalar
quintic lemma to every radius and sum it; the linear terms vanish.
The quadratic moment dual gives the displayed radial term, and
the normalized eight-row Gram inequality adds $-7X/16$.

At the band endpoint $s=31/60$, put $v=A^2=217/480$. Write
\[
 \Phi_8(\sqrt v)=r_0+r_1\sqrt v,\qquad
 r_0=v+\frac{43}{392}v^2,\quad
 r_1=\frac{10}{7}v-\frac{465}{343}v^2.
\]
Exact rational arithmetic proves $r_1>0$,
$181/250-r_0>0$, and
\[
                 (181/250-r_0)^2>r_1^2v.
\]
Consequently $\Phi_8(\sqrt{217/480})<181/250$.
Monotonicity and the previously certified strict logarithm gap
$\log(C_8/8^8)>181/250$ prove the whole-angle exclusion.
Collisions have zero value. The attained central-heptagon
competitor gives the final global necessary condition.
\end{proof}

This sharper band still leaves a substantial large-variance region.
No conclusion that all remaining global maximizers satisfy its
variance hypothesis or its sufficient value bound is claimed.

% Fragment for main.tex: uses its theorem/proof environments and booktabs.
% This replacement distinguishes the published coverage input from the
% independently replayable finite arithmetic certificate.
\section{Earlier classification-dependent six- and seven-point cross-checks}
\label{sec:six-seven-improved-proof}

Put $N_n=n(n-1)/2$ and
\[
 C_6=\frac{6^{15}}{5^{10}},\qquad
 C_7=\frac{7^{21}}{6^{15}}.
\]
This retained proof is a historical cross-check. None of its classification
coverage is used in the preceding non-classification arguments.
The proof has two separate inputs. The published Faug\`ere--Svartz
computation supplies coverage of all physical stationary shapes.
The rational certificates supplied here cover and compare every algebraic
candidate in the retained charts. The latter verification does not
recompute the former decomposition.

\begin{lemma}[Reduction to normalized coefficient charts]
\label{lem:six-seven-chart-reduction}
A global maximizer of $\Delta_n$ on $\sum_i|z_i|^2=n$ exists, is
collision-free, and has $\sum_i z_i=0$. If $w_i=\alpha z_i$ for
$\alpha\ne0$, and $P(t)=\prod_i(t-w_i)$, $Q=\sum_i|w_i|^2$, then
\begin{equation}\label{eq:six-seven-uniform-residual}
 QP''(w_i)-2N_n\overline{w_i}P'(w_i)=0\quad(1\le i\le n),
\end{equation}
and its value at moment $n$ is
\begin{equation}\label{eq:six-seven-uniform-value}
 \mathcal V_n(w)=\left(\frac nQ\right)^{N_n}
                    \prod_{i<j}|w_i-w_j|^2.
\end{equation}
Every such shape can be put in a chart in which the first nonzero
elementary symmetric coefficient is $1$; complex scaling loses no
physical shape and preserves \eqref{eq:six-seven-uniform-value}.
\end{lemma}
\begin{proof}
The constraint sphere is compact and the objective is continuous.
A regular $n$-gon gives a positive value, whereas a collision gives
zero, so a maximizer has distinct points. Let $c=n^{-1}\sum_i z_i$.
Translation by $-c$ preserves all pair differences and changes the
moment to $n-n|c|^2$. If $c\ne0$, rescaling the translated points to
moment $n$ strictly increases the positive objective, a contradiction.

The Lagrange equations for $\log\Delta_n$ at a maximizer are
\[
 F_i(z):=\sum_{j\ne i}\frac1{z_i-z_j}=\lambda\overline{z_i}.
\]
Pairwise summation gives
$\sum_i z_iF_i(z)=\sum_{i<j}(z_i-z_j)/(z_i-z_j)=N_n$.
Consequently $\lambda=N_n/n>0$. Under $w_i=\alpha z_i$,
$F_i(w)=\lambda|\alpha|^{-2}\overline{w_i}$, and the positive
force constant $\kappa=\lambda|\alpha|^{-2}$ obeys $\kappa Q=N_n$.
The identity $F_i(w)=P''(w_i)/(2P'(w_i))$ proves the residual equation.
Scaling a shape of moment $Q$ to moment $n$ multiplies every squared
pair difference by $n/Q$, proving the value formula.

The centered root polynomial is
$P(t)=t^n+\sum_{k=2}^n(-1)^ke_k t^{n-k}$.
Some $e_k$ is nonzero, since otherwise every root would be zero.
Choosing $\alpha$ with $\alpha^ke_k=1$ normalizes its first nonzero
coefficient, with all earlier coefficients still zero.
\end{proof}

\paragraph{The exact external coverage statement.}
We use the collision-saturated coefficient-chart computation reported in
Faug\`ere--Svartz~\cite[Sections~5.3--5.7]{FS}, with the archived
author bases~\cite{FSdata,FSdataSeven}, in the following precise form:
every collision-free physical stationary shape for $n=6$ or $n=7$,
up to complex scaling and permutation, is represented in the retained
charts specified below. This is a necessary coverage statement; the
charts may also contain nonphysical algebraic candidates. Theorem~5
of~\cite{FS} gives the symmetric equations, whereas the chart
classification is the subsequent published computation.
This coverage concerns solutions of the original physical equilibrium
equations, including unstable and collinear configurations; it is not a
classification restricted to stable equilibria. Sections~5.3--5.4 and
5.7 occur on proceedings pages~176 and~177, respectively.

For $n=6$ the author file supplies eleven components, of parameter
degrees
\[
 22,28,12,4,2,1,1,6,2,1,1,
\]
with total $80$. Its header $\texttt{N=7}$ is a metadata typo: the
linked file is \path{vortices_N6.txt} and its variables are
$e_2,\ldots,e_6$. Exact substitution into the six-point coefficient
equations supplies an additional integrity check. It proves forward
containment, and does not by itself prove completeness of these charts.

For $n=7$, Section~7 and Figure~4 of~\cite[p.~178]{FS} are also used. Their
$e_2=1$ component has degree $526$, decomposed into degrees
$6,18,356,82,64$. Figure~4 marks the degree-$356$ branch physically
empty. Its polynomial and its independent real-feasibility certificate
are absent from our supplied files. The retained file has seven charts,
of parameter degrees
\[
 1,1,2,6,18,82,64,
\]
with total $174$. The degree-$18$ branch, also marked physically
empty in the paper, is retained and independently excluded below.
Thus our seven-point verification includes all $174$ retained
coefficient candidates, and does not claim to include the missing
$356$ complex candidates. Their physical exclusion remains a cited
external input.

\begin{proposition}[Exact comparison on every retained chart]
\label{prop:six-seven-retained-comparison}
For every retained six-point candidate that satisfies
\eqref{eq:six-seven-uniform-residual},
$\mathcal V_6\le C_6$. Equality can occur only for
$P(t)=t(t^5-1)$.
For every retained seven-point candidate that satisfies the same
residual with $n=7$, $\mathcal V_7\le C_7$. Equality can occur only
for $P(t)=t(t^6+1)$.
\end{proposition}
\begin{proof}
The coefficient conventions are
\[
 \begin{split}
 P_6(t)&=t^6+e_2t^4-e_3t^3+e_4t^2-e_5t+e_6,\\
 P_7(t)&=t^7+e_2t^5-e_3t^4+e_4t^3-e_5t^2+e_6t-e_7.
 \end{split}
\]
The standard-library programs
\path{certify_fs6_disks.py} and
\path{certify_fs7_disks.py} parse the exact author polynomials
without executing them. They isolate every parameter root and every
root of the associated point polynomial in disjoint Gaussian-rational
disks. Each disk contains one root by a strict rational Rouch\'e
inequality, and the number of disks equals the polynomial degree.
Thus no parameter root or point root in any retained chart is omitted.
Coefficient uncertainty from the parameter disk is included throughout.

Here are the replayed six-point results. ``Bounded'' means a strict
upper bound on the normalized objective; it does not assert that the
candidate satisfies the physical equations.
\begin{center}
\begin{tabular}{rrrrl}
\toprule
Component&Parameters&Force excluded&Bounded&Upper endpoint\\
\midrule
1&22&20&2&$3170$\\
2&28&28&0&---\\
3&12&12&0&---\\
4&4&2&2&$42454$\\
5&2&2&0&---\\
6&1&0&1&$18699$\\
7&1&0&1&$18699$\\
8&6&4&2&$22400$, $133$\\
9&2&1&1&$7225$\\
10&1&0&0&$C_6$, attaining\\
11&1&0&1&$6^6$, exact\\
\bottomrule
\end{tabular}
\end{center}
There are $80$ parameter roots and $480$ point-root disks.
The $69$ force exclusions and $10$ strictly smaller candidates leave
one attaining chart. Every displayed upper endpoint is at most $6^6$,
and
\[
 C_6>6^6\quad\Longleftrightarrow\quad
 6^9=10077696>9765625=5^{10}.
\]

The seven-point results are
\begin{center}
\begin{tabular}{rrrrl}
\toprule
Chart&Parameters&Force excluded&Bounded&Upper endpoint\\
\midrule
1&1&0&1&$7^7=823543$, exact\\
2&1&0&0&$C_7$, attaining\\
3&2&1&1&$40253$\\
4&6&4&2&$159666$\\
5&18&18&0&---\\
6&82&78&4&$268962$\\
7&64&54&10&$822629$\\
\bottomrule
\end{tabular}
\end{center}
There are $174$ parameter roots and $1218$ point-root disks.
The $155$ force exclusions and $18$ strictly smaller candidates leave
one attaining chart. Every displayed upper endpoint is at most $7^7$,
and
\[
 C_7>7^7\quad\Longleftrightarrow\quad
 7^{14}=678223072849>470184984576=6^{15}.
\]
The force exclusions and objective bounds are justified uniformly on
the exact root disks by the estimates given next. All deciding checks
use explicit exceptions and remain active under Python $\texttt{-O}$.
\end{proof}

\paragraph{Root coverage, physical exclusions and value bounds.}
Write $\|c\|_1=|\Re c|+|\Im c|$ and
$\|c\|_\infty=\max\{|\Re c|,|\Im c|\}$ for Gaussian rationals.
Let $p_0$ be a nominal polynomial, let its coefficient of $t^j$ have
uncertainty at most $\epsilon_j$, and put
$a_k=p_0^{(k)}(c)/k!$. The test
\[
 \|a_0\|_1+\sum_{k\ge2}\|a_k\|_1r^k
       +\sum_j\epsilon_j(\|c\|_1+r)^j
 <r\|a_1\|_\infty
\]
compares the allowed polynomial on $|t-c|=r$ with its nominal
linear term $a_1(t-c)$. Rouch\'e's theorem therefore gives exactly
one root, counted with multiplicity. Pairwise disk disjointness and
the degree count prove exhaustive coverage in each retained chart.

Let $c_i$ be the point-root centers, of common radius $r$, and set
\[
 Q_0=\sum_i|c_i|^2,\qquad
 E_Q=\sum_i(2\|c_i\|_1r+r^2).
\]
Then $|Q-Q_0|\le E_Q$. Let $v_j=p_0^{(j)}(c_i)$ and let $E_j$
bound $|P^{(j)}(w_i)-v_j|$, including coefficient and root variation.
The actual physical residual differs from
$Q_0v_2-2N_n\overline{c_i}v_1$ by at most
\[
 E_Q\|v_2\|_1+(Q_0+E_Q)E_2
 +2N_n\{r\|v_1\|_1+(\|c_i\|_1+r)E_1\}.
\]
If the nominal residual's $\|\cdot\|_\infty$ exceeds this bound
at even one point, the entire coefficient candidate is nonphysical.
This tests the actual pointwise conjugate pairing, rather than only
equality of two unordered sets of roots.

For each surviving candidate, $Q_0>E_Q$ and
\begin{equation}\label{eq:six-seven-rational-upper}
 U_n=\left(\frac n{Q_0-E_Q}\right)^{N_n}
 \prod_{i<j}\bigl(|c_i-c_j|^2+4r\|c_i-c_j\|_1+4r^2\bigr)
\end{equation}
is a rational upper bound on \eqref{eq:six-seven-uniform-value}.
The tables give upward-rounded integer endpoints of these rational
bounds, or the exact regular-polygon values. Hence their strict
comparison with $C_n$ proves the proposition without needing to
establish physical sufficiency for any bounded candidate.

\begin{theorem}[Global values using the published physical coverage]
\label{thm:six-seven-global-improved}
For $\sum_i|z_i|^2=n$,
\[
 M_6=\frac{6^{15}}{5^{10}},\qquad
 M_7=\frac{7^{21}}{6^{15}}.
\]
Equality for $n=6$ holds exactly for the origin and a regular
pentagon of radius $\sqrt{6/5}$; equality for $n=7$ holds exactly
for the origin and a regular hexagon of radius $\sqrt{7/6}$,
in each case up to rotation and permutation.
\end{theorem}
\begin{proof}
Every global maximizer is a collision-free centered physical stationary
shape by Lemma~\ref{lem:six-seven-chart-reduction}. The stated external
coverage places it in a retained chart, and
Proposition~\ref{prop:six-seven-retained-comparison} gives the upper
bound with a strict gap for every chart other than the attaining one.

To verify attainment independently of the classification, put
$m=n-1$ and consider the roots of $P(t)=t(t^m-a)$, $|a|=1$.
They are the origin and a rotated unit regular $m$-gon, with moment
$m$ and squared pair-distance product $m^m$. At every nonzero root,
$P''(w)/(2P'(w))=n/(2w)$, and the force at the origin is zero.
Thus the residual equation holds directly. Scaling to moment $n$
gives
\[
 (n/m)^{N_n}m^m
 =\frac{n^{N_n}}{(n-1)^{(n-1)(n-2)/2}}.
\]
For $n=6,7$ this is respectively $C_6,C_7$.
The strict comparisons already prove that an attaining shape must
belong to this chart. Moment normalization fixes its radius and
the remaining complex phase is exactly an overall rotation, proving
both equality classifications.
\end{proof}

\paragraph{Reproducibility boundary.}
The exact replay counts are $80=69+10+1$ and $174=155+18+1$.
The six-point component-identity checker has also been changed from
$\texttt{assert}$ statements to explicit arithmetic checks, so its
decisions remain active under $\texttt{python -S -O}$.
An independent reconstruction of the saturated Gr\"obner decompositions
is not included, and the missing seven-point degree-$356$ physical
exclusion is not claimed to have been reproduced. A proof using the
published classification and an entirely standalone computational
classification are different reproducibility claims; the theorem here
uses the former.

% New fragment: independent completion of the n=6 e_2=0 coefficient charts.
% Uses the main document's theorem/proof environments; no new bibliography.
\section{Independent completion of the exceptional six-point charts}
\label{sec:six-special-independent}

The following argument replaces the published classification input for
the entire six-point region $e_2=0$. It does not reconstruct the remaining
$e_2\ne0$ region or the missing seven-point degree-$356$ branch.

For a centered monic polynomial $p$ of degree six with distinct roots
$x_i$, put $\delta=\operatorname{Res}(p,p')$,
$g_i=p''(x_i)/p'(x_i)$, $S_k=\sum_i g_i^k$ and
$m_k=\sum_i x_i g_i^k$. Define
\[
 R_0(Y)=\operatorname{Res}_X(p,Yp'-p''),\qquad
 R_t(Y)=\left.\partial_T\operatorname{Res}_X
 (p,(Y+TX)p'-p'')\right|_{T=0}.
\]
Write $R_0/\delta=Y^6+q_2Y^4+q_3Y^3+q_4Y^2+q_5Y+q_6$ and let
\begin{align*}
 E_1&=[Y^3]R_t,&E_2&=[Y^2]R_t-12[Y^4]R_0,\\
 E_3&=[Y]R_t-6[Y^3]R_0,&E_4&=[1]R_t-2[Y^2]R_0.
\end{align*}
These are necessary equations for a physical stationary shape after
any nonzero complex scaling. Indeed, $g_i=\kappa\overline{x_i}$ with
$\kappa=30/\sum_i|x_i|^2>0$ implies
$x_i=2\sum_{j\ne i}(g_i-g_j)^{-1}$.
Pairwise summation and Newton identities give
\begin{align*}
 E_1/\delta&=m_2,&E_2/\delta&=m_3+18q_2,\\
 E_3/\delta&=m_4+q_2m_2+24q_3,
 &E_4/\delta&=m_5+q_2m_3+q_3m_2+28q_4,
\end{align*}
and the four expressions vanish. Noncollision is precisely
$\delta\ne0$. Since $p$ is monic,
$\operatorname{Disc}(p)=(-1)^{15}\delta=-\delta$; hence the squared
pair-distance product is $|\delta|$, as used throughout.

\begin{proposition}[Complete necessary coverage when $e_2=0$]
\label{prop:six-e2-zero-independent}
Every collision-free six-point physical stationary shape with $e_2=0$,
up to complex scaling and permutation, has one of the following root
polynomials:
\[
 X(X^5-1),\qquad X^6+1,\qquad
 X^6-X^3+u,\quad 25u^2-17u+1=0.
\]
This coverage follows from explicit polynomial identities, independently
of the Faug\`ere--Svartz component catalogue.
\end{proposition}
\begin{proof}
First suppose $e_3\ne0$ and normalize $e_3=1$. Write
$p=X^6-X^3+cX^2+dX+u$. Direct expansion of the resultants gives four
polynomials $E_i(c,d,u)$ and $\delta(c,d,u)$. The attached exact
certificate supplies rational polynomials $A_i,B_i,C_i$, each of total
degree at most six, with the identities
\begin{equation}\label{eq:six-b1-combinations}
 c\delta=\sum_{i=1}^4 A_iE_i,\qquad
 d\delta=\sum_{i=1}^4 B_iE_i,\qquad
 (25u^2-17u+1)\delta=\sum_{i=1}^4 C_iE_i.
\end{equation}
The files \path{n6_n7_next_b1_equations.json} and
\path{n6_n7_next_b1_combinations.json} contain these full
polynomials (with $u$ named \texttt{e} in those files); the
standard-library verifier described below independently
reconstructs the equations and multiplies out all three identities.
Thus $E_i=0$ and $\delta\ne0$ force $c=d=0$ and
$25u^2-17u+1=0$. As a direct check on the last implication, after
$c=d=0$ one has
\[
 \delta=729u^2(4u-1)^3,\qquad
 E_3=-104976u(4u-1)(25u^2-17u+1).
\]

Next suppose $e_3=0$ but $e_4\ne0$. Normalize $e_4=1$ and write
$p=X^6+X^2+dX+u$. The necessary polynomials are explicitly
\begin{align*}
 E_1&=-4d(9375d^4-72000d^2u-291600u^4+120960u^2-256),\\
 E_2&=64(5625d^4-21600d^2u+58320u^4+10368u^2+256),\\
 E_3&=-720d(25d^2-96u)(135u^2+4),\\
 E_4&=96(84375d^4u-378000d^2u^2+3200d^2+259200u^3-10752u),
\end{align*}
and
\[
 \delta=-3125d^6+22500d^4u-43200d^2u^2-256d^2
        +46656u^5+13824u^3+1024u.
\]
The short polynomial identity
\begin{equation}\label{eq:six-e4-collision-identity}
 \delta=\frac d{96}E_1
 -\frac{175d^2-768u}{23040}E_2
 -\frac{du}{96}E_3-\frac{27u^2+4}{8640}E_4
\end{equation}
contradicts $\delta\ne0$ whenever the necessary equations hold.
Consequently this entire first-nonzero-$e_4$ chart is empty on the
noncollision locus.

Finally, if $e_3=e_4=0$, write $p=X^6+dX+u$. Direct expansion gives
\[
 \delta=-3125d^6+46656u^5,\qquad
 (E_1,E_2,E_3,E_4)
 =(1166400du^4,0,-2430000d^3u^2,8100000d^4u).
\]
Thus $E_1=0$ implies $d=0$ or $u=0$. They cannot both vanish because
$\delta\ne0$. Complex scaling normalizes these two alternatives to
$X^6+1$ or $X(X^5-1)$, respectively. The cases exhaust $e_2=0$.
\end{proof}

\begin{proposition}[Independent value comparison in the $e_2=0$ region]
\label{prop:six-special-value-independent}
Among physical stationary six-point shapes with $e_2=0$, normalized to
$\sum_i|z_i|^2=6$, the largest value is
$6^{15}/5^{10}$, attained only by the origin and a regular pentagon of
radius $\sqrt{6/5}$, up to rotation and permutation.
\end{proposition}
\begin{proof}
The two regular-polynomial charts give $6^{15}/5^{10}$ and $6^6$,
respectively. It remains to compare the two roots
$u_\pm=(17\pm3\sqrt{21})/50$ of the quadratic in the proposition.
For $p=X^6-X^3+u$, let $a,b$ be the roots of $t^2-t+u$.
The squared pair-distance product is
$|\delta|=729u^2|4u-1|^3$.

For $u_+>1/4$, $a,b$ are conjugate with modulus $\sqrt{u_+}$.
All six point roots therefore have squared radius $u_+^{1/3}$ and
moment $Q=6u_+^{1/3}$. The objective rescaled to moment six is
\[
 \left(\frac6Q\right)^{15}|\delta|
 =729\left(4-\frac1{u_+}\right)^3<729\cdot4^3=6^6.
\]
This comparison does not require this coefficient candidate to be
physical.

For the smaller root, evaluation of $25u^2-17u+1$ at $13/200$ and
$33/500$ gives $1/1600$ and $-131/10000$, respectively. The quadratic
is strictly decreasing on this interval, so
$13/200<u_-<33/500<1/4$. Here $a,b\in(0,1)$ and $a+b=1$.
Thus $Q=3(a^{2/3}+b^{2/3})>3$, and
\[
 \left(\frac6Q\right)^{15}|\delta|
 <2^{15}\,729\left(\frac{33}{500}\right)^2
                    \left(\frac{37}{50}\right)^3<6^6.
\]
The final strict comparison is the integer inequality
\[
 2^9\cdot33^2\cdot37^3=28242491904
 <31250000000=500^2\cdot50^3.
\]
Finally $6^{15}/5^{10}>6^6$ because
$6^9=10077696>9765625=5^{10}$. All alternatives except the central
pentagon are strictly smaller; its exact attainment follows directly
from its roots and discriminant. This proves the equality statement.
\end{proof}

\paragraph{Independent exact reconstruction of the certificate.}
The program \path{n6_n7_next_special_verify.py} uses only Python's
standard library and explicit checks, also under
\texttt{python -B -S -O}. It verifies the displayed resultant formulas
independently of their symbolic generation. Modulo monic $p$, replace
$(Y+TX)p'-p''$ by $Yp'-p''+T(Xp'-6p)$, whose degree in $X$ is at most
five. Its Sylvester matrix has size eleven. The parameter degrees of
its determinant are at most eleven in each of $c,d$, and at most six
in $u$ for the first variation: $u$ occurs in five copies of $p$ and
in at most one first-variation entry. Thus a $12\times12\times7$
Cartesian grid determines all coefficient polynomials exactly.

At each grid point the verifier computes $\delta$ by the integer
Bareiss determinant, verifies $\delta\ne0$, and computes $g=p''/p'$ in
$\mathbb Q[X]/(p)$ by the extended Euclidean algorithm. Newton root
sums give the exact traces $S_k,m_k$ and hence the four $E_i$ by the
displayed formulas. This reconstructs the full $e_3=1$ chart at
$1008$ exact points. The two smaller charts need $12\times7=84$
points each. The total is $1176$ exact points. Finally sparse rational
polynomial multiplication verifies
\eqref{eq:six-b1-combinations} and
\eqref{eq:six-e4-collision-identity} as formal identities. No numerical
decision or original component catalogue enters this proof.

\paragraph{Remaining coverage boundary.}
Only the $e_2=0$ part of the six-point coverage has been independently
reconstructed here. The general $e_2\ne0$ six-point charts and the
absent seven-point degree-$356$ physical exclusion still use the
published coverage input. These new identities reduce that dependence;
they do not remove it from the unrestricted global theorems.

% New fragment: independent n=7 coverage in the e_2=e_3=0 coefficient region.
\section{Independent completion of the sparse seven-point charts}
\label{sec:seven-sparse-independent}

The same discriminant-combination method independently closes the
seven-point region $e_2=e_3=0$, including its omitted first-nonzero
$e_4$ and $e_5$ charts. This section uses no published classification
as a coverage input for that region.

For a centered monic polynomial $p$ of degree seven with distinct roots
$x_i$, put $g_i=p''(x_i)/p'(x_i)$ and
$\delta=\operatorname{Res}(p,p')$. With
$Q=\sum_i|x_i|^2$, physical stationarity gives
\[
 F_i(x):=\sum_{j\ne i}(x_i-x_j)^{-1}
 =\frac{21}{Q}\overline{x_i},\qquad
 g_i=\frac{42}{Q}\overline{x_i}.
\]
Writing $\kappa=42/Q$, one gets
$F_i(g)=\kappa^{-1}\overline{F_i(x)}=x_i/2$, so
$x_i=2F_i(g)$. This remains true after any nonzero complex scaling.

Define the two resultants
\[
 R_0(Y)=\operatorname{Res}_X(p,Yp'-p''),\qquad
 R_t(Y)=\left.\partial_T\operatorname{Res}_X
 (p,(Y+TX)p'-p'')\right|_{T=0}.
\]
The five necessary coefficient equations are
\begin{align*}
 E_1&=[Y^4]R_t,&E_2&=[Y^3]R_t-20[Y^5]R_0,\\
 E_3&=[Y^2]R_t-12[Y^4]R_0,
 &E_4&=[Y]R_t-6[Y^3]R_0,\\
 E_5&=[1]R_t-2[Y^2]R_0.
\end{align*}
For clarity, let $S_k=\sum_i g_i^k$ and $m_k=\sum_i x_i g_i^k$.
Then $S_1=m_0=0$, $m_1=42$ and the double-force relation gives
\[
 m_2=0,\quad m_3=11S_2,\quad m_4=10S_3,\quad
 m_5=9S_4+S_2^2,\quad m_6=8S_5+2S_2S_3.
\]
Writing $R_0/\delta=Y^7+q_2Y^5+q_3Y^4+q_4Y^3+q_5Y^2+\cdots$,
synthetic division gives
\begin{align*}
 E_1/\delta&=m_2,&E_2/\delta&=m_3+22q_2,\\
 E_3/\delta&=m_4+q_2m_2+30q_3,\\
 E_4/\delta&=m_5+q_2m_3+q_3m_2+36q_4,\\
 E_5/\delta&=m_6+q_2m_4+q_3m_3+q_4m_2+40q_5.
\end{align*}
Newton identities make all five expressions zero at every physical
stationary shape. Noncollision is $\delta\ne0$. Here
$\operatorname{Disc}(p)=(-1)^{21}\delta=-\delta$, so the objective is
$|\delta|$ before moment rescaling.

\begin{proposition}[Complete sparse seven-point coverage]
\label{prop:seven-sparse-coverage}
Every collision-free physical stationary seven-point shape with
$e_2=e_3=0$ is, up to complex scaling and permutation, either a regular
heptagon or the origin together with a regular hexagon.
\end{proposition}
\begin{proof}
If $e_4\ne0$, normalize $e_4=1$ and write
$p=X^7+X^3+bX^2+cX+d$. Expansion of the five original necessary
equations yields a certificate of the formal identity
\begin{equation}\label{eq:seven-first-e4-collision}
 \delta=\sum_{i=1}^5 A_i(b,c,d)E_i(b,c,d),
 \qquad \deg A_i\le6.
\end{equation}
The complete rational multipliers are supplied in
\path{n6_n7_next_seven_a1_combination.json}, with the original
equations in \path{n6_n7_next_seven_a1_equations.json}.
The independent standard-library reconstruction and formal multiplication
described below verify the identity. Hence $E_i=0$ implies
$\delta=0$, and this entire chart has no noncollision candidate.

If $e_4=0$ but $e_5\ne0$, normalize $e_5=1$ and write
$p=X^7-X^2+cX+d$. In this case the short identity is
\begin{align}
 \delta={}&\frac{139799c}{27080060}E_1
 -\frac{2625912c^2+13009045d}{406200900}E_2\notag\\
 &-\frac{c(585288c^2+1303127d)}{162480360}E_3
 -\frac{d(53208c^2+85547d)}{92845920}E_4\notag\\
 &+\frac{3725}{129984288}E_5.
 \label{eq:seven-first-e5-collision}
\end{align}
This again excludes the entire noncollision chart. The full original
coefficient formulas and the same multipliers are supplied as the
\texttt{seven\_b1} JSON certificate pair.

Finally, if $e_4=e_5=0$, write $p=X^7+cX+d$. The original resultant
equations simplify to
\[
 \delta=46656c^7+823543d^6,\qquad
 E_1=29647548cd^5,
\]
\[
 E_2=0,\quad E_3=-130691232c^3d^3,\quad
 E_4=821487744c^4d^2,\quad E_5=-2352442176c^5d.
\]
Thus $c=0$ or $d=0$. They cannot both vanish because $\delta\ne0$.
Complex scaling normalizes the two alternatives to $X^7+1$ or
$X(X^6+1)$, respectively. Their roots have the asserted geometries,
and all cases in $e_2=e_3=0$ have now been covered.
\end{proof}

\begin{corollary}[Value comparison in the sparse seven-point region]
\label{cor:seven-sparse-value}
Among physical stationary seven-point shapes with $e_2=e_3=0$, at
moment seven, the largest value is $7^{21}/6^{15}$, attained only by
the origin and a regular hexagon of radius $\sqrt{7/6}$, up to rotation
and permutation.
\end{corollary}
\begin{proof}
The unit heptagon has moment seven and objective $7^7$.
The origin plus a unit hexagon has moment six and objective $6^6$,
so rescaling it to moment seven gives
$(7/6)^{21}6^6=7^{21}/6^{15}$. The strict comparison is
\[
 7^{14}=678223072849>470184984576=6^{15}.
\]
Both shapes satisfy the original force equations directly, and the
coverage proposition supplies no other physical candidate in this region.
\end{proof}

\paragraph{Independent certificate reconstruction.}
The verifier \path{n6_n7_next_seven_verify.py} uses only exact
standard-library integer and rational arithmetic. It imports elementary
polynomial arithmetic from the six-point verifier without running that
verifier or reading any six-point classification file.

Modulo monic $p$, replace $(Y+TX)p'-p''$ by
$Yp'-p''+T(Xp'-7p)$, of degree at most six in $X$.
The Sylvester matrix has size thirteen. The parameter degree bounds
are at most thirteen in $b,c$ and at most seven in $d$ for the first
variation: the constant $d$ occurs in six copies of $p$ and at most
one first-variation entry. Therefore $14\times14\times8$ distinct
Cartesian samples reconstruct each first-$e_4$ coefficient formula,
and $14\times8$ suffice for each of the two smaller charts.

At every sample, an integer Bareiss determinant gives $\delta$ and
is checked to be nonzero. Extended Euclidean arithmetic constructs
$g=p''/p'$ in $\mathbb Q[X]/(p)$; Newton root sums recover the exact
traces and the displayed $E_i/\delta$. This checks all $1792$
reconstruction points ($1568+112+112$). Sparse rational polynomial
multiplication then verifies
\eqref{eq:seven-first-e4-collision} and
\eqref{eq:seven-first-e5-collision} as formal identities. Thus the
coverage claim does not depend on a symbolic solver's output.

The check is reproducible offline with
\texttt{python -B -S -O n6\_n7\_next\_seven\_verify.py}. It does not
use floating-point decisions or the author component catalogue.
Only the region $e_2=e_3=0$ is independently closed here; the other
seven-point coefficient regions retain their previously stated coverage
inputs.

% Fragment compatible with main.tex; not a standalone document.
\section{A Gram deficit and an eight-point origin-annulus certificate}
\label{sec:eight-origin-annulus}

The failure of concavity for $S_7$ does not discard the angular information
in the Vandermonde determinant. The following deficit retains that information.
\begin{lemma}[Polynomial Gram deficit]
\label{lem:polynomial-gram-deficit}
Let $m\ge2$, $u_1,\ldots,u_m\in\C$, and set
\[
 A_{ij}=u_i^j,\qquad
 S_i=\sum_{j=1}^m|u_i|^{2j},\qquad
 Q_{ij}=\sum_{j=1}^m(u_i\overline{u_j})^j,
\]
\[
 H=\prod_iS_i,\qquad
 E=\sum_{i<j}|Q_{ij}|^2\prod_{k\ne i,j}S_k.
\]
Then
\begin{equation}\label{eq:polynomial-gram-deficit}
 |\det A|^2\le H-\binom m2^{-1}E.
\end{equation}
For $m=7$, one also has the fully polynomial inequality
\begin{equation}\label{eq:eight-gram-polynomial-cut}
 |\det A|^4H^5\le(H-E/21)^7.
\end{equation}
\end{lemma}
\begin{proof}
Suppose first that all $S_i>0$ and normalize the rows of $A$.
Their Gram matrix $K$ is positive semidefinite with $K_{ii}=1$,
so its nonnegative eigenvalues have sum $m$. Put
$X=\sum_{i<j}|K_{ij}|^2=E/H$.
The second elementary symmetric function of these eigenvalues is
$e_2(K)=\binom m2-X$, by the principal-minor formula.
Maclaurin's inequality gives
\[
 (\det K)^{2/m}\le1-X/\binom m2.
\]
Here $0\le X/\binom m2\le1$. Since $m/2\ge1$,
$(1-x)^{m/2}\le1-x$ on $[0,1]$. Multiplication by $H$ proves
\eqref{eq:polynomial-gram-deficit}. For $m=7$, square the sharper
Maclaurin inequality and multiply by $H^7$ to obtain
\eqref{eq:eight-gram-polynomial-cut}. Zero-radius cases follow
by continuity.
\end{proof}

\begin{lemma}[Seven-row trace-log deficit]
\label{lem:seven-trace-log-deficit}
If $K$ is a positive definite $7\times7$ Hermitian matrix with diagonal
entries one, then
\[
 \log\det K\le-\frac29\sum_{i<j}|K_{ij}|^2.
\]
\end{lemma}
\begin{proof}
Every eigenvalue lies in $(0,7]$ because their sum is seven.
For such $\lambda$,
\[
 \log\lambda\le\lambda-1-\frac19(\lambda-1)^2.
\]
Indeed the derivative of the difference between the right and left
sides is $(\lambda-1)(9-2\lambda)/(9\lambda)$. On $(0,1)$ the
difference decreases to zero. On $[1,7]$ its minimum is at an
endpoint, where its values are zero and $2-\log7>0$.
The latter inequality follows from the positive exponential series,
whose terms of orders zero through five already have sum greater
than seven at the argument two. Sum the scalar inequality and use
\[
 \sum_i(\lambda_i-1)=0,\qquad
 \sum_i(\lambda_i-1)^2=\operatorname{Tr}(K-I)^2
 =2\sum_{i<j}|K_{ij}|^2.
\]
\end{proof}

\begin{theorem}[Eight-point origin annulus]
\label{thm:eight-origin-annulus}
Suppose $z_1=0$ and $\sum_{i=1}^8|z_i|^2=8$. Put
\[
 t_i=\frac78|z_{i+1}|^2\quad(1\le i\le7),\qquad
 D=\sum_{i=1}^7(t_i-1)^2.
\]
If $49/50\le t_i\le51/50$ for every $i$, then
\begin{equation}\label{eq:eight-origin-annulus}
 \boxed{\displaystyle
 \Delta_8(z)\le\frac{8^{28}}{7^{21}}
               \exp\!\left(-\frac{17}{2700}D\right).}
\end{equation}
Equality in $\Delta_8\le8^{28}/7^{21}$ holds precisely when the
nonzero points form a regular heptagon of radius $\sqrt{8/7}$,
up to rotation and permutation. No centering, stationarity,
reflection, or angular ordering is assumed.
\end{theorem}
\begin{proof}
Let $u_i=\sqrt{7/8}\,z_{i+1}$, $S(t)=\sum_{j=1}^7t^j$, and
$A_{ij}=u_i^j$. Then $\sum_i t_i=7$ and
\[
 F:=\prod_i|u_i|^2\prod_{i<j}|u_i-u_j|^2=|\det A|^2.
\]
The unnormalized Gram matrix is $G_{ij}=S(u_i\overline{u_j})$.
Set $H=\prod_iS(t_i)$ and normalize to the correlation matrix
\[
 K_{ij}=\frac{S(u_i\overline{u_j})}{\sqrt{S(t_i)S(t_j)}}.
\]
Thus $F=H\det K$. If $K$ is singular, $F=0$ and the result is
immediate. Suppose henceforth that it is positive definite.

Write $R=S'^2-SS''$. Exact expansion gives
\begin{align*}
 R(1+x)={}&196x+1568x^2+5488x^3+11452x^4+16058x^5
             +16016x^6\\
          &+11648x^7+6188x^8+2352x^9+608x^{10}
             +96x^{11}+7x^{12}.
\end{align*}
For $|x|\le1/50$, the sum of absolute shifted coefficients yields
\[
 |R(1+x)|\le
 \frac{1121335142991470524807}{4882812500000000000}|x|
 <250|x|.
\]
Also $S(1+x)\ge7(49/50)^7>6$, so
\[
 |(\log S)''(1+x)|<7|x|.
\]
As $S(1)=7$ and $(\log S)'(1)=4$, integrating twice gives
\[
 |\log S(1+x)-\log7-4x|\le\frac76|x|^3.
\]
Consequently
\begin{equation}\label{eq:eight-annulus-radial-excess}
 \log H\le7\log7+\frac7{300}D,
\end{equation}
because $\sum_i(t_i-1)=0$ and $|t_i-1|\le1/50$.

For a complex $q$ with $r=|q|\in[49/50,51/50]$, the geometric
sum identity and the reverse triangle inequality give
\[
 |S(q)|=r\frac{|q^7-1|}{|q-1|}
 \ge\frac{r|r^7-1|}{r+1}
 \ge\frac{7(49/50)^7}{101/50}|r-1|
 >3|r-1|,
\]
with a non-strict bound when $r=1$. At $q=1$ this bound follows
directly, without dividing by $q-1$. For $q=u_i\overline{u_j}$,
the modulus $r=\sqrt{t_it_j}$ is in the same interval. Furthermore
$S(t_i)\le7(51/50)^7<81/10$. Hence
\begin{equation}\label{eq:eight-annulus-pair-deficit}
 |K_{ij}|^2\ge
 \frac9{((81/10)(101/50))^2}(t_it_j-1)^2
 \ge\frac1{30}(t_it_j-1)^2.
\end{equation}
The first rational coefficient is $250000/7436529>1/30$.

Put $d_i=t_i-1$. On the vector space indexed by unordered
pairs let $a_{ij}=d_i+d_j$ and $b_{ij}=d_id_j$.
The zero-sum condition gives the exact identity
\[
 \|a\|^2=\sum_{i<j}(d_i+d_j)^2=5D,
\]
while
\[
 \|b\|^2=\sum_{i<j}d_i^2d_j^2
 \le3(1/50)^2D.
\]
Therefore
\[
 \|a+b\|\ge(\sqrt5-\sqrt3/50)\sqrt D\ge2\sqrt D,
\]
where the last inequality follows already from
$\sqrt5>11/5$ and $\sqrt3<2$.
With $X=\sum_{i<j}|K_{ij}|^2$, equation
\eqref{eq:eight-annulus-pair-deficit} now implies
\[
 X\ge\frac1{30}\sum_{i<j}(t_it_j-1)^2\ge\frac2{15}D.
\]
Apply Lemma~\ref{lem:seven-trace-log-deficit} and
\eqref{eq:eight-annulus-radial-excess}:
\[
 \log F\le7\log7-\left(\frac4{135}-\frac7{300}\right)D
            =7\log7-\frac{17}{2700}D.
\]
Restoring the scale multiplies $F$ by $(8/7)^{28}$,
proving \eqref{eq:eight-origin-annulus}.

If $D>0$, the value is strictly below $8^{28}/7^{21}$.
If $D=0$, equality requires $\det K=1$, and the eigenvalue
arithmetic--geometric mean inequality forces $K=I$.
For distinct unit-modulus $u_i,u_j$, the equation
$S(u_i\overline{u_j})=0$ then forces
$(u_i\overline{u_j})^7=1$ and $u_i\overline{u_j}\ne1$.
The seven points are therefore a rotated set of seventh roots
of unity. Conversely that configuration attains equality.
\end{proof}

All polynomial expansions and rational constants above are checked
by \path{n8_verify.py}, including with Python \texttt{-S -O}.
The analytic Gram and Taylor inequalities are proved explicitly above.
The theorem excludes every angular arrangement in this annulus,
but does not cover its radial complement. In particular, a configuration
with an origin point and value greater than $8^{28}/7^{21}$ must have
some nonzero squared radius outside $[28/25,204/175]$.

For genuine full stationary configurations, a separate coefficient
reduction is also available. If an actual point is at zero and
$P(t)=tQ(t)$ with $Q(0)\ne0$, its force equation gives
$P''(0)=0$, equivalently $Q'(0)=0$. Centering removes the
coefficient of $t^7$, so the eight-point root polynomial has the form
\[
 P(t)=t^8+a t^6+b t^5+c t^4+d t^3+f t,\qquad f\ne0.
\]
One complex scaling chart leaves four complex coefficient variables.
This reduction concerns genuine full stationary configurations;
the force equation at a fixed origin must not be imposed without
justification on maximizers of an origin-constrained subclass.

For comparison, the earlier annulus argument used the following bounds.
For a genuine global eight-point maximizer with an origin point,
the subset constraints give $1/2<t_i<49/32$. The annulus theorem
does not cover this entire larger interval. The presently certified
unrestricted bounds remain
\[
\frac{8^{28}}{7^{21}}\le M_8<\left(\frac87\right)^{28}6^8,
\]
using the central heptagon and the existing strict deletion theorem.
Theorem~\ref{thm:noclass-origin8} now covers the entire origin branch.
The unrestricted value still requires the no-origin branch.
\begin{remark}[A limitation of the radial correlation relaxation]
The following counterexample concerns the older spectral constants
$2/9$ and the Maclaurin exponent $7/2$. The stronger seven-row
lemma in Section~\ref{sec:noclass-origin8} now overcomes this
particular obstruction on the full required radial domain.
For positive $a,b$ put $r=\sqrt{ab}$ and
\[
\beta(a,b)=\frac{ab(r^7-1)^2}{(1+r)^2S(a)S(b)},\qquad
Y=\sum_{i<j}\beta(t_i,t_j).
\]
The reverse triangle inequality gives $X\ge Y$ for the actual
normalized Gram correlations. This lower bound alone does not extend
the annulus proof to $1/2<t_i<49/32$. Indeed take
$t=(51/100,649/600,\ldots,649/600)$ with six equal outer entries.
Its sum is seven and all entries are strictly inside that interval.
Exact rational arithmetic gives $H/7^7>26/25$. For mixed pairs,
$37/50<r<1$, so substituting $37/50$ in
$(1-r^7)^2/(1+r)^2$ gives an upper bound and proves
$\beta(51/100,649/600)<1/60$.
The equal outer pairs give $\beta(649/600,649/600)<1/500$ directly.
Thus $Y<13/100$, and
\[
e^{2Y/9}<\frac{450}{437}<\frac{26}{25}<\frac H{7^7}.
\]
Even keeping the exact Maclaurin power leaves
\[
\frac H{7^7}(1-Y/21)^{7/2}
>\frac{26}{25}(1-13/600)=\frac{7631}{7500}>1.
\]
This is a counterexample to the proposed radial sufficient estimates,
not a configuration exceeding the central-heptagon value. The actual
angular correlations can be larger. The checker
\path{n8_relaxation_check.py} verifies all decisive rational bounds.
\end{remark}

% English section fragment; compatible with the existing theorem preamble.
\section{Eight-point stationary classification through three radial levels}
\label{sec:n8-two-radial-levels}

The following earlier structural argument is retained as an independent
cross-check. Theorem~\ref{thm:noclass-origin8} now gives the sharp
bound for all origin-containing configurations, including unequal radii.
The exact force equations at an actual origin impose more than radial
moment constraints: they tie the angular centroids at different radii
together. The following result classifies every genuine eight-point
stationary configuration with at most three distinct nonzero radii.
It therefore applies to global maximizers, but does not assert that a
maximizer of an origin-constrained subclass obeys the force equation
at its fixed origin.

Use $u_i=\sqrt{7/8}\,z_i$. A true eight-point stationary configuration
then satisfies
\[
 \sum_{j\ne i}\frac1{u_i-u_j}=4\overline u_i,
 \qquad \sum_i u_i=0,\qquad \sum_i|u_i|^2=7.
\]
The moment follows by multiplying the force equations by $u_i$ and
summing. Suppose $u_1=0$, and denote the seven nonzero roots by
$u_2,\ldots,u_8$.

\begin{lemma}[Inverse-centroid balance]\label{lem:n8-inverse-centroid}
Every such collision-free stationary configuration satisfies
\[
 \sum_{i=2}^8u_i=0,\qquad \sum_{i=2}^8\frac1{u_i}=0,
 \qquad\sum_{i=2}^8|u_i|^2u_i=0.
\]
If $t_i=|u_{i+1}|^2$ and $u_{i+1}=\sqrt{t_i}e^{i\theta_i}$, then
\begin{equation}\label{eq:n8-inverse-polygon}
 \sum_{i=1}^7\frac{t_i-1}{\sqrt{t_i}}e^{i\theta_i}=0,
 \qquad
 2\max_i\frac{|t_i-1|}{\sqrt{t_i}}
       \le\sum_i\frac{|t_i-1|}{\sqrt{t_i}}.
\end{equation}
If there are at most three distinct nonzero squared radii, then
the points on each radius separately have zero centroid.
\end{lemma}
\begin{proof}
Centering gives the first identity. The force equation at the origin
gives the second. Multiplying the force equations by $u_i^2$ and
summing by unordered pairs gives
\[
 4\sum_i|u_i|^2u_i
 =\sum_{i<j}\frac{u_i^2-u_j^2}{u_i-u_j}
 =7\sum_i u_i=0.
\]
Conjugating the origin force equation yields
$\sum u_i/|u_i|^2=0$; subtract it from the centering equation to obtain
\eqref{eq:n8-inverse-polygon}. The scalar inequality follows from the
triangle inequality by isolating its largest vector.
For radial groups of squared radii $t_j$, let $A_j$ be their vector
sums. The three identities give
\[
                  \sum_j A_j=\sum_j t_j A_j=\sum_j A_j/t_j=0.
\]
For three distinct positive radii $a,b,c$, the coefficient matrix has
rows $(1,1,1)$, $(a,b,c)$, and $(1/a,1/b,1/c)$. Its determinant is
$(b-a)(c-a)(c-b)/(abc)\ne0$, after multiplying its columns by
$a,b,c$. Thus every $A_j=0$. The same conclusion for one or two
levels follows from the first one or two independent equations.
\end{proof}

\begin{theorem}[Classification through three nonzero radial levels]
\label{thm:n8-two-level-exclusion}
A collision-free true eight-point stationary configuration containing
the origin and having at most three distinct nonzero radii is the
origin with a regular heptagon. Consequently every other
origin-containing true stationary configuration has at least four
distinct nonzero radial levels.
\end{theorem}
\begin{proof}
Let $P(x)=\prod_{i=1}^8(x-u_i)$. At a nonzero root of squared radius
$b$, stationarity is equivalent to
\begin{equation}\label{eq:n8-radius-divisibility}
                   xP''(x)-8bP'(x)=0.
\end{equation}
Thus the polynomial of all roots on that radius divides the expression
on the left.

Suppose first that there are two radii. By
Lemma~\ref{lem:n8-inverse-centroid}, each group is separately centered,
so neither group can be a singleton. Their sizes therefore split as
$2+5$ or $3+4$. For a group on a circle, zero centroid also implies
zero sum of reciprocals. Its monic root polynomial accordingly has
both its next-to-leading and its linear coefficient equal to zero.

In the $2+5$ case rotate the centered pair to $\pm\sqrt a$, $a>0$.
The five-point group has polynomial
\[
 Q(x)=x^5+cx^3+dx^2+f,\qquad f\ne0,
\]
and $P=x(x^2-a)Q$. Its squared radius is $b>0$.
The remainder of $xP''-8bP'$ upon division by $Q$ is
\begin{align*}
 &12(2b-3)d\,x^4
 +c\bigl[(26-16b)c+(18-16b)a\bigr]x^3\\
 &\quad+\bigl[(40b-50)f+((24-24b)a+(26-16b)c)d\bigr]x^2\\
 &\quad+f\bigl[(30-40b)a+(26-16b)c\bigr].
\end{align*}
All four coefficients must vanish. If $b=3/2$, the constant
coefficient gives $c=15a$, and the cubic coefficient becomes
$360a^2\ne0$, a contradiction. If $b\ne3/2$, the quartic coefficient
gives $d=0$. The quadratic coefficient then gives $b=5/4$, because
$f\ne0$. The constant coefficient gives $c=10a/3$, and the cubic
coefficient becomes $60a^2\ne0$. This also contradicts divisibility.

In the $3+4$ case the centered three-point circle group is a regular
triangle. Equivalently its root polynomial is $x^3-A$, $A\ne0$:
its quadratic and linear coefficients both vanish. The four-point
group has polynomial
\[
 Q(x)=x^4+Cx^2+F,\qquad F\ne0,
\]
so $P=x(x^3-A)Q$. If its squared radius is $b$, the remainder of
$xP''-8bP'$ upon division by $Q$ is
\begin{align*}
 &\bigl[(-44+32b)F+(26-16b)C^2\bigr]x^3
 +A(14-16b)C\,x^2\\
 &\qquad +(26-16b)CF\,x+AF(20-32b).
\end{align*}
The constant coefficient gives $b=5/8$ since $AF\ne0$.
The linear coefficient then gives $C=0$, and the cubic coefficient
is $-24F\ne0$. Thus this case is impossible as well.

Suppose next that there are exactly three nonzero radial levels.
Every group is centered by Lemma~\ref{lem:n8-inverse-centroid} and
must contain at least two points. Their seven points therefore split
only as $2+2+3$. The two centered pairs and centered triangle have
root polynomials $x^2-A$, $x^2-B$, and $x^3-C$, with $ABC\ne0$.
Thus
\[
                  P=x(x^2-A)(x^2-B)(x^3-C).
\]
Let the triangle squared radius be $c$. Reducing $xP''-8cP'$ modulo
$x^3-C$ gives exactly
\[
 24(c-1)C(A+B)x^2+12(3-2c)C^2x+12(1-2c)CAB.
\]
The linear coefficient requires $c=3/2$, while the constant requires
$c=1/2$. They cannot both vanish. Thus three levels are also impossible.

Finally suppose all seven nonzero points have one squared radius.
Their total moment is seven, so that radius is one. Put $Q=P/x$.
Equation~\eqref{eq:n8-radius-divisibility} at the seven distinct roots
shows that $xP''-8P'$ is divisible by $Q$. Its leading coefficient is
$-8$, so
\[
                      xP''-8P'=-8P/x.
\]
For the coefficient of $x^k$ in $P$, the resulting scalar multiplier
is $k(k-1)-8k+8=(k-1)(k-8)$. This forces every coefficient of degrees
$2$ through $7$ to vanish. Thus $P=x^8+fx$, $f\ne0$, whose seven
nonzero roots form a regular heptagon. Their unit radius requires
$|f|=1$.
\end{proof}

\begin{remark}[A reusable two-level stationary reduction]
For general $n$, normalize the moment to $n-1$ by
$u_i=\sqrt{(n-1)/n}\,z_i$. The force constant is $n/2$.
If the origin is a true stationary point and the other points have two
distinct squared radii $a,b$, each radial group has zero centroid.
Let their monic root polynomials be $A(x),B(x)$, of degrees $p,q$
with $p+q=n-1$. They have nonzero constant coefficients and zero
next-to-leading and linear coefficients. The exact necessary system is
\[
 P=xAB,\qquad pa+qb=n-1,
\]
\[
 A\mid xP''-naP',\qquad B\mid xP''-nbP'.
\]
Equivalently,
\[
 A\mid xA''B+2xA'B'+(2-na)A'B,
\]
with the analogous equation interchanging $A,a$ and $B,b$.
If $A=x^p+\alpha_1x^{p-1}+\cdots+\alpha_p$, its common-radius
property additionally gives
\[
 |\alpha_p|^2=a^p,\qquad
 \alpha_k=\alpha_p a^{k-p}\overline{\alpha_{p-k}}.
\]
The corresponding identities hold for $B$. These divisibility and
self-inversive equations form a smaller exact chart. The eight-point
contradictions above are special to its possible group sizes; they
do not establish a general two-level exclusion.
\end{remark}

The exact standard-library checker \path{n8_structure_verify.py}
verifies all three displayed polynomial remainders and the decisive
nonzero witnesses, including under Python \texttt{-S -O}.
This theorem independently covers the stated stationary subclasses.
The new fixed-origin theorem removes the former gap at four or more
radial levels. The unrestricted no-origin branch remains unresolved.


\subsection{A general three-level centroid theorem}

\begin{proposition}[Radial group centroids through three levels]
\label{prop:three-level-centroids-general}
For any $n\ge3$, a collision-free true stationary configuration
containing the origin and having at most three distinct nonzero
radial levels has zero centroid on each radial level separately.
In particular every level contains at least two points. Consequently
an origin-containing stationary configuration cannot have precisely
three nonzero radial levels with a singleton level.
\end{proposition}
\begin{proof}
For any nonzero force constant $\kappa$, centering and pairwise
summation of the force equations multiplied by $u_i^2$ give
\[
 \sum_i u_i=0,\qquad
 \kappa\sum_i |u_i|^2u_i=(n-1)\sum_i u_i=0.
\]
The force equation at the actual origin adds $\sum_{i\ne1}1/u_i=0$.
Conjugating it gives the third radial group equation
$\sum_j A_j/t_j=0$. The same invertible three-by-three matrix used
in Lemma~\ref{lem:n8-inverse-centroid} forces every group sum to zero.
A singleton at nonzero radius cannot have zero centroid.
\end{proof}

\begin{corollary}[Eight-point radial multiplicity gap]
\label{cor:n8-five-equal-radii}
A true stationary eight-point configuration containing the origin,
in which at least five of the seven nonzero points have the same
modulus, must be the origin with a regular heptagon. Every other
origin-containing stationary eight-point configuration has at least
four distinct nonzero radial levels, with no level containing more
than four points.
\end{corollary}
\begin{proof}
Five points on one circle leave at most two other radial levels.
Theorem~\ref{thm:n8-two-level-exclusion} therefore applies. At least
four distinct levels among seven nonzero points also directly limits
each level's population to four.
\end{proof}


% Fragment compatible with the main manuscript preamble.
\section{Executed eight-point coefficient resultants and sparse-chart exclusions}
\label{sec:eight-executed-coefficient-resultants}

At force-one scale, a centered collision-free stationary eight-point
configuration has root polynomial
\[
 P(t)=t^8+c_2t^6+c_3t^5+c_4t^4+c_5t^3+c_6t^2+c_7t+c_8
\]
and satisfies $P''(w_i)/(2P'(w_i))=\overline{w_i}$ at each root.
Define the exact resultant
\begin{equation}\label{eq:eight-executed-resultant}
 \mathcal R(u)=\operatorname{Res}_t(P(t),P''(t)-2uP'(t))
              =\sum_{j=0}^8R_ju^j.
\end{equation}
Its leading coefficient is $R_8=256\operatorname{disc}P\ne0$.
Stationarity necessarily implies
\begin{equation}\label{eq:eight-executed-dual-equations}
 R_7=0,\qquad R_{8-k}=R_8\overline{c_k}\quad(2\le k\le8).
\end{equation}
Indeed, the roots of $\mathcal R$ are the eight force values
$P''(w_i)/(2P'(w_i))$, with multiplicities. Equality with the conjugate
root polynomial gives the displayed identities.
The coefficient $R_7$ vanishes identically, since the sum of all forces
is zero by pairwise cancellation.

These are necessary equations. Equality of the unordered dual-root
sets alone does not establish actual conjugate pairing at individual
roots. A physical trace test or the pointwise force equations must
still filter their solutions.

The entire resultant in~\eqref{eq:eight-executed-resultant} has now been
generated with exact integer arithmetic, without assuming reflection
or polygon symmetry. It contains $11843$ monomials. The generator
\path{n8_algebra_full.py} uses the determinant of multiplication
by $P''-2uP'$ on $\mathbb Z[c_2,\ldots,c_8,u][t]/(P)$.
The machine-readable coefficients are in
\path{n8_algebra_full_resultant.json}; a text expansion is also
provided. Every coefficient $R_j$ has weighted degree $48+j$ when
$c_k$ has weight $k$. An independent verifier evaluates ten direct
Sylvester matrices using exact integer Bareiss determinants, rather
than the multiplication-matrix determinant used to generate the data.

\begin{theorem}[Complete stationary classification of centered octic trinomials]
\label{thm:eight-trinomial-classification}
Suppose a collision-free stationary eight-point configuration has
centered root polynomial
\[
 P(t)=t^8+A t^{8-k}+B,\qquad B\ne0,\quad 2\le k\le7.
\]
If $A\ne0$, then $k=4$. Every configuration in the $k=4$ chart,
allowing $A=0$, consists of two square rings and, up to rotation
and interchange of the rings, has precisely one of the following three shapes:
the regular octagon, two aligned unequal squares, or two staggered
unequal squares. At moment eight their respective products are
\begin{equation}\label{eq:eight-two-square-values}
 2^{24},\qquad \frac{2^{36}3^6}{5^{10}},\qquad
 \frac{2^{54}3^6}{7^{14}}.
\end{equation}
All are strictly less than the attainable central-heptagon value
$C_8=2^{84}/7^{21}$. Consequently no global eight-point maximizer
belongs to this trinomial class.
\end{theorem}
\begin{proof}
Absent coefficients of $P$ force the corresponding coefficients of
$\mathcal R$ to vanish. The exact specialized resultant gives the
following equations, each divided by a nonzero integer content.
For $k=3,5,7$, respectively, an absent dual coefficient is a nonzero
constant times $A^3B^5$, $A^5B^3$, or $A^6B$.
Since $B\ne0$, these immediately force $A=0$.

For $k=2$, two absent dual coefficients give
\[
 A^2B^5(-1215A^4+12544B)=0,
\]
\[
 AB^4(2278125A^8+78865920A^4B-1101463552B^2)=0.
\]
If $A\ne0$, set $x=A^4/B=12544/1215$. The second equation would
require $2278125x^2+78865920x-1101463552=0$, but exact substitution
gives $-3596615680/81\ne0$.
For $k=6$, the corresponding equations are
\[
 A^2B^2(-5643A^4+62720B^3)=0,
\]
\[
 A(729A^8+235008A^4B^3-2293760B^6)=0.
\]
Now $y=A^4/B^3=62720/5643$ would have to satisfy
$729y^2+235008y-2293760=0$, whose exact value is
$17836277760/43681\ne0$. This excludes every $k\ne4$.

For $k=4$, factor $P(t)=(t^4-x)(t^4-y)$ with $x,y\ne0$ and
$x\ne y$. Rotate so $x=R^2>0$, where $R$ is the squared radius
of the first ring; write $|y|=S^2$. At a point $w$ of the first ring,
the within-ring force is $3/(2w)$ and the other-ring force is
$4w^3/(w^4-y)$. Multiplying the force equation by $w$ gives
\begin{equation}\label{eq:eight-two-square-force}
 R=\frac32+\frac{4x}{x-y}.
\end{equation}
The left side is real, so $y/x$ is real. Hence
$y=\varepsilon S^2$ with $\varepsilon=1$ for aligned rings and
$\varepsilon=-1$ for rings staggered by $\pi/4$.
The full force-one moment is $28$, so $R+S=7$.
Put $X=R/S>0$. Equation~\eqref{eq:eight-two-square-force} becomes
\[
 3X^3-11X^2-11\varepsilon X+3\varepsilon=0.
\]
For aligned rings it factors as
$(X+1)(3X^2-14X+3)$, and for staggered rings as
$(X-1)(3X^2-8X+3)$.
The aligned positive roots give $(R-S)^2=98/5$.
The staggered roots give $(R-S)^2=0$ or $7$.
The zero difference gives the regular octagon. Conversely these
three choices directly satisfy the pointwise force equations on
both rings, so the classification is physical, not merely setwise.

The discriminant identity
\[
 |\operatorname{disc}P|
 =2^{16}(RS)^6|R^2-\varepsilon S^2|^8
\]
follows either from the two quartic discriminants and their resultant,
or by direct expansion. Scaling moment $28$ to moment eight multiplies
the squared-distance product by $(2/7)^{28}$.
Substituting the three differences above gives
\eqref{eq:eight-two-square-values} exactly. Their strict ordering is
\[
 \frac{2^{36}3^6}{5^{10}}<2^{24}
 <\frac{2^{54}3^6}{7^{14}}<\frac{2^{84}}{7^{21}}.
\]
All these comparisons are integer inequalities after clearing
denominators; the ratio of the largest two-square value to $C_8$ is
$3^6 7^7/2^{30}<1$.
\end{proof}

An actual origin point gives the smaller chart
\[
 P(t)=tQ(t),\qquad
 Q(t)=t^7+at^5+bt^4+ct^3+dt^2+f,\qquad f\ne0.
\]
Here $Q'(0)=0$ comes from the force equation at the origin.
Modulo $Q$, the force polynomial reduces to
\begin{align*}
 H(t,u)={}&56t^6+4aut^5+(30a+6bu)t^4
 +(20b+8cu)t^3\\
 &+(12c+10du)t^2+6dt+14fu.
\end{align*}
The exact origin resultant $\operatorname{Res}_t(Q,H)$ has
$759$ monomials; the full resultant specializes to $-2uf$ times
this one when $c_6=c_8=0$.

\begin{theorem}[Origin charts with at most two active coefficients]
\label{thm:eight-origin-two-coefficient-exclusion}
Every collision-free genuine stationary eight-point configuration
with an origin point, other than the central regular heptagon, has
at least three nonzero coefficients among $a,b,c,d$ in the displayed
origin chart.
\end{theorem}
\begin{proof}
Write the origin resultant as $\sum_{j=0}^7S_ju^j$.
Collision-freeness and $f\ne0$ give
$S_7=(-2)^7\prod_{Q(w)=0}P'(w)\ne0$.
The absent coefficient of $u$ in the dual polynomial forces $S_1=0$;
absent $a,b,c,d$ additionally force $S_5,S_4,S_3,S_2=0$, respectively.

For $\alpha\ne0$, replace $Q(t)$ by
$Q_\alpha(t)=\alpha^7Q(t/\alpha)$. Its coefficients become
$(\alpha^2a,\alpha^3b,\alpha^4c,\alpha^5d,\alpha^7f)$,
and its resultant coefficient $S_j$ becomes $\alpha^{42+j}S_j$.
Choose $\alpha$ with $\alpha^7f=1$. Every equation $S_j=0$ and
the condition $S_7\ne0$ are preserved. This use of scaling
concerns only those zero equations;
it does not assert preservation of the force-one normalization.

Allow one of the six pairs of coefficients to be active and write
its two variables as $(x,y)$. Let $F_1,F_2,F_3$ be the specialized
absent-coefficient rows in the following table, each divided by its
integer content, and let $L$ denote the specialized leading row $S_7$.
The stored exact rational polynomial identities directly give
\[
 x^pL=\sum_{i=1}^3q_iF_i,\qquad
 y^qL=\sum_{i=1}^3r_iF_i.
\]
The exponents and largest multiplier degrees are:
\begin{center}
\begin{tabular}{cccc}
\toprule
Active pair&Absent rows $(S_j)$&$(p,q)$&Multiplier degree\\
\midrule
$(a,b)$&$(1,3,2)$&$(1,1)$&10\\
$(a,c)$&$(1,4,2)$&$(1,1)$&10\\
$(a,d)$&$(1,4,3)$&$(2,2)$&9\\
$(b,c)$&$(1,5,2)$&$(3,2)$&7\\
$(b,d)$&$(1,5,3)$&$(1,1)$&10\\
$(c,d)$&$(1,5,4)$&$(1,1)$&9\\
\bottomrule
\end{tabular}
\end{center}
All twelve identities are supplied in the files
\path{n8_algebra_direct_witness_ab.json}, and analogously for
the other five pairs. The independent standard-library auditor
\path{n8_algebra_witness_audit.py} reconstructs the three
input rows and $L$ from the original resultant data, then expands
each identity exactly. Thus the proof relies on direct membership
in the original equation ideal, rather than an unverified claim that
a reported Gr\"obner basis belongs to that ideal.

At a stationary point the $F_i$ vanish and $L\ne0$, so the identities
force $x=y=0$. Any set of at most two active coefficients is contained
in one of these six charts; hence $a=b=c=d=0$.
The remaining polynomial $t(t^7+f)$ has a regular heptagon of roots
around the origin. Its force-one nonzero squared radius is $4$, and
scaling to moment eight gives $C_8$.
\end{proof}

The exact generated equations and these sparse-chart exclusions are
concrete progress toward eight-point exhaustion. They do not exhaust
the remaining mixed coefficient charts, including asymmetric physical
configurations, and therefore do not prove $M_8=C_8$.
The verifier \path{n8_algebra_verify.py} has effective exception
checks under Python \texttt{-S -O}; it independently checks the
Sylvester evaluations, all displayed elimination residuals, cubic
factorizations, and exact branch-value comparisons. The additional
\path{n8_algebra_witness_audit.py} has passed both normal and
\texttt{-S -O} execution, replaying all twelve original-ideal
exclusion identities without a computer algebra package.

% New fragment: independent stationary coverage in n=8, e_2=e_3=e_4=0.
\section{Independent completion of the sparse eight-point charts}
\label{sec:eight-sparse-independent}

The discriminant-combination certificates extend to the entire stationary
region $e_2=e_3=e_4=0$ for eight points. They independently exclude its
first-nonzero $e_5$ and $e_6$ charts and cover its remaining two shapes.
The result below concerns physical stationary shapes of the original
moment-constrained problem; it does not assert unrestricted eight-point
coverage.

For a centered monic polynomial $p$ of degree eight with distinct roots
$x_i$, set $Q=\sum_i|x_i|^2$, $\delta=\operatorname{Res}(p,p')$ and
$g_i=p''(x_i)/p'(x_i)$. Physical stationarity implies
\[
 F_i(x)=\sum_{j\ne i}(x_i-x_j)^{-1}
 =\frac{28}{Q}\overline{x_i},\qquad
 g_i=\frac{56}{Q}\overline{x_i}.
\]
Thus, writing $\kappa=56/Q$, one has
$F_i(g)=\kappa^{-1}\overline{F_i(x)}=x_i/2$.
In particular, the double-force system $x_i=2F_i(g)$ and its necessary
coefficient equations remain valid after any nonzero complex scaling.

Define
\[
 R_0(Y)=\operatorname{Res}_X(p,Yp'-p''),\qquad
 R_t(Y)=\left.\partial_T\operatorname{Res}_X
 (p,(Y+TX)p'-p'')\right|_{T=0}.
\]
The six necessary equations are
\begin{equation}\label{eq:eight-sparse-E-definition}
 E_1=[Y^5]R_t,\qquad
 E_r=[Y^{6-r}]R_t-(8-r)(7-r)[Y^{8-r}]R_0
 \quad(2\le r\le6).
\end{equation}
To verify necessity directly, put $S_k=\sum_i g_i^k$ and
$m_k=\sum_i x_i g_i^k$. Pairwise summation gives
\[
 S_1=m_0=0,\quad m_1=56,\qquad
 m_k=\sum_{a=0}^{k-1}S_{k-1-a}S_a-kS_{k-1}.
\]
There is also a direct generating-polynomial proof. Let
$H(Y)=\prod_i(Y-g_i)$; the $g_i$ are distinct at a physical shape.
The product formulas and $x_i=2F_i(g)$ give
\[
 \frac{R_t(Y)}\delta
 =\sum_i\frac{x_iH(Y)}{Y-g_i}
 =2\sum_{i<j}\frac{H(Y)}{(Y-g_i)(Y-g_j)}
 =H''(Y).
\]
Consequently $R_t=R_0''$. Taking its coefficients gives
$E_r=0$ for $1\le r\le6$; the term omitted from $E_1$ is zero because
$S_1=0$. This generating identity works in every point count.
Noncollision means $\delta\ne0$. Since $(-1)^{28}=1$, the monic
discriminant convention is $\operatorname{Disc}(p)=\delta$ for eight
points, and the squared pair-distance product is $|\delta|$.

\begin{proposition}[Independent sparse eight-point stationary coverage]
\label{prop:eight-sparse-coverage}
Every collision-free physical stationary eight-point shape with
$e_2=e_3=e_4=0$ is, up to complex scaling and permutation, either a
regular octagon or the origin together with a regular heptagon.
\end{proposition}
\begin{proof}
If $e_5\ne0$, normalize $e_5=1$ and write
$p=X^8-X^3+bX^2+cX+d$. The attached full rational polynomials certify
\begin{equation}\label{eq:eight-first-e5-collision}
 \delta=\sum_{i=1}^6 A_i(b,c,d)E_i(b,c,d),\qquad \deg A_i\le6.
\end{equation}
They are supplied in
\path{n6_n7_next_eight_a1_combination.json}, with the six original
coefficient equations in the corresponding
\path{eight_a1_equations.json} file. Independent standard-library
reconstruction and formal multiplication verify this identity. It implies
$\delta=0$ whenever the necessary equations hold, excluding the
entire noncollision chart.

If $e_5=0$ but $e_6\ne0$, normalize $e_6=1$ and write
$p=X^8+X^2+cX+d$. This chart has the shorter identity
\begin{align}
 \delta={}&-\frac{68399c}{13711104}E_1
 -\frac{14085295c^2-68264064d}{2303465472}E_2\notag\\
 &+\frac{c(424879c^2-1206256d)}{109688832}E_3
 -\frac{d(75313c^2+144320d)}{383910912}E_4\notag\\
 &-\frac{203cd^2}{1713888}E_5-\frac3{1938944}E_6.
 \label{eq:eight-first-e6-collision}
\end{align}
The original equations and the same multipliers are supplied in the
\texttt{eight\_b1} JSON pair and independently verified. This entire
chart is likewise excluded on $\delta\ne0$.

If $e_5=e_6=0$, the polynomial is $p=X^8+cX+d$. Direct expansion of
the original equations gives
\[
 \delta=-823543c^8+16777216d^7,\qquad
 E_1=822083584cd^6,\qquad E_2=0,
\]
\[
 \begin{aligned}
 E_3&=-6294077440c^3d^4,& E_4&=57276104704c^4d^3,\\
 E_5&=-300699549696c^5d^2,& E_6&=863547424768c^6d.
 \end{aligned}
\]
Thus $c=0$ or $d=0$, and noncollision excludes their simultaneous
vanishing. The two alternatives normalize to $X^8+1$ and
$X(X^7-1)$ by complex scaling. This covers all the remaining cases.
\end{proof}

\begin{corollary}[Exact comparison of the sparse stationary shapes]
\label{cor:eight-sparse-value}
At moment eight, every physical stationary shape with
$e_2=e_3=e_4=0$ has objective at most $8^{28}/7^{21}$.
Equality occurs precisely for the origin and a regular heptagon of
radius $\sqrt{8/7}$, up to rotation and permutation.
In particular, any unrestricted global maximizer satisfying these
coefficient conditions must have that shape and value.
\end{corollary}
\begin{proof}
The regular unit octagon has objective $8^8$ and moment eight.
The origin with a unit heptagon has objective $7^7$ and moment seven,
so its rescaled value is $(8/7)^{28}7^7=8^{28}/7^{21}$.
These two shapes satisfy the original force equations directly, and
the strict comparison is
\[
 8^{20}=1152921504606846976
 >558545864083284007=7^{21}.
\]
The coverage proposition and this strict gap prove the equality statement.
The central heptagon is available as a competitor in the unrestricted
problem, so a global maximizer in this stationary region cannot be the
smaller octagon.
\end{proof}

\paragraph{Exact independent reconstruction.}
The program \path{n6_n7_next_eight_verify.py} imports only the
standard-library integer and rational polynomial routines from the
six-point checker. It does not execute that checker's main routine or
read any six-point classification data.

Modulo monic $p$, replace $(Y+TX)p'-p''$ by
$Yp'-p''+T(Xp'-8p)$, whose degree in $X$ is at most seven.
The Sylvester matrix has size fifteen. Its parameter degree bounds are
at most fifteen in $b,c$ and at most eight in the constant $d$ for
the first variation. Hence $16\times16\times9$ Cartesian samples
determine each first-$e_5$ coefficient formula, and $16\times9$ suffice
for each smaller chart.

At all $2592$ samples ($2304+144+144$), the checker obtains $\delta$
by an integer Bareiss determinant and verifies $\delta\ne0$.
Extended Euclidean arithmetic constructs $g=p''/p'$ in
$\mathbb Q[X]/(p)$; Newton root sums recover the exact traces and all
six original resultant coefficients. Sparse rational polynomial
multiplication then verifies
\eqref{eq:eight-first-e5-collision} and
\eqref{eq:eight-first-e6-collision} as formal identities.
The offline command is
\texttt{python -B -S -O n6\_n7\_next\_eight\_verify.py}.
No floating-point decision, symbolic solver output or author component
list is a trusted coverage input.

This result supplies a complete stationary classification only in
$e_2=e_3=e_4=0$. An unrestricted eight-point proof must still address
maximizers in the other coefficient regions.

% Compatible with the existing theorem preamble and n8 algebra fragment.
\section{No-origin eight-point candidates with at most two radial levels}
\label{sec:n8-noorigin-two-levels}

The origin classification leaves configurations with no actual origin
point. The following argument excludes all global eight-point
maximizers in this second class having at most two nonzero radial
levels. Its $2+6$ and $3+5$ exclusions apply to all true stationary
configurations; the $4+4$ rectangle exclusion explicitly uses global
subset spread and is not asserted for every stationary configuration.
Continue to use $u=\sqrt{7/8}\,z$, total moment seven and force
constant four.

\begin{lemma}[A general holomorphic radial moment identity]
\label{lem:holomorphic-radial-two}
At any true stationary configuration normalized to moment $n-1$
and force constant $n/2$, put $s_2=\sum_i u_i^2$. Then
\[
 \sum_i|u_i|^2u_i=0,\qquad
 \sum_i|u_i|^2u_i^2=\frac{2n-3}{n}s_2.
\]
Consequently a no-origin stationary configuration having exactly two
nonzero radial levels has zero centroid on each level.
\end{lemma}
\begin{proof}
The first identity was proved by summing the force equations times
$u_i^2$. For the second, multiply them by $u_i^3$ and sum by pairs:
\[
 \frac n2\sum_i|u_i|^2u_i^2
 =\sum_{i<j}(u_i^2+u_iu_j+u_j^2)
 =\left(n-\frac32\right)s_2,
\]
since $\sum_i u_i=0$ implies $\sum_{i<j}u_i u_j=-s_2/2$.
If the two radial group sums are $A,B$ at distinct squared radii
$a,b$, centering and the first identity give $A+B=aA+bB=0$.
Hence $A=B=0$.
\end{proof}

\begin{theorem}[Exclusion of no-origin global candidates with two levels]
\label{thm:n8-noorigin-two-levels}
A global moment-eight maximizer with no origin point has at least
three distinct nonzero radial levels. More precisely, true no-origin
stationary configurations cannot have two distinct radial levels
with populations $2+6$ or $3+5$. For population $4+4$, any global
candidate is forced into the two-square family, whose values are
strictly below the attainable center-heptagon value.
\end{theorem}
\begin{proof}
If there is one radial level, its squared radius in $u$ coordinates
is $7/8$. The root force equation gives
$xP''-7P'=0$ at all eight roots. This polynomial has degree below
eight and is therefore identically zero. It follows that
$P=x^8+f$, so the configuration is a regular octagon. At moment eight
its value $8^8$ is strictly below $C_8=8^{28}/7^{21}$.

For two levels, Lemma~\ref{lem:holomorphic-radial-two} centers each
group. Each group has at least two points, so the populations split
only as $2+6$, $3+5$, or $4+4$. On a common circle, zero centroid
also implies zero reciprocal sum, hence zero next-to-leading and
linear coefficients in the monic group polynomial.

\paragraph{The $2+6$ stationary class.}
Rotate the pair to $\pm\sqrt a$, $a>0$, and write
\[
 Q=x^6+cx^4+dx^3+ex^2+f,\qquad f\ne0,
 \qquad P=(x^2-a)Q.
\]
Let $b$ be the six-point squared radius, so $2a+6b=7$ and
$0<b<7/6$. The exact remainder of $xP''-8bP'$ modulo $Q$ is
\begin{align*}
 &\bigl[(16b-26)c-(30-48b)a\bigr]x^5+12(2b-3)dx^4\\
 &\quad+\bigl[(32b-44)e-(12-32b)ac\bigr]x^3
 +(24b-6)ad\,x^2\\
 &\quad+\bigl[(48b-54)f-(2-16b)ae\bigr]x.
\end{align*}
Its quartic and quadratic coefficients force $d=0$: otherwise they
would require both $b=3/2$ and $b=1/4$.
The first and cubic equations give
\[
 c=\frac{30-48b}{16b-26}a,\qquad
 e=\frac{12-32b}{32b-44}ac.
\]
The denominators are nonzero since $b<7/6$. If $c=0$, these equations
give $b=5/8$ and $e=0$, while the linear equation gives $f=0$,
a contradiction. Thus $c\ne0$. Both $c,e$ are real. The circle
self-inversive identity is $e=fc/b^2$, so $f$ is real as well.
Since $|f|=b^3$, write $f=\sigma b^3$, $\sigma\in\{1,-1\}$, and
$e=\sigma bc$.

Substitute $a=(7-6b)/2$ into the cubic equation. For the two signs
one obtains respectively
\[
 Q_+(b)=32b^2-52b+21=0,\qquad
 Q_-(b)=64b^2-96b+21=0.
\]
The linear equation, after eliminating $c$ and clearing its nonzero
denominator, gives
\[
 E(b)=(48b-54)b^2(16b-26)
       -(2-16b)\left(\frac{7-6b}{2}\right)^2(30-48b)=0.
\]
Exact Euclidean division yields
\[
 E\bmod Q_+=336-384b,\qquad
 E\bmod Q_-=462-432b.
\]
The positive-sign case forces $b=7/8$ and thus $a=b$, contradicting
two distinct levels. The negative-sign case forces $b=77/72$, but
$Q_-(77/72)=-686/81\ne0$. Thus the distinct-level $2+6$ class is
impossible for true stationary configurations.

\paragraph{The $3+5$ stationary class.}
The triangle has polynomial $x^3-A$, $A\ne0$, and the five-point
group has $Q=x^5+cx^3+dx^2+f$, $f\ne0$. Let its squared radius be
$b$. For $P=(x^3-A)Q$, the remainder of $xP''-8bP'$ modulo $Q$ is
\begin{align*}
 &\bigl[(24b-36)d-(20-40b)A\bigr]x^4+(26-16b)c^2x^3\\
 &\quad+\bigl[(40b-50)f-(6-24b)Ac+(26-16b)cd\bigr]x^2\\
 &\quad +(16b-2)Ad\,x+(26-16b)cf.
\end{align*}
The moment $3a+5b=7$, $a>0$, gives $b<7/5<13/8$.
The constant coefficient therefore gives $c=0$.
The quadratic coefficient then gives $b=5/4$, the linear coefficient
gives $d=0$, and the quartic coefficient is $30A\ne0$.
This excludes the true stationary $3+5$ class.

\paragraph{The $4+4$ global class.}
Write the two group polynomials as
\[
 A=x^4+Cx^2+F,\qquad B=x^4+Dx^2+G,
\]
with squared radii $a,b$ and $a+b=7/4$. The global subset-spread
bound gives $4a>12/8$ and $4b>12/8$, so
\begin{equation}\label{eq:n8-fourfour-global-domain}
                       \frac38<a,b<\frac{11}8.
\end{equation}
Since the group second holomorphic moments are $-2C$ and $-2D$,
Lemma~\ref{lem:holomorphic-radial-two} gives
\[
                 (a-13/8)C+(b-13/8)D=0.
\]
If one of $C,D$ is zero, so is the other. These groups are then two
regular squares; their strictly sub-$C_8$ values are already classified
by Theorem~\ref{thm:eight-trinomial-classification}.

Suppose instead $C,D\ne0$. By
\eqref{eq:n8-fourfour-global-domain}, they are opposite positive-real
multiples of one another. The circle identities
$C=(F/a^2)\overline C$, $D=(G/b^2)\overline D$ allow a common rotation
making $F=a^2$, $G=b^2$ and $C,D$ real. For $P=AB$, the remainder
of $xP''-8bP'$ modulo $B$ is
\begin{align*}
 &\bigl[(12-32b)F+(32b-44)G+(16b-18)CD+(26-16b)D^2\bigr]x^3\\
 &\quad+\bigl[(32b-28)CG+(2-16b)DF+(26-16b)DG\bigr]x.
\end{align*}
Use $b=7/4-a$ and $D=-(13-8a)C/(8a-1)$ in its linear coefficient.
After division by $C\ne0$ and multiplication by $8a-1>0$, it factors
exactly as
\[
                  (a-7/8)(64a^2-112a+7)=0.
\]
On the global interval \eqref{eq:n8-fourfour-global-domain},
\[
 64a^2-112a+7=64(a-7/8)^2-42<16-42<0.
\]
Thus $a=7/8=b$, contradicting distinct radial levels. This excludes
the nonsquare $4+4$ global candidate. The use of global subset spread
is essential here; the argument does not exclude every stationary
rectangle configuration outside that interval.
\end{proof}

The standard-library rational checker
\path{n8_structure_noorigin_verify.py} independently verifies
all three remainders, both $2+6$ Euclidean remainders, the $3+5$
contradiction, the $4+4$ factorization, and the already-known
strict two-square value comparison. It passes under Python
\texttt{-S -O} and imports only the local standard-library sparse
polynomial arithmetic in \path{n8_structure_verify.py}.
Together with Theorem~\ref{thm:noclass-origin8}, any strictly better
global eight-point candidate must have no origin and at least three
distinct nonzero radial levels. This remaining class has not been
exhausted, so $M_8=C_8$ is not yet established.

\subsection{An additional centered three-level stationary exclusion}

\begin{proposition}[Three no-origin radial group centroids]
\label{prop:noorigin-three-centroids}
At any true stationary configuration with exactly three distinct
nonzero squared radii $a,b,c$, the three radial group centroid sums
are collinear vectors. If one of those sums vanishes, they all vanish.
For $n=8$, the true stationary subclass in which all three sums vanish
cannot have radial populations $2+3+3$. Its remaining population split is $2+2+4$, whose even root polynomial
is treated in Proposition~\ref{prop:eight-centered-threelevel-224}.
\end{proposition}
\begin{proof}
The group sums $A_1,A_2,A_3$ satisfy
$\sum A_j=\sum t_jA_j=0$ by centering and
Lemma~\ref{lem:holomorphic-radial-two}. This two-by-three real matrix
has rank two and a one-dimensional real kernel, so the three complex
vectors are real multiples of one another. Distinct $t_j$ also imply
that a zero group sum makes the other two sums vanish.

When all group sums vanish, each radial group contains at least two
points. Eight points in three groups split only as $2+3+3$ or $2+2+4$.
In the first case the centered pair and the centered triangles have
root polynomial
\[
                   P=(x^2-A)(x^3-B)(x^3-C),\qquad ABC\ne0.
\]
Let $b$ be the squared radius of the $B$ triangle. The exact remainder
of $xP''-8bP'$ modulo $x^3-B$ is
\[
 6A\bigl[4(b-1)B+(1-4b)C\bigr]x^2
 +6B\bigl[2(3-2b)B+(4b-3)C\bigr]x.
\]
Both coefficients must vanish. The coefficient matrix for $(B,C)$
has determinant
\[
       4(b-1)(4b-3)-2(1-4b)(3-2b)=6,
\]
so $B=C=0$, a contradiction. For the $2+2+4$ split, centered pairs
and a centered common-circle quadrilateral each have even root
polynomials; their product is therefore even. The next proposition excludes this chart under the necessary global
radius bounds. General uncentered three-level configurations remain
outside this centered-group argument.
\end{proof}

\subsection{Excluding the centered three-level pattern \texorpdfstring{$2+2+4$}{2+2+4}}
\label{sec:eight-centered-threelevel-224}

\begin{proposition}[The $2+2+4$ three-level exclusion]
\label{prop:eight-centered-threelevel-224}
Let $u_1,\ldots,u_8$ be a collision-free physical stationary
configuration satisfying
\[
 \sum_{j\ne i}\frac1{u_i-u_j}=4\overline{u_i},\qquad
 \sum_i|u_i|^2=7.
\]
Suppose it has three distinct squared radii, with multiplicities
$2,2,4$, and each radius group has zero centroid. Write $c$ for the
squared radius of the four-point group. If
\[
                     \frac38<c<\frac{49}{32},
\]
then
\[
             \Delta_8(u)<620000<7^7.
\]
Consequently this pattern cannot furnish a global eight-point maximum
whenever the previously established global radius cuts apply.
\end{proposition}

\begin{proof}
The two-point groups are antipodal pairs. For the four-point group,
the vanishing root sum also makes its third elementary symmetric
coefficient zero: it is the root product times the sum of reciprocals,
and the latter is the conjugate root sum divided by $c$. Its polynomial
is therefore even, giving a rectangle on its circle. Rotate so that its root polynomial is
$x^4+Cx^2+c^2$, with $C\in\R$ and $C^2<4c^2$. Thus
\[
 P(x)=(x^2-A)(x^2-B)(x^4+Cx^2+c^2),\qquad
 q(y)=(y-A)(y-B)(y^2+Cy+c^2).
\]
Put $a=|A|$, $b=|B|$, $S=A+B$, $T=AB$.
The three radii $a,b,c$ are distinct and
$a+b+2c=7/2$.
Since $P(x)=q(x^2)$, define $H=q'+2yq''$.
The force equation at a root with $x^2=y$ reads
\[
                      H(y)=8|x|^2q'(y).
\]
For the rectangular layer, $H-8cq'$ has zero remainder modulo
$y^2+Cy+c^2$. Its two coefficients are
\begin{align*}
E_1={}&(13-8c)C^2+(16c-22)c^2+(9-8c)CS+(6-16c)T,\\
E_2={}&(13-8c)Cc^2+(14-16c)c^2S+(1-8c)CT.
\end{align*}
Both vanish. The imaginary parts are a real linear system in
$\operatorname{Im}S,\operatorname{Im}T$, with determinant
\[
 D=(9-8c)(1-8c)C^2-(6-16c)(14-16c)c^2.
\]
The decisive polynomial identity is
\begin{align*}
 &(9-8c)\bigl[(1-8c)CE_1-(6-16c)E_2\bigr]\\
 &\hspace{1cm}-D\bigl[(9-8c)S+(13-8c)C\bigr]=192Cc^2.
\end{align*}
If $D=0$, then $C=0$. The determinant formula then forces
$c=3/8$ or $7/8$. At $c=3/8$, $E_1=-16c^2\ne0$.
At $c=7/8$, $E_1=0$ gives $T=-c^2$, so $ab=c^2$; the moment
condition gives $a+b=2c$. Equality in the elementary arithmetic--
geometric mean inequality forces $a=b=c$, contrary to three distinct
radii. Hence $D\ne0$ and $S,T$ are real. As roots of the real
quadratic $y^2-Sy+T$, the numbers $A,B$ are either both real or a
conjugate pair. The latter would give $a=b$, so $A,B$ are real.
This conclusion was obtained without assuming reflection symmetry.

There are now two cases, according to the signs of $A,B$.
Multiplying all $u_i$ by $\mathrm i$ changes both signs and $C$'s
sign, so the equal-sign case may take $A,B>0$, and the opposite-sign
case may take $A>0>B$ after exchanging the pairs.
The mixed moment identity is
\begin{equation}\label{eq:eight-224-mixed-moment}
 (a-13/8)A+(b-13/8)B-(c-13/8)C=0.
\end{equation}
To derive it directly, multiply the force equation by $u_i^3$ and
sum. Pairwise cancellation gives
$4\sum_i|u_i|^2u_i^2=(13/2)\sum_i u_i^2$, since the full
configuration is centered. Summation over the three groups gives
\eqref{eq:eight-224-mixed-moment}.

For equal signs, put
\[
 S=7/2-2c,\qquad
 T=\frac{S^2-(13/8)S-(c-13/8)C}{2}.
\]
Substitute these expressions in $E_1,E_2$, clear constant
denominators, and take their resultant in $C$. It is
$288(8c-13)Q_s(c)$, where
\begin{align*}
Q_s(c)={}&2097152c^7-14942208c^6+42336256c^5\\
 &-59654144c^4+42795008c^3-14368256c^2\\
 &+1840440c-77175.
\end{align*}
The factor $8c-13$ cannot vanish in the stipulated range.
A linear combination canceling the $C^2$ coefficients then fixes
\begin{align*}
414946C={}&65536c^6-19800064c^5+106128384c^4\\
 &-212024448c^3+183825488c^2-61502182c+4096155
\end{align*}
at every relevant root of $Q_s$. The nonzero coefficient of this
linear combination is invertible modulo $Q_s$, as checked by the
extended Euclidean algorithm. Thus this necessary parametrization
does not rely on a reported Gr\"obner basis.

For opposite signs, put $d=A-B=7/2-2c$. Equation
\eqref{eq:eight-224-mixed-moment} becomes
\[
 (15-16c)S=(8c-13)C.
\]
The value $c=15/16$ first forces $C=0$ by the mixed moment equation;
then $E_2=-c^2S=0$ forces $S=0$. Substituting
$T=-d^2/4$ in $E_1$ gives $-27/128\ne0$.
Otherwise
\[
 S=\frac{(8c-13)C}{15-16c},\qquad T=\frac{S^2-d^2}{4}.
\]
Here $C=0$ would imply $S=0$ and $a=b$, so $C\ne0$.
After substitution and clearing powers of $15-16c$, the equations
$E_1=0,E_2/C=0$ are linear in $V=C^2$. Eliminating $V$ gives,
apart from nonvanishing factors, the polynomial
\begin{align*}
Q_m(c)={}&32768c^6-225280c^5+596480c^4-767040c^3\\
 &+491088c^2-139020c+11025.
\end{align*}
The corresponding value is
\begin{align*}
23716C^2={}&-4718592c^6+29196288c^5-69335040c^4\\
 &+81177216c^3-48520256c^2+13181280c-981225.
\end{align*}
Its equality with the value forced by an original linear row is
verified by inversion modulo $Q_m$, rather than ideal-membership
claims about a computed basis.

For either case, set
\[
 R=(T-c^2)^2+CS(T+c^2)+c^2S^2+C^2T.
\]
The elementary product formula for the discriminant gives
\begin{equation}\label{eq:eight-224-disc-product}
 \Delta_8(u)=256|T|c^2(C^2-4c^2)^2(S^2-4T)^2R^4.
\end{equation}
It is important that this is the absolute discriminant itself.
Reduce its right side modulo $Q_s$ or $Q_m$, using $|T|=T$
in the equal-sign case and $|T|=-T$ in the opposite-sign case.
The exact rational residue coefficients are supplied with the
independent certificate.

Rational Sturm counting finds exactly two roots of $Q_s$ and three
of $Q_m$ in $(3/8,49/32)$. The following table gives intervals in
units of $10^{-12}$ and exact rational interval conclusions:
\begin{center}
\begin{tabular}{lll}
\toprule
Case&$10^{12}c$ interval&Certified consequence\\
\midrule
Equal signs&$(552295886410,552295886411)$&$C^2>4c^2$\\
Equal signs&$(1060148224689,1060148224690)$&$0<\Delta_8(u)<600000$\\
Opposite signs&$(594183633638,594183633639)$&$C^2>4c^2$\\
Opposite signs&$(1012057168860,1012057168861)$&$C^2<0$\\
Opposite signs&$(1171372842247,1171372842248)$&$0<\Delta_8(u)<620000$\\
\bottomrule
\end{tabular}
\end{center}
The first and third rows contradict a genuine rectangle, and the
fourth contradicts real $C$. The two remaining values are strictly
below $7^7=823543$. The certificate
\path{n8_algebra_threelevel_224_audit.py} uses only rational
standard-library arithmetic. It verifies the Sylvester resultant,
the original rectangle rows, the determinant identity, the direct row
combinations and inverses, the discriminant reduction, all Sturm
coverage counts, and all table inequalities. Its 42 checks pass with
explicit exceptions both normally and under \texttt{-S -O}.

Finally $u=\sqrt{7/8}\,z$ gives
$\Delta_8(z)=(8/7)^{28}\Delta_8(u)$, whereas the normalized
central heptagon has value
$C_8=(8/7)^{28}7^7=2^{84}/7^{21}$.
This proves the stated exclusion.
\end{proof}

This proposition concerns only the centered $2+2+4$ pattern with
three distinct radii and the stated global radius range. It does not
classify every even eight-point polynomial or the remaining
noncentral layer patterns, and does not by itself prove $M_8=C_8$.

\begin{corollary}[All remaining three-level centroids are nonzero]
\label{cor:n8-threelevel-nonzero-centroids}
A no-origin global eight-point maximum with exactly three distinct
nonzero radial levels has three nonzero, collinear radial-group
centroid sums. In particular the case where each group is centered
has been completely excluded.
\end{corollary}
\begin{proof}
Proposition~\ref{prop:noorigin-three-centroids} gives collinearity
and shows that one zero group sum would make all three sums zero.
In that event the split $2+3+3$ is already impossible, leaving $2+2+4$.
For its four-point group, Theorem~\ref{thm:subset-spread-new} gives
$4c>4\cdot3/8=3/2$ in the moment-seven scale; the global point-radius
bound gives $c<49/32$. Thus Proposition~\ref{prop:eight-centered-threelevel-224}
applies and yields a value strictly below the central heptagon.
This contradicts global maximality.
\end{proof}

% October 7 fragment: genuinely new constraints on nonzero three-shell centroids.
\section{Three radial levels with nonzero centroid sums}
\label{sec:n8-oct7-three-levels}

We use the normalization $\sum_i|u_i|^2=n-1$ and force constant
$n/2$; for eight points these are moment seven and force constant four.
The results below do not assume reflection symmetry. They derive it
in specified population classes from true stationarity alone.

\begin{lemma}[The holomorphic radial moment recursion]
\label{lem:n8-oct7-moment-recursion}
At a collision-free true stationary configuration, let
$p_r=\sum_i u_i^r$, $p_0=n$, and $t_i=|u_i|^2$. For every integer
$r\ge1$,
\begin{equation}\label{eq:n8-oct7-moment-recursion}
 \sum_i t_i u_i^r
 =\frac{2n-r-1}{n}p_r+
   \frac1n\sum_{\ell=1}^{r-1}p_\ell p_{r-\ell}.
\end{equation}
In particular $p_1=0$ and, when $n=8$,
\begin{align*}
 \sum_i t_i u_i^2&=\frac{13}{8}p_2,
 &\sum_i t_i u_i^3&=\frac32p_3,\\
 \sum_i t_i u_i^4&=\frac{11}{8}p_4+\frac18p_2^2,
 &\sum_i t_i u_i^5&=\frac54p_5+\frac14p_2p_3.
\end{align*}
\end{lemma}
\begin{proof}
Summing the force equations gives $p_1=0$. Multiply each equation
by $u_i^{r+1}$, sum, and pair the two terms for $i,j$:
\[
 \frac n2\sum_i t_i u_i^r
 =\sum_{i<j}\sum_{\ell=0}^r u_i^{r-\ell}u_j^\ell
 =\frac12\left(\sum_{\ell=0}^r p_\ell p_{r-\ell}
                         -(r+1)p_r\right).
\]
Separating the two end terms proves the recursion.
\end{proof}

\begin{lemma}[Reflection forced by nonresonant moments]
\label{lem:n8-oct7-reflection}
Suppose all points outside a circle group of $m\ge2$ points are invariant
as a multiset under complex conjugation. Let the group's squared
radius be $t>0$ and its sum be real and nonzero. If
\[
            t\ne\frac{2n-r-1}{n}\qquad(2\le r\le m-1),
\]
then the circle group, and hence the entire stationary configuration,
is invariant under conjugation.
\end{lemma}
\begin{proof}
The group first power sum is real. Induct on $r=2,\ldots,m-1$ in
\eqref{eq:n8-oct7-moment-recursion}. All lower global power sums and
the outside group's power sums are real, so the imaginary part is
$(t-(2n-r-1)/n)$ times the circle group's $r$th power sum. The
nonresonance hypothesis makes that power sum real. Newton identities
therefore make all nonconstant coefficients of the monic circle
polynomial $G(x)=x^m+g_{m-1}x^{m-1}+\cdots+g_0$ real.
The circle identity gives
$g_1=g_0\overline{g_{m-1}}/t$. Since $g_{m-1}$ is real and nonzero,
$g_0$ is real too. A real polynomial has a conjugation-invariant
root multiset.
\end{proof}

\begin{corollary}[Two singleton levels for arbitrary point count]
\label{cor:n8-oct7-general-singletons}
For $n\ge4$, a collision-free no-origin true stationary configuration
with exactly three radial levels of populations $1+1+(n-2)$ must
be reflection invariant about the common centroid line. If $c$ is
the $(n-2)$-point squared radius, then
\[
 \frac{n-1}{n}<c<\frac{3(n-1)}{3n-4},\qquad c\ne1.
\]
After rotating the line and writing the singleton points as
$x,y\in\mathbb R$, $h=x+y$, $d=xy$, one has
\[
 h^2=3(n-1)-(3n-4)c,\qquad d=(n-1)(1-c).
\]
\end{corollary}
\begin{proof}
The two radial sum equations make the three group sums nonzero
collinear vectors, since a singleton cannot have zero sum.
Rotate so the singleton points are real. Their distinct squared
radii give $h\ne0$ and $x\ne y$.
The first radial moment gives $c=x^2-xy+y^2$; the total moment is
$(n-1)h^2-(3n-4)d=n-1$. These yield the displayed identities and
$h^2-4d=nc-(n-1)>0$, proving both strict bounds and $c\ne1$.
The largest possible $c$ is smaller than the smallest coefficient
in the reflection lemma for $2\le r\le n-3$, because
\[
 \frac{3(n-1)}{3n-4}<\frac{n+2}{n}\quad(n\ge4).
\]
The lemma proves reflection. For $n=4$ its moment range is empty;
the same conclusion follows from the real nonzero sum of the
two-point circle group.
\end{proof}

\begin{proposition}[Three population classes force reflection]
\label{prop:n8-oct7-reflected-classes}
Suppose a no-origin eight-point stationary configuration has exactly
three distinct squared radii and all three radial group sums are
nonzero. These sums determine a common line. With this line rotated
to the real axis, the entire configuration has real root polynomial
for each of the population splits $1+1+6$, $1+2+5$, and $2+2+4$.
All three assertions hold for every true stationary configuration.
\end{proposition}
\begin{proof}
The group sums satisfy $\sum_j A_j=\sum_j t_jA_j=0$.
The two-by-three real coefficient matrix has a one-dimensional kernel
whose three entries are nonzero for distinct $t_j$. Thus all three
complex sums lie on one line and are nonzero by hypothesis. Rotate
the line to be real. Singleton points are real. Two same-circle
points with nonzero real sum form a conjugate pair: writing their
sum as $A$ and their common squared modulus as $t$, their product
is $t$, so their polynomial is $x^2-Ax+t$.

For $1+1+6$, the six-point squared radius is less than $7/6$ by the
moment constraint, whereas the smallest coefficient in the reflection
lemma, for $2\le r\le5$, is $5/4$. The lemma applies.
For $1+2+5$, let the singleton be $v\in\mathbb R$, the pair have
sum $A\ne0$ and squared radius $b$, and the five-point circle have
real nonzero sum $s$ and squared radius $c<7/5$. Write its polynomial
using the circle identities as
\[
 G=X^5-sX^4+qX^3+\frac f{c^2}\overline q\,X^2
                    -\frac{fs}{c}X+f,\qquad |f|=c^{5/2}.
\]
The $r=2,3$ identities make the group's first three power sums real,
since $c<7/5<3/2<13/8$. Thus both $q$ and
$f\overline q/c^2$ are real. If $q\ne0$, then $f$ is real and the
group is reflected. Suppose $q=0$. Its fourth power sum is
$s^4+4fs/c$, so $r=4$ makes $f$ real unless $c=11/8$.
At this exceptional radius the $r=5$ identity still makes the fifth
group power sum $s^5+5f(s^2/c-1)$ real. If $\Im f\ne0$, then
$s^2=c$ and
\[
                         G=(X-s)(X^4-f/s).
\]
The global sixth power sum and its radial weighting are real.
But the $r=6$ recursion has imaginary part
$p_2\Im p_4/4$, since $p_1=0$ and $p_3$ is real. The moment gives
$v^2+2b=1/8$, whence
\[
 p_2=v^2+(A^2-2b)+c=2v^2+A^2+\frac54>0,
 \qquad \Im p_4=\frac{4s}{c}\Im f\ne0.
\]
This contradiction proves reflection without invoking global spread.

For $2+2+4$, denote the two pair sums by $A,B\in\mathbb R\setminus
\{0\}$ and their squared radii by $a,b$. The quadrilateral radius
$c$ satisfies $2a+2b+4c=7$. Write its polynomial, using the
nonzero real sum $s$, as
\[
 G=x^4-sx^3+q x^2-\frac{ds}{c}x+d,\qquad |d|=c^2,\qquad
                      q=\frac d{c^2}\overline q.
\]
If $c\ne13/8$, the $r=2$ moment identity makes $q$ real.
Suppose $c=13/8$ and $\Im q\ne0$. The $r=3$ identity gives
$\Im d=c\Im q$. The imaginary part of
$c^2q=d\overline q$ then gives $\Re d=c\Re q-c^2$, and therefore
\[
 d=cq-c^2,\qquad G=(x^2-sx+c)(x^2+q-c).
\]
The first factor is a reflected pair. Put
$L=A^2+B^2+s^2-2(a+b+c)$, the second power sum of all three
reflected pairs. The remaining antipodal pair gives
\[
 p_2=L+2(c-q),\qquad \Im p_4=-4(c-\Re q)\Im q.
\]
At $c=13/8$, the $r=4$ identity reads
$\Im p_4=(\Re p_2)\Im p_2$, hence $L=0$.
But the total moment gives $a+b=1/4$. Collision-freeness of the
two original pairs gives $A^2+B^2<4(a+b)=1$, and
$s^2=(A+B)^2\le2(A^2+B^2)$. Thus
\[
                     L<3-\frac{15}{4}<0,
\]
a contradiction. Therefore $q$ is real even at this resonance.
If $q\ne0$ the circle
identity immediately makes $d$ real. Suppose $q=0$.
The third group power sum is $s^3+3ds/c$. If $c\ne3/2$, the
$r=3$ identity makes $d$ real as well. It remains to exclude
$c=3/2$ and $\Im d\ne0$.

Here the global $p_2$ is real and the $r=4$ identity makes the fourth
group power sum real. Newton identities give that sum as
$s^4+4d(s^2/c-1)$, so $s^2=c$. Consequently
\[
                       G=(x-s)(x^3-d/s).
\]
This group's fifth power sum is $s^5$ and is real, as are all
outside fifth power sums. The $r=5$ identity therefore gives
$p_2\Im p_3=0$. But $a+b=1/2$ and
\[
 p_2=(A^2-2a)+(B^2-2b)+c=A^2+B^2+\frac12>0,
 \qquad \Im p_3=\frac{3s}{c}\Im d\ne0.
\]
This is a contradiction. In every case all coefficients of $G$ are
real, completing the reflection assertion.
\end{proof}

\begin{theorem}[Exclusion of the three-level $1+1+6$ stationary class]
\label{thm:n8-oct7-exclude-116}
There is no collision-free no-origin eight-point true stationary
configuration with exactly three radial levels of populations
$1+1+6$. This exclusion does not require global maximality.
\end{theorem}
\begin{proof}
All three group sums are nonzero, since a singleton sum is nonzero
and the two radial moment equations make a zero sum force all sums
to be zero. By Proposition~\ref{prop:n8-oct7-reflected-classes},
rotate so the two singleton points are real, say $x,y$, and the
six-point polynomial has real coefficients. Put $h=x+y\ne0$,
$d=xy\ne0$, and let $c$ be the six-point squared radius.
The first radial moment and the total moment give
\[
 c=x^2-xy+y^2=h^2-3d,\qquad
                    h^2=21-20c,\quad d=7(1-c).
\]
The two singleton radii are distinct, so $x\ne y$; their real
quadratic has strictly positive discriminant
$h^2-4d=(7-2h^2)/5$. Since $h\ne0$,
\begin{equation}\label{eq:n8-oct7-116-domain}
                         \frac78<c<\frac{21}{20},\qquad c\ne1.
\end{equation}
The six-circle self-inversive coefficient identities now give exactly
two possibilities:
\[
 Q_+=X^6+hX^5+CX^4+DX^3+cCX^2+c^2hX+c^3,
\]
\[
 Q_-=X^6+hX^5+CX^4-cCX^2-c^2hX-c^3,
 \qquad C,D\in\mathbb R.
\]
For $P=(X^2-hX+7(1-c))Q_\pm$, the circle force equations require
the remainder of $XP''-8cP'$ modulo $Q_\pm$ to vanish.
After using $h^2=21-20c$, its $X^5$ and $X$ coefficients are
$2\mathcal A$ and $\pm14c(c-1)\mathcal B$, where
\[
 \mathcal A=(8c-13)C+328c^2-701c+378,
 \qquad \mathcal B=(8c-1)C+72c^2-81c.
\]
Thus $\mathcal A=\mathcal B=0$, and
\[
 C=\frac{9c(9-8c)}{8c-1}>0,
\quad
 (8c-1)\mathcal A-(8c-13)\mathcal B
       =2(8c-7)(128c^2-160c+27)=0.
\]
By \eqref{eq:n8-oct7-116-domain},
$J(c)=128c^2-160c+27=0$.
For $Q_-$ the $X^3$ coefficient is
$12C(2c-1)(8c-7)$, which is strictly positive on this domain.
This excludes the negative sign.

For $Q_+$ the $X^4$ coefficient gives, with $E=D/h$,
\[
 E=C+\frac{35(2c-1)(c-1)}{3(3-2c)}.
\]
The denominator is nonzero on the domain. The $X^3$ coefficient gives
\[
 (64c^2-88c+21)C+(11-8c)(21-20c)E=0.
\]
Substitute the displayed $C,E$ and clear the nonzero denominator.
The resulting polynomial is
\[
 F(c)=-7(c-1)(26624c^4-90880c^3+102832c^2-39734c+1155)=0.
\]
Exact Euclidean division gives
\[
                 F\bmod J=\frac{105}{4}(848c-889).
\]
Hence $c=889/848$, but
$J(889/848)=-333/5618\ne0$. This excludes the positive sign and
finishes the proof.
\end{proof}

The standard-library program \path{n8_oct7_next_verify.py}
checks the relevant circle remainder coefficients, the radial
substitutions, the eliminant and Euclidean remainder, and the
quadrilateral and quintic Newton identities used in the exceptional reflection
case. These arguments leave the three-level population classes
$1+2+5$, $1+3+4$, $2+2+4$, and $2+3+3$ with nonzero collinear
centroid sums. Reflection is now necessary for every true stationary
configuration in the first and third classes. None of these remaining
classes, nor the configurations with at least four radial levels,
is declared exhausted here.

% Independent Oct 7 fragment: no symmetry assumption in the hypothesis.
\section{Origin configurations with exactly four nonzero radial levels}
\label{sec:n8-oct7-origin-fourlevels}

This earlier structural constraint is retained as an independent
check. Its entire class now satisfies the sharp origin inequality
by Theorem~\ref{thm:noclass-origin8}. Use $u=\sqrt{7/8}\,z$, moment seven, and force constant four.

\begin{lemma}[Four-level origin balances]
\label{lem:n8-oct7-origin-four-balances}
Let a collision-free true stationary eight-point configuration contain
the origin and have exactly four distinct positive squared radii
$t_1,\ldots,t_4$. Let $A_j$ be the sum of the points on the $j$th
circle. Then all $A_j$ are nonzero and lie on one common line.
More generally the kernel of the four-column balance matrix is
spanned over $\mathbb R$ by
\[
 \left(\frac{t_j}{\prod_{k\ne j}(t_j-t_k)}\right)_{j=1}^4.
\]
With $p_r=\sum_i u_i^r$, the stationary moment recursion is
\begin{equation}\label{eq:n8-oct7-origin-recursion}
 \sum_i|u_i|^2u_i^r
 =\frac{15-r}{8}p_r+\frac18\sum_{\ell=1}^{r-1}p_\ell p_{r-\ell}
 \qquad(r\ge1),\qquad p_1=0.
\end{equation}
\end{lemma}
\begin{proof}
Summing the force equations gives $p_1=0$. The force equation at
the actual origin gives $\sum_{u_i\ne0}1/u_i=0$. Its conjugate,
together with the radial moment for $r=1$, gives the three balances
\[
          \sum_j A_j=\sum_j t_jA_j=\sum_j A_j/t_j=0.
\]
The determinant of any three columns, with radii $a,b,c$, is
$(b-a)(c-a)(c-b)/(abc)\ne0$. The balance matrix therefore has rank
three, a one-dimensional real kernel, and every entry of its kernel
vector is nonzero. The displayed kernel vector follows from the
Lagrange interpolation identities for polynomials of degrees zero,
one and two. At least one of the seven nonzero points is a singleton
radial group, since there are four groups. Its sum is nonzero, so
none of the four sums can vanish. The complex solution vector is a
complex scalar times the real kernel vector, proving collinearity.

For the recursion, multiply the force equation by $u_i^{r+1}$ and
sum by pairs to obtain
\[
 4\sum_i|u_i|^2u_i^r
 =\frac12\left(\sum_{\ell=0}^r p_\ell p_{r-\ell}
                            -(r+1)p_r\right),\qquad p_0=8.
\]
Separating the end terms proves
\eqref{eq:n8-oct7-origin-recursion}.
\end{proof}

\begin{theorem}[Four-level origin stationary configurations force reflection]
\label{thm:n8-oct7-origin-four-reflection}
Every collision-free true stationary eight-point configuration
containing the origin and exactly four nonzero radial levels is
invariant under reflection about its radial-centroid line.
After rotating that line to the real axis, its root polynomial has
real coefficients. The assertion holds for all true stationary
configurations, without requiring global maximality or any radius cut.
\end{theorem}
\begin{proof}
Rotate the line from Lemma~\ref{lem:n8-oct7-origin-four-balances}
to be real. All group sums are real and nonzero. Singleton points
are real. A two-point circle group with real nonzero sum $A$ and
squared radius $b$ has product $b$, because its reciprocal sum is
$A/b$. Thus it is the conjugate pair with polynomial $X^2-AX+b$.
The seven nonzero points can split into four populations only as
$1+1+1+4$, $1+1+2+3$, or $1+2+2+2$. The last case is already
reflected. We handle the two larger-circle cases below.

\paragraph{Population $1+1+1+4$.}
Let the three singleton points be real, and let their total squared
modulus be $T$. The quadrilateral sum is real and nonzero, say $s$,
and its squared radius is $c$. Circle coefficient identities give
\[
 G=X^4-sX^3+qX^2-\frac{ds}{c}X+d,
 \qquad |d|=c^2,\qquad c^2q=d\overline q.
\]
All outside power sums are real. If $c\ne13/8$, the $r=2$ recursion
makes $q$ real. Suppose $c=13/8$ and $\Im q\ne0$.
The $r=3$ recursion makes the group third power sum
$s^3-3sq+3ds/c$ real, hence $\Im d=c\Im q$.
The imaginary part of the circle identity then gives
$\Re d=c\Re q-c^2$. Consequently
\[
 d=cq-c^2,\qquad G=(X^2-sX+c)(X^2+q-c).
\]
The first factor is a reflected pair. Put $L=T+s^2-2c$.
The remaining antipodal pair gives
$p_2=L+2(c-q)$ and $\Im p_4=-4(c-\Re q)\Im q$.
At this radius, the $r=4$ recursion reads
$\Im p_4=(\Re p_2)\Im p_2$, forcing $L=0$.
However $T=7-4c=1/2$, and centering implies $s^2\le3T=3/2$.
Thus $L\le3/2+1/2-13/4=-5/4<0$, a contradiction.
So $q$ is real in every case.

If $q\ne0$, the circle identity makes $d$ real. If $q=0$, the
third power sum is $s^3+3ds/c$, so the $r=3$ recursion makes $d$
real unless $c=3/2$. At that exceptional radius suppose $\Im d\ne0$.
The global $p_2$ is real. The fourth group power sum is
$s^4+4d(s^2/c-1)$; the $r=4$ recursion forces $s^2=c$ and hence
\[
                       G=(X-s)(X^3-d/s).
\]
Its fifth power sum is $s^5$, so both $p_5$ and its radial weighting
are real. The $r=5$ recursion therefore gives $p_2\Im p_3=0$.
But $T=7-4c=1$ and
\[
 p_2=T+c=\frac52>0,\qquad
                  \Im p_3=\frac{3s}{c}\Im d\ne0.
\]
This contradiction makes $d$ real, completing the quadrilateral case.

\paragraph{Population $1+1+2+3$.}
Write the two singleton points as $x,y\in\mathbb R$, the reflected
pair sum as $A\ne0$ and its squared radius as $b$. Let the triangle
sum be $s\in\mathbb R\setminus\{0\}$ and its squared radius be $c$.
Its circle polynomial is
\[
             G=X^3-sX^2-\frac{fs}{c}X+f,
             \qquad |f|=c^{3/2}.
\]
The triangle second and third power sums are
$s^2+2fs/c$ and $s^3+3f(s^2/c-1)$.
If $c\ne13/8$, the $r=2$ recursion makes $f$ real.
At $c=13/8$, suppose $\Im f\ne0$. The $r=3$ recursion makes the
third sum real, so $s^2=c$ and
\[
 G=(X-s)(X^2-f/s),\qquad \alpha=f/s.
\]
Here $\Im\alpha\ne0$. Put
$L=x^2+y^2+A^2-2b+c$. Then
$p_2=L+2\alpha$ and $\Im p_4=4\Re\alpha\Im\alpha$.
The $r=4$ recursion again reads
$\Im p_4=(\Re p_2)\Im p_2$, forcing $L=0$.
Together with the total moment
$x^2+y^2+2b+3c=7$, this yields
\[
                       2x^2+2y^2+A^2=7-4c=\frac12.
\]
Centering says $x+y+A=-s$. Weighted Cauchy--Schwarz now gives
\[
 s^2=(x+y+A)^2
 \le\left(\frac12+\frac12+1\right)(2x^2+2y^2+A^2)=1,
\]
contradicting $s^2=c=13/8$. Thus $f$ is real at the resonance too.
In all three population cases every circle group has real polynomial,
so the entire configuration is reflected as claimed.
\end{proof}

The standard-library certificate
\path{n8_oct7_next_origin_fourlevels_verify.py}
checks the three-column determinant, all four kernel cofactor
identities, the cubic and quartic Newton sums, both resonance
factorizations, the exact phase cancellations, and the final
weighted Cauchy--Schwarz identity and rational constants.
This is a symmetry reduction for the origin four-level class;
an objective comparison or instability certificate for all reflected
configurations in that class remains to be supplied.

\section{A neighborhood certificate allowing all eight points to move}
\label{sec:n8-full-neighborhood}
Let $z^*=(0,r,r\omega,\ldots,r\omega^6)$, where
$r^2=8/7$ and $\omega=e^{2\pi i/7}$, and let
$C_8=8^{28}/7^{21}$. This section gives an explicit neighborhood
exclusion in the full configuration space; the initially central
point is allowed to move.

\begin{theorem}\label{thm:n8-full-neighborhood}
Let $z$ be centered and satisfy $\sum_i|z_i|^2=8$. Define
\[
 d(z)=\min_{\theta,\sigma}
       \|z-e^{i\theta}\sigma z^*\|,
\]
where $\sigma$ ranges over permutations. If $d(z)\le1/24$, then
\[
       \Delta_8(z)\le C_8\exp\{-d(z)^2/2\}.
\]
Equality at $C_8$ occurs precisely when $d(z)=0$.
\end{theorem}
\begin{proof}
Choose the minimizing phase and permutation, relabel, and rotate
so that $\langle z^*,z\rangle=8a$ is positive real.
Then $d(z)^2=16(1-a)$. Put $h=z/a-z^*$ and $\rho=\|h\|$.
We have $\sum_i h_i=0$, $\langle z^*,h\rangle=0$, and
\[
 z=\frac{\sqrt8\,(z^*+h)}{\sqrt{8+\rho^2}},\qquad
 \rho^2=8(a^{-2}-1)\ge d(z)^2.
\]
For $d(z)\le1/24$, $a\ge9215/9216$, hence
\[
 \rho^2\le\frac{147448}{84916225}<\frac1{529}.
\]

Let $B$ be the inverse-square Laplacian at $z^*$. The central-polygon
Fourier calculation gives its singular values on the complex slice
$\sum h_i=\langle z^*,h\rangle=0$ as
\[
 \frac7{\sqrt8},\ \frac7{\sqrt8},\
 \frac78,\ \frac74,\ \frac74,\ \frac78.
\]
Indeed, in unnormalized Fourier vectors
$f_k=(0,1,\omega^k,\ldots,\omega^{6k})$ and
$u=(-7,1,\ldots,1)$, the unit-radius inverse-square Laplacian maps
$f_2\mapsto u$, $u\mapsto8f_5$, and
$f_3,f_4,f_5,f_6\mapsto-f_1,-2f_2,-2f_3,-f_4$,
respectively. Multiplication by $7/8$ gives the actual-radius
operator. These maps are checked independently in $\mathbb Q(\omega)$
by the supplied arithmetic checker.
Thus the quadratic term in the normalized logarithmic product is
\[
 -\Re(h^{\mathsf T}Bh)-\frac72\rho^2
 \le-\delta\rho^2,\qquad
 \delta=\frac72-\frac7{\sqrt8}>1.
\]

The minimum squared distance at $z^*$ is
$\frac{32}{7}\sin^2(\pi/7)>4/5$.
For an elementary rational verification, use $31/10<\pi<7/2$ and
$\sin x\ge x-x^3/6$ for $0\le x\le1/2$:
\[
 \frac{32}{7}\sin^2(\pi/7)
 >\frac{32}{7}\left\{\frac{31}{70}
              -\frac16\left(\frac{31}{70}\right)^3\right\}^2>\frac45.
\]
The lower bound for $\pi$ already follows from the inscribed
regular twelve-gon, $\pi>3(\sqrt6-\sqrt2)>31/10$;
the upper bound follows from the circumscribed regular hexagon.
The central distances have squared length $8/7$ and are larger.
For $v_{ij}=(h_i-h_j)/(z_i^*-z_j^*)$ we consequently have
\[
 \max|v_{ij}|\le b:=\frac85\rho,
 \qquad \sum_{i<j}|v_{ij}|^2\le10\rho^2.
\]
The logarithmic Taylor remainder and the scale correction yield
\[
 \log\frac{\Delta_8(z)}{C_8}
 \le\left\{-\delta+\frac{20b}{3(1-b)}
                   +\frac7{32}\rho^2\right\}\rho^2.
\]
Here $|\log(1+v)-v+v^2/2|\le |v|^3/[3(1-|v|)]$
and $\log(1+x)\ge x-x^2/2$ supply the two remainders.
For $\rho\le1/23$,
\[
 \frac{20b}{3(1-b)}+\frac7{32}\rho^2
 \le\frac{160}{321}+\frac7{16928}<\frac12.
\]
Therefore the logarithmic ratio is at most $-\rho^2/2$,
which is at most $-d(z)^2/2$, proving the assertion.
\end{proof}

This certificate excludes a full neighborhood from the search for a
better eight-point configuration, including perturbations that move
the central point. It gives no comparison outside that neighborhood.
For the central $(n-1)$-gon in other dimensions, the same Fourier
singular values, a minimum-distance bound, and a Taylor remainder
produce explicit neighborhoods whenever the slice has a strict
spectral gap. For the equal-strength central polygons this applies
to $5\le n\le9$. The degenerate case $n=10$ needs higher-order terms
and is already known here to have a larger competitor.

% Fragment: an exact smaller physical model for the centrally symmetric n8 class.
\section{A half-size physical model for centrally symmetric eight-point candidates}
\label{sec:n8-even-half-model}

This is a reduction of a still-unresolved subclass, not an assertion
that every eight-point maximizing shape is centrally symmetric.
At force-one scale a centrally symmetric collision-free no-origin
configuration has root polynomial
\[
 P(z)=q(z^2),\qquad q(y)=y^4+ay^3+by^2+cy+d,\qquad d\ne0.
\]
Write $q'=\partial_yq$ and
\[
 H_q(y)=q'(y)+2yq''(y)=28y^3+15ay^2+6by+c.
\]
For a root $z^2=y$,
\[
 \frac{P''(z)}{2P'(z)}=\frac{H_q(y)}{2zq'(y)}.
\]
Thus true stationarity is equivalent to
$H_q(y)/(2q'(y))=|y|$ at every root of $q$.

\begin{proposition}[Exact half-size physical trace formulation]
\label{prop:n8-half-trace}
Let $T,S$ be polynomials of degree below four, interpreted modulo
$q$, and impose
\begin{equation}\label{eq:n8-half-congruences}
 2Tq'-H_q\equiv0\pmod q,\qquad yS-T^2\equiv0\pmod q.
\end{equation}
In $\C[y]/(q)$ let
\[
 K_{jk}=\operatorname{Tr}(y^jS^k),\quad 0\le j,k\le4,
 \qquad G=(K_{jk})_{0\le j,k<4},
\]
\[
 L_{jk}=\operatorname{Tr}(y^jS^kT),\quad 0\le j,k<4.
\]
Then \eqref{eq:n8-half-congruences} together with
\begin{equation}\label{eq:n8-half-positive}
                  K=K^*,\qquad G\succ0,\qquad L\succ0
\end{equation}
are equivalent to actual collision-free force-one stationarity of
all eight roots of $P=q(z^2)$.
For every such physical configuration,
\[
 \operatorname{Tr}T=14,\qquad \sum_{P(z)=0}|z|^2=28,
\]
\begin{equation}\label{eq:n8-half-objective}
 \Delta_8(z)=4^4\det(M_T)(\det G)^2,
\end{equation}
where $M_T$ is multiplication by $T$ in the four-dimensional quotient.
Scaling the moment from $28$ to eight multiplies this expression by
$(2/7)^{28}$.
\end{proposition}
\begin{proof}
For roots $y_i$ of $q$ with multiplicities $m_i$, the quotient traces
are $K_{jk}=\sum_i m_i y_i^jS(y_i)^k$. Thus $G\succ0$ forces four
distinct roots and, by its factorization $G=VW^{\mathsf T}$, four
distinct values $S(y_i)$, where $V_{ji}=y_i^j$ and
$W_{ki}=S(y_i)^k$. Write $q(y)=y^4+\sum_{j=0}^3q_jy^j$.
The quotient relation gives $K_{4,k}=-\sum_jq_jG_{jk}$.
Hermitian symmetry of the border then gives
$K_{k,4}=-\sum_j\overline{q_j}G_{kj}$, or
$\sum_i y_i^k\overline q(S(y_i))=0$ for $0\le k<4$.
Invertibility of $V$ proves $\overline q(S(y_i))=0$.
Consequently the distinct $S(y_i)$ permute the conjugate roots, and
$G=V\Pi V^*$ for a permutation matrix $\Pi$. Hermitian positive
definiteness passes to $\Pi$; a Hermitian permutation is an involution,
and each nontrivial transposition has eigenvalue $-1$. Therefore
$\Pi=I$, proving actual conjugate pairing
$S(y_i)=\overline{y_i}$. The border is essential; positivity of the
leading four-by-four block alone would not suffice. Since $d\ne0$,
all $y_i\ne0$.
With $V_{ji}=y_i^j$, the localizer is
$L=V\operatorname{diag}(T(y_i))V^*$. Its positive definiteness makes
all $T(y_i)$ positive real numbers. The second congruence gives
$T(y_i)^2=|y_i|^2$, hence $T(y_i)=|y_i|$.
The first congruence now gives
$P''(z)/(2P'(z))=|y_i|/z=\overline z$ for both square roots
$z=\pm\sqrt{y_i}$. Conversely true stationarity gives these
congruences, $S(y_i)=\overline{y_i}$, and all stated positivity.

At distinct roots,
\[
 \operatorname{Tr}T
 =2+\sum_i\frac{y_iq''(y_i)}{q'(y_i)}
 =2+2\sum_{i<j}\frac{y_i-y_j}{y_i-y_j}=14.
\]
Therefore the eight-point moment is $2\sum_iT(y_i)=28$.
The Vandermonde factorization of the centrally symmetric points gives
\[
 \Delta_8=4^4\prod_i|y_i|\prod_{i<j}|y_i-y_j|^4.
\]
Here $\det M_T=\prod_i|y_i|$ and
$\det G=\prod_{i<j}|y_i-y_j|^2$, proving
\eqref{eq:n8-half-objective}.
\end{proof}

\subsection{Necessary global cuts and the symmetry-breaking stability block}
At a global eight-point maximum the previously proved force-one
radius and pair separation bounds give, on the quartic roots,
\[
 \frac12\le T(y_i)\le\frac{49}8,\qquad
 |y_i-y_j|^2=|z_i-z_j|^2|z_i+z_j|^2\ge4.
\]
The two-pair subset-spread bound also gives
$T(y_i)+T(y_j)\ge3$. The closed inequalities suffice in a compact
certification model. The first bounds become four-by-four Hermitian
localizers for $T-1/2$ and $49/8-T$. For the separation product,
\[
 |q'(y_i)|^2\ge4^3,
\]
so the localizer for
$q'(y)\overline{q'}(S(y))-64$ is positive semidefinite.
These are necessary conditions specifically for global candidates.

The inverse-square stability matrix of the eight roots splits by
parity of perturbations under $z\mapsto-z$. On perturbations taking
the same value at $z_i$ and $-z_i$, its four-by-four block is
\[
 C_{ij}=-\frac{2(y_i+y_j)}{(y_i-y_j)^2}\quad(i\ne j),
 \qquad C_{ii}=-\sum_{j\ne i}C_{ij}.
\]
This block tests perturbations breaking central symmetry. Every local
maximum requires $\|C\|\le1$. Indeed the unconstrained weighted
objective $F=\log\Delta_8-\sum|z|^2$ has the same local shape maxima
as its moment-$28$ restriction, because each nearby shape has best
scale at moment $28$. Its quadratic term for a perturbation $h$ is
\[
 -\sum_i|h_i|^2-\Re\sum_{i<j}
             \left(\frac{h_i-h_j}{z_i-z_j}\right)^2.
\]
Assigning $h=s_i$ to both $z_i,-z_i$ reduces this to
$-2\|s\|^2-2\Re(s^{\mathsf T}Cs)$. Varying the common phase of
$s$ gives $|s^{\mathsf T}Cs|\le\|s\|^2$; Takagi factorization of
the complex symmetric $C$ is equivalent to $\|C\|\le1$.
This condition must not be dropped merely because the candidate
configuration is centrally symmetric.
Writing $s_k=\sum_i y_i^k$, its exact polynomial actions are
\[
 C(y_i)_i=(2T_i-7)_i,
\]
\[
 C(y_i^2)_i=(4y_iT_i-18y_i-2s_1)_i,
\]
\[
 C(y_i^3)_i=(6y_i^2T_i-25y_i^2-6s_1y_i-2s_2)_i.
\]
The centered powers $y,y^2,y^3$ span $\mathbf1^\perp$ when the roots
are distinct. Their three-by-three Gram difference therefore expresses
$\|C\|\le1$ exactly in quotient traces of $q,T,S$. This supplies a
small polynomial stability cut for the full even chart, although
feasibility or an objective certificate for that chart is not solved
here.

\begin{proposition}[All stationary two-square configurations are unstable]
\label{prop:n8-two-square-unstable}
Every collision-free stationary eight-point configuration consisting
of two concentric square rings is unstable under perturbations that
break central symmetry. In force-one scale the largest singular value
of the same-pair block $C$ has square
\[
 \frac{64}{49}\quad\hbox{(regular octagon)},\qquad
 \frac{532}{245}\quad\hbox{(aligned unequal squares)},\qquad
 \frac87\quad\hbox{(staggered unequal squares)}.
\]
In particular none is a local maximum among unrestricted eight-point
configurations.
\end{proposition}
\begin{proof}
Let $R,S>0$ be the squared ring radii. Rotate the first square so its
fourth powers are $R^2$. Multiplying its pointwise force equation by
the point gives
\[
 R=\frac32+\frac{4R^2}{R^2-X},
\]
where $|X|=S^2$ is the fourth power on the other square. Reality of
the left side forces $X$ real, hence $X=\varepsilon S^2$,
$\varepsilon=1$ or $-1$. The moment gives $R+S=7$.
For $\varepsilon=1$, collision-freeness requires $R\ne S$ and the
force equation gives $(R-S)^2=98/5$. For $\varepsilon=-1$, subtract
the force equations of the two rings to obtain either $R=S=7/2$ or
$R^2+S^2=28$, equivalently $(R-S)^2=7$.

The quartic roots are $\{R,-R,S,-S\}$ in the aligned case and
$\{R,-R,iS,-iS\}$ in the staggered case. In this order the block is
\[
 C=\begin{pmatrix}
 A+B&0&-A&-B\\
 0&-A-B&B&A\\
 -A&B&A-B&0\\
 -B&A&0&B-A
 \end{pmatrix}.
\]
Here aligned rings have
$A=2(R+S)/(R-S)^2$, $B=2(R-S)/(R+S)^2$; staggered rings have
$A=2(R+iS)/(R-iS)^2$, $B=\overline A$.
Let $J$ interchange the first two and last two coordinates. Then
$JCJ=-C$ and $C\mathbf1=0$. The even and odd eigenspaces of $J$
both have dimension two. In an orthonormal real basis of these
spaces, with the first even vector proportional to $\mathbf1$,
$C$ therefore has the form
\[
 \begin{pmatrix}
 0&0&0&0\\0&0&\alpha&\beta\\
 0&\alpha&0&0\\0&\beta&0&0
 \end{pmatrix}.
\]
Its two nonzero singular values are equal, with square
$|\alpha|^2+|\beta|^2=\tfrac12\operatorname{Tr}(CC^*)
=4(|A|^2+|B|^2)$.
For aligned rings, $A=5/7$ and $B^2=8/245$, giving $532/245$.
For staggered rings, $|A|^2=|B|^2=4/(R^2+S^2)$, giving
$32/(R^2+S^2)$. The equal and unequal stationary choices yield
$64/49$ and $8/7$, respectively. Each exceeds one and hence violates
the necessary Hessian contraction proved above.
\end{proof}

\begin{remark}[Why the geometric cuts do not alone close the even chart]
The force equations are essential to any extension of this argument.
For example take the staggered-square configuration with
$R=109/20$, $S=31/20$. Its full moment is $28$. Its physical moduli
$|y_i|$ satisfy $1/2\le |y_i|\le49/8$, all pairwise
$|y_i-y_j|^2\ge4$, every $|y_i|+|y_j|\ge3$, and
$|q'(y_i)|^2\ge64$. Nevertheless the block above has
\[
 \|C\|^2=\frac{32}{R^2+S^2}=\frac{6400}{6421}<1.
\]
This configuration is not stationary: the force interpolation has
values $3/2+4R^2/(R^2+S^2)$ on $y=\pm R$ and
$3/2+4S^2/(R^2+S^2)$ on $y=\pm iS$, which differ from $R,S$.
Thus these explicit geometric inequalities together with the
same-pair Hessian block cannot supply a universal instability
certificate without using physical stationarity. This example makes
no assertion about feasibility of the full stationary trace model,
nor does it exclude or exhibit a global competitor.
\end{remark}

\subsection{A smaller exact dual resultant for this chart}
The necessary squared-force dual polynomial is
\[
 \mathcal E(v)=\operatorname{Res}_y
 \bigl(q(y),4vyq'(y)^2-H_q(y)^2\bigr)=\sum_{k=0}^4E_kv^k.
\]
Its leading coefficient is
\[
                     E_4=256d(\operatorname{disc}q)^2\ne0.
\]
Physical stationarity necessarily implies
\[
 E_3=E_4\overline a,\quad E_2=E_4\overline b,\quad
 E_1=E_4\overline c,\quad E_0=E_4\overline d.
\]
The full eight-point force resultant specializes exactly to
$256\mathcal E(u^2)$. These squared setwise conditions remain only
necessary; the bordered trace and positive $T$ test above resolve
both actual pairing and the sign lost when the force is squared.
The script \path{n8_next_even_resultant.py} generates the
four-by-four multiplication determinant and checks its leading
coefficient and independent scalar Sylvester specializations.
An execution or an eliminant alone is not an exhaustion certificate
for the unresolved centrally symmetric eight-point class.

\section{A fixed finite moment determines the exact maximum}
\label{sec:fixed-moment-reconstruction}

The moment limit can be turned into an explicitly prescribed finite
algebraic selection, without constructing the stationary-value polynomial.
The price is a prohibitive uniform height bound and a very large moment
order. This result is an effective finiteness theorem, rather than a
simple product expression or a practical computation of previously
unknown point counts.

For $n\ge3$, put $r=2n$, $d=n(n-1)=2N$, and define the integers
\begin{equation}\label{eq:uniform-height-constants}
 h_n=\max\{8^N,4n+2\},\qquad
 \mathfrak{D}_n=rd(d-1)^{r-1},\qquad
 \mathfrak{H}_n=(16h_nd^r)^{r2^rd^r}.
\end{equation}
These constants are deliberately generous.

\begin{lemma}[An explicit annihilator height bound]
\label{lem:uniform-height-bound}
There is a nonzero $q_n\in\mathbb Z[T]$ satisfying
\[
 q_n(M_n)=0,\qquad \deg q_n\le \mathfrak{D}_n,\qquad
 \max_j|[T^j]q_n|\le \mathfrak{H}_n.
\]
The bound is for an annihilating polynomial, without a claim that its
primitive irreducible factor has the same coefficient bound.
\end{lemma}
\begin{proof}
In the $r$ real coordinates $(x_i,y_i)$, the constraint
\[
 S=\left\{(x,y):\sum_i(x_i^2+y_i^2)=n\right\}
\]
is a compact connected sphere defined by one integer polynomial.
Each factor $(x_i-x_j)^2+(y_i-y_j)^2$ has coefficient $\ell^1$ norm $8$.
Consequently $\Delta_n$ has degree $d$ and coefficient height at most
$8^N$. The defining sphere polynomial has degree $2$ and coefficient
height $n\le8^N$.

Apply the deformation and resultant construction of
Jeronimo--Perrucci--Tsigaridas~\cite[Proposition~2.7 and the proof of
Theorem~1.1]{JPT} to the polynomial $-\Delta_n$ on this sphere.
Their construction supplies a nonzero integer polynomial vanishing at
$-M_n$, including when the minimizing set has positive dimension.
There are either $s=0$ or $s=1$ active constraints in that construction;
the degree bound is
\[
 \binom rs d^s(d-1)^{r-s}\le rd(d-1)^{r-1}.
\]
The coefficient estimate explicitly derived in their proof is
\[
 (2^{4-r/2}\widetilde H d^r)^{r2^rd^r},\qquad
 \widetilde H=\max\{H_{\rm input},2r+2m_{\rm input}\}.
\]
Here $H_{\rm input}=8^N$ and $m_{\rm input}=1$, so
$\widetilde H=h_n$. Replacing $2^{4-r/2}$ by the larger integer $16$
gives $\mathfrak{H}_n$. Finally replace $T$ by $-T$ in the annihilator, which
does not change its coefficient magnitudes.
\end{proof}

For positive integers $D,H$, let $\mathcal A(D,H)$ be the finite set
of all distinct positive real roots of nonzero integer polynomials of
degree at most $D$ and coefficient height at most $H$.
Multiplicity and repetition across different polynomials are ignored.

\begin{lemma}[Uniform separation without a factor-height estimate]
\label{lem:universal-resultant-separation}
Set
\[
 B(D,H)=(2D)!\bigl[(D+1)H2^D\bigr]^{2D},\qquad
 \delta(D,H)=\frac1{B(D,H)+1}.
\]
Any two distinct elements of $\mathcal A(D,H)$ are separated by at
least $\delta(D,H)$.
\end{lemma}
\begin{proof}
Choose integer polynomials $f,g$ of positive degrees at most $D$ and
coefficient heights at most $H$, with $f(\alpha)=g(\beta)=0$.
The polynomial
\[
 R(u)=\operatorname{Res}_t(f(t),g(t+u))\in\mathbb Z[u]
\]
is nonzero: over $\mathbb C$ its roots, with multiplicities, are the
finitely many differences between roots of $g$ and roots of $f$.
Thus $\gamma=\beta-\alpha$ is a nonzero root of $R$.

The coefficient $\ell^1$ norm of $g(t+u)$ is at most
$(D+1)H2^D$. Every polynomial entry of the Sylvester matrix has
$\ell^1$ norm bounded by this number, including the entries from $f$.
The matrix has order at most $2D$. Expanding its determinant and using
submultiplicativity of the coefficient $\ell^1$ norm gives
\[
 \|R\|_1\le(2D)![(D+1)H2^D]^{2D}=B(D,H).
\]
Write $R(u)=u^aQ(u)$ with $Q(0)\ne0$. The integer $Q(0)$ has
absolute value at least one, while every coefficient of $Q$ has
absolute value at most $B=B(D,H)$. If $0<|\gamma|<1/(B+1)$, then
\[
 \left|\sum_{j\ge1}[u^j]Q\,\gamma^j\right|
 \le\frac{B|\gamma|}{1-|\gamma|}<1\le|Q(0)|,
\]
contradicting $Q(\gamma)=0$.
This argument permits repeated roots and requires no factorization.
\end{proof}

Use $\mathfrak{B}_n=B(\mathfrak{D}_n,\mathfrak{H}_n)$ and $\delta_n=1/(\mathfrak{B}_n+1)$, and put
\begin{equation}\label{eq:explicit-fixed-moment-order}
 \mathfrak{U}_n=\left(\frac{2n}{n-1}\right)^N,\qquad
 m=n-2,\qquad \mathfrak{L}_n=m(N+m+1),\qquad
 \mathfrak{K}_n=\left\lceil4\mathfrak{U}_n\mathfrak{L}_n(\mathfrak{B}_n+1)\right\rceil^2.
\end{equation}
The complete coefficient expansion in
\eqref{eq:coefficient-moment}, evaluated at this prescribed $\mathfrak{K}_n$,
gives the rational number $A_{n,\mathfrak{K}_n}$. Let
\[
 E_n=\binom{\mathfrak{K}_nN+n-2}{n-2}.
\]

\begin{theorem}[A finite coefficient-only algebraic selection]
\label{thm:fixed-moment-reconstruction}
For every $n\ge3$, the exact value $M_n$ is the unique element
$\xi\in\mathcal A(\mathfrak{D}_n,\mathfrak{H}_n)$ satisfying
\begin{equation}\label{eq:fixed-moment-root-selector}
 \boxed{\qquad A_{n,\mathfrak{K}_n}\le\xi^{\mathfrak{K}_n}\le E_nA_{n,\mathfrak{K}_n}.\qquad}
\end{equation}
For $n=2$, $M_2=4$ directly.
Thus a single finite moment, whose index is fixed from $n$ before any
geometric calculation, determines the exact maximum.
\end{theorem}
\begin{proof}
The identity
\[
 \sum_{i<j}|z_i-z_j|^2
 =n\sum_i|z_i|^2-\left|\sum_i z_i\right|^2\le n^2
\]
and AM--GM give $M_n\le \mathfrak{U}_n$.
The moment enclosure~\eqref{eq:moment-enclosure} places $M_n$ in
\[
 I_n=[A_{n,\mathfrak{K}_n}^{1/\mathfrak{K}_n},(E_nA_{n,\mathfrak{K}_n})^{1/\mathfrak{K}_n}].
\]
This interval has length at most $\delta_n/2$. Indeed, abbreviating
$K=\mathfrak{K}_n$, the elementary bounds
$E_n\le[K(N+m)]^m$, $\log K\le\sqrt K$ for $K\ge1$, and
$\log(N+m)\le N+m$ give
\[
 \frac{\log E_n}{K}
 \le\frac{m[\log K+\log(N+m)]}{K}
 \le\frac{\mathfrak{L}_n}{\sqrt K}\le\frac{\delta_n}{4\mathfrak{U}_n}.
\]
The last quantity is at most $1/2$, so $e^x-1\le2x$ implies
$E_n^{1/K}-1\le\delta_n/(2\mathfrak{U}_n)$. The lower endpoint of $I_n$
is at most $M_n\le \mathfrak{U}_n$, whence
\[
 |I_n|=A_{n,K}^{1/K}(E_n^{1/K}-1)\le\delta_n/2.
\]
By Lemma~\ref{lem:uniform-height-bound}, $M_n$ belongs to
$\mathcal A(\mathfrak{D}_n,\mathfrak{H}_n)$. By
Lemma~\ref{lem:universal-resultant-separation}, no two distinct
elements of that finite set can lie in $I_n$. This proves existence
and uniqueness in~\eqref{eq:fixed-moment-root-selector}.
\end{proof}

To implement the theorem, enumerate the finitely many bounded integer
coefficient arrays, isolate and deduplicate their positive real roots,
and compute the complete expansion of $V_n^{\mathfrak{K}_n}$. Real algebraic
arithmetic decides the two inequalities in
\eqref{eq:fixed-moment-root-selector} exactly. Factoring a polynomial
containing the selected root and supplying a rational isolating
interval produces the canonical exact-output representation described
earlier. No enumeration of stationary shapes or construction of their
critical-value polynomial is required by this procedure.

The constants make this procedure prohibitively expensive, even for
small $n$. The sharper degree estimate $n^{2n-4}$ from
the earlier critical-value construction remains useful; the larger $\mathfrak{D}_n$ here
is used only to obtain an explicit height bound. The theorem replaces
an unspecified limiting process by finite algebraic selection, and
does not supply a short closed expression or an executed new maximum
for $n\ge8$. The standard-library checker
\path{algebraic_verify.py} verifies small exact moment examples,
translation resultants including repeated roots, and the prescribed
moment-order inequality on manageable constants. It does not execute
the universal polynomial enumeration or the worst-case moment order.



\section{Exact arithmetic and reproducibility}\label{ref5-app:certificate}
\subsection{Interval operations}
All endpoints in the program are instances of Python's
\texttt{fractions.Fraction}. Addition, subtraction, multiplication,
positive-part evaluation and squaring use their elementary outward
interval formulas. Division is performed only when the denominator
interval has positive lower endpoint.

For a nonnegative rational $v=a/b$, an outward square-root interval
with denominator $2^{40}$ is obtained from
\[
 m=\operatorname{isqrt}\left(\left\lfloor
              a\,2^{80}/b\right\rfloor\right).
\]
Its lower endpoint is $m/2^{40}$; the upper endpoint is this value
when it squares exactly to $v$, and otherwise $(m+1)/2^{40}$.
For a nondegenerate interval the lower and upper endpoint calculations
are applied separately. Thus no floating-point approximation enters
a mathematical decision. Wall-clock timings in the JSON output have
no role in certification.

The factor $20-12s-4q/5$ is nonnegative in the target domain.
The other factor is also positive: in A,
\[
 D_*=15+x^2(-180+5760x-15600x^2)\ge15-180/16>0,
\]
since $0\le x\le1/4$.
In B, put $T=5+s\in[25/4,20/3]$; then
$D_*=T(-12T^2+144T-405)/5>0$ by checking the endpoints of the
concave quadratic. Clipping interval factors to their nonnegative
parts is therefore valid, including rectangles extending outside the
curved domain boundary.

\subsection{Transcendental and algebraic constants}
The only transcendental checks concern logarithms.
The coefficient $c$ has the positive-series representation
\[
 c=\frac15-\frac4{25}\sum_{j=0}^{\infty}
                \frac1{(2j+1)5^j}.
\]
After $N$ terms the series tail is at most
$5/(4(2N+1)5^N)$. The program uses $N=24$.
For rational $x\ge1$, with $y=(x-1)/(x+1)$,
\[
 \log x=2\sum_{j=0}^{\infty}\frac{y^{2j+1}}{2j+1}.
\]
After $N$ terms, its tail is at most
$2y^{2N+1}/((2N+1)(1-y^2))$.
Using $N=120$ at $x=5$ and $x=41/6$ certifies
\eqref{ref5-eq:c-bounds} and
$H(R)-\log R-k(R^2-5R+5)^2>0$.

The checks in \eqref{ref5-eq:two-steps}--\eqref{ref5-eq:spectral-gap}
use the algebraic constants \eqref{ref5-eq:algebraic-trig}, each enclosed by
the integer-square-root procedure above.
No numerical value of $\pi$ or trigonometric series is required.
The underlying bounds $\sin x\le x$ and $\cos x\ge1-x^2/2$ follow
by integration of $\cos x\le1$ and $\sin x\le x$ for $x\ge0$.

\subsection{Polynomial expressions and source files}
To make the interval encoding explicit, in A the functions in
\eqref{ref5-eq:scalar-definitions} simplify to
\begin{equation}\label{ref5-eq:branch-A}
\begin{aligned}
 f(20x^2)&=x^2\bigl(-10+450C-500k
             +x(-300C+3000k-8500kx)\bigr),\\
 D_*(20x^2)&=15+x^2(-180+5760x-15600x^2).
\end{aligned}
\end{equation}
In B, with $T=5+s$,
\begin{equation}\label{ref5-eq:branch-B}
\begin{aligned}
 f(s)&=\left(-\frac12+\frac{15C}{2}\right)s
       +5C(T^2/5-5)
       -k(T^3/5+24T^2-275T+750),\\
 D_*(s)&=T(-12T^2+144T-405)/5.
\end{aligned}
\end{equation}

\noindent\begin{minipage}{\textwidth}
The ancillary files supplied with this manuscript are:
\begin{itemize}
\item \path{verify_mordell.py};
\item \path{verification_result.json};
\item \path{algebra_identity_check.py};
\item \path{algebra_identity_result.json};
\item \path{generalization_check.py};
\item \path{generalization_result.json}.
\end{itemize}
\end{minipage}
From their directory, the full certificate is reproduced by
\begin{verbatim}
python verify_mordell.py --output verification_result.json
python algebra_identity_check.py
python generalization_check.py
\end{verbatim}
The first command must be run with both regions, its default.
The program exits unsuccessfully if any scalar region is unresolved
or a constant comparison fails. Explicit exception checks remain active
under Python's optimization option \texttt{-O}.

The second script independently expands \eqref{ref5-eq:moment-expansion}
as formal polynomials over $\Q$, substitutes the mixed-moment
relations, and checks all four identities in
\eqref{ref5-eq:distance-moments}. It also derives
\eqref{ref5-eq:branch-A}--\eqref{ref5-eq:branch-B} from
\eqref{ref5-eq:scalar-definitions} and checks exact point encodings.
Its point checks are an encoding audit; the covering proof for region A
of \eqref{ref5-eq:certificate} is performed by the interval program, while
region B is handled by \eqref{ref5-eq:branch-b-analytic}.
The third script checks the general Gram coefficients and supporting
polynomial identities, the six-point rational constants, and the
exact six-point value in \eqref{ref5-eq:Chalk-competitor}.
The recursive identities and the spectral and variational arguments
are proved analytically above; finite script checks support those
proofs and do not replace them.


\section{An independent algebraic cross-check for five points}
\label{sec:n5-complete-improved}

The five-point argument below reconstructs the exact classification from
the locally supplied verification script. It does not invoke the numerical
list of vortex shapes, assume reflection symmetry, or use the sharp
inequality of Wang~\cite{Wang} as a premise. The resultant elimination is
the same method developed in that preprint; the attribution and the
computational dependence should both be retained.

\begin{theorem}[Sharp five-point inequality]
\label{thm:n5-reconstructed-sharp}
If $z_1,\ldots,z_5\in\C$ and $\sum_i|z_i|^2=5$, then
\[
 \Delta_5(z)=\prod_{i<j}|z_i-z_j|^2\le3125.
\]
Equality holds precisely for a rotated regular unit pentagon, with
arbitrary permutation of its vertices.
\end{theorem}

\subsection{Variational reduction and the necessary resultant equations}
The constraint set is compact, and a regular unit pentagon has positive
product. Thus a maximizer exists and is collision-free. At such a point,
the Lagrange equations for $\log\Delta_5$ give
\[
 \sum_{j\ne i}\frac1{z_i-z_j}=\lambda\overline{z_i}.
\]
Multiplication by $z_i$ and summation give $5\lambda=10$, so
$\lambda=2$. Summation without weights gives $\sum_i z_i=0$.
It therefore suffices to classify the collision-free physical solutions
of this equation.

Introduce $x_i,y_i$ satisfying
\begin{equation}\label{eq:n5-reconstructed-complexified}
 \sum_{j\ne i}\frac1{x_i-x_j}=2y_i,
 \qquad
 \sum_{j\ne i}\frac1{y_i-y_j}=2x_i.
\end{equation}
Physical solutions have $x_i=z_i$, $y_i=\overline{z_i}$.
The transformation $(x_i,y_i)\mapsto(t x_i,t^{-1}y_i)$ preserves
\eqref{eq:n5-reconstructed-complexified}. After this transformation the
physical relation becomes
\begin{equation}\label{eq:n5-reconstructed-kappa}
 y_i=\kappa\overline{x_i},\qquad \kappa=|t|^{-2}>0.
\end{equation}
This positive factor must be retained: arbitrary complex coefficient
normalization is not a rotation and does not preserve the moment five.

Let
\[
 p(X)=\prod_i(X-x_i)=X^5+aX^3+bX^2+cX+d,
 \qquad \delta=\operatorname{Disc}(p)\ne0,
 \qquad \xi_i=\frac{p''(x_i)}{p'(x_i)}=4y_i.
\]
Define
\[
 R(Y,T)=\operatorname{Res}_X\bigl(p,(Y+TX)p'-p''\bigr)
        =\delta\prod_i(Y+Tx_i-\xi_i).
\]
Write
\[
 R(Y,0)=\delta Y^5+r_3Y^3+r_2Y^2+r_1Y+r_0,
 \qquad
 R_T(Y,0)=20\delta Y^3+s_2Y^2+s_1Y+s_0.
\]
Indeed, $\sum_i x_i=\sum_i\xi_i=0$ and
$\sum_i x_i\xi_i=20$. The second equation in
\eqref{eq:n5-reconstructed-complexified} gives
$\sum_{j\ne i}(\xi_i-\xi_j)^{-1}=x_i/2$.
For $m=2,3,4$, pair the summands in
\[
 \frac12\sum_i x_i\xi_i^m
 =\sum_{i<j}\frac{\xi_i^m-\xi_j^m}{\xi_i-\xi_j}.
\]
Using $\sum_i\xi_i=0$ gives respectively
\[
 \sum_i x_i\xi_i^2=0,\quad
 \sum_i x_i\xi_i^3=7\sum_i\xi_i^2,\quad
 \sum_i x_i\xi_i^4=6\sum_i\xi_i^3.
\]
Expansion of the product defining $R$ then yields the three necessary
polynomial equations
\begin{equation}\label{eq:n5-reconstructed-moment-equations}
 s_2=0,\qquad s_1-6r_3=0,\qquad s_0-2r_2=0.
\end{equation}
The monic polynomial of the dual coordinates is
\begin{equation}\label{eq:n5-reconstructed-dual}
 p_y(Y)=\frac{R(4Y,0)}{4^5\delta}
 =Y^5+\frac{r_3}{16\delta}Y^3+
       \frac{r_2}{64\delta}Y^2+
       \frac{r_1}{256\delta}Y+
       \frac{r_0}{1024\delta}.
\end{equation}
Only necessary equations are used to prove coverage. We do not infer
physical pairing merely from equality of two unordered root sets.

\subsection{Reflection symmetry, including all exceptional charts}
Suppose first $ab\ne0$. Choose $t$ with $t^2a=1$ and use the transformed
coefficients, still denoted by $a,b,c,d$. Put $B=b^2$ and $k=d/b$.
In \eqref{eq:n5-reconstructed-moment-equations}, substitute $a=1$,
$d=bk$, divide the first and third rows by their factor $b$, and replace
$b^2$ by $B$. Denote the resulting primitive rational polynomials by
$F_1,F_2,F_3$, and denote the transformed discriminant by
$\widetilde\delta(B,c,k)$.

The exact computation starts with these defined inputs, rather than a
preselected list of shapes. It computes a Gr\"obner basis of
\[
 \langle F_1,F_2,F_3,sB\widetilde\delta-1\rangle
       \subset\Q[s,k,c,B]
\]
in graded reverse lexicographic order, followed by exact FGLM conversion
to lexicographic order. Its $s$-free part is
\begin{equation}\label{eq:n5-reconstructed-generic-ideal}
 \langle P_1,c-C_1,k-K_1\rangle
 \cap\langle P_2,c-C_2,k-K_2\rangle,
\end{equation}
where
\begin{align*}
P_1&=125B^4+7900B^3-37040B^2+44608B+512,\\
P_2&=2500B^4-114000B^3+338925B^2-11000B+6912,\\
C_1&=\frac{3(125B^3+8020B^2-16496B+4224)}{57088},\\
K_1&=\frac{75B^3+4812B^2-16320B+13952}{28544},\\
C_2&=-\frac{1250B^3-69425B^2+327305B-145296}{443535},\\
K_2&=-\frac{28500B^3-1287200B^2+3884705B-1798816}{5913800}.
\end{align*}
The two quartics are coprime and their product is square-free.
Reduction in that same computed saturated basis gives
\begin{equation}\label{eq:n5-reconstructed-dual-invariant}
 \delta r_2^2=b^2r_3^3.
\end{equation}
The supplied script also checks the two further identities
$\delta r_1=cr_3^2$ and $\delta r_0=kr_3r_2$; they are not needed
for the following reality argument.

Let $a_y,b_y$ be the corresponding coefficients in
\eqref{eq:n5-reconstructed-dual}. Since $a=1$, physical conjugacy gives
$a_y=\kappa^2>0$ and $b_y=\kappa^3\overline b$. Consequently
\[
 \overline B=\frac{b_y^2}{a_y^3}
 =\frac{\delta r_2^2}{r_3^3}=B.
\]
The denominator is nonzero because $a_y>0$ and $\delta\ne0$.
Thus $B$ is real, and \eqref{eq:n5-reconstructed-generic-ideal} makes
$c,k$ real. If $B>0$, the polynomial has real coefficients. If $B<0$,
$b,d$ are purely imaginary; multiplying its roots by $\mathrm i$
makes all coefficients real. Undoing the original complex scale gives
reflection symmetry in a line through the origin.

The exceptional charts exhaust the complement of $ab\ne0$. For each,
use \eqref{eq:n5-reconstructed-moment-equations}, add the vanished
dual coefficients from \eqref{eq:n5-reconstructed-dual}, and saturate
by $\delta$. Exact lexicographic computation gives:
\begin{center}
\begin{tabular}{lll}
\toprule
Chart after normalization&Added dual rows&Collision-free consequence\\
\midrule
$a=1,b=0$&$r_2=0$&$d=0,\ c=3/20$\\
$a=0,b=1$&$r_3=0$&unit ideal; no physical solution\\
$a=b=0$&$r_3=r_2=0$&$cd=0$\\
\bottomrule
\end{tabular}
\end{center}
The first polynomial is $X(X^4+X^2+3/20)$; both nonzero values of
$X^2$ are negative, so its roots lie on a line. In the last chart,
if $c=0$, then $d\ne0$ and the roots form a regular pentagon; if
$d=0$, then $c\ne0$ and the roots form a square together with the
origin. These sets are reflection symmetric. This proves reflection
symmetry for every collision-free physical equilibrium; no
non-symmetric branch remains outside the chart partition.

\subsection{All physical branches after reflection}
Rotate back to moment five and make the reflection axis the real axis.
A real monic polynomial of odd degree has either one, three, or five
real roots, with the other roots in genuine nonreal conjugate pairs.
These three cases cover every reflection-symmetric configuration.

For one real point and two conjugate pairs, write
\[
 p(z)=(z-h)q_1(z)q_2(z),\quad
 q_j(z)=z^2+u_jz+v_j,
\]
where
\[
 u_1=\frac{h+d}{2},\quad u_2=\frac{h-d}{2},\qquad
 v_1=\frac{5-h^2+2e}{4},\quad v_2=\frac{5-h^2-2e}{4}.
\]
The parameters $h,d,e$ are real. The formulas impose zero centroid
and moment five, because $|z|^2=v_j$ for a root of a genuine
nonreal pair. At the root $h$, the force condition is
$p''(h)-4hp'(h)=0$. At a root of $q_j$ its conjugate is $-u_j-z$,
so the condition is that
\begin{equation}\label{eq:n5-reconstructed-remainders}
 p''+4(u_j+z)p'
\end{equation}
have zero remainder on division by $q_j$. These are five real
polynomial equations, obtained directly from the stationary equations.

Their exact lexicographic basis gives
\[
 h(h^2-1)Q(h^2)=0,\qquad 2de=5h(h^2-1),\qquad
 Q(H)=10H^4-65H^3+140H^2-97H+6.
\]
The full basis is retained in the verification log. Its specializations
give the following exhaustive cases:
\begin{enumerate}
\item $h=0$: $d^2+4e^2=10$ and $de=0$. If $d=0$, then
      $e^2=5/2$ and the four nonzero points lie on the imaginary
      axis, giving the collinear configuration. If $e=0$, then
      $d^2=10$ and the four points form a square around the origin.
      The two quadratic factors in both cases have distinct roots
      and no common root.
\item $h^2=1$: $d^2=5$ and $e=0$, giving the regular unit pentagon.
\item $Q(H)=0$, $H=h^2$: the remaining equations give
\[
 d^2=J(H)=\frac{-350H^3+1775H^2-2430H+975}{92},\qquad
 e=\frac{5h(H-1)}{2d}.
\]
Sturm counting gives exactly two real roots of $Q$, both positive.
Their rational isolating intervals are
\[
 I_- =\left(\frac{68395128919}{10^{12}},
             \frac{68395128920}{10^{12}}\right),\qquad
 I_+ =\left(\frac{1214112777579}{10^{12}},
             \frac{1214112777580}{10^{12}}\right).
\]
On each interval $J(H)>0$. The sum of the two quadratic
discriminants is $(5H+J(H)-20)/2<0$ and their product is positive;
thus both pairs are genuinely nonreal. Also $q_1(h)q_2(h)>0$ and
$\operatorname{Res}(q_1,q_2)>0$, so no collision occurs. Direct
substitution, followed by reduction modulo $Q(H)$ and
$d^2-J(H),h^2-H$, verifies all five original equations.
Changing the sign of $d$ exchanges the two pairs, while changing
the sign of $h$ gives the same shape after a rotation and an exchange
of pairs. Hence these two roots contribute exactly two shapes.
\end{enumerate}

For three real points and one nonreal pair, write
\[
 p(z)=\bigl(z^3-uz^2+bz+c\bigr)(z^2+uz+v),
 \qquad b=\frac{u^2+2v-5}{2}.
\]
Use zero remainders of $p''-4zp'$ modulo the cubic and of
$p''+4(u+z)p'$ modulo the quadratic. Exact elimination yields
\[
 uQ_{32}(u^2)=0,\qquad
 Q_{32}(U)=10U^4-55U^3+115U^2-37U-9.
\]
If $u\ne0$, then
\[
 v=V_{32}(u^2),\qquad
 V_{32}(U)=-\frac{130U^3-625U^2+1074U-211}{304}.
\]
Sturm counting gives exactly one positive root of $Q_{32}$, in
$(607600/10^6,607601/10^6)$, and rational interval arithmetic gives
$V_{32}<0$ there. This contradicts the necessary condition
$v>u^2/4\ge0$ for a nonreal pair. For $u=0$, the full basis gives
$c=0,v=5/4$, which is the square and its center already listed.

For five real points, each root satisfies $p''-4zp'=0$.
The polynomial $p''-4zp'+20p$ has degree at most four and vanishes
at five distinct roots; therefore it is identically zero. Comparing
coefficients in its unique monic degree-five solution gives
\[
 p(z)=z^5-\frac52z^3+\frac{15}{16}z.
\]
Its roots are $0$ and
$\pm\sqrt{(5+\sqrt{10})/4},\ \pm\sqrt{(5-\sqrt{10})/4}$.
They are distinct and real, and their total squared modulus is five.
Thus this case has exactly the collinear shape already listed.

\subsection{Exact comparison and computational scope}
For a monic polynomial, the objective is
$|\operatorname{Disc}(p)|$, rather than its square. The three
classical shapes give respectively
\[
 3125,\qquad \frac{5^{10}}{2^{12}},\qquad \frac{84375}{1024}.
\]
For the two other shapes, the product identity
\[
 \operatorname{Disc}(p)=
 \operatorname{Disc}(q_1)\operatorname{Disc}(q_2)
 [q_1(h)q_2(h)]^2\operatorname{Res}(q_1,q_2)^2
\]
reduces modulo $Q$ to
\[
 D(H)=-\frac{84375}{385875968}
       (7632290H^3-30197665H^2+32980990H-12854393).
\]
Rational Horner enclosures on $I_-,I_+$ prove
\[
 2347<D(H_-)<2348,\qquad 801<D(H_+)<802.
\]
All four competitors are strictly below $3125$. Every maximizer
belongs to one of the five covered shapes, which proves
Theorem~\ref{thm:n5-reconstructed-sharp} and its equality statement.

The computation supporting the generic and exceptional ideal statements
and the two reflection branch eliminations is exact rational
Gr\"obner-basis computation in the supplied
\path{n5_complete_verification.py}. The wrapper
\path{n5_root_runner.py} requires enabled assertions, reruns that
local script without fetching code, saves its complete log, and exports
the exact input and basis coefficients. The independent standard-library
checker \path{n5_physical_audit.py} verifies Sturm isolation,
physical signs, discriminant factor reduction, all classical values,
and the strict comparisons. This smaller checker alone does not prove
the ideal coverage; the preceding exact reconstruction supplies it.
The full local reconstruction was rerun successfully on 6 October 2026
with assertions enabled, exporting ten Gr\"obner computations. The
independent physical checker passed all 28 effective checks both normally
and under \texttt{-B -S -O}.

\section{Verification and unresolved work}
The accompanying standard-library checkers verify formal quadratic-form
identities, exact trace constants, the rational annulus estimates, small
complete coefficient moments and translated Sylvester resultants.
They use explicit exception checks, including under Python optimization.
These finite arithmetic audits do not replace the general analytic proofs.
The five-point full symbolic classification was rerun with assertions
enabled, reconstructing ten exact rational Gr\"obner computations.
Its separate physical/value checker passed all $28$ deciding tests.
The six- and seven-point disk certificates were independently replayed
under Python optimization: $80$ and $174$ parameter candidates and $480$
and $1218$ point roots respectively. The six-point component containment
checker now uses explicit arithmetic exceptions. The external coverage
dependencies, including the missing degree-$356$ seven-point exclusion,
are retained.
The three standard-library checkers from the primary five-point manuscript
were replayed successfully: the first scalar branch has a complete
$183$-node, $92$-leaf cover, while the second has a uniform rational bound.
The new six-point exceptional-chart certificate reconstructs $1176$ exact
grid values with proved degree bounds and verifies original-equation
combination identities. The new eight-point three-level exclusions use
formal polynomial and rational arithmetic checks.

The new fixed-origin theorem replaces all remaining origin-branch
exclusions: every origin-containing eight-point configuration is now
bounded sharply, including arbitrary radial and angular arrangements.
Its scalar step has a complete exact cover of 1102 leaf rectangles,
independently replayed, and no unresolved boxes. The original annulus,
radial-level, and origin coefficient results remain useful cross-checks.
They are no longer required to cover the origin branch.

The unresolved unrestricted branch has no point at the origin.
Earlier subset and radial arguments show that a strictly better global
competitor has at least three radial levels. With exactly three levels,
all radial-group centroid sums are nonzero and collinear, and the
$1+1+6$ split is impossible. Those structural exclusions are not a
proof for arbitrary no-origin configurations. A classification-free
alternative is to prove the centroid compensation inequality
\eqref{eq:noclass-centroid-target}. The new frame, finite-rank, and
high-order Gram identities preserve angular information for this task,
but the required uniform scalar deficit has not been established.
For six and seven points, the complete no-classification replacement
of the retained published coverage is likewise not yet proved.

For general $n$, the fixed moment
selection is an exact finite specification but its universal polynomial
enumeration and moment expansion were not executed. A useful computational
breakthrough must reduce these costs. A simple finite expression for every
$M_n$ and a classification of every attaining shape remain unresolved here.

\begin{thebibliography}{9}
\bibitem{Wang}
J.~Wang, \emph{The Five-Point Case of Mordell's Discriminant Inequality},
arXiv:2610.00358 (2026).
\url{https://arxiv.org/abs/2610.00358}.
\bibitem{JiangFive}
F.~Jiang, \emph{A second-order and moment proof of the five-point Mordell
inequality}, user-supplied local manuscript, 3 October 2026.
The source and exact ancillary certificate are included in the verification
package; no public submission or priority claim is inferred.
\bibitem{Mordell}
L.~J.~Mordell, \emph{On a discriminant inequality},
Canadian Journal of Mathematics \textbf{12} (1960), 699--704.
\href{https://doi.org/10.4153/CJM-1960-063-5}{doi:10.4153/CJM-1960-063-5}.
\bibitem{Chalk}
J.~H.~H.~Chalk, \emph{A note on a discriminant inequality},
Canadian Mathematical Bulletin \textbf{3} (1960), no.~1, 7--9.
\href{https://doi.org/10.4153/CMB-1960-002-x}{doi:10.4153/CMB-1960-002-x}.
\bibitem{FS}
J.-C.~Faug\`ere and J.~Svartz,
\emph{Solving polynomial systems globally invariant under an action of
the symmetric group and application to the equilibria of $N$ vortices
in the plane}, Proceedings of ISSAC 2012, 170--178.
\href{https://doi.org/10.1145/2442829.2442856}{doi:10.1145/2442829.2442856}.
\bibitem{FSdata}
J.-C.~Faug\`ere and J.~Svartz, author computational archive,
\emph{Vortices for $N=6$}, exact coefficient bases supplied with the
manuscript as \path{archivepolsys_vortices_N6.txt}.
\bibitem{FSdataSeven}
J.-C.~Faug\`ere and J.~Svartz, author computational archive,
\emph{Vortices for $N=7$}, retained exact coefficient bases supplied as
\path{archivepolsys_vortices_N7.txt}.
\bibitem{JPT}
G.~Jeronimo, D.~Perrucci and E.~Tsigaridas,
\emph{On the minimum of a polynomial function on a basic closed
semialgebraic set and applications},
SIAM Journal on Optimization \textbf{23} (2013), no.~1, 241--255.
\href{https://doi.org/10.1137/110857751}{doi:10.1137/110857751}.
\end{thebibliography}
\end{document}
